\documentclass[11pt]{article}
\usepackage[a4paper,margin=24mm]{geometry}
\usepackage{fontspec}
\usepackage[AutoFallBack=true]{xeCJK}
\setmainfont{Noto Serif}
\setCJKmainfont{Noto Serif CJK SC}
\setCJKfallbackfamilyfont{rm}{SimSun-ExtB}
\xeCJKDeclareCharClass{CJK}{"25A1, "25CB}
\usepackage{xcolor}
\usepackage[hidelinks]{hyperref}
\setlength{\parindent}{0pt}
\setlength{\parskip}{0.4em}
\newcommand{\zhsmall}{\footnotesize\color{black!70}}
\newcommand{\unitlabel}[1]{\marginpar{\raggedleft\tiny\color{gray}#1}}
\newenvironment{unitzh}{\par\CJKfontspec{Noto Serif CJK SC}[AutoFallBack=true]\large\linespread{1.25}\selectfont}{\par}
\newenvironment{uniten}{\par\small\color{black!85}}{\par\medskip}
\sloppy
\begin{document}
\begin{titlepage}\centering
{\fontsize{34}{40}\selectfont 測圓海鏡}\par\vspace{1em}
{\LARGE Sea Mirror of Circle Measurements}\par\vspace{0.6em}
{\large by Li Ye (李冶), 1248}\par\vspace{2em}
{\normalsize Chinese text with modern punctuation, and a complete English translation}\par\vfill
\begin{minipage}{0.86\textwidth}\small
Chinese text: Siku quanshu (四庫全書) edition, as transcribed on Chinese Wikisource (the pages and revisions used are listed at the end).
Punctuation, division into units and English translation: AI-produced by Claude Opus 5.5 (Anthropic), 24--25 October 2026.
The punctuated Chinese was checked by program against the Wikisource text: every character is reproduced, in order,
and small print (commentary) stays small print. In the translation, commentary is given in square brackets with its
author. The labels in the margin identify the units of the machine-readable file units.jsonl.
\end{minipage}
\end{titlepage}
\tableofcontents\clearpage
\section*{測圓海鏡}\addcontentsline{toc}{section}{測圓海鏡}

\begin{unitzh}\unitlabel{ceyuan-001.1}欽定四庫全書·子部六\end{unitzh}
\begin{uniten}Imperially Commissioned Complete Library of the Four Treasuries (Qinding Siku quanshu). Masters section (zibu), 6.\end{uniten}

\begin{unitzh}\unitlabel{ceyuan-001.2}測圓海鏡·天文算法類二{\zhsmall 算書之屬}·提要\end{unitzh}
\begin{uniten}Sea Mirror of Circle Measurements (Ceyuan haijing). Category of astronomy and mathematical methods (tianwen suanfa lei), 2 [Note: class of mathematical books]. Summary (tiyao).\end{uniten}

\begin{unitzh}\unitlabel{ceyuan-001.3}{\zhsmall 臣}等謹案：《測圓海鏡》十二卷，元李冶撰。冶字鏡齋，欒城人，金末登進士，入元官翰林學士，事蹟具《元史》本傳。其書以勾股容圓為題，自圓心、圓外縱横取之，得大小十五形，皆無竒零。次列識别雜記數百條，以窮其理；次設問一百七十則，以盡其用。探賾索隱，參伍錯綜，雖習其法者不能驟解。而其草多言立天元一。按：立天元一法見扵宋秦九韶《九章》大衍術中，厥後《授時草》及《四元玉鑑》等書皆屢見之，而此書言之獨詳，其法關乎數學者甚大。\end{unitzh}
\begin{uniten}[Your servant] et al. respectfully note: The Sea Mirror of Circle Measurements (Ceyuan haijing), in 12 juan, was composed by Li Ye of the Yuan. Ye's style name was Jingzhai; he was a man of Luancheng. At the end of the Jin he passed the jinshi examination; after entering the Yuan he served as Hanlin Academician; his deeds are fully recorded in his biography in the History of the Yuan (Yuanshi). His book takes "the circle contained in base and altitude" (gougu rong yuan) as its subject: taking from it lengthwise and crosswise, from the center of the circle and from outside the circle, it obtains 15 figures, large and small, all without odd remainders (jiling). Next it lists several hundred items of identifications and miscellaneous notes (shibie zaji), in order to exhaust its principles; next it sets 170 problems, in order to exhaust its uses. It probes the abstruse and searches out the hidden, combining by threes and fives and interweaving, so that even those versed in its methods cannot understand it at once. And its workings (cao) mostly speak of setting up the celestial unknown (li tianyuan yi). Note that the method of setting up the celestial unknown appears in the great extension method (dayan shu) of the Nine Chapters (Jiuzhang) by Qin Jiushao of the Song; thereafter it appears repeatedly in books such as the Shoushi cao (Workings of the Season-Granting System) and the Siyuan yujian (Jade Mirror of the Four Unknowns), but this book alone discusses it in detail; its method is of very great consequence for mathematics.\end{uniten}

\begin{unitzh}\unitlabel{ceyuan-001.4}然自元以來，疇人皆株守立成，習而不察，至遂無知其法者。故唐順之與頋應祥書稱「立天元一漫不省為何語」；頋應祥演是書為《分類釋術》，其自序亦云「立天元一無下手之術」，則是書雖存而其傳已泯矣。明萬厯中，利瑪竇與徐光啟、李之𧁑等譯為《同文算指》諸書，扵古《九章》皆有辨訂，獨於立天元一法闕而不言。徐光啟扵《勾股義》序中引此書，又謂「欲説其義而未遑」，是此書已為利瑪竇所見，而猶未得其解也。\end{unitzh}
\begin{uniten}However, since the Yuan the mathematical specialists (chouren) have all clung stubbornly to ready-made tables (licheng), practicing them without examining them, until in the end there was no one who knew this method. Thus Tang Shunzhi, in a letter to Gu Yingxiang, says: "Setting up the celestial unknown — I cannot at all make out what these words mean"; Gu Yingxiang expanded this book into the Fenlei shishu (Methods Explained by Category), and in his own preface he too says: "For setting up the celestial unknown there is no procedure by which to set to work"; so, although this book survived, its transmission had already died out. In the Wanli period of the Ming, Matteo Ricci, with Xu Guangqi, Li Zhizao and others, produced in translation the Tongwen suanzhi (Guide to Calculation in the Common Script) and other books; on the ancient Nine Chapters they in every case gave critical discussion and emendation, but the method of setting up the celestial unknown alone they left out and did not speak of. Xu Guangqi, in his preface to the Gougu yi (Meaning of Base and Altitude), cites this book, and further says that he "wished to explain its meaning but had not had the leisure"; this shows that this book had already been seen by Matteo Ricci, yet its explanation had still not been obtained.\end{uniten}

\begin{unitzh}\unitlabel{ceyuan-001.5}迨我國朝醲化翔洽，梯航鱗萃，歐邏巴人始以借根方法進呈聖祖仁皇帝，授蒙養齋諸臣習之。梅㲄成乃悟即古立天元一法，扵《赤水遺珍》中詳解之，且載西名阿爾熱巴拉{\zhsmall 案：原本作阿爾熱巴逹，謹據西洋借根法改正。}，即華言東來法，知即冶之遺書流入西域，又轉而還入中原也。今用以勘騐西法，一一脗合，瑴成所説，信而有徴。特錄存之，以為算法之秘鑰，且以見中法西法互發益明，無容設畛域之見焉。乾隆四十六年二月恭校上。\end{unitzh}
\begin{uniten}When it came to our dynasty, whose rich civilizing influence soared and spread everywhere, while those who came by ladder and by boat thronged together like fish scales, the Europeans (Ouluoba ren) for the first time presented the method of borrowed roots and powers (jiegenfang fa) to the Emperor Shengzu Ren, who gave it to the officials of the Mengyang Studio (Mengyang zhai) to study. Mei Juecheng then realized that it was precisely the ancient method of setting up the celestial unknown; he explained it in detail in the Chishui yizhen (Pearls Left at the Red Water), and moreover recorded its Western name, A'errebala [Siku editors: the original text writes A'errebada; respectfully corrected on the basis of the Western borrowed-root method (jiegen fa).], which in Chinese means "the method that came from the East"; from this it is known that it is Ye's surviving book, which flowed into the Western Regions and then turned round and came back into the Central Plain. Now, when it is used to check the Western method, they agree point by point: what Juecheng said is trustworthy and has proof. We have specially copied and preserved it as a secret key to the methods of calculation, and also to show that the Chinese method and the Western method illuminate each other and become ever clearer, so that there is no room for views that set up boundaries. Respectfully collated and submitted in the 2nd month of the 46th year of Qianlong.\end{uniten}

\begin{unitzh}\unitlabel{ceyuan-001.6}總纂官{\zhsmall 臣}紀昀、{\zhsmall 臣}陸錫熊、{\zhsmall 臣}孫士毅。總校官{\zhsmall 臣}陸費墀。\end{unitzh}
\begin{uniten}Chief compilers (zongzuan guan): [your servant] Ji Yun, [your servant] Lu Xixiong, [your servant] Sun Shiyi. Chief collator (zongjiao guan): [your servant] Lufei Chi.\end{uniten}

\begin{unitzh}\unitlabel{ceyuan-001.7}原序\end{unitzh}
\begin{uniten}Original preface\end{uniten}

\begin{unitzh}\unitlabel{ceyuan-001.8}數本難窮，吾欲以力强窮之，彼其數不惟不能得其凡，而吾之力且憊矣。然則數果不可以窮耶？既已名之數矣，則又何為而不可窮也？故謂數為難窮，斯可；謂數為不可窮，斯不可。何則？彼其冥冥之中，固有昭昭者存。夫昭昭者，其自然之數也；非自然之數，其自然之理也。數一出於自然，吾欲以力强窮之，使隸首復生，亦末如之何也已。苟能推自然之理，以明自然之數，則雖逺而乾端坤倪，幽而神情鬼状，未有不合者矣。\end{unitzh}
\begin{uniten}Numbers are by nature hard to exhaust. If I wish to exhaust them by forcing my strength, then not only shall I be unable to grasp their general outline (fan), but my strength will moreover be worn out. That being so, can numbers truly not be exhausted? Since they have already been named "numbers", why then should they be inexhaustible? Therefore, to say that numbers are hard to exhaust is acceptable; to say that numbers cannot be exhausted is not acceptable. Why so? Within their darkness there is certainly something bright. What is bright is their natural numbers; what is not natural numbers is their natural principles. Numbers issue entirely from nature: if I wish to exhaust them by force, then even if Li Shou were to live again, he too could do nothing about it. If one can extend the natural principles so as to elucidate the natural numbers, then even what is as remote as the ends of Heaven and the bounds of Earth, or as hidden as the feelings of spirits and the forms of ghosts, there is nothing that will not accord.\end{uniten}

\begin{unitzh}\unitlabel{ceyuan-001.9}予自幼喜算數，恒病夫考圓之術，例出於牽强，殊乖于自然：如古率、徽率、密率之不同，截弧、截矢、截背之互見，内外諸角，析剖支條，莫不各自名家，與世作法。及反覆研究，卒無以當吾心焉。老大以來，得洞淵九容之說，日夕玩繹，而向之病我者，使爆然落去而無遺餘。山中多暇，客有從余求其說者，於是乎又為衍之，遂累一百七十問。既成編，客復目之《測圓海鏡》，蓋取夫天臨海鏡之義也。\end{unitzh}
\begin{uniten}From my youth I have delighted in calculation, and I was constantly vexed that the methods for investigating the circle as a rule proceed from forced contrivance and run quite contrary to what is natural: such as the differences among the ancient rate (gu lü), Hui's rate (Hui lü) and the close rate (mi lü); the cut arc (jiehu), the cut sagitta (jieshi) and the cut back (jiebei) appearing alongside one another; the various inner and outer corners (jiao), analyzed and dissected into branches and strands — none but set itself up as a school of its own and made methods for the world. When I came to study them over and over, in the end none could satisfy my mind. Since growing old I have obtained Dongyuan's theory of the nine containments (jiu rong); day and night I savored and unraveled it, and what had previously vexed me was made to drop away in a burst, leaving nothing behind. In the mountains there was much leisure, and among my guests there were those who sought this theory from me; thereupon I further expanded it for them, and so accumulated 170 problems. When it had been completed as a volume, a guest further titled it Sea Mirror of Circle Measurements, taking, it seems, the meaning of "heaven looking down on the mirror of the sea".\end{uniten}

\begin{unitzh}\unitlabel{ceyuan-001.10}昔半山老人集《唐百家詩選》，自謂「廢日力于此，良可惜」。明道先生以上蔡謝君記誦為「玩物喪志」。夫文史尚矣，猶之為不足貴，況九九賤技能乎？嗜好酸鹹，平生每痛自戒勅，竟莫能已，類有物慿之者，吾亦不知其然而然也。故嘗私為之解曰：由技兼于事者言之，夷之禮、夔之樂，亦不免為一技；由技進乎道者言之，石之斤、扁之輪，非聖人之所與乎？覽吾之編，察吾苦心，其憫我者當百數，其笑我者當千數。乃若吾之所自得，則自得焉耳，寧復為人憫笑計哉！李冶序。\end{unitzh}
\begin{uniten}Formerly the Old Man of Banshan compiled the Selection of Poems by a Hundred Tang Poets (Tang baijia shixuan) and said of himself: "Wasting the strength of my days on this is truly a pity." Master Mingdao regarded the memorizing and reciting of Mr. Xie of Shangcai as "trifling with things and losing one's purpose". Now, literature and history are esteemed indeed, and yet even they were regarded as not worth valuing; how much more the lowly skill of the nine nines! My predilection is like a taste for the sour or the salty: all my life I have often bitterly admonished myself, yet in the end I could not stop; it is as if something possessed me, and I myself do not know why it is so, yet so it is. Therefore I once privately offered an explanation of it, saying: speaking of skill as comprehended in affairs, even the rites of Yi and the music of Kui cannot escape being a skill; speaking of skill as advancing into the Way, are not Shi's axe and Bian's wheel what the sages approve of? Of those who read my compilation and examine my painstaking efforts, those who pity me will number in the hundreds, and those who laugh at me will number in the thousands. As for what I have gained for myself, I have gained it for myself and that is all; why should I further take account of people's pity or laughter! Preface by Li Ye.\end{uniten}

\begin{unitzh}\unitlabel{ceyuan-001.11}全覽。卷一、卷二、卷三、卷四、卷五、卷六、卷七、卷八、卷九、卷十、卷十一、卷十二。\end{unitzh}
\begin{uniten}Full view. Juan 1; Juan 2; Juan 3; Juan 4; Juan 5; Juan 6; Juan 7; Juan 8; Juan 9; Juan 10; Juan 11; Juan 12.\end{uniten}

\section*{卷01}\addcontentsline{toc}{section}{卷01}

\begin{unitzh}\unitlabel{ceyuan-002.1}總率名號\end{unitzh}
\begin{uniten}Names of the General Rates (zonglü)\end{uniten}

\begin{unitzh}\unitlabel{ceyuan-002.2}天之地為通弦，天之乾為通股，乾之地為通勾。天之川為邊弦，天之西為邊股，西之川為邊勾。日之地為底弦，日之北為底股，北之地為底勾。天之山為黄廣弦，天之金為股，金之山為勾。月之地為黄長弦，月之泉為股，泉之地為勾。天之日為上髙弦，天之旦為股，旦之日為勾。日之山為下髙弦，日之朱為股，朱之山為勾。月之川為上平弦，月之青為股，青之川為勾。川之地為下平弦，川之夕為股，夕之地為勾。天之月為大差弦，天之坤為股，坤之月為勾。山之地為小差弦，山之艮為股，艮之地為勾。日之川為皇極弦，日之心為股，心之川為勾。月之山為太虚弦，月之水為股，水之山為勾。日之月為明弦，日之南為股，南之月為勾。山之川為□弦，山之東為股，東之川為勾。\end{unitzh}
\begin{uniten}Tian (天) to Di (地) is the through hypotenuse (tongxian); Tian to Qian (乾) is the through altitude (tonggu); Qian to Di is the through base (tonggou). Tian to Chuan (川) is the side hypotenuse (bianxian); Tian to Xi (西) is the side altitude (biangu); Xi (西) to Chuan is the side base (biangou). Ri (日) to Di is the bottom hypotenuse (dixian); Ri to Bei (北) is the bottom altitude (digu); Bei to Di is the bottom base (digou). Tian to Shan (山) is the yellow-breadth hypotenuse (huangguang xian); Tian to Jin (金) is the altitude (gu); Jin to Shan is the base (gou). Yue (月) to Di is the yellow-length hypotenuse (huangchang xian); Yue to Quan (泉) is the altitude; Quan to Di is the base. Tian to Ri is the upper height hypotenuse (shanggao xian); Tian to Dan (旦) is the altitude; Dan to Ri is the base. Ri to Shan is the lower height hypotenuse (xiagao xian); Ri to Zhu (朱) is the altitude; Zhu to Shan is the base. Yue to Chuan is the upper level hypotenuse (shangping xian); Yue to Qing (青) is the altitude; Qing to Chuan is the base. Chuan to Di is the lower level hypotenuse (xiaping xian); Chuan to Xi (夕) is the altitude; Xi (夕) to Di is the base. Tian to Yue is the great-difference hypotenuse (dacha xian); Tian to Kun (坤) is the altitude; Kun to Yue is the base. Shan to Di is the small-difference hypotenuse (xiaocha xian); Shan to Gen (艮) is the altitude; Gen to Di is the base. Ri to Chuan is the imperial-pole hypotenuse (huangji xian); Ri to Xin (心) is the altitude; Xin to Chuan is the base. Yue to Shan is the great-void hypotenuse (taixu xian); Yue to Shui (水) is the altitude; Shui to Shan is the base. Ri to Yue is the bright hypotenuse (mingxian); Ri to Nan (南) is the altitude; Nan to Yue is the base. Shan to Chuan is the □ hypotenuse (□ xian) [Translator's note: □ stands for a character left blank in the transcription]; Shan to Dong (東) is the altitude; Dong to Chuan is the base.\end{uniten}

\begin{unitzh}\unitlabel{ceyuan-002.3}今問正數\end{unitzh}
\begin{uniten}Correct numbers of the present question\end{uniten}

\begin{unitzh}\unitlabel{ceyuan-002.4}通弦六百八十，勾三百二十，股六百。勾股和九百二十，較二百八十；勾弦和一千，較三百六十；股弦和一千二百八十，較八十；弦較和九百六十，較四百；弦和和一千六百，較二百四十。\end{unitzh}
\begin{uniten}Through hypotenuse 680; base 320; altitude 600. Sum (he) of base and altitude 920, difference (jiao) 280; sum of base and hypotenuse 1000, difference 360; sum of altitude and hypotenuse 1280, difference 80; sum of hypotenuse and difference (xianjiao he) 960, difference 400; sum of hypotenuse and sum (xianhe he) 1600, difference 240.\end{uniten}

\begin{unitzh}\unitlabel{ceyuan-002.5}邊弦五百四十四，勾二百五十六，股四百八十。勾股和七百三十六，較二百二十四；勾弦和八百，較二百八十八；股弦和一千零二十四，較六十四；弦較和七百六十八，較三百二十；弦和和一千二百八十，較一百九十二。\end{unitzh}
\begin{uniten}Side hypotenuse 544; base 256; altitude 480. Sum of base and altitude 736, difference 224; sum of base and hypotenuse 800, difference 288; sum of altitude and hypotenuse 1024, difference 64; sum of hypotenuse and difference 768, difference 320; sum of hypotenuse and sum 1280, difference 192.\end{uniten}

\begin{unitzh}\unitlabel{ceyuan-002.6}底弦四百二十五，勾二百，股三百七十五。勾股和五百七十五，較一百七十五；勾弦和六百二十五，較二百二十五；股弦和八百，較五十；弦較和六百，較二百五十；弦和和一千，較一百五十。\end{unitzh}
\begin{uniten}Bottom hypotenuse 425; base 200; altitude 375. Sum of base and altitude 575, difference 175; sum of base and hypotenuse 625, difference 225; sum of altitude and hypotenuse 800, difference 50; sum of hypotenuse and difference 600, difference 250; sum of hypotenuse and sum 1000, difference 150.\end{uniten}

\begin{unitzh}\unitlabel{ceyuan-002.7}黄廣弦五百一十，勾二百四十{\zhsmall 即城徑也}，股四百五十。勾股和六百九十，較二百一十；勾弦和七百五十，較二百七十；股弦和九百六十，較六十；弦較和七百二十，較三百；弦和和一千二百，較一百八十。\end{unitzh}
\begin{uniten}Yellow-breadth hypotenuse 510; base 240 [Note: this is the city diameter (chengjing)]; altitude 450. Sum of base and altitude 690, difference 210; sum of base and hypotenuse 750, difference 270; sum of altitude and hypotenuse 960, difference 60; sum of hypotenuse and difference 720, difference 300; sum of hypotenuse and sum 1200, difference 180.\end{uniten}

\begin{unitzh}\unitlabel{ceyuan-002.8}黄長弦二百七十二，勾一百二十八，股二百四十{\zhsmall 即城徑也}。勾股和三百六十八，較一百一十二；勾弦和四百，較一百四十四；股弦和五百一十二，較三十二；弦較和三百八十四，較一百六十；弦和和六百四十，較九十六。\end{unitzh}
\begin{uniten}Yellow-length hypotenuse 272; base 128; altitude 240 [Note: this is the city diameter]. Sum of base and altitude 368, difference 112; sum of base and hypotenuse 400, difference 144; sum of altitude and hypotenuse 512, difference 32; sum of hypotenuse and difference 384, difference 160; sum of hypotenuse and sum 640, difference 96.\end{uniten}

\begin{unitzh}\unitlabel{ceyuan-002.9}髙弦二百五十五{\zhsmall 上下同}，勾一百二十{\zhsmall 即半徑}，股二百二十五。勾股和三百四十五，較一百零五；勾弦和三百七十五，較一百三十五；股弦和四百八十，較三十；弦較和三百六十，較一百五十；弦和和六百，較九十。\end{unitzh}
\begin{uniten}Height hypotenuse (gaoxian) 255 [Note: the upper and the lower are the same]; base 120 [Note: this is the radius (banjing)]; altitude 225. Sum of base and altitude 345, difference 105; sum of base and hypotenuse 375, difference 135; sum of altitude and hypotenuse 480, difference 30; sum of hypotenuse and difference 360, difference 150; sum of hypotenuse and sum 600, difference 90.\end{uniten}

\begin{unitzh}\unitlabel{ceyuan-002.10}平弦一百三十六{\zhsmall 上下同}，勾六十四，股一百二十{\zhsmall 即半徑也}。勾股和一百八十四，較五十六；勾弦和二百，較七十二；股弦和二百五十六，較十六；弦較和一百九十二，較八十；弦和和三百二十，較四十八。\end{unitzh}
\begin{uniten}Level hypotenuse (pingxian) 136 [Note: the upper and the lower are the same]; base 64; altitude 120 [Note: this is the radius]. Sum of base and altitude 184, difference 56; sum of base and hypotenuse 200, difference 72; sum of altitude and hypotenuse 256, difference 16; sum of hypotenuse and difference 192, difference 80; sum of hypotenuse and sum 320, difference 48.\end{uniten}

\begin{unitzh}\unitlabel{ceyuan-002.11}大差弦四百零八，勾一百九十二，股三百六十。勾股和五百五十二，較一百六十八；勾弦和六百，較二百一十六；股弦和七百六十八，較四十八；弦較和五百七十六，較二百四十；弦和和九百六十，較一百四十四。\end{unitzh}
\begin{uniten}Great-difference hypotenuse 408; base 192; altitude 360. Sum of base and altitude 552, difference 168; sum of base and hypotenuse 600, difference 216; sum of altitude and hypotenuse 768, difference 48; sum of hypotenuse and difference 576, difference 240; sum of hypotenuse and sum 960, difference 144.\end{uniten}

\begin{unitzh}\unitlabel{ceyuan-002.12}小差弦一百七十，勾八十，股一百五十。勾股和二百三十，較七十；勾弦和二百五十，較九十；股弦和三百二十，較二十；弦較和二百四十，較一百；弦和和四百，較六十。\end{unitzh}
\begin{uniten}Small-difference hypotenuse 170; base 80; altitude 150. Sum of base and altitude 230, difference 70; sum of base and hypotenuse 250, difference 90; sum of altitude and hypotenuse 320, difference 20; sum of hypotenuse and difference 240, difference 100; sum of hypotenuse and sum 400, difference 60.\end{uniten}

\begin{unitzh}\unitlabel{ceyuan-002.13}皇極弦二百八十九，勾一百三十六，股二百五十五。勾股和三百九十一，較一百一十九；勾弦和四百二十五，較一百五十三；股弦和五百四十四，較三十四；弦較和四百零八，較一百七十；弦和和六百八十，較一百零二。\end{unitzh}
\begin{uniten}Imperial-pole hypotenuse 289; base 136; altitude 255. Sum of base and altitude 391, difference 119; sum of base and hypotenuse 425, difference 153; sum of altitude and hypotenuse 544, difference 34; sum of hypotenuse and difference 408, difference 170; sum of hypotenuse and sum 680, difference 102.\end{uniten}

\begin{unitzh}\unitlabel{ceyuan-002.14}太虚弦一百零二，勾四十八，股九十。勾股和一百三十八，較四十二；勾弦和一百五十，較五十四；股弦和一百九十二，較一十二；弦較和一百四十四，較六十；弦和和二百四十，較三十六。\end{unitzh}
\begin{uniten}Great-void hypotenuse 102; base 48; altitude 90. Sum of base and altitude 138, difference 42; sum of base and hypotenuse 150, difference 54; sum of altitude and hypotenuse 192, difference 12; sum of hypotenuse and difference 144, difference 60; sum of hypotenuse and sum 240, difference 36.\end{uniten}

\begin{unitzh}\unitlabel{ceyuan-002.15}明弦一百五十三，勾七十二，股一百三十五。勾股和二百零七，較六十三；勾弦和二百二十五，較八十一；股弦和二百八十八，較一十八；弦較和二百一十六，較九十；弦和和三百六十，較五十四。\end{unitzh}
\begin{uniten}Bright hypotenuse 153; base 72; altitude 135. Sum of base and altitude 207, difference 63; sum of base and hypotenuse 225, difference 81; sum of altitude and hypotenuse 288, difference 18; sum of hypotenuse and difference 216, difference 90; sum of hypotenuse and sum 360, difference 54.\end{uniten}

\begin{unitzh}\unitlabel{ceyuan-002.16}□弦三十四，勾十六，股三十。勾股和四十六，較一十四；勾弦和五十，較一十八；股弦和六十四，較四；弦較和四十八，較二十；弦和和八十，較十二。\end{unitzh}
\begin{uniten}□ hypotenuse 34; base 16; altitude 30. Sum of base and altitude 46, difference 14; sum of base and hypotenuse 50, difference 18; sum of altitude and hypotenuse 64, difference 4; sum of hypotenuse and difference 48, difference 20; sum of hypotenuse and sum 80, difference 12.\end{uniten}

\begin{unitzh}\unitlabel{ceyuan-002.17}識别雜記\end{unitzh}
\begin{uniten}Miscellaneous Notes on Identification\end{uniten}

\begin{unitzh}\unitlabel{ceyuan-002.18}天之于日與日之於心同，心之于川與川之于地同，日之于心與日之于山同，故以山之川為小差。川之于心與川之于月同，故以月之日為大差。\end{unitzh}
\begin{uniten}Tian to Ri is the same as Ri to Xin; Xin to Chuan is the same as Chuan to Di; Ri to Xin is the same as Ri to Shan; therefore Shan to Chuan is taken as the small difference (xiaocha). Chuan to Xin is the same as Chuan to Yue; therefore Yue to Ri is taken as the great difference (dacha).\end{uniten}

\begin{unitzh}\unitlabel{ceyuan-002.19}明勾□股相得，名為内率，求虚積；明股□勾相得，名為外率，求虛積；虛勾虚股相得，名為虚率，求虚積。\end{unitzh}
\begin{uniten}The bright base (minggou) and the □ altitude brought together (xiangde) are named the inner rate (neilü), for finding the void area (xuji); the bright altitude (minggu) and the □ base brought together are named the outer rate (wailü), for finding the void area; the void base (xugou) and the void altitude (xugu) brought together are named the void rate (xulü), for finding the void area.\end{uniten}

\begin{unitzh}\unitlabel{ceyuan-002.20}凡勾股和即弦黄和；凡大差即股黄較；凡小差即勾黄較。\end{unitzh}
\begin{uniten}Every sum of base and altitude is the sum of the hypotenuse and the yellow (xianhuang he); every great difference is the difference of the altitude and the yellow (guhuang jiao); every small difference is the difference of the base and the yellow (gouhuang jiao). [Translator's note: "the yellow" (huang) is the diameter of the circle inscribed in the triangle.]\end{uniten}

\begin{unitzh}\unitlabel{ceyuan-002.21}髙股平勾差名角差，{\zhsmall 又}名逺差。此數即髙平二差共也，又為明和□和較也，{\zhsmall 又}為通差内去極差，{\zhsmall 又}為極差虚差共。明□二差共名次差，{\zhsmall 又}名近差，{\zhsmall 又}名戾{\zhsmall 音列}和。此數{\zhsmall 又}為明大差□小差較也。勾圓差之股、股圓差之勾相併，名混同和。此數{\zhsmall 又}為一徑一虛弦共也。\end{unitzh}
\begin{uniten}The difference of the height altitude (gaogu) and the level base (pinggou) is named the corner difference (jiaocha), [Note: also] named the far difference (yuancha). This number is the two differences, height and level, together; it is also the difference of the bright sum (minghe) and the □ sum, [Note: also] it is the through difference (tongcha) minus the pole difference (jicha), [Note: also] it is the pole difference and the void difference (xucha) together. The two differences, bright and □, together are named the secondary difference (cicha), [Note: also] named the near difference (jincha), [Note: also] named the contrary [Note: pronounced 列 (lie)] sum (liehe). This number [Note: also] is the difference of the bright great difference and the □ small difference. The altitude of the base-circle difference (gouyuan cha) and the base of the altitude-circle difference (guyuan cha) added together are named the amalgamated sum (huntong he). This number [Note: also] is one diameter (jing) and one void hypotenuse together.\end{uniten}

\begin{unitzh}\unitlabel{ceyuan-002.22}明□二差較名傍差。此數又為髙平二差較，{\zhsmall 又}為極雙差内減虚和，{\zhsmall 又}為極和内減城徑也。虚差不及傍差名蓌差。此數又為大差差内去角差，{\zhsmall 又}為極差内去二之平差，{\zhsmall 又}為次差内去小差差，{\zhsmall 又}為明股□勾共内去二之明勾也。虚差傍差共為蓌和{\zhsmall 蓌音剉}。\end{unitzh}
\begin{uniten}The difference of the two differences, bright and □, is named the adjacent difference (bangcha). This number is also the difference of the two differences, height and level, [Note: also] it is the pole double difference (ji shuangcha) less the void sum (xuhe), [Note: also] it is the pole sum (jihe) less the city diameter. That by which the void difference falls short of the adjacent difference is named the cuo difference (cuocha). This number is also the great-difference difference minus the corner difference, [Note: also] it is the pole difference minus twice the level difference, [Note: also] it is the secondary difference minus the small-difference difference, [Note: also] it is the bright altitude and the □ base together minus twice the bright base. The void difference and the adjacent difference together make the cuo sum (cuohe) [Note: 蓌 is pronounced 剉 (cuo)].\end{uniten}

\begin{unitzh}\unitlabel{ceyuan-002.23}凡大差股小差勾相乘，為半段徑冪；大差勾小差股相乘，亦同上。虚勾乗大股，得半段徑冪；虚股乘大勾，亦同上。邊股□股相乘，得半徑冪；明勾底勾相乘，亦同上。黄廣股黄長勾相乗，得徑冪。髙股平勾相乗，得半徑冪。明弦明股併，與□弦□勾併相乘，得半徑冪；明弦明勾併，與□弦□股併相乘，亦同上。髙弦平弦相乘，為一段皇極積。明勾□股相乘，倍之，為一段太虚積；明股□勾相乘，亦同。\end{unitzh}
\begin{uniten}In every case, the great-difference altitude and the small-difference base multiplied together make half a portion (duan) of the diameter area (mi); the great-difference base and the small-difference altitude multiplied together are also the same as above. The void base multiplied by the great altitude (dagu) gives half a portion of the diameter area; the void altitude multiplied by the great base (dagou) is also the same as above. The side altitude and the □ altitude multiplied together give the radius area; the bright base and the bottom base multiplied together are also the same as above. The yellow-breadth altitude and the yellow-length base multiplied together give the diameter area. The height altitude and the level base multiplied together give the radius area. The bright hypotenuse and the bright altitude added together, and the □ hypotenuse and the □ base added together, multiplied together, give the radius area; the bright hypotenuse and the bright base added together, and the □ hypotenuse and the □ altitude added together, multiplied together, are also the same as above. The height hypotenuse and the level hypotenuse multiplied together make one portion of the imperial-pole area. The bright base and the □ altitude multiplied together and doubled make one portion of the great-void area; the bright altitude and the □ base multiplied together are also the same.\end{uniten}

\begin{unitzh}\unitlabel{ceyuan-002.24}右諸雜名目。\end{unitzh}
\begin{uniten}The above are the various miscellaneous names.\end{uniten}

\begin{unitzh}\unitlabel{ceyuan-002.25}通弦上勾股和即一城徑一通弦也，其較即勾圓差之勾、股圓差之股相較也。勾弦和即二勾一大差，其較則大差也。股弦和即二股一小差，其較則小差也。弦較和為一徑三差共，其較則大勾小差共也。三事和即邊弦三事和上帶大勾也，{\zhsmall 又}為底弦三事和上帶大股也，其較則城徑也。\end{unitzh}
\begin{uniten}The sum of base and altitude on the through hypotenuse is one city diameter and one through hypotenuse; its difference is the difference between the base of the base-circle difference and the altitude of the altitude-circle difference. The sum of base and hypotenuse is two bases and one great difference; its difference is the great difference. The sum of altitude and hypotenuse is two altitudes and one small difference; its difference is the small difference. The sum of hypotenuse and difference is one diameter and three differences together; its difference is the great base and the small difference together. The sum of the three things (sanshi he) is the sum of the three things on the side hypotenuse with the great base added on (shangdai), [Note: also] it is the sum of the three things on the bottom hypotenuse with the great altitude added on; its difference is the city diameter.\end{uniten}

\begin{unitzh}\unitlabel{ceyuan-002.26}邊弦上勾股和為通股平弦共，其較則大差股内去平弦也。勾弦和即通股底勾共，其較則明股明弦共也。股弦和即通股通弦和内少个邊勾也，其較則平勾也。弦較和為大差上股弦和，其較則大勾也。三事和即通弦上股弦和，{\zhsmall 又}為黄廣三事和上帶勾圓差也，其較則大差勾也，{\zhsmall 又}為平弦上弦較和，{\zhsmall 又}為太虛弦上股弦和也。\end{unitzh}
\begin{uniten}The sum of base and altitude on the side hypotenuse is the through altitude and the level hypotenuse together; its difference is the great-difference altitude minus the level hypotenuse. The sum of base and hypotenuse is the through altitude and the bottom base together; its difference is the bright altitude and the bright hypotenuse together. The sum of altitude and hypotenuse is the sum of the through altitude and the through hypotenuse lacking one side base; its difference is the level base. The sum of hypotenuse and difference is the sum of altitude and hypotenuse on the great difference; its difference is the great base. The sum of the three things is the sum of altitude and hypotenuse on the through hypotenuse, [Note: also] it is the sum of the three things of the yellow breadth with the base-circle difference added on; its difference is the great-difference base, [Note: also] it is the sum of hypotenuse and difference on the level hypotenuse, [Note: also] it is the sum of altitude and hypotenuse on the great-void hypotenuse.\end{uniten}

\begin{unitzh}\unitlabel{ceyuan-002.27}底弦上勾股和為通勾髙弦共，其較則髙弦内去小差勾也。勾弦和為通弦上弦較較與髙股共，其較則髙股也。股弦和為半个通弦上三事和，其較則□弦上勾弦和也。弦較和為大差上勾弦和也，其較則小差上勾弦和也。三事和即通弦上勾弦和，{\zhsmall 又}為黄長三事和上帶股圓差，其較則小差股也，{\zhsmall 又}為髙弦上弦較較，{\zhsmall 又}為太虚弦上勾弦和。\end{unitzh}
\begin{uniten}The sum of base and altitude on the bottom hypotenuse is the through base and the height hypotenuse together; its difference is the height hypotenuse minus the small-difference base. The sum of base and hypotenuse is the difference of hypotenuse and difference (xianjiao jiao) on the through hypotenuse together with the height altitude; its difference is the height altitude. The sum of altitude and hypotenuse is half of the sum of the three things on the through hypotenuse; its difference is the sum of base and hypotenuse on the □ hypotenuse. The sum of hypotenuse and difference is the sum of base and hypotenuse on the great difference; its difference is the sum of base and hypotenuse on the small difference. The sum of the three things is the sum of base and hypotenuse on the through hypotenuse, [Note: also] it is the sum of the three things of the yellow length with the altitude-circle difference added on; its difference is the small-difference altitude, [Note: also] it is the difference of hypotenuse and difference on the height hypotenuse, [Note: also] it is the sum of base and hypotenuse on the great-void hypotenuse.\end{uniten}

\begin{unitzh}\unitlabel{ceyuan-002.28}黄廣弦上勾股和為大股虚股共，{\zhsmall 又}為通勾通股共内少个小差上勾股和，其較則兩个髙差也。勾弦和為二髙弦一圓徑共，其較則二明股也。股弦和為通弦上弦較和，其較則二□股也。弦較和即兩个大差股也，其較即兩个小差股也。三事和兩大股也，其較則兩虚股也。\end{unitzh}
\begin{uniten}The sum of base and altitude on the yellow-breadth hypotenuse is the great altitude and the void altitude together, [Note: also] it is the through base and the through altitude together lacking one sum of base and altitude on the small difference; its difference is two height differences (gaocha). The sum of base and hypotenuse is two height hypotenuses and one circle diameter (yuanjing) together; its difference is two bright altitudes. The sum of altitude and hypotenuse is the sum of hypotenuse and difference on the through hypotenuse; its difference is two □ altitudes. The sum of hypotenuse and difference is two great-difference altitudes; its difference is two small-difference altitudes. The sum of the three things is two great altitudes; its difference is two void altitudes.\end{uniten}

\begin{unitzh}\unitlabel{ceyuan-002.29}黄長弦上勾股和為大勾虚勾共，{\zhsmall 又}為通和内少个大差上勾股和也，其較則兩个平差也。勾弦和為通弦上弦較較，其較則兩个明勾也。股弦和為二圓徑二□勾，其較則二□勾也。弦較和為兩个大差勾也，其較則兩个小差勾也。三事和為兩大勾，其較則兩虚勾也。\end{unitzh}
\begin{uniten}The sum of base and altitude on the yellow-length hypotenuse is the great base and the void base together, [Note: also] it is the through sum (tonghe) lacking one sum of base and altitude on the great difference; its difference is two level differences (pingcha). The sum of base and hypotenuse is the difference of hypotenuse and difference on the through hypotenuse; its difference is two bright bases. The sum of altitude and hypotenuse is two circle diameters and two □ bases; its difference is two □ bases. The sum of hypotenuse and difference is two great-difference bases; its difference is two small-difference bases. The sum of the three things is two great bases; its difference is two void bases.\end{uniten}

\begin{unitzh}\unitlabel{ceyuan-002.30}髙弦上勾股和為髙弦虚股共，{\zhsmall 又}為一徑及髙勾髙股差也，其較則底弦内減大勾也，{\zhsmall 又}為邊股内減底股也。勾弦共則底股，其較則明股也。股弦共即邊股，其差則□股也。弦較共則大差股，其較則小差股也。三事和即大股，其較則虚股也，{\zhsmall 又}為小差上勾弦較，{\zhsmall 又}為明弦上弦較較。\end{unitzh}
\begin{uniten}The sum of base and altitude on the height hypotenuse is the height hypotenuse and the void altitude together, [Note: also] it is one diameter and the difference of the height base and the height altitude; its difference is the bottom hypotenuse less the great base, [Note: also] it is the side altitude less the bottom altitude. Base and hypotenuse together are the bottom altitude; their difference is the bright altitude. Altitude and hypotenuse together are the side altitude; their difference is the □ altitude. Hypotenuse and difference together are the great-difference altitude; their difference is the small-difference altitude. The sum of the three things is the great altitude; its difference is the void altitude, [Note: also] it is the difference of base and hypotenuse on the small difference, [Note: also] it is the difference of hypotenuse and difference on the bright hypotenuse.\end{uniten}

\begin{unitzh}\unitlabel{ceyuan-003.1}平弦上勾股共即平弦虚勾共也，其較則大股内減邊弦也。勾弦共即底勾，其差則明勾也。股弦共即邊勾，其較則□勾也。弦較共即大差勾，其較則小差勾也。三事和即大勾，其較則虚勾也，{\zhsmall 又}為大差上股弦較，{\zhsmall 又}為□弦上弦較和。\end{unitzh}
\begin{uniten}On the Level (ping) hypotenuse (xian), base (gou) and altitude (gu) together are the Level hypotenuse and the Void (xu) base together; their difference (jiao) is the Great (da) altitude with the Edge (bian) hypotenuse subtracted. Base and hypotenuse together are the Bottom (di) base; their difference (cha) is the Bright (ming) base. Altitude and hypotenuse together are the Edge base; their difference is the □ base [Translator's note: □ marks a character not reproduced in the transcription; it names one of the figures, like "Bright"; so throughout]. Hypotenuse and difference together are the Great Difference (dacha) base; their difference is the Small Difference (xiaocha) base. The sum of the three things (sanshi he) is the Great base; its difference is the Void base, [Note: also] is the difference of altitude and hypotenuse on the Great Difference, [Note: also] is the sum of hypotenuse and difference on the □ hypotenuse.\end{uniten}

\begin{unitzh}\unitlabel{ceyuan-003.2}大差上勾股和即大股内去虚勾，其差則大差弦内去圓徑也。弦勾共即大股，其差則大差股内去二之明勾也。股弦和為大股上加个大中差也{\zhsmall 按大中差乃明股弦和與半徑之較}，其較則虚勾也。弦較和為兩个邊弦上勾弦較，其較即城徑也。三事和即大股與股圓差共，{\zhsmall 又}為大弦大較共，{\zhsmall 又}為二邊股，其較則太虚上弦較和也。\end{unitzh}
\begin{uniten}On the Great Difference, the sum of base and altitude is the Great altitude with the Void base removed; their difference is the Great Difference hypotenuse with the circle diameter removed. Hypotenuse and base together are the Great altitude; their difference is the Great Difference altitude with 2 times the Bright base removed. The sum of altitude and hypotenuse is the Great altitude plus a great middle difference (dazhongcha) [Note: the great middle difference is the difference between the sum of the Bright altitude and hypotenuse and the half-diameter (banjing)]; their difference is the Void base. The sum of hypotenuse and difference is 2 differences of base and hypotenuse on the Edge hypotenuse; their difference is the city diameter (chengjing). The sum of the three things is the Great altitude together with the altitude–circle difference (guyuancha), [Note: also] is the Great hypotenuse and the Great difference (dajiao) together, [Note: also] is 2 Edge altitudes; its difference is the sum of hypotenuse and difference on the Supreme Void (taixu).\end{uniten}

\begin{unitzh}\unitlabel{ceyuan-003.3}小差上勾股和即大勾内去虚股也，其較則圓徑内去小差弦也。勾弦和為大勾上減个小中差也{\zhsmall 按小中差乃□勾弦和與半徑之較}，其較則虚股也。股弦共即大勾，其較則小差勾内去兩个□股也。弦較和為圓徑，其較則為兩个底弦上股弦較，{\zhsmall 又}為兩个□弦上勾弦和也。三事和即大勾與勾圓差共也，又為大弦大較較{\zhsmall 按即通弦又上弦較較}，為二底勾，其較則太虚上弦較較也。\end{unitzh}
\begin{uniten}On the Small Difference, the sum of base and altitude is the Great base with the Void altitude removed; its difference is the circle diameter with the Small Difference hypotenuse removed. The sum of base and hypotenuse is the Great base minus a small middle difference (xiaozhongcha) [Note: the small middle difference is the difference between the sum of the □ base and hypotenuse and the half-diameter]; their difference is the Void altitude. Altitude and hypotenuse together are the Great base; their difference is the Small Difference base with 2 □ altitudes removed. The sum of hypotenuse and difference is the circle diameter; its difference is 2 differences of altitude and hypotenuse on the Bottom hypotenuse, [Note: also] is 2 sums of base and hypotenuse on the □ hypotenuse. The sum of the three things is the Great base together with the base–circle difference (gouyuancha), also is the difference of the Great hypotenuse and the Great difference [Note: that is, on the Full (tong) hypotenuse, also, the difference of hypotenuse and difference] [Translator's note: the 又 "also" inside this note appears misplaced], is 2 Bottom bases; its difference is the difference of hypotenuse and difference on the Supreme Void.\end{uniten}

\begin{unitzh}\unitlabel{ceyuan-003.4}皇極勾股和即髙弦平弦共，其較則明股内去□勾也。勾弦共即底弦，其較則明弦也。股弦共則邊弦，其較則□弦也。弦較和為髙弦明弦共，{\zhsmall 又}為大股内減大差勾，{\zhsmall 又}為大差弦，其較則小差弦也。三事和即通弦，其較則太虛弦也，{\zhsmall 又}為明勾□股共，{\zhsmall 又}為髙弦内減明弦，{\zhsmall 又}為平弦内減□弦，{\zhsmall 又}為大差勾上減虚股，{\zhsmall 又}為小差股上減虚勾也。\end{unitzh}
\begin{uniten}The Imperial Pole (huangji) sum of base and altitude is the Height (gao) hypotenuse and the Level hypotenuse together; its difference is the Bright altitude with the □ base removed. Base and hypotenuse together are the Bottom hypotenuse; their difference is the Bright hypotenuse. Altitude and hypotenuse together are the Edge hypotenuse; their difference is the □ hypotenuse. The sum of hypotenuse and difference is the Height hypotenuse and the Bright hypotenuse together, [Note: also] is the Great altitude with the Great Difference base subtracted, [Note: also] is the Great Difference hypotenuse; its difference is the Small Difference hypotenuse. The sum of the three things is the Full hypotenuse; its difference is the Supreme Void hypotenuse, [Note: also] is the Bright base and the □ altitude together, [Note: also] is the Height hypotenuse with the Bright hypotenuse subtracted, [Note: also] is the Level hypotenuse with the □ hypotenuse subtracted, [Note: also] is the Great Difference base minus the Void altitude, [Note: also] is the Small Difference altitude minus the Void base.\end{uniten}

\begin{unitzh}\unitlabel{ceyuan-003.5}太虚勾股和即圓徑内減虚弦，{\zhsmall 又}為虚弦虚黄方共，{\zhsmall 又}為皇極弦内去明股□勾共，其差則大差勾内減个小差股也。勾弦共即小差股也，其較則虚股内減个小黄方也。股弦共即大差勾，其較則虚勾内減个小黄方也。弦較和為大差弦上弦和較，{\zhsmall 又}黄長弦上勾弦較，{\zhsmall 又}為兩个明勾，其較小差弦上黄方面也。三事和即大黄方，其較則為兩个明弦上股弦較，{\zhsmall 又}為□弦上兩个勾弦較，{\zhsmall 又}為明弦上小差與□弦上大差共也。\end{unitzh}
\begin{uniten}The Supreme Void sum of base and altitude is the circle diameter with the Void hypotenuse subtracted, [Note: also] is the Void hypotenuse and the Void yellow square (xu huangfang) together, [Note: also] is the Imperial Pole hypotenuse with the Bright altitude and the □ base together removed; its difference is the Great Difference base with a Small Difference altitude subtracted. Base and hypotenuse together are the Small Difference altitude; their difference is the Void altitude with a small yellow square (xiao huangfang) subtracted. Altitude and hypotenuse together are the Great Difference base; their difference is the Void base with a small yellow square subtracted. The sum of hypotenuse and difference is the difference of hypotenuse and sum on the Great Difference hypotenuse, [Note: also] the difference of base and hypotenuse on the Yellow Length (huangchang) hypotenuse, [Note: also] is 2 Bright bases; its difference is the side of the yellow square on the Small Difference hypotenuse. The sum of the three things is the great yellow square (da huangfang); its difference is 2 differences of altitude and hypotenuse on the Bright hypotenuse, [Note: also] is 2 differences of base and hypotenuse on the □ hypotenuse, [Note: also] is the small difference on the Bright hypotenuse together with the great difference on the □ hypotenuse.\end{uniten}

\begin{unitzh}\unitlabel{ceyuan-003.6}明弦勾股和即大差股内減明弦，其較則明弦内減虚股也。勾弦併即髙股，其較則髙股内少二之明勾也。股弦和即邊股内減大差勾，{\zhsmall 又}為邊勾邊弦差，其較則半个虚黄方也。弦較和即大差上勾弦較，其較則虚股也。三事和即股圓差，其較則太虚上勾弦較，{\zhsmall 又}為虚股内減虛黄方也。\end{unitzh}
\begin{uniten}The Bright hypotenuse sum of base and altitude is the Great Difference altitude with the Bright hypotenuse subtracted; its difference is the Bright hypotenuse with the Void altitude subtracted. Base and hypotenuse combined are the Height altitude; their difference is the Height altitude less 2 times the Bright base. The sum of altitude and hypotenuse is the Edge altitude with the Great Difference base subtracted, [Note: also] is the difference of the Edge base and the Edge hypotenuse; their difference is 1/2 of a Void yellow square. The sum of hypotenuse and difference is the difference of base and hypotenuse on the Great Difference; its difference is the Void altitude. The sum of the three things is the altitude–circle difference; its difference is the difference of base and hypotenuse on the Supreme Void, [Note: also] is the Void altitude with the Void yellow square subtracted.\end{uniten}

\begin{unitzh}\unitlabel{ceyuan-003.7}□弦上勾股和即小差内減□弦，其較則虚勾内減□弦也。勾弦和即底勾内減小差股，{\zhsmall 又}為底股底弦差，其較則半个虚黄方也。股弦和即平勾，其較則平勾内少二个□股也。弦較和即虛勾，其較則小差上股弦較也。三事和即勾圓差，其較則太虚上股弦較，{\zhsmall 又}為虚勾内減虚黄方也。\end{unitzh}
\begin{uniten}On the □ hypotenuse, the sum of base and altitude is the Small Difference with the □ hypotenuse subtracted; its difference is the Void base with the □ hypotenuse subtracted. The sum of base and hypotenuse is the Bottom base with the Small Difference altitude subtracted, [Note: also] is the difference of the Bottom altitude and the Bottom hypotenuse; their difference is 1/2 of a Void yellow square. The sum of altitude and hypotenuse is the Level base; their difference is the Level base less 2 □ altitudes. The sum of hypotenuse and difference is the Void base; its difference is the difference of altitude and hypotenuse on the Small Difference. The sum of the three things is the base–circle difference; its difference is the difference of altitude and hypotenuse on the Supreme Void, [Note: also] is the Void base with the Void yellow square subtracted.\end{uniten}

\begin{unitzh}\unitlabel{ceyuan-003.8}前黄廣勾股下：其勾股較，{\zhsmall 又}為大差股上少个小差股，{\zhsmall 又}為中差{\zhsmall 按中差係通勾股較}内少个小差較，{\zhsmall 又}為黄廣股内少一徑。勾弦共，{\zhsmall 又}為兩个底股，{\zhsmall 又}為大股與小差股共。股弦和，{\zhsmall 又}為大弦中差共，{\zhsmall 又}為兩个邊股。股弦差，{\zhsmall 又}為小差上黄方面。\end{unitzh}
\begin{uniten}Under the preceding Yellow Breadth (huangguang) right triangle (gougu): its difference of base and altitude [Note: also] is the Great Difference altitude less a Small Difference altitude, [Note: also] is the middle difference (zhongcha) [Note: the middle difference is the difference of the Full base and altitude] less a Small Difference difference (xiaocha jiao), [Note: also] is the Yellow Breadth altitude less 1 diameter. Base and hypotenuse together [Note: also] are 2 Bottom altitudes, [Note: also] are the Great altitude together with the Small Difference altitude. The sum of altitude and hypotenuse [Note: also] is the Great hypotenuse and the middle difference together, [Note: also] is 2 Edge altitudes. The difference of altitude and hypotenuse [Note: also] is the side of the yellow square on the Small Difference.\end{uniten}

\begin{unitzh}\unitlabel{ceyuan-003.9}前黄長勾股下：其勾股較，{\zhsmall 又}為大差勾上少个小差勾也，{\zhsmall 又}為圓徑内少个黄長勾。勾弦共，{\zhsmall 又}為兩个底勾，{\zhsmall 又}為大勾與小差勾共。勾弦較，{\zhsmall 又}為大差上黄方靣。股弦共，{\zhsmall 又}為兩个邊勾。\end{unitzh}
\begin{uniten}Under the preceding Yellow Length right triangle: its difference of base and altitude [Note: also] is the Great Difference base less a Small Difference base, [Note: also] is the circle diameter less a Yellow Length base. Base and hypotenuse together [Note: also] are 2 Bottom bases, [Note: also] are the Great base together with the Small Difference base. The difference of base and hypotenuse [Note: also] is the side of the yellow square on the Great Difference. Altitude and hypotenuse together [Note: also] are 2 Edge bases.\end{uniten}

\begin{unitzh}\unitlabel{ceyuan-003.10}右五和五較\end{unitzh}
\begin{uniten}The above: the five sums and five differences.\end{uniten}

\begin{unitzh}\unitlabel{ceyuan-003.11}大弦為大勾與股圓差共，{\zhsmall 又}為大股與勾圓差共。邊弦乃邊股平勾共，{\zhsmall 又}為大股内減平弦上勾股較。底弦乃底勾髙股共，{\zhsmall 又}為大勾内加一个髙差。黄廣弦為大股内減虚股，{\zhsmall 又}為邊股□股共。黄長弦乃大勾内減虚勾，{\zhsmall 又}為底勾明勾共。\end{unitzh}
\begin{uniten}The Great hypotenuse is the Great base together with the altitude–circle difference, [Note: also] is the Great altitude together with the base–circle difference. The Edge hypotenuse is the Edge altitude and the Level base together, [Note: also] is the Great altitude with the difference of base and altitude on the Level hypotenuse subtracted. The Bottom hypotenuse is the Bottom base and the Height altitude together, [Note: also] is the Great base with 1 Height difference (gaocha) added. The Yellow Breadth hypotenuse is the Great altitude with the Void altitude subtracted, [Note: also] is the Edge altitude and the □ altitude together. The Yellow Length hypotenuse is the Great base with the Void base subtracted, [Note: also] is the Bottom base and the Bright base together.\end{uniten}

\begin{unitzh}\unitlabel{ceyuan-003.12}髙弦乃大差弦内減明弦，{\zhsmall 又}為明弦虚弦共。平弦乃小差弦内減□弦，{\zhsmall 又}為□弦虚弦共。大差弦乃大股内減大差勾，{\zhsmall 又}為髙弦明弦共，{\zhsmall 又}大弦内去黄長弦。小差弦為大勾内減小差股，{\zhsmall 又}為平弦□弦共，{\zhsmall 又}為大弦内去黄廣弦。極弦乃髙股平勾共，{\zhsmall 又}為平弦明弦共，{\zhsmall 又}為髙弦□弦共，{\zhsmall 又}為大差弦内減髙平二弦較，{\zhsmall 又}為小差弦内加髙平二弦較。虚弦乃皇極黄方靣，{\zhsmall 又}為明勾□股共，{\zhsmall 又}為髙弦内減明弦，{\zhsmall 又}為平弦内減□弦。明弦乃髙弦内減虚弦。□弦乃平弦内減虚弦。\end{unitzh}
\begin{uniten}The Height hypotenuse is the Great Difference hypotenuse with the Bright hypotenuse subtracted, [Note: also] is the Bright hypotenuse and the Void hypotenuse together. The Level hypotenuse is the Small Difference hypotenuse with the □ hypotenuse subtracted, [Note: also] is the □ hypotenuse and the Void hypotenuse together. The Great Difference hypotenuse is the Great altitude with the Great Difference base subtracted, [Note: also] is the Height hypotenuse and the Bright hypotenuse together, [Note: also] the Great hypotenuse with the Yellow Length hypotenuse removed. The Small Difference hypotenuse is the Great base with the Small Difference altitude subtracted, [Note: also] is the Level hypotenuse and the □ hypotenuse together, [Note: also] is the Great hypotenuse with the Yellow Breadth hypotenuse removed. The Pole (ji) hypotenuse is the Height altitude and the Level base together, [Note: also] is the Level hypotenuse and the Bright hypotenuse together, [Note: also] is the Height hypotenuse and the □ hypotenuse together, [Note: also] is the Great Difference hypotenuse with the difference of the 2 hypotenuses Height and Level subtracted, [Note: also] is the Small Difference hypotenuse with the difference of the 2 hypotenuses Height and Level added. The Void hypotenuse is the side of the Imperial Pole yellow square, [Note: also] is the Bright base and the □ altitude together, [Note: also] is the Height hypotenuse with the Bright hypotenuse subtracted, [Note: also] is the Level hypotenuse with the □ hypotenuse subtracted. The Bright hypotenuse is the Height hypotenuse with the Void hypotenuse subtracted. The □ hypotenuse is the Level hypotenuse with the Void hypotenuse subtracted.\end{uniten}

\begin{unitzh}\unitlabel{ceyuan-003.13}黄廣弦、黄長弦相併為大弦虚弦共也。以此數減于大和，餘即虚和。若以二弦相減，餘即虚弦平弦共也{\zhsmall 按「虚弦平弦共」，此題數偶合，當云「二極差」}。黄廣弦{\zhsmall 又}為大差弦虚弦共。黄長弦{\zhsmall 又}為小差弦虚弦共。以黄長弦減于大勾，餘即虚勾。以黄廣弦減于大股，餘即虚股。\end{unitzh}
\begin{uniten}The Yellow Breadth hypotenuse and the Yellow Length hypotenuse combined make the Great hypotenuse and the Void hypotenuse together. Subtract this number from the Great sum (dahe); the remainder is the Void sum (xuhe). If the 2 hypotenuses are subtracted from each other, the remainder is the Void hypotenuse and the Level hypotenuse together [Note: as for "the Void hypotenuse and the Level hypotenuse together", the numbers of this problem coincide by chance; it should say "2 Pole differences (jicha)"]. The Yellow Breadth hypotenuse [Note: also] is the Great Difference hypotenuse and the Void hypotenuse together. The Yellow Length hypotenuse [Note: also] is the Small Difference hypotenuse and the Void hypotenuse together. Subtract the Yellow Length hypotenuse from the Great base; the remainder is the Void base. Subtract the Yellow Breadth hypotenuse from the Great altitude; the remainder is the Void altitude.\end{uniten}

\begin{unitzh}\unitlabel{ceyuan-003.14}邊弦、底弦相併為大弦皇極弦共也。于此併數内減大和，餘為皇極弦内減圓徑也。若以二弦相減，餘即皇極差也。此數同者最多，故{\zhsmall 又}為皇極弦内少个小差弦，{\zhsmall 又}為髙弦平弦較，{\zhsmall 又}為明股内少□勾，{\zhsmall 又}為大差弦内少皇極弦，{\zhsmall 又}為次差虚差共也。邊弦{\zhsmall 又}為皇極股弦共，{\zhsmall 又}為黄廣弦□弦共。\end{unitzh}
\begin{uniten}The Edge hypotenuse and the Bottom hypotenuse combined make the Great hypotenuse and the Imperial Pole hypotenuse together. From this combined number subtract the Great sum; the remainder is the Imperial Pole hypotenuse with the circle diameter subtracted. If the 2 hypotenuses are subtracted from each other, the remainder is the Imperial Pole difference (huangjicha). Those equal to this number are the most numerous; therefore it [Note: also] is the Imperial Pole hypotenuse less a Small Difference hypotenuse, [Note: also] is the difference of the Height hypotenuse and the Level hypotenuse, [Note: also] is the Bright altitude less the □ base, [Note: also] is the Great Difference hypotenuse less the Imperial Pole hypotenuse, [Note: also] is the secondary difference (cicha) and the Void difference (xucha) together. The Edge hypotenuse [Note: also] is the Imperial Pole altitude and hypotenuse together, [Note: also] is the Yellow Breadth hypotenuse and the □ hypotenuse together.\end{uniten}

\begin{unitzh}\unitlabel{ceyuan-003.15}底弦{\zhsmall 又}為皇極勾弦共，{\zhsmall 又}為黄長弦明弦共也。以邊弦減大股，餘為半徑内減平勾，{\zhsmall 又}為平弦内減小差勾也。底弦内減大勾，餘為髙股内減半徑，{\zhsmall 又}為大差股内減髙弦也。\end{unitzh}
\begin{uniten}The Bottom hypotenuse [Note: also] is the Imperial Pole base and hypotenuse together, [Note: also] is the Yellow Length hypotenuse and the Bright hypotenuse together. Subtract the Edge hypotenuse from the Great altitude; the remainder is the half-diameter with the Level base subtracted, [Note: also] is the Level hypotenuse with the Small Difference base subtracted. From the Bottom hypotenuse subtract the Great base; the remainder is the Height altitude with the half-diameter subtracted, [Note: also] is the Great Difference altitude with the Height hypotenuse subtracted.\end{uniten}

\begin{unitzh}\unitlabel{ceyuan-003.16}黄廣弦内減邊股即□股，黄長弦内減底勾即明勾也。\end{unitzh}
\begin{uniten}The Yellow Breadth hypotenuse with the Edge altitude subtracted is the □ altitude; the Yellow Length hypotenuse with the Bottom base subtracted is the Bright base.\end{uniten}

\begin{unitzh}\unitlabel{ceyuan-003.17}髙弦髙股共即邊股，平弦平勾共即底勾，髙弦髙勾共即底股，平弦平股共即邊勾。\end{unitzh}
\begin{uniten}The Height hypotenuse and the Height altitude together are the Edge altitude; the Level hypotenuse and the Level base together are the Bottom base; the Height hypotenuse and the Height base together are the Bottom altitude; the Level hypotenuse and the Level altitude together are the Edge base.\end{uniten}

\begin{unitzh}\unitlabel{ceyuan-003.18}上髙弦減于通股，餘即邊股内減□股也。下平弦減于通勾，餘即邊勾内減明勾也。髙弦、平弦相併，即大弦内少个皇極弦也；若以相併數減於大和，餘為皇極弦圓徑共也。髙弦、平弦相減，餘即皇極差也，{\zhsmall 又}為皇極弦上減小差弦也；若以相減數却加于相併數，即黄廣弦也。\end{unitzh}
\begin{uniten}Subtract the Upper Height (shanggao) hypotenuse from the Full altitude; the remainder is the Edge altitude with the □ altitude subtracted. Subtract the Lower Level (xiaping) hypotenuse from the Full base; the remainder is the Edge base with the Bright base subtracted. The Height hypotenuse and the Level hypotenuse combined are the Great hypotenuse less an Imperial Pole hypotenuse; if the combined number is subtracted from the Great sum, the remainder is the Imperial Pole hypotenuse and the circle diameter together. Subtract the Height hypotenuse and the Level hypotenuse from each other; the remainder is the Imperial Pole difference, [Note: also] is the Imperial Pole hypotenuse minus the Small Difference hypotenuse; if the number from the subtraction is in turn added to the combined number, it is the Yellow Breadth hypotenuse.\end{uniten}

\begin{unitzh}\unitlabel{ceyuan-003.19}髙弦内減明股得半徑，平弦内減□勾亦同上。皇極勾上加明弦為皇極弦，皇極股上加□弦亦同上。\end{unitzh}
\begin{uniten}The Height hypotenuse with the Bright altitude subtracted gives the half-diameter; the Level hypotenuse with the □ base subtracted is also the same as above. The Imperial Pole base plus the Bright hypotenuse is the Imperial Pole hypotenuse; the Imperial Pole altitude plus the □ hypotenuse is also the same as above.\end{uniten}

\begin{unitzh}\unitlabel{ceyuan-003.20}皇極弦：得極勾即底弦；得極股即邊弦；内去極勾即明弦；去極股即□弦；減于通弦即極和；得虚弦亦同上；内去虚弦即明弦□弦共；去虚黄即明和□和共也；去城徑即傍差；内加極差即大差弦；去極差即小差弦；加角差即兩个髙股；減角差即二平勾。\end{unitzh}
\begin{uniten}The Imperial Pole hypotenuse: gaining the Pole base, it is the Bottom hypotenuse; gaining the Pole altitude, it is the Edge hypotenuse; with the Pole base removed, it is the Bright hypotenuse; removing the Pole altitude, it is the □ hypotenuse; subtracted from the Full hypotenuse, it is the Pole sum (jihe); gaining the Void hypotenuse, it is also the same as above; with the Void hypotenuse removed, it is the Bright hypotenuse and the □ hypotenuse together; removing the Void yellow (xu huang), it is the Bright sum (minghe) and the □ sum together; removing the city diameter, it is the side difference (bangcha); with the Pole difference added, it is the Great Difference hypotenuse; removing the Pole difference, it is the Small Difference hypotenuse; adding the angle difference (jiaocha), it is 2 Height altitudes; subtracting the angle difference, it is 2 Level bases.\end{uniten}

\begin{unitzh}\unitlabel{ceyuan-003.21}太虛弦：加入極弦為極和；極弦内去之即明□二弦共；再去之則明大差□小差併也；加于大差弦即黄廣弦；加于小差弦即黄長弦；内去明勾則□勾；加明勾為圓徑内少虛黄□股共；加入明股為明和□股共；減于明股即明較内去□股；加入明弦為極股；減于明弦為明大差□小差内少个□弦；加于明和即兩个虚弦一个髙差共也；減于明和即髙差也。\end{unitzh}
\begin{uniten}The Supreme Void hypotenuse: added into the Pole hypotenuse, it makes the Pole sum; the Pole hypotenuse with it removed is the 2 hypotenuses Bright and □ together; removing it once more gives the Bright great difference and the □ small difference combined; added to the Great Difference hypotenuse, it is the Yellow Breadth hypotenuse; added to the Small Difference hypotenuse, it is the Yellow Length hypotenuse; with the Bright base removed, it is the □ base; adding the Bright base, it makes the circle diameter less the Void yellow and the □ altitude together; added into the Bright altitude, it makes the Bright sum and the □ altitude together; subtracted from the Bright altitude, it is the Bright difference (mingjiao) with the □ altitude removed; added into the Bright hypotenuse, it makes the Pole altitude; subtracted from the Bright hypotenuse, it makes the Bright great difference and the □ small difference less a □ hypotenuse; added to the Bright sum, it is 2 Void hypotenuses and 1 Height difference together; subtracted from the Bright sum, it is the Height difference.\end{uniten}

\begin{unitzh}\unitlabel{ceyuan-003.22}内去□勾即明勾□較共，{\zhsmall 又}為□股平差共；加于□勾即□和明勾共；加于□股為二虚弦内少明勾，{\zhsmall 又}為圓徑内少虚黄明勾共；内減□股即明勾；内加□弦即極勾；減于□弦為明勾内少个□小差；加入□和即兩个虚弦内少个平差也；内減□和即平差也；加入明□二和共即極和内少个虚黄也；若減於明□二和共即明股□勾共也；減于髙弦即明弦；減于平弦即□弦；加于角差即二明勾一極差也；減于角差即一極差二□股較也；得傍差即明股□勾共；内減傍差即太虚三事和内去了極雙差也{\zhsmall 按雙差係勾弦差股弦差}；内加虚差即二明勾；内減虚差即二□股；内加虚黄方即虚和；内減虚黄方即太虚大小差併也。\end{unitzh}
\begin{uniten}With the □ base removed, it is the Bright base and the □ difference (jiao) together, [Note: also] is the □ altitude and the Level difference (pingcha) together; added to the □ base, it is the □ sum and the Bright base together; added to the □ altitude, it makes 2 Void hypotenuses less the Bright base, [Note: also] is the circle diameter less the Void yellow and the Bright base together; with the □ altitude subtracted, it is the Bright base; with the □ hypotenuse added, it is the Pole base; subtracted from the □ hypotenuse, it makes the Bright base less a □ small difference; added into the □ sum, it is 2 Void hypotenuses less a Level difference; with the □ sum subtracted, it is the Level difference; added into the 2 sums Bright and □ together, it is the Pole sum less a Void yellow; if subtracted from the 2 sums Bright and □ together, it is the Bright altitude and the □ base together; subtracted from the Height hypotenuse, it is the Bright hypotenuse; subtracted from the Level hypotenuse, it is the □ hypotenuse; added to the angle difference, it is 2 Bright bases and 1 Pole difference; subtracted from the angle difference, it is the difference of 1 Pole difference and 2 □ altitudes; gaining the side difference, it is the Bright altitude and the □ base together; with the side difference subtracted, it is the Supreme Void sum of the three things with the Pole double difference (ji shuangcha) removed [Note: the double difference is the difference of base and hypotenuse and the difference of altitude and hypotenuse]; with the Void difference added, it is 2 Bright bases; with the Void difference subtracted, it is 2 □ altitudes; with the Void yellow square added, it is the Void sum; with the Void yellow square subtracted, it is the Supreme Void great and small differences combined.\end{uniten}

\begin{unitzh}\unitlabel{ceyuan-003.23}右諸弦\end{unitzh}
\begin{uniten}The above: the various hypotenuses.\end{uniten}

\begin{unitzh}\unitlabel{ceyuan-003.24}大差弦、小差弦共即兩个極弦也，以兩个極差為之較。大差差、小差差共即兩个極差也，以兩个傍差為之較。大差上大差、小差上大差共即兩个明弦也，以兩个明差為之較。大差上小差、小差上小差共即兩个□弦也，以兩个□差為之較。大差黄{\zhsmall 按即二明勾}、小差黄{\zhsmall 按即二□股}數共即兩个極黄{\zhsmall 按即二虚弦}也，以兩个虚差為之較。大差勾、小差勾共即兩个極勾也，以兩个平差為之較。大差股、小差股共即兩个極股也，以兩个髙差為之較。二和共為二極和，以二角差為之較。\end{unitzh}
\begin{uniten}The Great Difference hypotenuse and the Small Difference hypotenuse together are 2 Pole hypotenuses, with 2 Pole differences as their difference. The Great Difference difference (dacha cha) and the Small Difference difference (xiaocha cha) together are 2 Pole differences, with 2 side differences as their difference. The great difference on the Great Difference and the great difference on the Small Difference together are 2 Bright hypotenuses, with 2 Bright differences (mingcha) as their difference. The small difference on the Great Difference and the small difference on the Small Difference together are 2 □ hypotenuses, with 2 □ differences (cha) as their difference. The Great Difference yellow [Note: that is, 2 Bright bases] and the Small Difference yellow [Note: that is, 2 □ altitudes], their numbers together, are 2 Pole yellows (ji huang) [Note: that is, 2 Void hypotenuses], with 2 Void differences as their difference. The Great Difference base and the Small Difference base together are 2 Pole bases, with 2 Level differences as their difference. The Great Difference altitude and the Small Difference altitude together are 2 Pole altitudes, with 2 Height differences as their difference. The 2 sums together make 2 Pole sums, with 2 angle differences as their difference.\end{uniten}

\begin{unitzh}\unitlabel{ceyuan-004.1}大差上弦較較即圓徑，小差上弦較和亦同上。大差上小差即虚勾，小差上大差即虚股也。大差弦與明勾共即邊股，小差弦與□股共即底勾也。大差弦内減中差即黄長勾{\zhsmall 按：「勾」應作「股」}；小差弦内加中差即黄廣股也{\zhsmall 按：「股」應作「勾」}。大股内減小差股即黄廣股，大勾内減大差勾即黄長勾也。虚弦得虛股即大差勾，虚弦得虚勾即小差股也。明段弦較和即大差上勾弦較，明段弦較較即小差上勾弦較也；□段弦較和即大差上股弦較，□段弦較較即小差上股弦較也。\end{unitzh}
\begin{uniten}The hypotenuse-difference difference (xianjiaojiao) on the great difference (dacha) is the circle diameter (yuanjing); the hypotenuse-difference sum (xianjiaohe) on the small difference (xiaocha) is also the same as the above. The small difference on the great difference is the void base (xugou); the great difference on the small difference is the void altitude. The great-difference hypotenuse together with the bright base (minggou) is the side altitude (biangu); the small-difference hypotenuse together with the □ altitude is the bottom base (digou). [Translator's note: □ marks a lost character, the name of one of the triangles; it is kept as □.] From the great-difference hypotenuse subtract the middle difference (zhongcha): this is the yellow-length base (huangchang gou) [Siku editors: "base" should read "altitude"]; to the small-difference hypotenuse add the middle difference: this is the yellow-breadth altitude (huangguang gu) [Siku editors: "altitude" should read "base"]. From the great altitude (dagu) subtract the small-difference altitude: this is the yellow-breadth altitude; from the great base subtract the great-difference base: this is the yellow-length base. The void hypotenuse joined with (de) the void altitude is the great-difference base; the void hypotenuse joined with the void base is the small-difference altitude. The hypotenuse-difference sum of the bright segment (duan) is the base-hypotenuse difference (gouxianjiao) on the great difference; the hypotenuse-difference difference of the bright segment is the base-hypotenuse difference on the small difference; the hypotenuse-difference sum of the □ segment is the altitude-hypotenuse difference (guxianjiao) on the great difference; the hypotenuse-difference difference of the □ segment is the altitude-hypotenuse difference on the small difference.\end{uniten}

\begin{unitzh}\unitlabel{ceyuan-004.2}大差勾内減虚弦，餘即虚股；小差股内減虚弦，餘即虚勾也。以大差和減大股，即虚勾；以小差和減大勾，即虚股也。以大差差減圓徑，即明勾；此差若多於圓徑，則内減圓徑，餘即虚勾也{\zhsmall 按：此條因題數偶合而誤，若勾股差甚大甚小者，皆不能合}。以小差差減圓徑，即小差弦也。大差弦上加一徑，即大股上加虚勾也；小差弦上加一徑，即大勾上加虚股也。大差股内減髙弦，餘即髙股内減半徑；平弦内減小差勾，餘即半徑内減平勾也。大差内減虚差，即二明差；小差内減虚差，即二□差也。\end{unitzh}
\begin{uniten}From the great-difference base subtract the void hypotenuse: the remainder is the void altitude; from the small-difference altitude subtract the void hypotenuse: the remainder is the void base. Subtract the great-difference sum (he) from the great altitude: this is the void base; subtract the small-difference sum from the great base: this is the void altitude. Subtract the great-difference difference from the circle diameter: this is the bright base; if this difference is more than the circle diameter, then subtract the circle diameter from it: the remainder is the void base [Siku editors: This item is mistaken because the numbers of the problem happen to agree; where the base-altitude difference is very large or very small, none of them can agree]. Subtract the small-difference difference from the circle diameter: this is the small-difference hypotenuse. The great-difference hypotenuse with 1 diameter added onto it is the great altitude with the void base added onto it; the small-difference hypotenuse with 1 diameter added onto it is the great base with the void altitude added onto it. From the great-difference altitude subtract the high hypotenuse (gaoxian): the remainder is the high altitude less the half-diameter (banjing); from the level hypotenuse (pingxian) subtract the small-difference base: the remainder is the half-diameter less the level base. From the great difference subtract the void difference (xucha): this is 2 bright differences; from the small difference subtract the void difference: this is 2 □ differences.\end{uniten}

\begin{unitzh}\unitlabel{ceyuan-004.3}大弦内減大差股、小差勾共，即圓徑；三事和内減二之大差股、小差勾共，即三个圓徑也。\end{unitzh}
\begin{uniten}From the great hypotenuse subtract the great-difference altitude and the small-difference base together: this is the circle diameter; from the sum of the three things (sanshihe) subtract 2 times the great-difference altitude and the small-difference base together: this is 3 circle diameters.\end{uniten}

\begin{unitzh}\unitlabel{ceyuan-004.4}大差勾、小差股相併，名「混同」，即一圓徑、一虚弦也；若以相減，即虚差也。\end{unitzh}
\begin{uniten}The great-difference base and the small-difference altitude combined are named "merged-together" (huntong): this is 1 circle diameter and 1 void hypotenuse; if they are subtracted from each other, this is the void difference.\end{uniten}

\begin{unitzh}\unitlabel{ceyuan-004.5}大差和、小差和二數相併，即大弦、虚弦共也；二數相減，即中差、虚差共也。{\zhsmall 又}半之併數，即為極弦、虚弦共也；{\zhsmall 又}為髙弦、平弦共；{\zhsmall 又}為皇極勾股共也。\end{unitzh}
\begin{uniten}The great-difference sum and the small-difference sum: the 2 numbers combined are the great hypotenuse and the void hypotenuse together; the 2 numbers subtracted from each other are the middle difference and the void difference together. [Note: also] Half the combined number is the pole hypotenuse (jixian) and the void hypotenuse together; [Note: also] it is the high hypotenuse and the level hypotenuse together; [Note: also] it is the base and altitude of the august pole (huangji) together.\end{uniten}

\begin{unitzh}\unitlabel{ceyuan-004.6}大差差、小差差二數相併，即兩个皇極差，{\zhsmall 又}為大差弦内減小差弦也；二數相減而半之，即是皇極弦上減圓徑也{\zhsmall 即傍差}。\end{unitzh}
\begin{uniten}The great-difference difference and the small-difference difference: the 2 numbers combined are 2 august-pole differences, [Note: also] it is the great-difference hypotenuse less the small-difference hypotenuse; the 2 numbers subtracted from each other and halved are the august-pole hypotenuse with the circle diameter subtracted from it [Note: that is, the adjacent difference (pangcha)].\end{uniten}

\begin{unitzh}\unitlabel{ceyuan-004.7}右大小差\end{unitzh}
\begin{uniten}The above: the great and small differences (daxiaocha).\end{uniten}

\begin{unitzh}\unitlabel{ceyuan-004.8}大差差、小差差、虚差共為一个通差；髙、平、極三差共，亦同上；明、□、虚三差共為一个極差也。諸黄方面亦倣此。\end{unitzh}
\begin{uniten}The great-difference difference, the small-difference difference and the void difference together make 1 overall difference (tongcha); the 3 differences, high, level and pole, together are also the same as the above; the 3 differences, bright, □ and void, together make 1 pole difference. The sides of the various yellow squares (huangfang) also follow this pattern.\end{uniten}

\begin{unitzh}\unitlabel{ceyuan-004.9}邊黄内減底黄，即虚差；黄廣黄内減黄長黄，即二虚差；髙黄内減平黄，即虚差，蓋髙黄即虚股，平黄即虚勾也；大差黄内減小差黄，即二虚差，蓋大差黄即二明勾，小差黄即二□股也；明黄内減□黄，餘即虚差。□弦上三差合成一个虚黄方。\end{unitzh}
\begin{uniten}From the side yellow (bianhuang) subtract the bottom yellow: this is the void difference; from the yellow-breadth yellow subtract the yellow-length yellow: this is 2 void differences; from the high yellow subtract the level yellow: this is the void difference, for the high yellow is the void altitude and the level yellow is the void base; from the great-difference yellow subtract the small-difference yellow: this is 2 void differences, for the great-difference yellow is 2 bright bases and the small-difference yellow is 2 □ altitudes; from the bright yellow subtract the □ yellow: the remainder is the void difference. The 3 differences on the □ hypotenuse combine to form 1 void yellow square.\end{uniten}

\begin{unitzh}\unitlabel{ceyuan-004.10}髙差内減平差為傍差，邊差内減底差亦同上，明差内減□差亦同上；大差差内減小差差為二旁差，黄廣差内減黄長差亦同上。\end{unitzh}
\begin{uniten}From the high difference subtract the level difference: this makes the adjacent difference; from the side difference subtract the bottom difference: also the same as the above; from the bright difference subtract the □ difference: also the same as the above. From the great-difference difference subtract the small-difference difference: this makes 2 adjacent differences; from the yellow-breadth difference subtract the yellow-length difference: also the same as the above.\end{uniten}

\begin{unitzh}\unitlabel{ceyuan-004.11}極雙差即明、□二弦共。内加虚雙差，即明、□二和共；内減虚雙差，即明雙差、□雙差共也；内加旁差，即極弦内少个虚弦旁差差；内減旁差，即虚和也；内加虚差，即極弦内少二□股；内減虚差，則極弦内少二明勾也。\end{unitzh}
\begin{uniten}The pole double difference (ji shuangcha) is the 2 hypotenuses, bright and □, together. Add to it the void double difference: this is the 2 sums, bright and □, together; subtract from it the void double difference: this is the bright double difference and the □ double difference together; add to it the adjacent difference: this is the pole hypotenuse lacking a difference of the void hypotenuse and the adjacent difference; subtract from it the adjacent difference: this is the void sum; add to it the void difference: this is the pole hypotenuse lacking 2 □ altitudes; subtract from it the void difference: then it is the pole hypotenuse lacking 2 bright bases.\end{uniten}

\begin{unitzh}\unitlabel{ceyuan-004.12}極差内加旁差為大差差，内減旁差為小差差也；内加虚差即角差，内減虚差即次差也。倍極差為大差差、小差差共，則倍旁差為之較；倍極弦為大差弦、小差弦共，倍極差為之較。以極差為明差、平差共，則以蓌差為之較；以極差為髙差、□差共，則以蓌和為之較。副置蓌和，上加蓌差而半之，即旁差也；減蓌差而半之，則虛差也。極差内減二之平差，得蓌差。\end{unitzh}
\begin{uniten}To the pole difference add the adjacent difference: this makes the great-difference difference; subtract from it the adjacent difference: this makes the small-difference difference; add to it the void difference: this is the corner difference (jiaocha); subtract from it the void difference: this is the secondary difference (cicha). Twice the pole difference makes the great-difference difference and the small-difference difference together, and then twice the adjacent difference is their difference; twice the pole hypotenuse makes the great-difference hypotenuse and the small-difference hypotenuse together, and twice the pole difference is their difference. If the pole difference is taken as the bright difference and the level difference together, then the askew difference (cuocha) is taken as their difference; if the pole difference is taken as the high difference and the □ difference together, then the askew sum (cuohe) is taken as their difference. Set up the askew sum in a second position; add the askew difference onto it and halve: this is the adjacent difference; subtract the askew difference and halve: then it is the void difference. From the pole difference subtract 2 times the level difference: this gives the askew difference.\end{uniten}

\begin{unitzh}\unitlabel{ceyuan-004.13}角差内加旁差為二髙差，内減旁差即二平差也；内加明、□二差併而半之得極差，内減明、□二差而半之則虚差也；内加極差則通差，内減極差則虚差也。\end{unitzh}
\begin{uniten}To the corner difference add the adjacent difference: this makes 2 high differences; subtract from it the adjacent difference: this is 2 level differences; add to it the 2 differences, bright and □, combined, and halve: this gives the pole difference; subtract from it the 2 differences, bright and □, and halve: then it is the void difference; add to it the pole difference: then it is the overall difference; subtract from it the pole difference: then it is the void difference.\end{uniten}

\begin{unitzh}\unitlabel{ceyuan-004.14}以虚差減於明和，為明、□二股共；以虚差加於□和，為明、□二勾共也。又副置二和共，上加次差而半之，即明、□二股共；減次差而半之，即明、□二勾共也。明、□二股共，以髙差為之較；明、□二勾共，以平差為之較。\end{unitzh}
\begin{uniten}Subtract the void difference from the bright sum: this makes the 2 altitudes, bright and □, together; add the void difference to the □ sum: this makes the 2 bases, bright and □, together. Again, set up the 2 sums together in a second position; add the secondary difference onto it and halve: this is the 2 altitudes, bright and □, together; subtract the secondary difference and halve: this is the 2 bases, bright and □, together. The 2 altitudes, bright and □, together have the high difference as their difference; the 2 bases, bright and □, together have the level difference as their difference.\end{uniten}

\begin{unitzh}\unitlabel{ceyuan-004.15}以髙差減明和即虚弦，以平差加□和亦同上；以髙差減髙股即半徑，以平差加平勾亦同上；以髙差減大差差即明差，以平差減小差差即□差也；以髙差減大差即髙弦，以平差加小差即平弦也。二之平差内去虚差，餘即小差差；去二虚差，即兩个□差。\end{unitzh}
\begin{uniten}Subtract the high difference from the bright sum: this is the void hypotenuse; add the level difference to the □ sum: also the same as the above. Subtract the high difference from the high altitude: this is the half-diameter; add the level difference to the level base: also the same as the above. Subtract the high difference from the great-difference difference: this is the bright difference; subtract the level difference from the small-difference difference: this is the □ difference. Subtract the high difference from the great difference: this is the high hypotenuse; add the level difference to the small difference: this is the level hypotenuse. From 2 times the level difference remove the void difference: the remainder is the small-difference difference; remove 2 void differences: this is 2 □ differences.\end{uniten}

\begin{unitzh}\unitlabel{ceyuan-004.16}髙股即半徑上股方差，平勾即半徑上勾方差，故髙勾、平股共為全徑也。黄廣股即全徑上股方差，黄長勾即全徑上勾方差，故黄廣勾、黄長股共數為兩个全徑也。\end{unitzh}
\begin{uniten}The high altitude is the altitude-square difference (gu fangcha) on the half-diameter; the level base is the base-square difference on the half-diameter; therefore the high base and the level altitude together make the full diameter (quanjing). The yellow-breadth altitude is the altitude-square difference on the full diameter; the yellow-length base is the base-square difference on the full diameter; therefore the number of the yellow-breadth base and the yellow-length altitude together makes 2 full diameters.\end{uniten}

\begin{unitzh}\unitlabel{ceyuan-004.17}邊弦内減底弦，即皇極差。邊股内減底股，即髙差，{\zhsmall 又}為底弦内減大勾。邊勾内減底勾，即平差，{\zhsmall 又}為大股内減邊弦也。\end{unitzh}
\begin{uniten}From the side hypotenuse subtract the bottom hypotenuse: this is the august-pole difference. From the side altitude subtract the bottom altitude: this is the high difference, [Note: also] it is the bottom hypotenuse less the great base. From the side base subtract the bottom base: this is the level difference, [Note: also] it is the great altitude less the side hypotenuse.\end{uniten}

\begin{unitzh}\unitlabel{ceyuan-004.18}大勾減底弦，餘即半徑，為勾之中差也；大股内減邊弦，餘即半徑，為股之中差也。邊股、底勾相併，即大弦；若以相減，即通中差也。\end{unitzh}
\begin{uniten}From the great base subtract the bottom hypotenuse: the remainder is the half-diameter, which is the middle difference of the base; from the great altitude subtract the side hypotenuse: the remainder is the half-diameter, which is the middle difference of the altitude. The side altitude and the bottom base combined are the great hypotenuse; if they are subtracted from each other, this is the overall middle difference (tong zhongcha).\end{uniten}

\begin{unitzh}\unitlabel{ceyuan-004.19}二髙股、一虚差合成一个股圓差，二平勾、一虚差合成一个勾圓差{\zhsmall 按：此二條誤，當云「二明股、一虚股合成一个股圓差；二□勾、一虚勾合成一个勾圓差」也}。\end{unitzh}
\begin{uniten}2 high altitudes and 1 void difference combine to form 1 altitude-circle difference (guyuancha); 2 level bases and 1 void difference combine to form 1 base-circle difference (gouyuancha) [Siku editors: These 2 items are mistaken; they should say: "2 bright altitudes and 1 void altitude combine to form 1 altitude-circle difference; 2 □ bases and 1 void base combine to form 1 base-circle difference"].\end{uniten}

\begin{unitzh}\unitlabel{ceyuan-004.20}明雙差亦為明、□二大差，其較則明差也；□雙差亦為明、□二小差，其較則□差也。明雙差内減明差，即虚黄；□雙差上加□差，亦同上。以明雙差加明和，即兩明弦；以□雙差加□和，則兩□弦也。以明雙差減明和而半之，即明黄，{\zhsmall 又}為虚大差；以□雙差減於□和而半之，即□黄，{\zhsmall 又}為虚小差也。以虚大差減明和，即為明弦；以虚小差減□和，即□弦也。明雙差、□雙差相較，則次差也。明雙差、□雙差相併，加於明、□二和共，則為兩个極雙差；若以減於明、□二和共，則為兩个虚雙差也。明雙差上加虚雙差，即明、□二股共；□雙差上加虚雙，即明、□二勾共也。\end{unitzh}
\begin{uniten}The bright double difference is also the 2 great differences, bright and □; their difference is the bright difference. The □ double difference is also the 2 small differences, bright and □; their difference is the □ difference. From the bright double difference subtract the bright difference: this is the void yellow; onto the □ double difference add the □ difference: also the same as the above. Add the bright double difference to the bright sum: this is 2 bright hypotenuses; add the □ double difference to the □ sum: then it is 2 □ hypotenuses. Subtract the bright double difference from the bright sum and halve: this is the bright yellow, [Note: also] it is the void great difference; subtract the □ double difference from the □ sum and halve: this is the □ yellow, [Note: also] it is the void small difference. Subtract the void great difference from the bright sum: this is the bright hypotenuse; subtract the void small difference from the □ sum: this is the □ hypotenuse. The bright double difference and the □ double difference compared with each other: then it is the secondary difference. The bright double difference and the □ double difference combined and added to the 2 sums, bright and □, together make 2 pole double differences; if they are subtracted from the 2 sums, bright and □, together, they make 2 void double differences. Onto the bright double difference add the void double difference: this is the 2 altitudes, bright and □, together; onto the □ double difference add the void double: this is the 2 bases, bright and □, together.\end{uniten}

\begin{unitzh}\unitlabel{ceyuan-004.21}以明、□二股共為明弦、□黄共，則髙差、虚黄共為之較{\zhsmall 按：明弦又□黄較}；為明大小差、虚大小差共，則明、□二股共内去兩个虚雙差為之較也{\zhsmall 按：明大小差、虚大小差之較}。以明、□二勾共為□弦、明黄共，則以平差、虚黄較為之較；{\zhsmall 又}為□大小差、虚大小差共，則明、□二勾共内減兩个虚大小差為之較也{\zhsmall 按：虚大小差、□大小差之較}。\end{unitzh}
\begin{uniten}If the 2 altitudes, bright and □, together are taken as the bright hypotenuse and the □ yellow together, then the high difference and the void yellow together are their difference [Siku editors: bright hypotenuse also □ yellow difference] [Translator's note: the small-print 又 "also" seems displaced into this note; cf. the parallel construction below]; if taken as the bright great and small differences and the void great and small differences together, then the 2 altitudes, bright and □, together with 2 void double differences removed are their difference [Siku editors: the difference of the bright great and small differences and the void great and small differences]. If the 2 bases, bright and □, together are taken as the □ hypotenuse and the bright yellow together, then the difference of the level difference and the void yellow is taken as their difference; [Note: also] if taken as the □ great and small differences and the void great and small differences together, then the 2 bases, bright and □, together less 2 void great and small differences are their difference [Siku editors: the difference of the void great and small differences and the □ great and small differences].\end{uniten}

\begin{unitzh}\unitlabel{ceyuan-004.22}明、□二和共内減旁差，即二虚弦；虚弦内加旁差，明股、□勾共也。\end{unitzh}
\begin{uniten}From the 2 sums, bright and □, together subtract the adjacent difference: this is 2 void hypotenuses; to the void hypotenuse add the adjacent difference: the bright altitude and the □ base together.\end{uniten}

\begin{unitzh}\unitlabel{ceyuan-004.23}明和内去平差，即明股、□勾共；□和上加髙差，亦同上也。明和内去髙差，即虚弦；□和上加平差，亦同上。明弦内去髙差，即虚勾；□弦上加平差，即虚股也。明股内去□股，即髙差；去□勾，則極差也。明勾内去□股，即虚差；去□勾，則平差也。\end{unitzh}
\begin{uniten}From the bright sum remove the level difference: this is the bright altitude and the □ base together; onto the □ sum add the high difference: also the same as the above. From the bright sum remove the high difference: this is the void hypotenuse; onto the □ sum add the level difference: also the same as the above. From the bright hypotenuse remove the high difference: this is the void base; onto the □ hypotenuse add the level difference: this is the void altitude. From the bright altitude remove the □ altitude: this is the high difference; remove the □ base: then it is the pole difference. From the bright base remove the □ altitude: this is the void difference; remove the □ base: then it is the level difference.\end{uniten}

\begin{unitzh}\unitlabel{ceyuan-004.24}明、□二股併内減虚弦，即明差；明、□二勾併減於虚弦，即□差。\end{unitzh}
\begin{uniten}From the 2 altitudes, bright and □, combined subtract the void hypotenuse: this is the bright difference; subtract the 2 bases, bright and □, combined from the void hypotenuse: this is the □ difference.\end{uniten}

\begin{unitzh}\unitlabel{ceyuan-004.25}明、□二和共，{\zhsmall 又}為明、□二弦共與明、□二黄共數也，其較則明雙差、□雙差共數也。其明、□二和共數内減旁差，即二虚弦也；若内減虚雙差，即明、□二弦共也。\end{unitzh}
\begin{uniten}The 2 sums, bright and □, together [Note: also] are the number of the 2 hypotenuses, bright and □, together and the 2 yellows, bright and □, together; their difference is the number of the bright double difference and the □ double difference together. From that number of the 2 sums, bright and □, together subtract the adjacent difference: this is 2 void hypotenuses; if the void double difference is subtracted from it, this is the 2 hypotenuses, bright and □, together.\end{uniten}

\begin{unitzh}\unitlabel{ceyuan-004.26}極弦得極差為大差弦，大差弦内減明和，則髙弦内減虚大差也；内減極差則為小差弦，小差弦内減□和，則是平弦内減虚小差也。又大差弦内減明和與髙股共，餘則為虚勾不及明勾數；小差弦内減□和與平勾共，餘則為□股不及虚股數也。\end{unitzh}
\begin{uniten}The pole hypotenuse joined with the pole difference makes the great-difference hypotenuse; from the great-difference hypotenuse subtract the bright sum: then it is the high hypotenuse less the void great difference. Subtract from it the pole difference: then it makes the small-difference hypotenuse; from the small-difference hypotenuse subtract the □ sum: then it is the level hypotenuse less the void small difference. Again, from the great-difference hypotenuse subtract the bright sum and the high altitude together: the remainder then is the amount by which the void base falls short of the bright base; from the small-difference hypotenuse subtract the □ sum and the level base together: the remainder then is the amount by which the □ altitude falls short of the void altitude.\end{uniten}

\begin{unitzh}\unitlabel{ceyuan-004.27}右諸差\end{unitzh}
\begin{uniten}The above: the various differences.\end{uniten}

\begin{unitzh}\unitlabel{ceyuan-004.28}邊勾邊股差，{\zhsmall 又}為皇極差與髙差共也，{\zhsmall 又}為邊弦内去大勾也。邊勾邊弦共，{\zhsmall 又}為大勾邊股共。邊勾邊弦較，{\zhsmall 又}為大差弦内減半徑也。邊股邊弦較，{\zhsmall 又}為□股弦和。\end{unitzh}
\begin{uniten}The difference of the side base and the side altitude [Note: also] is the august-pole difference and the high difference together, [Note: also] it is the side hypotenuse with the great base removed. The side base and the side hypotenuse together [Note: also] are the great base and the side altitude together. The difference of the side base and the side hypotenuse [Note: also] is the great-difference hypotenuse less the half-diameter. The difference of the side altitude and the side hypotenuse [Note: also] is the □ altitude-hypotenuse sum (guxianhe).\end{uniten}

\begin{unitzh}\unitlabel{ceyuan-004.29}底勾底股差，{\zhsmall 又}為皇極差、平差共，{\zhsmall 又}為大股内去底弦，{\zhsmall 又}為髙股内去底小差。底勾底弦共，為大弦内少个底股大勾差。底勾底弦較，{\zhsmall 又}為明弦上勾弦和。底股底弦共，與邊勾邊弦共同。底股底弦較，{\zhsmall 又}為底勾内少小差股也。\end{unitzh}
\begin{uniten}The difference of the bottom base and the bottom altitude [Note: also] is the august-pole difference and the level difference together, [Note: also] it is the great altitude with the bottom hypotenuse removed, [Note: also] it is the high altitude with the bottom small difference removed. The bottom base and the bottom hypotenuse together are the great hypotenuse lacking a difference of the bottom altitude and the great base. The difference of the bottom base and the bottom hypotenuse [Note: also] is the base-hypotenuse sum (gouxianhe) on the bright hypotenuse. The bottom altitude and the bottom hypotenuse together are the same as the side base and the side hypotenuse together. The difference of the bottom altitude and the bottom hypotenuse [Note: also] is the bottom base lacking the small-difference altitude.\end{uniten}

\begin{unitzh}\unitlabel{ceyuan-004.30}邊股内減髙弦，餘則髙股；内減大差弦，餘則明勾；内減底弦，即底股内減大勾也，{\zhsmall 又}為髙弦内減底勾也。\end{unitzh}
\begin{uniten}From the side altitude subtract the high hypotenuse: the remainder then is the high altitude; subtract from it the great-difference hypotenuse: the remainder then is the bright base; subtract from it the bottom hypotenuse: this is the bottom altitude less the great base, [Note: also] it is the high hypotenuse less the bottom base.\end{uniten}

\begin{unitzh}\unitlabel{ceyuan-004.31}底勾内減平弦，餘即平勾；内減小差弦，餘即□股。以底勾減於邊弦，餘即大股内減邊勾也，{\zhsmall 又}為邊股内減平弦也。\end{unitzh}
\begin{uniten}From the bottom base subtract the level hypotenuse: the remainder is the level base; subtract from it the small-difference hypotenuse: the remainder is the □ altitude. Subtract the bottom base from the side hypotenuse: the remainder is the great altitude less the side base, [Note: also] it is the side altitude less the level hypotenuse.\end{uniten}

\begin{unitzh}\unitlabel{ceyuan-004.32}邊弦内減底股與底弦内減邊勾，同為皇極弦内減半徑也。\end{unitzh}
\begin{uniten}The side hypotenuse less the bottom altitude and the bottom hypotenuse less the side base are alike the august-pole hypotenuse less the half-diameter.\end{uniten}

\begin{unitzh}\unitlabel{ceyuan-004.33}皇極勾内減明勾，餘即平勾也；若減□勾，即半徑也，倍之則為底勾、明勾共。皇極股内減□股，餘即髙股也；若減明股，餘即半徑也，倍之則為邊股、□股共也。\end{unitzh}
\begin{uniten}From the august-pole base subtract the bright base: the remainder is the level base; if the □ base is subtracted, this is the half-diameter; doubling it then makes the bottom base and the bright base together. From the august-pole altitude subtract the □ altitude: the remainder is the high altitude; if the bright altitude is subtracted, the remainder is the half-diameter; doubling it then makes the side altitude and the □ altitude together.\end{uniten}

\begin{unitzh}\unitlabel{ceyuan-004.34}明股得虚股即髙股，明勾得虚勾即半徑，□股得虚股即半徑，□勾得虚勾即平勾也。髙弦内減髙股，即□股；平弦内減平勾，即明勾也。明弦内減明差，即虚股；□弦内加□差，即虚勾也。髙股即虚、明二股共，平勾即虚、□二勾共也。明弦、明勾併數與髙股同，□弦、□股併數與平勾同也。\end{unitzh}
\begin{uniten}The bright altitude joined with the void altitude is the high altitude; the bright base joined with the void base is the half-diameter; the □ altitude joined with the void altitude is the half-diameter; the □ base joined with the void base is the level base. From the high hypotenuse subtract the high altitude: this is the □ altitude; from the level hypotenuse subtract the level base: this is the bright base. From the bright hypotenuse subtract the bright difference: this is the void altitude; to the □ hypotenuse add the □ difference: this is the void base. The high altitude is the 2 altitudes, void and bright, together; the level base is the 2 bases, void and □, together. The combined number of the bright hypotenuse and the bright base is the same as the high altitude; the combined number of the □ hypotenuse and the □ altitude is the same as the level base.\end{uniten}

\begin{unitzh}\unitlabel{ceyuan-004.35}明股、□勾相倂減於極弦，即虚和，{\zhsmall 又}為極黄、虚黄共數也。\end{unitzh}
\begin{uniten}The bright altitude and the □ base combined and subtracted from the pole hypotenuse: this is the void sum, [Note: also] it is the number of the pole yellow and the void yellow together.\end{uniten}

\begin{unitzh}\unitlabel{ceyuan-005.1}明□二弦併，内減□雙差，即明□二股併；内減明雙差，即明□二勾併；内加虚弦，即極弦；内減虚弦，即明大差、□小差併也。\end{unitzh}
\begin{uniten}The two hypotenuses (xian) of ming (明) and □ combined (bing): subtract the □ double difference (shuangcha) from it, and it is the two altitudes (gu) of ming and □ combined; subtract the ming double difference from it, and it is the two bases (gou) of ming and □ combined; add the xu (虚) hypotenuse to it, and it is the ji (極) hypotenuse; subtract the xu hypotenuse from it, and it is the ming great difference (da cha) and the □ small difference (xiao cha) combined. [Translator's note: □ is so in the transcription and stands for a character it does not reproduce. It is the name of one of the triangles and is used like ming.]\end{uniten}

\begin{unitzh}\unitlabel{ceyuan-005.2}以明和為明弦、明黄共，則明雙差為之較；以□和為□弦、□黄共，則□雙差為之較也。明和{\zhsmall 又}為髙差、虚弦共，{\zhsmall 又}為極差與明□二勾共數。□和{\zhsmall 又}為平差少於虚弦數，{\zhsmall 又}為極差少於明□二股數。\end{unitzh}
\begin{uniten}Taking the ming sum (he) as the ming hypotenuse and the ming yellow (huang) together (gong), then the ming double difference is their difference (jiao); taking the □ sum as the □ hypotenuse and the □ yellow together, then the □ double difference is their difference. The ming sum [also] is the gao (髙) difference (cha) and the xu hypotenuse together, [also] is the amount of the ji difference and the two bases of ming and □ together. The □ sum [also] is the amount by which the ping (平) difference is less than the xu hypotenuse, [also] is the amount by which the ji difference is less than the two altitudes of ming and □.\end{uniten}

\begin{unitzh}\unitlabel{ceyuan-005.3}半之三事和，内加半黄方，即勾股共；若減之，則弦也。半圓徑内加半虚黄，即虚和；減半虚黄，即虚弦也。{\zhsmall 又}以半虚黄加明和，即髙股；以半虚黄加□和，即平勾也。加明股，則明弦；加□股，則□弦也。減明勾，則明黄；減□股，則□黄也。以虚黄加明黄，則為虚股；以加□黄，則虚勾也。\end{unitzh}
\begin{uniten}The sum of the three things (sanshi he), halved: add half the yellow square (huangfang) to it, and it is the base and altitude together; if it is subtracted, then it is the hypotenuse. Half the circle diameter (yuanjing): add half the xu yellow to it, and it is the xu sum; subtracting half the xu yellow, it is the xu hypotenuse. [Also] adding half the xu yellow to the ming sum, it is the gao altitude; adding half the xu yellow to the □ sum, it is the ping base. Adding the ming altitude, then it is the ming hypotenuse; adding the □ altitude, then it is the □ hypotenuse. Subtracting it from the ming base, then it is the ming yellow; subtracting it from the □ altitude, then it is the □ yellow. Adding the xu yellow to the ming yellow, then it is the xu altitude; adding it to the □ yellow, then it is the xu base.\end{uniten}

\begin{unitzh}\unitlabel{ceyuan-005.4}右諸率弦見。\end{unitzh}
\begin{uniten}The above: the various rates (lü), hypotenuse seen. [Translator's note: 弦見 "hypotenuse seen" is evidently corrupt, perhaps for 互見 "seen in relation to one another".]\end{uniten}

\begin{unitzh}\unitlabel{ceyuan-005.5}髙弦、□弦共為極弦，其差即虚弦、極差共也。髙股、□股共為髙弦，其差即虚股、髙差共也。髙勾、□勾共為平弦，其差即半徑内減□勾也。髙和、□和共為極和，其差即極和内少二□和也。髙差、□差共為極差，其差即虚差、旁差共也。髙黄、□黄共為虚弦，其差即□黄不及虚股數也。{\zhsmall 髙黄即虚股。}髙大差、□大差共即明弦，其差即半虚黄不及明股數也。此髙大差即明股，此□大差即半虚黄也。髙小差{\zhsmall 即□股}、□小差共即□弦，其差即□小差不及□股數也。\end{unitzh}
\begin{uniten}The gao hypotenuse and the □ hypotenuse together are the ji hypotenuse; their difference is the xu hypotenuse and the ji difference together. The gao altitude and the □ altitude together are the gao hypotenuse; their difference is the xu altitude and the gao difference together. The gao base and the □ base together are the ping hypotenuse; their difference is the half-diameter (banjing) with the □ base subtracted from it. The gao sum and the □ sum together are the ji sum; their difference is the ji sum less 2 □ sums. The gao difference and the □ difference together are the ji difference; their difference is the xu difference and the lateral difference (pangcha) together. The gao yellow and the □ yellow together are the xu hypotenuse; their difference is the amount by which the □ yellow falls short of the xu altitude. [Note: The gao yellow is the xu altitude.] The gao great difference and the □ great difference together are the ming hypotenuse; their difference is the amount by which half the xu yellow falls short of the ming altitude. This gao great difference is the ming altitude; this □ great difference is half the xu yellow. The gao small difference [Note: that is, the □ altitude] and the □ small difference together are the □ hypotenuse; their difference is the amount by which the □ small difference falls short of the □ altitude.\end{uniten}

\begin{unitzh}\unitlabel{ceyuan-005.6}明平二弦共亦為極弦，其較即虚弦不及極差數也。明平二股共亦為髙弦，其較即明股内減半徑也。明平二勾共亦為平弦，其較即平差内去虚勾也。明平二和共亦為極和，其較即極和内少二之平和也。明平二差共亦為極差，其較即虚差不及旁差數也。明平二黄共亦為虚弦，其較則虚勾{\zhsmall 按：虚勾即平黄。}不及明黄數也。明平二大差共亦為明弦，其較即明勾不及明大差數。{\zhsmall 平大差即明勾。}明平二小差共亦為□弦，其較則□勾不及半虚黄數也。此明小差即半虚黄，此平小差即□勾。\end{unitzh}
\begin{uniten}The two hypotenuses of ming and ping together are also the ji hypotenuse; their difference is the amount by which the xu hypotenuse falls short of the ji difference. The two altitudes of ming and ping together are also the gao hypotenuse; their difference is the ming altitude with the half-diameter subtracted from it. The two bases of ming and ping together are also the ping hypotenuse; their difference is the ping difference with the xu base removed from it. The two sums of ming and ping together are also the ji sum; their difference is the ji sum less 2 of the ping sum. The two differences of ming and ping together are also the ji difference; their difference is the amount by which the xu difference falls short of the lateral difference. The two yellows of ming and ping together are also the xu hypotenuse; their difference is then the amount by which the xu base [Siku editors: The xu base is the ping yellow.] falls short of the ming yellow. The two great differences of ming and ping together are also the ming hypotenuse; their difference is the amount by which the ming base falls short of the ming great difference. [Note: The ping great difference is the ming base.] The two small differences of ming and ping together are also the □ hypotenuse; their difference is then the amount by which the □ base falls short of half the xu yellow. This ming small difference is half the xu yellow; this ping small difference is the □ base.\end{uniten}

\begin{unitzh}\unitlabel{ceyuan-005.7}右四位相套。\end{unitzh}
\begin{uniten}The above: four positions nested in one another (siwei xiangtao).\end{uniten}

\begin{unitzh}\unitlabel{ceyuan-005.8}邊弦：自減其股，為平勾；自減其勾，為明股、明弦併；減於通弦，餘平弦；減於通股，餘平差；内減通勾，餘邊差；内減底弦，餘極差；内減底股，為半徑、旁差共，{\zhsmall 又}為極弦内少半徑；内減底勾，即大股内去邊勾也；内減黄廣弦，餘□弦；内減黄廣股，即小差股内去平差；内減黄廣勾，即大差股内去平差；内減黄長弦，{\zhsmall 又}得黄長弦{\zhsmall 按：此條誤}；内減黄長股，與内減黄廣勾同；内減黄長勾，即大股内去極勾、虚勾共；内減皇極弦，餘髙弦。\end{unitzh}
\begin{uniten}The bian (邊) hypotenuse: subtracting its own altitude from itself, it is the ping base; subtracting its own base from itself, it is the ming altitude and the ming hypotenuse combined; subtracted from the tong (通) hypotenuse, the remainder is the ping hypotenuse; subtracted from the tong altitude, the remainder is the ping difference; subtracting the tong base from it, the remainder is the bian difference; subtracting the di (底) hypotenuse from it, the remainder is the ji difference; subtracting the di altitude from it, it is the half-diameter and the lateral difference together, [also] it is the ji hypotenuse less the half-diameter; subtracting the di base from it, it is the great altitude (dagu) with the bian base removed; subtracting the huangguang (黄廣) hypotenuse from it, the remainder is the □ hypotenuse; subtracting the huangguang altitude from it, it is the xiaocha (小差) altitude with the ping difference removed; subtracting the huangguang base from it, it is the dacha (大差) altitude with the ping difference removed; subtracting the huangchang (黄長) hypotenuse from it, [also] one obtains the huangchang hypotenuse [Siku editors: This item is in error.]; subtracting the huangchang altitude from it is the same as subtracting the huangguang base from it; subtracting the huangchang base from it, it is the great altitude with the ji base and the xu base together removed; subtracting the huangji (皇極) hypotenuse from it, the remainder is the gao hypotenuse.\end{uniten}

\begin{unitzh}\unitlabel{ceyuan-005.9}底弦：自減其股，為□勾、□弦併；自減其勾，為髙股；減於通弦，餘髙弦；減於通股，餘底差；内減通勾，餘髙差；減於邊弦，餘極差；減於邊股，即底差内去半徑；内減邊勾，即髙差、平勾共；減於黄廣弦，餘為明大差、□小差併{\zhsmall 按：此條亦係數偶合}；減於黄廣股，即底差内去小差股；内減黄廣勾，即一个明弦、一个黄長股弦較；内減去黄長弦，餘明弦；内減黄長股，與内減黄廣勾同；内減黄長勾，餘為髙股、明勾共；内減極弦，為平弦；減於邊股，{\zhsmall 又}為底股内去大勾。\end{unitzh}
\begin{uniten}The di hypotenuse: subtracting its own altitude from itself, it is the □ base and the □ hypotenuse combined; subtracting its own base from itself, it is the gao altitude; subtracted from the tong hypotenuse, the remainder is the gao hypotenuse; subtracted from the tong altitude, the remainder is the di difference; subtracting the tong base from it, the remainder is the gao difference; subtracted from the bian hypotenuse, the remainder is the ji difference; subtracted from the bian altitude, it is the di difference with the half-diameter removed; subtracting the bian base from it, it is the gao difference and the ping base together; subtracted from the huangguang hypotenuse, the remainder is the ming great difference and the □ small difference combined [Siku editors: This item, too, is a chance coincidence of the numbers.]; subtracted from the huangguang altitude, it is the di difference with the xiaocha altitude removed; subtracting the huangguang base from it, it is 1 ming hypotenuse and 1 huangchang altitude–hypotenuse difference; subtracting away the huangchang hypotenuse from it, the remainder is the ming hypotenuse; subtracting the huangchang altitude from it is the same as subtracting the huangguang base from it; subtracting the huangchang base from it, the remainder is the gao altitude and the ming base together; subtracting the ji hypotenuse from it, it is the ping hypotenuse; subtracted from the bian altitude, [also] it is the di altitude with the great base (dagou) removed.\end{uniten}

\begin{unitzh}\unitlabel{ceyuan-005.10}髙差、平差共，{\zhsmall 又}為平勾髙股差。以半徑減髙股，即髙差；半徑内減平勾，即平差。明勾内減□勾，與平差同；明股内減□股，與髙差同。股圓差内減極股，即髙差也；勾圓差減於極勾，即平差。正股内去邊弦，即平差也；底弦内去正勾，即髙差也。大差勾内去極勾，即平差也；極股内去小差股，即髙差也。極差内去□差，即髙差也；内去明差，即平差也。\end{unitzh}
\begin{uniten}The gao difference and the ping difference together [also] are the difference of the ping base and the gao altitude. Subtracting the half-diameter from the gao altitude, it is the gao difference; the half-diameter with the ping base subtracted from it is the ping difference. The ming base with the □ base subtracted from it is the same as the ping difference; the ming altitude with the □ altitude subtracted from it is the same as the gao difference. The altitude–circle difference (gu yuan cha) with the ji altitude subtracted from it is the gao difference; the base–circle difference (gou yuan cha) subtracted from the ji base is the ping difference. The proper altitude (zhenggu) with the bian hypotenuse removed from it is the ping difference; the di hypotenuse with the proper base (zhenggou) removed from it is the gao difference. The dacha base with the ji base removed from it is the ping difference; the ji altitude with the xiaocha altitude removed from it is the gao difference. The ji difference with the □ difference removed from it is the gao difference; with the ming difference removed from it, it is the ping difference.\end{uniten}

\begin{unitzh}\unitlabel{ceyuan-005.11}旁差即城徑極弦較也，{\zhsmall 又}為明差□差較，{\zhsmall 又}為髙差平差較。極差得之為大差差也，去之則為小差差也。\end{unitzh}
\begin{uniten}The lateral difference is the difference of the city diameter (chengjing) and the ji hypotenuse; [also] it is the difference of the ming difference and the □ difference; [also] it is the difference of the gao difference and the ping difference. The ji difference gaining it is the dacha difference; removing it, then it is the xiaocha difference.\end{uniten}

\begin{unitzh}\unitlabel{ceyuan-005.12}又髙差平差下：明和内去虚弦，即髙差；虚弦内去□和，即平差。\end{unitzh}
\begin{uniten}Further, under the gao difference and the ping difference: the ming sum with the xu hypotenuse removed from it is the gao difference; the xu hypotenuse with the □ sum removed from it is the ping difference.\end{uniten}

\begin{unitzh}\unitlabel{ceyuan-005.13}大差弦内加虚差，即黄廣股；小差股内減虚差，即黄長勾。\end{unitzh}
\begin{uniten}The dacha hypotenuse with the xu difference added to it is the huangguang altitude; the xiaocha altitude with the xu difference subtracted from it is the huangchang base.\end{uniten}

\begin{unitzh}\unitlabel{ceyuan-005.14}通差内去髙差，即底差；内去平差，即邊差也。虚大差得二虚勾，即勾圓差之股；虚小差得二虚股，即股圓差之勾也。\end{unitzh}
\begin{uniten}The tong difference with the gao difference removed from it is the di difference; with the ping difference removed from it, it is the bian difference. The xu great difference gaining 2 xu bases is the altitude of the base–circle difference; the xu small difference gaining 2 xu altitudes is the base of the altitude–circle difference.\end{uniten}

\begin{unitzh}\unitlabel{ceyuan-005.15}明股弦較與勾共，即虚股也；□勾弦較與股共，即虚勾也。\end{unitzh}
\begin{uniten}The ming altitude–hypotenuse difference together with the base is the xu altitude; the □ base–hypotenuse difference together with the altitude is the xu base.\end{uniten}

\begin{unitzh}\unitlabel{ceyuan-005.16}半虚黄：□勾得之，即□弦也；減於此數，即虚黄内去□弦也。□股得之，虚勾也；去之，即□黄方也。□弦得之，即平勾内去□黄也；去之，則□勾也。明勾内得之，即虚股也；去之，則明黄方也。明股得之，即明弦也；去之，則明弦内去个虚黄方也。明弦得之，即髙股内去明黄也；去之，則明股也。\end{unitzh}
\begin{uniten}Half the xu yellow: the □ base gaining it is the □ hypotenuse; subtracted from this amount, it is the xu yellow with the □ hypotenuse removed. The □ altitude gaining it is the xu base; removing it, it is the □ yellow square. The □ hypotenuse gaining it is the ping base with the □ yellow removed; removing it, then it is the □ base. The ming base gaining it within is the xu altitude; removing it, then it is the ming yellow square. The ming altitude gaining it is the ming hypotenuse; removing it, then it is the ming hypotenuse with a xu yellow square removed. The ming hypotenuse gaining it is the gao altitude with the ming yellow removed; removing it, then it is the ming altitude.\end{uniten}

\begin{unitzh}\unitlabel{ceyuan-005.17}右拾遺。\end{unitzh}
\begin{uniten}The above: gleanings (shiyi).\end{uniten}

\begin{unitzh}\unitlabel{ceyuan-005.18}按：「識别雜記」約五百條，皆隨時録其所得未經審定者，故難易淺深不拘先後，要皆精思妙義，足以開示數理之蘊奥者。徐光啟亟𫝊新法，而於《勾股義》中獨推是書，其必有所見矣。\end{unitzh}
\begin{uniten}Note: The "Miscellaneous Notes on Identification" (Shibie zaji), about 500 items, are all records made from time to time of what was found, not yet examined and settled; therefore difficult and easy, shallow and deep, are not bound to an order of before and after. In essence, all are refined thought and subtle meaning, sufficient to open up and reveal the hidden depths of the principles of number. Xu Guangqi eagerly transmitted the new methods, yet in the Gougu yi (The Meaning of Base and Altitude) he singled out this book alone for esteem; he must surely have seen something in it.\end{uniten}

\begin{unitzh}\unitlabel{ceyuan-005.19}測圓海鏡卷一\end{unitzh}
\begin{uniten}Sea Mirror of Circle Measurements (Ceyuan haijing), juan 1.\end{uniten}

\section*{卷02}\addcontentsline{toc}{section}{卷02}

\begin{unitzh}\unitlabel{ceyuan-006.1}欽定四庫全書\end{unitzh}
\begin{uniten}Imperially Authorized Complete Library of the Four Treasuries (Qinding Siku quanshu)\end{uniten}

\begin{unitzh}\unitlabel{ceyuan-006.2}測圓海鏡卷二\end{unitzh}
\begin{uniten}Sea Mirror of Circle Measurements, juan 2\end{uniten}

\begin{unitzh}\unitlabel{ceyuan-006.3}元·李冶撰\end{unitzh}
\begin{uniten}Composed by Li Ye of the Yuan\end{uniten}

\begin{unitzh}\unitlabel{ceyuan-006.4}正率一十四問\end{unitzh}
\begin{uniten}Standard rates (zhenglü): 14 problems\end{uniten}

\begin{unitzh}\unitlabel{ceyuan-006.5}假令有圓城一所，不知周徑。四面開門，門外縱横各有十字大道。其西北十字道頭定為乾地，其東北十字道頭定為艮地，其東南十字道頭定為㢲地，其西南十字道頭定為坤地。所有測望雜法，一一設問如後。\end{unitzh}
\begin{uniten}Suppose there is a circular city (yuancheng) whose circumference and diameter are not known. Gates are opened on its four sides, and outside the gates, lengthwise and crosswise, there are great crossroads (shizi dadao) at each. The head of the crossroads at its northwest is designated point Qian (Qian di); the head of the crossroads at its northeast is designated point Gen; the head of the crossroads at its southeast is designated point Xun; the head of the crossroads at its southwest is designated point Kun. All the various methods of measuring by sighting are set out as problems one by one, as follows.\end{uniten}

\begin{unitzh}\unitlabel{ceyuan-006.6}或問：甲乙二人俱在乾地，乙東行三百二十步而立，甲南行六百步望見乙。問徑幾里？答曰：城徑二百四十步。\end{unitzh}
\begin{uniten}Someone asks: two persons, Jia and Yi, are both at point Qian. Yi walks east 320 bu (步) and stands; Jia walks south 600 bu and sights Yi. Question: how many li (里) is the diameter? Answer: The diameter of the city is 240 bu.\end{uniten}

\begin{unitzh}\unitlabel{ceyuan-006.7}法曰：此為勾股容圓也。以勾股相乗倍之為實，併勾股冪以求弦，復加入勾股共以為法。\end{unitzh}
\begin{uniten}Rule: This is a circle contained in a right triangle (gougu rong yuan). Multiply the base (gou) and the altitude (gu) together and double it to make the dividend (shi); combine the areas (mi) of the base and the altitude in order to find the hypotenuse (xian); further add into it the base and altitude together, to make the divisor (fa).\end{uniten}

\begin{unitzh}\unitlabel{ceyuan-006.8}草曰：置甲南行六百步在地，以乙東行三百二十步乘之，得一十九萬二千步，倍之，得三十八萬四千步為實。以乙東行步自之，得一十萬零二千四百步為勾冪；以甲南行步自之，得三十六萬步為股冪。二冪相併，得四十六萬二千四百步為弦方實，以平方開之，得六百八十步，則弦也。以弦加勾股共，共得一千六百步，以為法。如法而一，得二百四十步，則城徑也。合問。\end{unitzh}
\begin{uniten}Working: Set down Jia's southward walk of 600 bu on the ground; multiply it by Yi's eastward walk of 320 bu, obtaining 192000 bu; double it, obtaining 384000 bu, as the dividend. Multiply Yi's eastward bu by itself, obtaining 102400 bu, as the area of the base (gou mi); multiply Jia's southward bu by itself, obtaining 360000 bu, as the area of the altitude (gu mi). Add the two areas together, obtaining 462400 bu, as the square dividend of the hypotenuse (xian fang shi); extract its square root, obtaining 680 bu, which is the hypotenuse. Add the hypotenuse to the base and altitude together; altogether one obtains 1600 bu, which is taken as the divisor. Divide by the divisor, obtaining 240 bu, which is the diameter of the city. This agrees with the question.\end{uniten}

\begin{unitzh}\unitlabel{ceyuan-006.9}或問：甲乙二人俱在西門，乙東行二百五十六步，甲南行四百八十步望見乙。問答同前。\end{unitzh}
\begin{uniten}Someone asks: two persons, Jia and Yi, are both at the west gate. Yi walks east 256 bu; Jia walks south 480 bu and sights Yi. The question and answer are the same as before.\end{uniten}

\begin{unitzh}\unitlabel{ceyuan-006.10}法曰：此為勾上容圓也。以勾股相乘倍之為實，併勾股冪以求弦，加入股以為法。\end{unitzh}
\begin{uniten}Rule: This is a circle contained on the base (gou shang rong yuan). Multiply the base and altitude together and double it to make the dividend; combine the areas of the base and altitude in order to find the hypotenuse; add the altitude into it to make the divisor.\end{uniten}

\begin{unitzh}\unitlabel{ceyuan-006.11}草曰：置甲南行四百八十步在地，以乙東行二百五十六步乘之，得一十二萬二千八百八十步，倍之，得二十四萬五千七百六十步為實。以乙東行步自之，得六萬五千五百三十六步為勾冪；以甲南行步自之，得二十三萬零四百步為股冪。勾股冪相併，得二十九萬五千九百三十六步為弦方實，以平方開之，得五百四十四步為弦也。以加入南行步，共得一千零二十四步，以為法而一，得二百四十步，則城徑。合問。\end{unitzh}
\begin{uniten}Working: Set down Jia's southward walk of 480 bu on the ground; multiply it by Yi's eastward walk of 256 bu, obtaining 122880 bu; double it, obtaining 245760 bu, as the dividend. Multiply Yi's eastward bu by itself, obtaining 65536 bu, as the area of the base; multiply Jia's southward bu by itself, obtaining 230400 bu, as the area of the altitude. Add the areas of the base and altitude together, obtaining 295936 bu, as the square dividend of the hypotenuse; extract its square root, obtaining 544 bu, as the hypotenuse. Add it into the southward bu, altogether obtaining 1024 bu; take this as the divisor and divide, obtaining 240 bu, which is the diameter of the city. This agrees with the question.\end{uniten}

\begin{unitzh}\unitlabel{ceyuan-006.12}或問：甲乙二人俱在北門，乙東行二百步而止，甲南行三百七十五步望見乙。問答同前。\end{unitzh}
\begin{uniten}Someone asks: two persons, Jia and Yi, are both at the north gate. Yi walks east 200 bu and stops; Jia walks south 375 bu and sights Yi. The question and answer are the same as before.\end{uniten}

\begin{unitzh}\unitlabel{ceyuan-006.13}法曰：此為股上容圓也。以勾股相乘倍之為實，以勾股冪求弦，加入勾以為法。\end{unitzh}
\begin{uniten}Rule: This is a circle contained on the altitude (gu shang rong yuan). Multiply the base and altitude together and double it to make the dividend; find the hypotenuse from the areas of the base and altitude, and add the base into it to make the divisor.\end{uniten}

\begin{unitzh}\unitlabel{ceyuan-006.14}草曰：置甲南行三百七十五步，以乙東行二百步乘之，得七萬五千步，倍之，得一十五萬步為實。以乙東行自之，得四萬步為勾冪；以甲南行自之，得一十四萬零六百二十五步為股冪。勾股冪相併，得一十八萬零六百二十五步為弦方實，如平方而一，得四百二十五步，則弦也。加入乙東行二百步，共得六百二十五步，以為法，以法除之，得二百四十步，則城徑也。合問。\end{unitzh}
\begin{uniten}Working: Set down Jia's southward walk of 375 bu; multiply it by Yi's eastward walk of 200 bu, obtaining 75000 bu; double it, obtaining 150000 bu, as the dividend. Multiply Yi's eastward walk by itself, obtaining 40000 bu, as the area of the base; multiply Jia's southward walk by itself, obtaining 140625 bu, as the area of the altitude. Add the areas of the base and altitude together, obtaining 180625 bu, as the square dividend of the hypotenuse; divide it by the square-root method (ru pingfang er yi), obtaining 425 bu, which is the hypotenuse. Add in Yi's eastward walk of 200 bu, altogether obtaining 625 bu, which is taken as the divisor; divide it by the divisor, obtaining 240 bu, which is the diameter of the city. This agrees with the question.\end{uniten}

\begin{unitzh}\unitlabel{ceyuan-006.15}或問：甲乙二人俱在圓城中心而立，乙穿城向東行一百三十六步而止，甲穿城南行二百五十五步望見乙。問答同前。\end{unitzh}
\begin{uniten}Someone asks: two persons, Jia and Yi, both stand at the centre of the circular city. Yi goes through the city, walking toward the east 136 bu, and stops; Jia goes through the city, walking south 255 bu, and sights Yi. The question and answer are the same as before.\end{uniten}

\begin{unitzh}\unitlabel{ceyuan-006.16}法曰：此為勾股上容圓也。以勾股相乘倍之為實，併勾股冪，如法求弦，以為法。\end{unitzh}
\begin{uniten}Rule: This is a circle contained on the base and altitude (gougu shang rong yuan). Multiply the base and altitude together and double it to make the dividend; combine the areas of the base and altitude, find the hypotenuse according to the method, and take it as the divisor.\end{uniten}

\begin{unitzh}\unitlabel{ceyuan-006.17}草曰：以二行步相乘，得三萬四千六百八十步，倍之，得六萬九千三百六十步為實。置乙東行自之，得一萬八千四百九十六步為勾冪；又以甲南行自之，得六萬五千零二十五步為股冪。二冪相併，得八萬三千五百二十一步為弦方實，以平方開之，得二百八十九步，即弦也。便以為法，如法除實，得二百四十步，即圓城之徑也。合問。\end{unitzh}
\begin{uniten}Working: Multiply the bu of the two walks together, obtaining 34680 bu; double it, obtaining 69360 bu, as the dividend. Set down Yi's eastward walk and multiply it by itself, obtaining 18496 bu, as the area of the base; further multiply Jia's southward walk by itself, obtaining 65025 bu, as the area of the altitude. Add the two areas together, obtaining 83521 bu, as the square dividend of the hypotenuse; extract its square root, obtaining 289 bu, which is the hypotenuse. Thereupon take it as the divisor; divide the dividend by the divisor, obtaining 240 bu, which is the diameter of the circular city. This agrees with the question.\end{uniten}

\begin{unitzh}\unitlabel{ceyuan-006.18}或問：甲乙二人同立於乾地，乙東行一百八十步遇塔而止，甲南行三百六十步，回望其塔，正居城徑之半。問答同前。\end{unitzh}
\begin{uniten}Someone asks: two persons, Jia and Yi, stand together at point Qian. Yi walks east 180 bu, comes upon a pagoda and stops; Jia walks south 360 bu and looks back at the pagoda: it lies exactly at the half of the city's diameter. The question and answer are the same as before.\end{uniten}

\begin{unitzh}\unitlabel{ceyuan-006.19}法曰：此為弦上容圓也。以勾股相乘倍之為實，以勾股和為法。\end{unitzh}
\begin{uniten}Rule: This is a circle contained on the hypotenuse (xian shang rong yuan). Multiply the base and altitude together and double it to make the dividend; take the sum of base and altitude (gougu he) as the divisor.\end{uniten}

\begin{unitzh}\unitlabel{ceyuan-006.20}草曰：以二行步相乘，得六萬四千八百步，倍之，得一十二萬九千六百步為實。併二行步，得五百四十步，以為法除實，得二百四十步，即城徑也。合問。\end{unitzh}
\begin{uniten}Working: Multiply the bu of the two walks together, obtaining 64800 bu; double it, obtaining 129600 bu, as the dividend. Add the bu of the two walks together, obtaining 540 bu; take it as the divisor and divide the dividend, obtaining 240 bu, which is the diameter of the city. This agrees with the question.\end{uniten}

\begin{unitzh}\unitlabel{ceyuan-006.21}或問：甲乙二人俱在坤地，乙東行一百九十二步而止，甲南行三百六十步，望乙與城㕘相直。問答同前。\end{unitzh}
\begin{uniten}Someone asks: two persons, Jia and Yi, are both at point Kun. Yi walks east 192 bu and stops; Jia walks south 360 bu and looks toward Yi in line with the city (can xiang zhi). The question and answer are the same as before.\end{uniten}

\begin{unitzh}\unitlabel{ceyuan-006.22}法曰：此為勾外容圓也。以勾股相乘倍之為實，以弦較和為法。\end{unitzh}
\begin{uniten}Rule: This is a circle contained outside the base (gou wai rong yuan). Multiply the base and altitude together and double it to make the dividend; take the sum of the hypotenuse and the difference (xian jiao he) as the divisor.\end{uniten}

\begin{unitzh}\unitlabel{ceyuan-006.23}草曰：以二行步相乘，得六萬九千一百二十步，倍之，得一十三萬八千二百四十步為實。置乙東行自之，得三萬六千八百六十四步為勾冪；又置甲南行自之，得一十二萬九千六百步為股冪。二冪相併，得一十六萬六千四百六十四步為弦方實，以平方開之，得四百零八，即弦也。又置甲南行步，内減乙東行步，餘一百六十八步，即較也。以較加弦，共得五百七十六步，以為法。實如法而一，得二百四十步，為城徑也。合問。\end{unitzh}
\begin{uniten}Working: Multiply the bu of the two walks together, obtaining 69120 bu; double it, obtaining 138240 bu, as the dividend. Set down Yi's eastward walk and multiply it by itself, obtaining 36864 bu, as the area of the base; further set down Jia's southward walk and multiply it by itself, obtaining 129600 bu, as the area of the altitude. Add the two areas together, obtaining 166464 bu, as the square dividend of the hypotenuse; extract its square root, obtaining 408, which is the hypotenuse. Further set down Jia's southward bu and subtract Yi's eastward bu from it; the remainder, 168 bu, is the difference (jiao). Add the difference to the hypotenuse, altogether obtaining 576 bu, which is taken as the divisor. Divide the dividend by the divisor, obtaining 240 bu, as the diameter of the city. This agrees with the question.\end{uniten}

\begin{unitzh}\unitlabel{ceyuan-006.24}按：此題用勾股求得弦，即可加減得弦較較為城徑。今必以勾股相乘倍積為實，求得弦，加減得弦較和為法，而後始得弦較較為城徑者，蓋欲因此並明勾股相乘之倍積為弦較較、弦較和相乘之積，非故為紆廻也。\end{unitzh}
\begin{uniten}Note: In this problem, once the hypotenuse has been found by means of the base and altitude, one could at once add and subtract to obtain the difference of the hypotenuse and the difference (xian jiao jiao) as the diameter of the city. Now, that one must take the doubled area (bei ji) of base and altitude multiplied together as the dividend, find the hypotenuse, add and subtract to obtain the sum of the hypotenuse and the difference as the divisor, and only then obtain the difference of the hypotenuse and the difference as the diameter of the city, is presumably because one wishes thereby also to make clear that the doubled area of base and altitude multiplied together is the area of the difference of the hypotenuse and the difference and the sum of the hypotenuse and the difference multiplied together; it is not deliberately taking a roundabout course.\end{uniten}

\begin{unitzh}\unitlabel{ceyuan-006.25}或問：甲乙二人同立於艮地，甲南行一百五十步而止，乙東行八十步，望乙與城㕘相直。問答同前。\end{unitzh}
\begin{uniten}Someone asks: two persons, Jia and Yi, stand together at point Gen. Jia walks south 150 bu and stops; Yi walks east 80 bu, and Yi is sighted in line with the city. The question and answer are the same as before.\end{uniten}

\begin{unitzh}\unitlabel{ceyuan-006.26}法曰：此為股外容圓也。以勾股相乘倍之為實，以弦較較為法。\end{unitzh}
\begin{uniten}Rule: This is a circle contained outside the altitude (gu wai rong yuan). Multiply the base and altitude together and double it to make the dividend; take the difference of the hypotenuse and the difference as the divisor.\end{uniten}

\begin{unitzh}\unitlabel{ceyuan-006.27}草曰：二行步相乘，得一萬二千，倍之，得二萬四千步為實。以甲南行自之，得二萬二千五百步為股冪；又以乙東行步自之，得六千四百步為勾冪。勾股冪相併，得二萬八千九百步為弦方實，以平方開之，得一百七十步，即弦也。以二行步相減，餘七十步，為勾股較也。以此較又減弦，餘一百步，即弦較較也。便以為法，實如法而一，得二百四十步，即城徑也。合問。\end{unitzh}
\begin{uniten}Working: Multiply the bu of the two walks together, obtaining 12000; double it, obtaining 24000 bu, as the dividend. Multiply Jia's southward walk by itself, obtaining 22500 bu, as the area of the altitude; further multiply Yi's eastward bu by itself, obtaining 6400 bu, as the area of the base. Add the areas of the base and altitude together, obtaining 28900 bu, as the square dividend of the hypotenuse; extract its square root, obtaining 170 bu, which is the hypotenuse. Subtract the bu of the two walks from each other; the remainder, 70 bu, is the difference of base and altitude (gougu jiao). Further subtract this difference from the hypotenuse; the remainder, 100 bu, is the difference of the hypotenuse and the difference. Thereupon take it as the divisor; divide the dividend by the divisor, obtaining 240 bu, which is the diameter of the city. This agrees with the question.\end{uniten}

\begin{unitzh}\unitlabel{ceyuan-006.28}按：此題係弦較和為城徑，其用法實以較取和之意，與上題同。\end{unitzh}
\begin{uniten}Note: In this problem the sum of the hypotenuse and the difference is the diameter of the city; its idea of using the divisor and the dividend to obtain the sum by means of the difference is the same as in the preceding problem.\end{uniten}

\begin{unitzh}\unitlabel{ceyuan-006.29}或問：甲乙二人同立於㢲地，乙西行四十八步而止，甲北行九十步，望乙與城㕘相直。問答同前。\end{unitzh}
\begin{uniten}Someone asks: two persons, Jia and Yi, stand together at point Xun. Yi walks west 48 bu and stops; Jia walks north 90 bu and looks toward Yi in line with the city. The question and answer are the same as before.\end{uniten}

\begin{unitzh}\unitlabel{ceyuan-006.30}法曰：此為弦外容圓也。勾股相乘倍之為實，以弦和較為法。\end{unitzh}
\begin{uniten}Rule: This is a circle contained outside the hypotenuse (xian wai rong yuan). Base and altitude are multiplied together and doubled to make the dividend; take the difference of the hypotenuse and the sum (xian he jiao) as the divisor.\end{uniten}

\begin{unitzh}\unitlabel{ceyuan-006.31}草曰：以二行步相乘，得四千三百二十步，倍之，得八千六百四十步為實。以甲北行自之，得八千一百步為股冪；又以乙西行自之，得二千三百零四步為勾冪。併二冪，得一萬零四百零四步為弦方實，以平方開之，得一百零二步，為弦也。又併二行步，得一百三十八步，為和。以弦減和，餘三十六步，得黄方，以為法。實如法而一，得二百四十步，即城徑也。合問。\end{unitzh}
\begin{uniten}Working: Multiply the bu of the two walks together, obtaining 4320 bu; double it, obtaining 8640 bu, as the dividend. Multiply Jia's northward walk by itself, obtaining 8100 bu, as the area of the altitude; further multiply Yi's westward walk by itself, obtaining 2304 bu, as the area of the base. Combine the two areas, obtaining 10404 bu, as the square dividend of the hypotenuse; extract its square root, obtaining 102 bu, as the hypotenuse. Further add the bu of the two walks together, obtaining 138 bu, as the sum (he). Subtract the hypotenuse from the sum; the remainder, 36 bu, gives the yellow square (huang fang), which is taken as the divisor. Divide the dividend by the divisor, obtaining 240 bu, which is the diameter of the city. This agrees with the question.\end{uniten}

\begin{unitzh}\unitlabel{ceyuan-006.32}按：此題弦和和即城徑。其以勾股相乘倍積為實、黄方為法者，亦以明弦和和、黄方相乘之積，與勾股相乘之倍積為相等也。\end{unitzh}
\begin{uniten}Note: In this problem the sum of the hypotenuse and the sum (xian he he) is the diameter of the city. That it takes the doubled area of base and altitude multiplied together as the dividend and the yellow square as the divisor is likewise in order to make clear that the area of the sum of the hypotenuse and the sum and the yellow square multiplied together is equal to the doubled area of base and altitude multiplied together.\end{uniten}

\begin{unitzh}\unitlabel{ceyuan-006.33}或問：甲乙二人俱在南門，乙東行七十二步而止，甲南行一百三十五步，望乙與城㕘相直。問答同前。\end{unitzh}
\begin{uniten}Someone asks: two persons, Jia and Yi, are both at the south gate. Yi walks east 72 bu and stops; Jia walks south 135 bu and looks toward Yi in line with the city. The question and answer are the same as before.\end{uniten}

\begin{unitzh}\unitlabel{ceyuan-006.34}法曰：此為勾外容圓半也。以勾股相乘倍之為實，以大差為法。\end{unitzh}
\begin{uniten}Rule: This is a half circle contained outside the base (gou wai rong yuan ban). Multiply the base and altitude together and double it to make the dividend; take the greater difference (da cha) as the divisor.\end{uniten}

\begin{unitzh}\unitlabel{ceyuan-006.35}草曰：以二行步相乘，得九千七百二十步，倍之，得一萬九千四百四十步為實。又以乙東行自之，得五千一百八十四步為勾冪；又以南行自之，得一萬八千二百二十五步為股冪。二冪相併，得二萬三千四百零九步為弦方實，以平方開之，得一百五十三步，即弦也。以乙東行七十二步為勾，以減弦，餘八十一步，即勾弦差也。便以為法，實如法而一，得二百四十步，即城徑也。合問。\end{unitzh}
\begin{uniten}Working: Multiply the bu of the two walks together, obtaining 9720 bu; double it, obtaining 19440 bu, as the dividend. Further multiply Yi's eastward walk by itself, obtaining 5184 bu, as the area of the base; further multiply the southward walk by itself, obtaining 18225 bu, as the area of the altitude. Add the two areas together, obtaining 23409 bu, as the square dividend of the hypotenuse; extract its square root, obtaining 153 bu, which is the hypotenuse. Take Yi's eastward walk of 72 bu as the base and subtract it from the hypotenuse; the remainder, 81 bu, is the difference of base and hypotenuse (gou xian cha). Thereupon take it as the divisor; divide the dividend by the divisor, obtaining 240 bu, which is the diameter of the city. This agrees with the question.\end{uniten}

\begin{unitzh}\unitlabel{ceyuan-006.36}或問：甲乙二人俱在東門，甲南行三十步而止，乙東行一十六步，回望甲與城㕘相直。問答同前。\end{unitzh}
\begin{uniten}Someone asks: two persons, Jia and Yi, are both at the east gate. Jia walks south 30 bu and stops; Yi walks east 16 bu and looks back toward Jia in line with the city. The question and answer are the same as before.\end{uniten}

\begin{unitzh}\unitlabel{ceyuan-006.37}法曰：此為股外容圓半也。以勾股相乘倍之為實，以小差為法。\end{unitzh}
\begin{uniten}Rule: This is a half circle contained outside the altitude (gu wai rong yuan ban). Multiply the base and altitude together and double it to make the dividend; take the lesser difference (xiao cha) as the divisor.\end{uniten}

\begin{unitzh}\unitlabel{ceyuan-006.38}草曰：以二行步相乘，得四百八十步，倍之，得九百六十步為實。又以乙東行自之，得二百五十六步為勾冪；又以甲南行自之，得九百步為股冪。二冪相併，得一千一百五十六步為弦方實，以平方開之，得三十四步，即弦也。以甲南行三十步為股，以減弦，餘四步，以為法。以法除實，得二百四十步，即城徑也。合問。\end{unitzh}
\begin{uniten}Working: Multiply the bu of the two walks together, obtaining 480 bu; double it, obtaining 960 bu, as the dividend. Further multiply Yi's eastward walk by itself, obtaining 256 bu, as the area of the base; further multiply Jia's southward walk by itself, obtaining 900 bu, as the area of the altitude. Add the two areas together, obtaining 1156 bu, as the square dividend of the hypotenuse; extract its square root, obtaining 34 bu, which is the hypotenuse. Take Jia's southward walk of 30 bu as the altitude and subtract it from the hypotenuse; the remainder, 4 bu, is taken as the divisor. Divide the dividend by the divisor, obtaining 240 bu, which is the diameter of the city. This agrees with the question.\end{uniten}

\begin{unitzh}\unitlabel{ceyuan-006.39}或問：甲出西門南行四百八十步而止，乙出東門南行三十步望見甲。問答同前。\end{unitzh}
\begin{uniten}Someone asks: Jia goes out of the west gate, walks south 480 bu and stops; Yi goes out of the east gate, walks south 30 bu and sights Jia. The question and answer are the same as before.\end{uniten}

\begin{unitzh}\unitlabel{ceyuan-006.40}法曰：此為半矮梯也。以二行步相乘為實，如平方而一，得半徑。\end{unitzh}
\begin{uniten}Rule: This is a half low trapezoid (ban aiti). Multiply the bu of the two walks together to make the dividend; divide it by the square-root method, obtaining the radius (banjing).\end{uniten}

\begin{unitzh}\unitlabel{ceyuan-006.41}草曰：以二行步相乘，得一萬四千四百步為實，以平方開之，得一百二十步，倍之即城徑也。合問。\end{unitzh}
\begin{uniten}Working: Multiply the bu of the two walks together, obtaining 14400 bu, as the dividend; extract its square root, obtaining 120 bu; double it: this is the diameter of the city. This agrees with the question.\end{uniten}

\begin{unitzh}\unitlabel{ceyuan-006.42}又問：甲乙二人，乙出南門折而東行七十二步而止，甲出北門折而東行二百，望見乙。問答同前。\end{unitzh}
\begin{uniten}It is further asked: two persons, Jia and Yi: Yi goes out of the south gate, turns and walks east 72 bu and stops; Jia goes out of the north gate, turns and walks east 200, and sights Yi. The question and answer are the same as before.\end{uniten}

\begin{unitzh}\unitlabel{ceyuan-007.1}法曰：以二行步相乘，得數四之為實，如平方而一，得城徑。\end{unitzh}
\begin{uniten}Rule: Multiply the bu (步) of the two walks together; quadruple the number obtained to make the dividend (shi); divide by the square-root procedure (ru pingfang er yi), and you obtain the diameter of the city.\end{uniten}

\begin{unitzh}\unitlabel{ceyuan-007.2}草曰：二行步相乘，得一萬四千四百步，又四之，得五萬七千六百步為實，以平方開之，得二百四十步，即城徑也。合問。\end{unitzh}
\begin{uniten}Working: Multiplying the bu of the two walks together gives 14400 bu; quadrupling it then gives 57600 bu, which makes the dividend (shi). Extracting its square root gives 240 bu, which is the diameter of the city. This fits the question.\end{uniten}

\begin{unitzh}\unitlabel{ceyuan-007.3}又假令乙出南門折東行二十步，甲出北門折東行七百二十步，如此之類，亦同上法。{\zhsmall 以上三問是以半矮梯求之。}按：右三題通為一問。\end{unitzh}
\begin{uniten}Again, suppose B goes out of the south gate and turns to walk east 20 bu, and A goes out of the north gate and turns to walk east 720 bu: cases of this kind likewise follow the rule above. [Note: The three questions above are all solved by means of the half low trapezoid (ban aiti).] Note: the three problems above are all one question.\end{uniten}

\begin{unitzh}\unitlabel{ceyuan-007.4}或問：甲乙二人，乙在艮地東行八十步而立，甲在坤地南行三百六十步望見乙。問答同前。\end{unitzh}
\begin{uniten}Someone asks: There are two persons, A and B. B, at the position Gen (艮), walks east 80 bu and stands; A, at the position Kun (坤), walks south 360 bu and catches sight of B. The question and answer are the same as before.\end{uniten}

\begin{unitzh}\unitlabel{ceyuan-007.5}法曰：此為兩差求黄方也。以二行步相乘，倍之為實，以平方開，得城徑。\end{unitzh}
\begin{uniten}Rule: This is finding the yellow square (huangfang) from the two differences. Multiply the bu of the two walks together and double it to make the dividend (shi); extract the square root, and you obtain the diameter of the city.\end{uniten}

\begin{unitzh}\unitlabel{ceyuan-007.6}草曰：二行步相乘，得二萬八千八百步，倍之，得五萬七千六百步為實，以平方開之，得二百四十步，即城徑也。合問。别得甲南行即股圓差也，乙東行即勾圓差也。\end{unitzh}
\begin{uniten}Working: Multiplying the bu of the two walks together gives 28800 bu; doubling it gives 57600 bu, which makes the dividend (shi). Extracting its square root gives 240 bu, which is the diameter of the city. This fits the question. It is separately obtained that A's walk south is the altitude–circle difference (gu yuan cha), and B's walk east is the base–circle difference (gou yuan cha).\end{uniten}

\begin{unitzh}\unitlabel{ceyuan-007.7}或問：甲出東門四十八步而立，乙出南門四十八步見之。問答同前。\end{unitzh}
\begin{uniten}Someone asks: A goes out of the east gate 48 bu and stands; B goes out of the south gate 48 bu and sees him. The question and answer are the same as before.\end{uniten}

\begin{unitzh}\unitlabel{ceyuan-007.8}法曰：此當以方五斜七求之，每出門二步管徑十步。\end{unitzh}
\begin{uniten}Rule: This should be solved with "square 5, diagonal 7" (fang wu xie qi): every 2 bu of going out of the gate accounts for 10 bu of the diameter.\end{uniten}

\begin{unitzh}\unitlabel{ceyuan-007.9}草曰：置出門步在地，以五之，得二百四十步，即城徑也。據此法，合置出門步在地，以十之，二而一，以二數相折，故五因便是。合問。\end{unitzh}
\begin{uniten}Working: Set down the bu out of the gate on the ground and multiply it by 5, obtaining 240 bu, which is the diameter of the city. According to this rule, one ought to set down the bu out of the gate on the ground, multiply it by 10 and divide by 2; because the two numbers offset each other (xiangzhe), multiplying by 5 (wu yin) is just right. This fits the question.\end{uniten}

\begin{unitzh}\unitlabel{ceyuan-007.10}按：方五斜七，古率，非密率也。設問以盡此題之變，故率之踈密勿論。\end{unitzh}
\begin{uniten}Note: "Square 5, diagonal 7" is an ancient rate (lü), not a precise rate (milü). The question is set in order to exhaust the variations of this problem; therefore the coarseness or precision of the rate is not considered.\end{uniten}

\begin{unitzh}\unitlabel{ceyuan-007.11}或問：出西門南行四百八十步有樹，出北門東行二百步見之。問答同前。\end{unitzh}
\begin{uniten}Someone asks: Going out of the west gate and walking south 480 bu there is a tree; going out of the north gate and walking east 200 bu one sees it. The question and answer are the same as before.\end{uniten}

\begin{unitzh}\unitlabel{ceyuan-007.12}法曰：以二行步相乗為實，二行步相併為從，二步常法，得半徑。\end{unitzh}
\begin{uniten}Rule: Multiply the bu of the two walks together to make the dividend (shi); add the bu of the two walks together to make the following (cong); 2 bu is the constant divisor (changfa); you obtain the radius. [Translator's note: so the text; the equation in the working below has 1 square.]\end{uniten}

\begin{unitzh}\unitlabel{ceyuan-007.13}草曰：立天元一為半徑，置南行步在地，内減天元半徑，得□□為股圓差。{\zhsmall 按：斜畫者，少之記也。□□是為四百八十步少一元也。下倣此。}又置乙東行步在地，内減天元，得下式□□為勾圓差。以勾圓差乘股圓差，得丨□□{\zhsmall 按：丨□□為一平方少六百八十元多九萬六千步。}為半段黄方冪，即城冪之半也。{\zhsmall 寄左。}又置天元冪，以倍之，得□□，亦為半段黄方冪。與左相消，得丨□□，如帶縱法之，得半徑。合問。{\zhsmall 按：相消者，取上兩相等之數，同加減相等之數，使一為步數，一為方元數，仍相等也。如寄數内減一平方，加六百八十元，則得九萬六千步；又數内亦減一平方，加六百八十元，則得一平方六百八十元。是為一平方六百八十元與九萬六千步等，故其式為丨□□。舊稿方元數皆作斜畫以别之，然遇方元數有多少異號者，殊混人目，今不用。}\end{unitzh}
\begin{uniten}Working: Set up the celestial unknown (tianyuan) as the radius. Place the bu walked south on the ground and from it subtract the celestial-unknown radius, obtaining □□ as the altitude–circle difference. [Translator's note: each □ stands for a counting-rod figure not reproduced in the transcription.] [Siku editors: The slanting stroke is the mark of "less". □□ is 480 bu less 1 yuan. The same applies below.] Again place B's bu walked east on the ground and from it subtract the celestial unknown, obtaining the following form, □□, as the base–circle difference. Multiplying the altitude–circle difference by the base–circle difference gives 丨□□ [Siku editors: 丨□□ is 1 square, less 680 yuan, more 96000 bu.], which is a half-portion of the area (mi) of the yellow square, that is, half of the area (mi) of the city. [Note: Set this aside on the left.] Again place the area (mi) of the celestial unknown and double it, obtaining □□, which is also a half-portion of the area of the yellow square. Cancelling this against the left (xiangxiao) gives 丨□□; treating it by the method with an attached extension (daizong), you obtain the radius. This fits the question. [Siku editors: "Cancelling against each other" means taking the two equal quantities above and adding or subtracting equal quantities to both alike, so that one becomes a number of bu and the other a number of squares and yuan, and they are still equal. Thus, if from the number set aside one subtracts 1 square and adds 680 yuan, one obtains 96000 bu; and if from the other number one likewise subtracts 1 square and adds 680 yuan, one obtains 1 square and 680 yuan. This means that 1 square and 680 yuan are equal to 96000 bu; therefore its form is 丨□□. In the old draft the numbers of squares and yuan were all written with slanting strokes to distinguish them; but where numbers of squares and yuan carry the different signs "more" and "less", this greatly confuses people's eyes, so it is not used now.]\end{uniten}

\begin{unitzh}\unitlabel{ceyuan-007.14}又法：識别得二行併即大弦也。立天元一為半徑，置甲南行步，加天元一，得□□為大股；又置乙東行步，加天元，得□□為大勾也。勾股相乘，得丨□□為一个大直積，以天元除之，得下式□□□為三事和。{\zhsmall 寄左。黄方除倍積得三事和，今以半黄方除直積，亦為三事和也。}然後併二行步，又併入勾股，共得□□為同數，與左相消，得□□□，以帶縱平方開之，得一百二十步，倍之得全徑也。合問。\end{unitzh}
\begin{uniten}Another method: By identification (shibie) one obtains that the two walks added together are the great hypotenuse (da xian). Set up the celestial unknown as the radius. Place A's bu walked south and add the celestial unknown, obtaining □□ as the great altitude (da gu); again place B's bu walked east and add the celestial unknown, obtaining □□ as the great base (da gou). Multiplying base and altitude together gives 丨□□, which is one great rectangular area (da zhi ji). Dividing it by the celestial unknown gives the following form, □□□, as the sum of the three things (sanshi he). [Note: Set this aside on the left. Dividing twice the area by the yellow square gives the sum of the three things; now dividing the rectangular area by half the yellow square also gives the sum of the three things.] After that, add together the bu of the two walks, and further add in the base and the altitude; altogether one obtains □□ as the same number (tong shu). Cancelling it against the left gives □□□. Extracting its square root with an attached extension (daizong pingfang) gives 120 bu; doubling it gives the whole diameter. This fits the question.\end{uniten}

\begin{unitzh}\unitlabel{ceyuan-007.15}按：是書皆先法後草。草者，以立天元一推衍而得其方元積數者也；法者，又取推衍中之支節條目，融㑹而歸於簡約者也。草者，法之本；法者，草之用。法使人易於推步，而草則存其義以俟知者。二者相須，不可偏廢。顧應祥僅演其開方乘除之數，而去其細草，蓋亦不得其理矣。\end{unitzh}
\begin{uniten}Note: This book always gives the rule first and the working after. The working is that which, by setting up the celestial unknown and deriving from it (tuiyan), obtains the numbers of its square, yuan and area; the rule, moreover, takes the branches and joints, articles and items within the derivation, fuses them together and brings them back to brevity and concision. The working is the root of the rule; the rule is the use of the working. The rule makes it easy for people to calculate (tuibu), while the working preserves its meaning to await those who understand. The two require each other, and neither may be neglected. Gu Yingxiang (顧應祥) only expounded its numbers for root extraction, multiplication and division and removed its detailed working (xicao); presumably he, too, did not grasp its principle.\end{uniten}

\begin{unitzh}\unitlabel{ceyuan-007.16}按：元時未有筆算，故加減乘除之式不能詳載，觀者遂以為無下手處。今借根方法既明，視此則渙如氷釋矣。\end{unitzh}
\begin{uniten}Note: In Yuan times there was not yet written calculation (bisuan); therefore the forms of addition, subtraction, multiplication and division could not be recorded in detail, and readers consequently thought there was no place to set their hands. Now that the method of borrowed roots and powers (jiegenfang) has become clear, when one looks at this, it dissolves like melting ice.\end{uniten}

\begin{unitzh}\unitlabel{ceyuan-007.17}《測圓海鏡》卷二\end{unitzh}
\begin{uniten}Sea Mirror of Circle Measurements, juan 2\end{uniten}

\section*{卷03}\addcontentsline{toc}{section}{卷03}

\begin{unitzh}\unitlabel{ceyuan-008.1}欽定四庫全書\end{unitzh}
\begin{uniten}Imperially Commissioned Complete Library of the Four Treasuries (Qinding Siku quanshu)\end{uniten}

\begin{unitzh}\unitlabel{ceyuan-008.2}測圓海鏡卷三\end{unitzh}
\begin{uniten}Sea Mirror of Circle Measurements (Ceyuan haijing), juan 3\end{uniten}

\begin{unitzh}\unitlabel{ceyuan-008.3}元 李冶 撰\end{unitzh}
\begin{uniten}Composed by Li Ye of the Yuan\end{uniten}

\begin{unitzh}\unitlabel{ceyuan-008.4}邊股一十七問\end{unitzh}
\begin{uniten}Edge altitude (biangu): 17 problems\end{uniten}

\begin{unitzh}\unitlabel{ceyuan-008.5}或問：乙出東門南行，不知歩數而止。甲出西門南行四百八十歩，望見乙，復就乙行五百一十步，與乙相㑹。問答同前。\end{unitzh}
\begin{uniten}Someone asks: Yi goes out of the east gate and walks south an unknown number of bu (步), and stops. Jia goes out of the west gate and walks south 480 bu, sees Yi, then walks toward Yi 510 bu and meets Yi. The question and the answer are the same as before.\end{uniten}

\begin{unitzh}\unitlabel{ceyuan-008.6}法曰：倍相減步，以乘二之甲南行步，為平方實，得城徑。\end{unitzh}
\begin{uniten}Rule: Double the bu of the difference, and with it multiply twice Jia's southward bu, as the dividend (shi) for the square root (pingfang); one obtains the city diameter.\end{uniten}

\begin{unitzh}\unitlabel{ceyuan-008.7}草曰：識别得二行相減餘三十步，即乙出東門南行步也。倍相減步，得六十步，以乘二之甲南行步九百六十步，得五萬七千六百步，為平方實。如法開之，得二百四十步，即城徑也。合問。\end{unitzh}
\begin{uniten}Working: By discernment (shibie) one obtains: the two walks subtracted from each other leave 30 bu, which is the bu Yi walks south after going out of the east gate. Doubling the bu of the difference gives 60 bu; with it multiply twice Jia's southward bu, 960 bu, obtaining 57600 bu as the dividend (shi) for the square root. Extracting the root according to the rule gives 240 bu, which is the city diameter. This agrees with the question.\end{uniten}

\begin{unitzh}\unitlabel{ceyuan-008.8}或問：甲出西門南行四百八十步而止，乙從艮隅東行八十步，望見甲。問答同前。\end{unitzh}
\begin{uniten}Someone asks: Jia goes out of the west gate, walks south 480 bu and stops; Yi, from the Gen corner [Translator's note: Gen 艮 is the northeast.], walks east 80 bu and sees Jia. The question and the answer are the same as before.\end{uniten}

\begin{unitzh}\unitlabel{ceyuan-008.9}法曰：倍南行步，以東行步乘之，為實；東行歩為從方；一步常法；得全徑。\end{unitzh}
\begin{uniten}Rule: Double the southward bu and multiply it by the eastward bu, as the dividend (shi); the eastward bu is the following square (congfang); 1 bu as the constant divisor (changfa); one obtains the full diameter.\end{uniten}

\begin{unitzh}\unitlabel{ceyuan-008.10}草曰：立天元一為全徑，以減於二之甲南行步，得□□，為兩个大差也。以乙東行步乘之，得□□，為圓徑冪{\zhsmall 寄左}。然後以天元冪與左相消，得丨□□，以帶縱平方開之，得二百四十步，即城徑也。合問。\end{unitzh}
\begin{uniten}Working: Set up the celestial unknown (tianyuan) as the full diameter; subtract it from twice Jia's southward bu, obtaining □□ [Translator's note: here and below, □ marks array content not reproduced in the transcription; 丨 and ○ are the rod numerals 1 and 0.] as two great differences (dacha). Multiply it by Yi's eastward bu, obtaining □□ as the area (mi) of the circle diameter [Note: deposit on the left (jizuo).]. Then eliminate (xiangxiao) the celestial-unknown area (mi) against the left, obtaining 丨□□; extract the square root with attached lengthwise extension (daizong pingfang), obtaining 240 bu, which is the city diameter. This agrees with the question.\end{uniten}

\begin{unitzh}\unitlabel{ceyuan-008.11}又法：半之乙東行步乘南行步，為實；半之乙東行步為從；一步常法；得半徑。\end{unitzh}
\begin{uniten}Another rule: Half of Yi's eastward bu multiplied by the southward bu, as the dividend (shi); half of Yi's eastward bu as the following (cong); 1 bu as the constant divisor; one obtains the half-diameter.\end{uniten}

\begin{unitzh}\unitlabel{ceyuan-008.12}草曰：立天元一為半城徑，減甲南行步，得□□，為大差也。以半之東行步乘之，得□□，即半徑冪{\zhsmall 寄左}。然後以天元冪為同數，與左相消，得丨□□，開帶縱平方，得一百二十步，倍之，即城徑也。合問。\end{unitzh}
\begin{uniten}Working: Set up the celestial unknown as half the city diameter; with it reduce Jia's southward bu, obtaining □□ as the great difference. Multiply it by half the eastward bu, obtaining □□, which is the area (mi) of the half-diameter [Note: deposit on the left.]. Then take the celestial-unknown area (mi) as the equal quantity (tongshu) and eliminate it against the left, obtaining 丨□□; extract the square root with attached lengthwise extension, obtaining 120 bu; doubling it gives the city diameter. This agrees with the question.\end{uniten}

\begin{unitzh}\unitlabel{ceyuan-008.13}或問：甲出西門南行四百八十步而止，乙從艮隅亦南行一百五十步，望見甲。問答同前。\end{unitzh}
\begin{uniten}Someone asks: Jia goes out of the west gate, walks south 480 bu and stops; Yi, from the Gen corner, also walks south 150 bu and sees Jia. The question and the answer are the same as before.\end{uniten}

\begin{unitzh}\unitlabel{ceyuan-008.14}法曰：兩行步相乘為實，南行步為從方，一為隅，得半徑。\end{unitzh}
\begin{uniten}Rule: Multiply the bu of the two walks together as the dividend (shi); the southward bu is the following square; 1 is the corner (yu); one obtains the half-diameter.\end{uniten}

\begin{unitzh}\unitlabel{ceyuan-008.15}草曰：立天元一為半城徑，以減乙南行步，得□□，為半梯頭。以甲行步為梯底，以乘之，得□□，為半徑冪{\zhsmall 寄左}。然後以天元冪與左相消，得丨□□，開帶縱平方，得一百二十步，倍之，即城徑也。合問。\end{unitzh}
\begin{uniten}Working: Set up the celestial unknown as half the city diameter; subtract it from Yi's southward bu, obtaining □□ as half the trapezoid head (titou). Take Jia's walk in bu as the trapezoid bottom (tidi) and multiply by it, obtaining □□ as the area (mi) of the half-diameter [Note: deposit on the left.]. Then eliminate the celestial-unknown area (mi) against the left, obtaining 丨□□; extract the square root with attached lengthwise extension, obtaining 120 bu; doubling it gives the city diameter. This agrees with the question.\end{uniten}

\begin{unitzh}\unitlabel{ceyuan-008.16}或問：甲出西門南行四百八十步，乙出東門直行一十六步，望見甲。問答同前。\end{unitzh}
\begin{uniten}Someone asks: Jia goes out of the west gate and walks south 480 bu; Yi goes out of the east gate, walks straight 16 bu and sees Jia. The question and the answer are the same as before.\end{uniten}

\begin{unitzh}\unitlabel{ceyuan-008.17}法曰：以四之東行步乗南行冪為實，從空，東行為亷，一步為隅法，得全徑。\end{unitzh}
\begin{uniten}Rule: With 4 times the eastward bu multiply the area (mi) of the southward walk, as the dividend (shi); the following is empty; the eastward walk is the edge (lian); 1 bu is the corner divisor (yufa); one obtains the full diameter.\end{uniten}

\begin{unitzh}\unitlabel{ceyuan-008.18}草曰：立天元一為圓徑，加乙東行步，得□□，為中勾；其甲南行即中股也。置東行步為小勾，以中股乘之，得□，合以中勾除；今不受除，便以為小股也{\zhsmall 内寄中勾分母}。乃復以中股乗之，得三百六十八萬六千四百；又四之，得一千四百七十四萬五千六百，為一段圓徑冪{\zhsmall 寄中勾分母，寄左}。然後以天元徑自之，又以中勾乘之，得□□，為同數，與左相消，得丨□□□，以𢃄縱立方開之，得二百四十步，為城徑也。合問。\end{unitzh}
\begin{uniten}Working: Set up the celestial unknown as the circle diameter; add Yi's eastward bu, obtaining □□ as the middle base (zhonggou); Jia's southward walk is the middle altitude (zhonggu). Set up the eastward bu as the small base (xiaogou) and multiply it by the middle altitude, obtaining □; it should be divided by the middle base; now it does not accept division (bushouchu), so simply take it as the small altitude (xiaogu) [Note: within it the middle base is deposited as denominator (fenmu).]. Then again multiply by the middle altitude, obtaining 3686400; further multiply by 4, obtaining 14745600 as one piece of the area (mi) of the circle diameter [Note: the middle base is deposited as denominator; deposit on the left.]. Then multiply the celestial-unknown diameter by itself and further multiply by the middle base, obtaining □□ as the equal quantity; eliminate it against the left, obtaining 丨□□□; extract the cube root with attached lengthwise extension (daizong lifang), obtaining 240 bu as the city diameter. This agrees with the question.\end{uniten}

\begin{unitzh}\unitlabel{ceyuan-008.19}按：不受除者，無可除之理也。凡二數，此數於彼數有可除之理則受除，無可除之理則不受除也。蓋除有法有實，實可二，法不可二。此題以中勾為法，而中勾内有一元，又有十六步，其為數已二矣，又何以均分不一之數乎？故曰不受也。寄分者，姑寄其應除之數也。俟求得兩相等數，而此數内尚少一除，不除此而轉乘彼，則兩數仍相等，猶之受除者也。此所謂以乘代除也。\end{unitzh}
\begin{uniten}Remark: "Does not accept division" means there is no principle by which it can be divided. In general, of two numbers, if this number has, with respect to that number, a principle by which it can be divided, then it accepts division; if there is no principle by which it can be divided, then it does not accept division. For division has a divisor (fa) and a dividend (shi): the dividend may be twofold, the divisor may not be twofold. In this problem the middle base is taken as the divisor, and within the middle base there is one yuan and also 16 bu, so that as a number it is already twofold; how, then, could one divide equally by a number that is not single? Therefore it is said that it "does not accept". "Depositing the division" (jifen) means provisionally depositing the number by which one ought to divide. When the two mutually equal numbers have been found, this number still lacks one division within it; not dividing this but instead multiplying that, the two numbers are still equal, just as with what has accepted division. This is what is called "using multiplication in place of division".\end{uniten}

\begin{unitzh}\unitlabel{ceyuan-008.20}或問：乙出南門東行七十二步而止，甲出西門南行四百八十步，望乙與城㕘相直。問答同前。\end{unitzh}
\begin{uniten}Someone asks: Yi goes out of the south gate, walks east 72 bu and stops; Jia goes out of the west gate, walks south 480 bu, and looks toward Yi in alignment with the city (canxiangzhi). The question and the answer are the same as before.\end{uniten}

\begin{unitzh}\unitlabel{ceyuan-008.21}法曰：以乙東行冪乗甲南行為實，乙東行冪為從方，甲南行步内減二之東行步為益亷，一步常法，得半徑。\end{unitzh}
\begin{uniten}Rule: With the area (mi) of Yi's eastward walk multiply Jia's southward walk, as the dividend (shi); the area (mi) of Yi's eastward walk is the following square; Jia's southward bu minus twice the eastward bu is the augmenting edge (yilian); 1 bu as the constant divisor; one obtains the half-diameter.\end{uniten}

\begin{unitzh}\unitlabel{ceyuan-008.22}草曰：立天元一為半城徑，以減南行步，得□□，為小股；又以天元加乙東行步，得□□，為小勾；又以天元加南行步，得□□，為大股。乃置大股在地，以小勾乘之，得下式丨□□，合以小股除之；今不受除，便以為大勾{\zhsmall 内寄小股分母}。又置天元半徑，以分母小股乘之，得□□，以減大勾，得□□□，為半个梯底，於上。以乙東行七十二步為半个梯頭，以乘上位，得□□□，為半徑冪{\zhsmall 内寄小股分母}，寄左。然後置天元冪，又以分母小股乘之，得□□□，為同數，與左相消，得□□□□，以立方開之，得一百二十步，倍之，即城徑也。合問。\end{unitzh}
\begin{uniten}Working: Set up the celestial unknown as half the city diameter; subtract it from the southward bu, obtaining □□ as the small altitude. Further add the celestial unknown to Yi's eastward bu, obtaining □□ as the small base. Further add the celestial unknown to the southward bu, obtaining □□ as the great altitude (dagu). Then set the great altitude on the ground and multiply it by the small base, obtaining the following form (xiashi) 丨□□; it should be divided by the small altitude; now it does not accept division, so simply take it as the great base (dagou) [Note: within it the small altitude is deposited as denominator.]. Further set up the celestial-unknown half-diameter and multiply it by the denominator, the small altitude, obtaining □□; subtract it from the great base, obtaining □□□ as half a trapezoid bottom, at the top. Take Yi's eastward 72 bu as half a trapezoid head and with it multiply the upper position (shangwei), obtaining □□□ as the area (mi) of the half-diameter [Note: within it the small altitude is deposited as denominator.]; deposit on the left. Then set up the celestial-unknown area (mi) and further multiply it by the denominator, the small altitude, obtaining □□□ as the equal quantity; eliminate it against the left, obtaining □□□□; extract the cube root, obtaining 120 bu; doubling it gives the city diameter. This agrees with the question.\end{uniten}

\begin{unitzh}\unitlabel{ceyuan-008.23}又法曰：以二數相乘為實，相減為從，一虚法，平開，得半徑。\end{unitzh}
\begin{uniten}Another rule: Multiply the two numbers together as the dividend (shi); their difference as the following; 1 as the void divisor (xufa); extract the square root (pingkai), obtaining the half-diameter.\end{uniten}

\begin{unitzh}\unitlabel{ceyuan-008.24}草曰：别得二數相併為大股内少一虚勾，其二數相減為大差弦也。立天元一為半徑，副置之：上位減於四百八十，得□□，為股圓差{\zhsmall 即大差股也}；下位加七十二，得□□，與股圓差相乘，得下式□□□，為一大差積{\zhsmall 寄左}。再以大差勾減於大差股，餘□□，為較；又加入大差弦四百單八，共得□□，為弦較共也。以天元乘之，得□□，為同數，與左相消，得□□□，以平方開之，得一百二十步，即半徑。合問。前法太煩，故又立此法以就簡也。\end{unitzh}
\begin{uniten}Working: By distinguishing one obtains: the two numbers added together make the great altitude less one void base (xugou); the two numbers subtracted from each other make the hypotenuse of the great difference. Set up the celestial unknown as the half-diameter and set it up in duplicate (fuzhi): subtract the upper position from 480, obtaining □□ as the altitude–circle difference (guyuancha) [Note: that is, the altitude of the great difference.]; add 72 to the lower position, obtaining □□, and multiply it with the altitude–circle difference, obtaining the following form □□□ as one great-difference area [Note: deposit on the left.]. Again subtract the base of the great difference from the altitude of the great difference; the remainder □□ is the difference (jiao); further add in the hypotenuse of the great difference, 408, together obtaining □□ as the sum of hypotenuse and difference (xianjiaogong). Multiply it by the celestial unknown, obtaining □□ as the equal quantity; eliminate it against the left, obtaining □□□; extract the square root, obtaining 120 bu, which is the half-diameter. This agrees with the question. The previous rule was too cumbersome; therefore this rule was set up in addition, so as to approach simplicity.\end{uniten}

\begin{unitzh}\unitlabel{ceyuan-008.25}或問：乙出南門東行，不知步數而立。甲出西門南行四百八十步，望見乙與城㕘相直，又就乙行四百零八步，與乙相㑹。問答同前。\end{unitzh}
\begin{uniten}Someone asks: Yi goes out of the south gate and walks east an unknown number of bu, and stands. Jia goes out of the west gate, walks south 480 bu, and sees Yi in alignment with the city; further he walks toward Yi 408 bu and meets Yi. The question and the answer are the same as before.\end{uniten}

\begin{unitzh}\unitlabel{ceyuan-008.26}法曰：二行步相減，以乘甲南行步，為實；甲東行步内減相減步，為益方；一步常法；得半徑。\end{unitzh}
\begin{uniten}Rule: Subtract the bu of the two walks from each other, and with this multiply Jia's southward bu, as the dividend (shi); Jia's eastward bu [Translator's note: so the text; in this problem Jia walks south.] minus the bu of the difference, as the augmenting square (yifang); 1 bu as the constant divisor; one obtains the half-diameter.\end{uniten}

\begin{unitzh}\unitlabel{ceyuan-008.27}草曰：識别得二行相減餘七十二步，即是乙出南門東行數也，更不湏用弦。遂立天元一為半城徑，加乙東行，得□□，為小勾也。副置南行步：上減天乙，得□□，為小股；下加天元，得□□，為大股。乃置大股，以小勾乘之，得下式丨□□，合以小股除之；今不受除，便以此為大勾也{\zhsmall 内帶小股分母}。又倍天元，以小股乘之，得下式□□，以減於大勾，得□□□，為勾圓差也。合以股圓差乘之，縁此勾圓差内已帶小股分母{\zhsmall 小股即股圓差也}，更不湏乘，便以此為半段黄方冪{\zhsmall 更無分母也}，寄左。乃以天元自之，又倍之，得□□，為同數，與左相消，得□□□，以平方開之，得一百二十步，倍之，即城徑也。合問。\end{unitzh}
\begin{uniten}Working: By discernment one obtains: the two walks subtracted from each other leave 72 bu, which is the number of Yi's eastward walk out of the south gate; there is no further need to use the hypotenuse (xian). Then set up the celestial unknown as half the city diameter; add Yi's eastward walk, obtaining □□ as the small base. Set up the southward bu in duplicate: from the upper subtract tianyi [Translator's note: 天乙, evidently standing for the celestial unknown.], obtaining □□ as the small altitude; to the lower add the celestial unknown, obtaining □□ as the great altitude. Then set up the great altitude and multiply it by the small base, obtaining the following form 丨□□; it should be divided by the small altitude; now it does not accept division, so simply take this as the great base [Note: within it the small altitude is carried as denominator.]. Further double the celestial unknown and multiply it by the small altitude, obtaining the following form □□; subtract it from the great base, obtaining □□□ as the base–circle difference (gouyuancha). It should be multiplied by the altitude–circle difference; because this base–circle difference already carries within it the small altitude as denominator [Note: the small altitude is the altitude–circle difference.], there is no further need to multiply; simply take this as half a piece of the area (mi) of the yellow square (huangfang) [Note: there is no longer any denominator.]; deposit on the left. Then multiply the celestial unknown by itself and further double it, obtaining □□ as the equal quantity; eliminate it against the left, obtaining □□□; extract the square root, obtaining 120 bu; doubling it gives the city diameter. This agrees with the question.\end{uniten}

\begin{unitzh}\unitlabel{ceyuan-008.28}或問：乙出東門直行，不知歩數而止。甲出西門南行四百八十步，望見乙，復就乙斜行五百四十四步，與乙相會。問答同前。\end{unitzh}
\begin{uniten}Someone asks: Yi goes out of the east gate and walks straight an unknown number of bu, and stops. Jia goes out of the west gate, walks south 480 bu and sees Yi; then he walks obliquely toward Yi 544 bu and meets Yi. The question and the answer are the same as before.\end{uniten}

\begin{unitzh}\unitlabel{ceyuan-008.29}法曰：半南行步減半斜行步，以乘南行步，為實；從方空；半斜行、半南行相減得數，加入南行步，為隅法；得半徑。\end{unitzh}
\begin{uniten}Rule: Half the southward bu subtracted from half the oblique bu, and with this multiply the southward bu, as the dividend (shi); the following square is empty; half the oblique walk and half the southward walk subtracted from each other, and the number obtained added to the southward bu, as the corner divisor; one obtains the half-diameter.\end{uniten}

\begin{unitzh}\unitlabel{ceyuan-008.30}草曰：識别得二行相減餘六十四步，即半徑為股之勾也。立天元為半徑，就以為小股；其二行相減餘六十四步，即小勾也。乃置甲南行步，加天元，得下式□□，為大股；以小勾乘之，得□□；又以小股除之，得□□，為大勾。又倍天元一減之，得下式□□□，為勾圓差也。半之，得□□□，於上。乃以天元減甲南行步，得□□，為股圓差，以乘上位，得丨□○□，為半徑冪{\zhsmall 寄左}。然後以天元冪與左相消，得下式□□□，以平方開之，得一百二十步，倍之，即城徑也。合問。\end{unitzh}
\begin{uniten}Working: By discernment one obtains: the two walks subtracted from each other leave 64 bu, which is the base (gou) of which the half-diameter is the altitude (gu). Set up the celestial unknown as the half-diameter and directly take it as the small altitude; the remainder of the two walks subtracted from each other, 64 bu, is the small base. Then set up Jia's southward bu and add the celestial unknown, obtaining the following form □□ as the great altitude; multiply it by the small base, obtaining □□; further divide it by the small altitude, obtaining □□ as the great base. Further double the celestial unknown and subtract it from that, obtaining the following form □□□ as the base–circle difference. Halve it, obtaining □□□, at the top. Then subtract the celestial unknown from Jia's southward bu, obtaining □□ as the altitude–circle difference; multiply the upper position by it, obtaining 丨□○□ as the area (mi) of the half-diameter [Note: deposit on the left.]. Then eliminate the celestial-unknown area (mi) against the left, obtaining the following form □□□; extract the square root, obtaining 120 bu; doubling it gives the city diameter. This agrees with the question.\end{uniten}

\begin{unitzh}\unitlabel{ceyuan-008.31}按：此問以小股為除法，蓋因小股只一天元，其數不二，猶有可除之理也。然得數降於實數之下者，皆不可以命名，至開方時，仍湏各升一位以計之，是兩邊各加一乘，猶是寄分之理也。\end{unitzh}
\begin{uniten}Remark: In this problem the small altitude is taken as the divisor (chufa), because the small altitude is only one celestial unknown; its number is not twofold, so there is still a principle by which it can be divided. However, the numbers obtained that have descended below the number of the dividend (shi) can none of them be named; when it comes to root extraction, each must still be raised one position in order to reckon them. This is adding one multiplication to each of the two sides, which is still the principle of depositing the division.\end{uniten}

\begin{unitzh}\unitlabel{ceyuan-008.32}又法：以二數差乘二數併，開方，得邊勾；復以邊股乘之，為實；併二數而半之，為法；實如法，得二百四十步，即城徑{\zhsmall 此蓋用前勾上容圓法也}。\end{unitzh}
\begin{uniten}Another rule: With the difference of the two numbers multiply the sum of the two numbers, and extract the square root, obtaining the edge base (biangou); again multiply by the edge altitude, as the dividend (shi); add the two numbers and halve, as the divisor (fa); dividing the dividend by the divisor gives 240 bu, which is the city diameter [Note: This presumably uses the earlier rule of the circle contained on the base (goushang rongyuan).].\end{uniten}

\begin{unitzh}\unitlabel{ceyuan-008.33}或問：乙從乾地東行，不知幾步而止。甲出西門南行四百八十步，望見乙，復就乙斜行六百八十步，與乙相㑹。問答同前。\end{unitzh}
\begin{uniten}Someone asks: Yi, from the Qian ground [Translator's note: Qian 乾 is the northwest.], walks east, how many bu not being known, and stops. Jia goes out of the west gate, walks south 480 bu and sees Yi; then he walks obliquely toward Yi 680 bu and meets Yi. The question and the answer are the same as before.\end{uniten}

\begin{unitzh}\unitlabel{ceyuan-008.34}法曰：併二行數，以二行差乘之，内減二行差冪，為實；併二行步及二行差，為從方；二步常法；得半徑。\end{unitzh}
\begin{uniten}Rule: Add the numbers of the two walks, multiply by the difference of the two walks, and from this subtract the area (mi) of the difference of the two walks, as the dividend (shi); add the bu of the two walks and the difference of the two walks, as the following square; 2 bu as the constant divisor; one obtains the half-diameter.\end{uniten}

\begin{unitzh}\unitlabel{ceyuan-008.35}草曰：識别得二行相減餘二百步，即半圓徑與小差勾之共數也。立天元一為半城徑，加於二百步，得□□，為大勾也。又以天元加於甲南行步四百八十，得□□，即大股也。乃以大勾自之，得丨□□，為勾冪{\zhsmall 寄左}。乃置乙斜行六百八十步為大弦，加入大股，共得□□，於上。再置二行差，内減天元，得□□，為小差勾，即股弦較，以乘上位，得□□□，為同數，與左相消，得□□□，以平方開之，得一百二十步，倍之，即城徑也。合問。\end{unitzh}
\begin{uniten}Working: By discernment one obtains: the two walks subtracted from each other leave 200 bu, which is the sum (gongshu) of half the circle diameter and the base of the small difference (xiaocha). Set up the celestial unknown as half the city diameter; add it to 200 bu, obtaining □□ as the great base. Further add the celestial unknown to Jia's southward bu, 480, obtaining □□, which is the great altitude. Then multiply the great base by itself, obtaining 丨□□ as the area (mi) of the base [Note: deposit on the left.]. Then set up Yi's oblique walk of 680 bu as the great hypotenuse (daxian); add in the great altitude, together obtaining □□, at the top. Again set up the difference of the two walks and subtract the celestial unknown from it, obtaining □□ as the base of the small difference, which is the altitude–hypotenuse difference (guxianjiao); multiply the upper position by it, obtaining □□□ as the equal quantity; eliminate it against the left, obtaining □□□; extract the square root, obtaining 120 bu; doubling it gives the city diameter. This agrees with the question.\end{uniten}

\begin{unitzh}\unitlabel{ceyuan-008.36}又法求小差：二行相減，以自之，又四之，為實；二行相減，八之，於上；二之南行步内減二之二行相減數，又以加上位，為益方；二步常法。\end{unitzh}
\begin{uniten}Another rule, to find the small difference: subtract the two walks from each other, multiply it by itself, further multiply by 4, as the dividend (shi); subtract the two walks from each other, multiply by 8, at the top; twice the southward bu minus twice the number of the two walks subtracted from each other, further added to the upper position, as the augmenting square; 2 bu as the constant divisor.\end{uniten}

\begin{unitzh}\unitlabel{ceyuan-008.37}草曰：立天元一為小差，減二行差，得□□，為半城徑；以自之，得丨□□；又四之，得□□□，為圓徑冪{\zhsmall 寄左}。然後以半城徑減於甲南行，得□□；又倍之，得□□，為兩个大差也；又以天元乘之，得□□○，為同數，與左相消，得下式□□□，以平方開之，得八十步，為小差也。\end{unitzh}
\begin{uniten}Working: Set up the celestial unknown as the small difference; with it reduce the difference of the two walks, obtaining □□ as half the city diameter; multiply it by itself, obtaining 丨□□; further multiply by 4, obtaining □□□ as the area (mi) of the circle diameter [Note: deposit on the left.]. Then subtract half the city diameter from Jia's southward walk, obtaining □□; further double it, obtaining □□ as two great differences; further multiply by the celestial unknown, obtaining □□○ as the equal quantity; eliminate it against the left, obtaining the following form □□□; extract the square root, obtaining 80 bu as the small difference.\end{uniten}

\begin{unitzh}\unitlabel{ceyuan-008.38}或問：乙出南門南行，不知步數而立。甲出西門南行四百八十步，望乙與城㕘相直，復就乙斜行二百五十五步，與乙相㑹。問答同前。\end{unitzh}
\begin{uniten}Someone asks: Yi goes out of the south gate and walks south an unknown number of bu, and stands. Jia goes out of the west gate, walks south 480 bu, and looks toward Yi in alignment with the city; then he walks obliquely toward Yi 255 bu and meets Yi. The question and the answer are the same as before.\end{uniten}

\begin{unitzh}\unitlabel{ceyuan-008.39}法曰：甲南行内減二之兩行差，餘以乘甲南行，又倍之，為實；二步為隅；得半徑。\end{unitzh}
\begin{uniten}Rule: From Jia's southward walk subtract twice the difference of the two walks; with the remainder multiply Jia's southward walk, and further double it, as the dividend (shi); 2 bu as the corner; one obtains the half-diameter.\end{uniten}

\begin{unitzh}\unitlabel{ceyuan-009.1}草曰：别得二行步相減，餘二百二十五步，乃是半徑為勾之股也。立天元一為半城徑，就以為小勾率；其二行差二百二十五步，即為小股率。乃置甲南行步，加入天元，得□□，為大股。以天元小勾乘之，得丨□，合以小股除，今不受除，{\zhsmall 按：此所謂不受除，乃其數竒零不能盡，非無可除之理也，與前辭同而意異。}便以此為大勾。{\zhsmall 内寄小股分母。}乃倍天元，以小股乘之，得□，以減大勾，餘丨□，為一个小差，於上。{\zhsmall 内寄小股分母。}乃以天元減甲南行步，得□□，為大差也。以乘上位，得□□□，又倍之，得□□□，為圓徑冪，{\zhsmall 内寄小股分母。}寄左。然後倍天元以自之，又以小股乘之，得□□，為同數。與左相消，得□○□。以平方開之，得一百二十步，倍之即城徑也。合問。\end{unitzh}
\begin{uniten}Working: By distinction (bie) one obtains: the bu of the two walks subtracted from each other leave 225 bu (步); this is the altitude (gu) [of the triangle] whose base (gou) is the radius. Set up the celestial unknown (tianyuan) as the half city-diameter, and accordingly take it as the rate (lü) of the small base; the difference of the two walks, 225 bu, is then the rate of the small altitude. Then put the bu of Jia's southward walk and add the celestial unknown into it, obtaining □□ [Translator's note: here and below, □ stands for a counting-rod array, or part of one, not reproduced in the transcription; 丨 is the rod numeral 1 and ○ is 0.], as the great altitude. Multiply it by the celestial unknown, the small base, obtaining 丨□; it should be divided by the small altitude, but now it does not admit division; [Siku editors: What is here called "does not admit division" means that the number has an odd remainder and cannot be divided out; it is not that there is no principle by which it could be divided. The wording is the same as before, but the meaning is different.] simply take this as the great base. [Note: Within it is lodged the small altitude as denominator (fen mu).] Then double the celestial unknown and multiply it by the small altitude, obtaining □; subtract it from the great base, leaving 丨□, as one small difference (xiao cha); [place it] above. [Note: Within it is lodged the small altitude as denominator.] Then subtract the celestial unknown from the bu of Jia's southward walk, obtaining □□, as the great difference (da cha). Multiply the upper position by it, obtaining □□□; double it, obtaining □□□, as the area (mi) of the circle's diameter; [Note: Within it is lodged the small altitude as denominator.] lodge it on the left (ji zuo). Afterwards double the celestial unknown and multiply it by itself; multiply again by the small altitude, obtaining □□, as the equal number (tong shu). Eliminate (xiang xiao) it against the left, obtaining □○□. Extract the square root, obtaining 120 bu; doubling it gives the city diameter. This agrees with the question.\end{uniten}

\begin{unitzh}\unitlabel{ceyuan-009.2}按：此題止用股弦求勾法，即得城半徑。其必展轉數次而後始得者，益見其為發明立天元一之術，使人易曉也。後多有倣此者。\end{unitzh}
\begin{uniten}Note: In this problem, merely by using the method of finding the base from the altitude and the hypotenuse one immediately obtains the city's radius; that it must go roundabout several times before it is first obtained shows all the more that it is meant to elucidate the method of setting up the celestial unknown and to make it easy for people to understand. Later there are many [problems] that imitate this.\end{uniten}

\begin{unitzh}\unitlabel{ceyuan-009.3}或問：乙出南門直行一百三十五步而止，甲出西門南行四百八十步，望乙與城㕘相直。問答同前。\end{unitzh}
\begin{uniten}Someone asks: Yi goes out of the south gate, walks straight on 135 bu and stops; Jia goes out of the west gate, walks south 480 bu, and looks toward Yi, [Yi] being in line with the city. The question and answer are the same as before.\end{uniten}

\begin{unitzh}\unitlabel{ceyuan-009.4}法曰：二行步相減，餘以自乘，内減乙行冪，為實；二之甲南行，為益從；一步常法。得半徑。\end{unitzh}
\begin{uniten}Rule: Subtract the bu of the two walks from each other; multiply the remainder by itself, and from it subtract the area of Yi's walk; this is the dividend (shi). Twice Jia's southward walk is the negative follower (yi cong); 1 bu is the constant divisor (chang fa). One obtains the radius.\end{uniten}

\begin{unitzh}\unitlabel{ceyuan-009.5}草曰：立天元一以為半徑，便以為勾率。又以天元加乙行步，併以減於甲行步，得□□，為股率。乃置乙南行步一百三十五步為小股，以勾率乘之，得□，合以股率除之，今不受除，乃便以此為小勾。{\zhsmall 内寄股率分母。}又置乙南行步，加二天元，得□□，為大股，以勾率乘之，得□□，合以股率除之，今不受除，便以此為大勾。{\zhsmall 内寄股率分母。}以小勾、大勾相乘，得□□□，為半徑冪，{\zhsmall 内帶股率冪為分母。}寄左。然後置天元，以自乘，又以股率冪乘之，得丨□□□，為同數。與左相消，得□□□。以平方開之，得一百二十步，倍之即城徑也。合問。\end{unitzh}
\begin{uniten}Working: Set up the celestial unknown as the radius, and simply take it as the rate of the base. Again add the celestial unknown to the bu of Yi's walk, and subtract the sum from the bu of Jia's walk, obtaining □□, as the rate of the altitude. Then put the bu of Yi's southward walk, 135 bu, as the small altitude; multiply it by the rate of the base, obtaining □; it should be divided by the rate of the altitude, but now it does not admit division, so simply take this as the small base. [Note: Within it is lodged the rate of the altitude as denominator.] Again put the bu of Yi's southward walk and add 2 celestial unknowns, obtaining □□, as the great altitude; multiply it by the rate of the base, obtaining □□; it should be divided by the rate of the altitude, but now it does not admit division; simply take this as the great base. [Note: Within it is lodged the rate of the altitude as denominator.] Multiply the small base and the great base together, obtaining □□□, as the area of the radius; [Note: Within it is carried the area of the rate of the altitude as denominator.] lodge it on the left. Afterwards put the celestial unknown, multiply it by itself, and multiply again by the area of the rate of the altitude, obtaining 丨□□□, as the equal number. Eliminate it against the left, obtaining □□□. Extract the square root, obtaining 120 bu; doubling it gives the city diameter. This agrees with the question.\end{uniten}

\begin{unitzh}\unitlabel{ceyuan-009.6}按：此草得數為九百六十立方少一三乘方，與十萬零八百平方等，皆虚數也。各降二位，即如各以平方除之，乃為九百六十元少一平方，與十萬零八百步等。兩數等，所降之位又等，則兩數仍相等，而實積步數乃出矣，故可以帶縱平方開之也。此係降位而得實數者，與前升位而得實數者，其理互相發明。草中不言，蓋以為不待於言也。\end{unitzh}
\begin{uniten}Note: The number obtained in this working is 960 cubes less 1 third-multiplication power (san cheng fang), equal to 100800 squares; all are empty numbers (xu shu). Lowering each by two places is the same as dividing each by the square: this gives 960 yuan less 1 square, equal to 100800 bu. The two numbers are equal and the places lowered are also equal, so the two numbers remain equal, and the bu of the actual area then come out; therefore it can be extracted by the square [root] with an accompanying lengthwise term (dai zong). This is a case of obtaining the actual number (shi shu) by lowering places; it and the earlier case of obtaining the actual number by raising places illuminate each other in principle. The working does not mention it, presumably because it was thought to need no mention.\end{uniten}

\begin{unitzh}\unitlabel{ceyuan-009.7}或問：甲乙二人同出西門，向南行，至西南十字道口分路。乙折東行一百九十二步而立；甲又南行，甲通行四百八十步，望乙與城㕘相直。問答同前。\end{unitzh}
\begin{uniten}Someone asks: Two persons, Jia and Yi, go out of the west gate together and walk toward the south; at the southwest crossroads they part ways. Yi turns and walks east 192 bu and stands; Jia walks further south, Jia's walk in all being 480 bu, and looks toward Yi, [Yi] being in line with the city. The question and answer are the same as before.\end{uniten}

\begin{unitzh}\unitlabel{ceyuan-009.8}法曰：兩行相乘得數，又以乙東行乘之，為實。二行相乘，於上位；又置乙東行，以二行相減數乘之，得數，加上位，為法。\end{unitzh}
\begin{uniten}Rule: Multiply the two walks together, and multiply the number obtained again by Yi's eastward walk; this is the dividend. Multiply the two walks together [and put it] in the upper position; again put Yi's eastward walk and multiply it by the difference of the two walks; add the number obtained to the upper position; this is the divisor (fa).\end{uniten}

\begin{unitzh}\unitlabel{ceyuan-009.9}草曰：立天元一為半城徑，副置。上位加南行步，得□□，為大股也；下位減於甲行步，得□□，為小股也。其乙東行即小勾也。置大股，以小勾乘之，得□□，内寄小股□□為母，便以為大勾也。置天元，以母通之，得□□，減於大勾，得丨□□，為半个矮梯底，於上。再置乙東行，内減天元，得下式□□，為半个矮梯頭，以乘上位，得下式□□□□，為半徑冪，寄左。再置天元，以自之為冪，又以分母乘之，得□□□，為如積。與左相消，得□□，上法下實，得一百二十步，即城之半徑也。合問。\end{unitzh}
\begin{uniten}Working: Set up the celestial unknown as the half city-diameter and put it in duplicate (fu zhi). To the upper position add the bu of the southward walk, obtaining □□, as the great altitude; the lower position is subtracted from the bu of Jia's walk, obtaining □□, as the small altitude. Yi's eastward walk is the small base. Put the great altitude and multiply it by the small base, obtaining □□, within which is lodged the small altitude □□ as denominator; simply take this as the great base. Put the celestial unknown and convert (tong) it with the denominator, obtaining □□; subtract it from the great base, obtaining 丨□□, as half of the base of the low trapezoid (ai ti di); [place it] above. Again put Yi's eastward walk and subtract the celestial unknown from it, obtaining the form below (xia shi) □□, as half of the top of the low trapezoid (ai ti tou); multiply the upper position by it, obtaining the form below □□□□, as the area of the radius; lodge it on the left. Again put the celestial unknown, multiply it by itself to make the area, and multiply again by the denominator, obtaining □□□, as the like area (ru ji). Eliminate it against the left, obtaining □□: the upper is the divisor, the lower the dividend; one obtains 120 bu, which is the radius of the city. This agrees with the question.\end{uniten}

\begin{unitzh}\unitlabel{ceyuan-009.10}按：草中相消法皆得兩邊數，此獨得一邊二數。蓋此條共數比彼條共數少一數，又多一數，為相等，則多少二數其必為相等無疑矣。多少數多者亦倣此。此又相消法中之一變也。\end{unitzh}
\begin{uniten}Note: In the workings the method of elimination always yields numbers of both sides; this alone yields two numbers of one side. For [when] the total of this item, compared with the total of that item, is less by one number and also more by one number, and [they] are equal, then the two numbers, the less and the more, must certainly be equal, beyond doubt. Where the numbers of less and more are many, imitate this too. This is yet another variation within the method of elimination.\end{uniten}

\begin{unitzh}\unitlabel{ceyuan-009.11}又法：二行步相乘，為實；倍甲南行，内減乙東行，為法。\end{unitzh}
\begin{uniten}Another rule: Multiply the bu of the two walks together; this is the dividend. Double Jia's southward walk and subtract Yi's eastward walk from it; this is the divisor.\end{uniten}

\begin{unitzh}\unitlabel{ceyuan-009.12}草曰：立天元一為半城徑，副置。上位加甲南行，得□□，為大股；下位減甲行步，得□□，為小股，便是股圓差也。其乙東行即小勾也。置大股，以小勾乘之，得□□，内寄小股□□為母，便以為大勾也。再置天元，以二之，又以分母乘之，得□□，為全徑，以減於大勾，餘□□□，為勾圓差也。合以股圓差乘之，縁内已有小股分母，不湏乘，便以此為兩段之半徑冪也，更無分母。{\zhsmall 寄左。}然後置天元冪，以二之，得□□，為如積。以左相消，得□□，上法下實，得一百二十步，即半城徑也。合問。\end{unitzh}
\begin{uniten}Working: Set up the celestial unknown as the half city-diameter and put it in duplicate. To the upper position add Jia's southward walk, obtaining □□, as the great altitude; subtract the lower position from the bu of Jia's walk, obtaining □□, as the small altitude, which is the difference of the altitude and the circle (gu yuan cha). Yi's eastward walk is the small base. Put the great altitude and multiply it by the small base, obtaining □□, within which is lodged the small altitude □□ as denominator; simply take this as the great base. Again put the celestial unknown, double it, and multiply again by the denominator, obtaining □□, as the whole diameter; subtract it from the great base, leaving □□□, as the difference of the base and the circle (gou yuan cha). It should be multiplied by the difference of the altitude and the circle, but because within it there is already the small altitude as denominator, it need not be multiplied; simply take this as two portions (liang duan) of the area of the radius; there is no further denominator. [Note: Lodge it on the left.] Afterwards put the area of the celestial unknown and double it, obtaining □□, as the like area. Eliminate it against the left, obtaining □□: the upper is the divisor, the lower the dividend; one obtains 120 bu, which is the half city-diameter. This agrees with the question.\end{uniten}

\begin{unitzh}\unitlabel{ceyuan-009.13}或問：見邊股四百八十步，□弦三十四步。問答同前。{\zhsmall 此題在「甲乙二人同出西門南行至十字道，乙折東行一百九十二步而立，甲又南行，甲通行四百八十步，望見乙與城㕘相直」之後。}\end{unitzh}
\begin{uniten}Someone asks: Given the edge altitude (bian gu), 480 bu, and the □ hypotenuse (xian) [Translator's note: □ before a term stands for a character not reproduced in the transcription.], 34 bu. The question and answer are the same as before. [Note: This problem comes after "Jia and Yi, two persons, go out of the west gate together and walk south to the crossroads; Yi turns and walks east 192 bu and stands; Jia walks further south, Jia's walk in all being 480 bu, and he looks and sees Yi, [Yi] being in line with the city".]\end{uniten}

\begin{unitzh}\unitlabel{ceyuan-009.14}法曰：□弦乘邊股，半之，為實；半□弦、半邊股相併，為從；半步隅法。平方得□股。\end{unitzh}
\begin{uniten}Rule: Multiply the edge altitude by the □ hypotenuse and halve it; this is the dividend. Half the □ hypotenuse and half the edge altitude added together are the follower (cong); 1/2 bu is the corner divisor (yu fa). [Extracting by] the square [root] one obtains the □ altitude.\end{uniten}

\begin{unitzh}\unitlabel{ceyuan-009.15}草曰：立天元一為□股，加□弦，得□□，為平勾也。又以天元減邊股而半之，得□□，為髙股也。平勾、髙股相乘，得□□□，為半徑冪，{\zhsmall 寄左。}然後以天元乘邊股，得□，為同數。與左相消，得下式□□□，開平方，得□股三十步，以乘邊股，開平方，倍之，即圓城徑也。合問。\end{unitzh}
\begin{uniten}Working: Set up the celestial unknown as the □ altitude; add the □ hypotenuse, obtaining □□, as the level base (ping gou). Again subtract the celestial unknown from the edge altitude and halve it, obtaining □□, as the high altitude (gao gu). Multiply the level base and the high altitude together, obtaining □□□, as the area of the radius; [Note: Lodge it on the left.] afterwards multiply the edge altitude by the celestial unknown, obtaining □, as the equal number. Eliminate it against the left, obtaining the form below □□□; extract the square root, obtaining the □ altitude, 30 bu; multiply the edge altitude by it, extract the square root, and double it: this is the diameter of the round city. This agrees with the question.\end{uniten}

\begin{unitzh}\unitlabel{ceyuan-009.16}按：此問原稿在三卷末。\end{unitzh}
\begin{uniten}Note: In the original manuscript this question was at the end of juan 3.\end{uniten}

\begin{unitzh}\unitlabel{ceyuan-009.17}或問：見邊股四百八十，明弦一百五十三。問答同前。\end{unitzh}
\begin{uniten}Someone asks: Given the edge altitude, 480, and the bright hypotenuse (ming xian), 153. The question and answer are the same as before.\end{uniten}

\begin{unitzh}\unitlabel{ceyuan-009.18}法曰：二云數相減，復倍之，内減邊股，復以邊股乘之，於上；又以明弦冪乘上位，為實。以邊股乘明弦冪，又二之，為從。二云數相減，餘以自之，為第一亷。二云數相減，又倍之，為第二益亷。一常法。開三乘方，得明勾。\end{unitzh}
\begin{uniten}Rule: Subtract the two stated numbers (yun shu) from each other, double [the remainder], subtract the edge altitude from it, and again multiply by the edge altitude; [place it] above. Again multiply the upper position by the area of the bright hypotenuse; this is the dividend. Multiply the area of the bright hypotenuse by the edge altitude and double it; this is the follower. Subtract the two stated numbers from each other and multiply the remainder by itself; this is the first edge (lian). Subtract the two stated numbers from each other and double [the remainder]; this is the second, negative edge. 1 is the constant divisor. Extract the root of the third-multiplication power, obtaining the bright base (ming gou).\end{uniten}

\begin{unitzh}\unitlabel{ceyuan-009.19}草曰：立天元一為明勾，加明弦，得□□，為髙股也。以髙股減邊股，餘□□，為髙弦。以倍之，得□□，為黄廣弦也。内減邊股，得□□，為□股。復以邊股乘之，得□□，於上。又以明弦自乘，得二萬三千四百零九，為分母，以乘上位，得□□，為𢃄分半徑冪，{\zhsmall 寄左。}然後置黄廣弦，以天元乘之，得□□，復合以明弦除之，不除，寄為母，便以此為全徑。又半之，得□□，為半徑，以自之，得□□□，為同數。與左相消，得下式丨□□□□，開三乘方，得七十二步，即明勾也。餘各依法求之。合問。\end{unitzh}
\begin{uniten}Working: Set up the celestial unknown as the bright base; add the bright hypotenuse, obtaining □□, as the high altitude. Subtract the high altitude from the edge altitude, leaving □□, as the high hypotenuse. Double it, obtaining □□, as the yellow-breadth hypotenuse (huang guang xian). Subtract the edge altitude from it, obtaining □□, as the □ altitude. Again multiply it by the edge altitude, obtaining □□; [place it] above. Again multiply the bright hypotenuse by itself, obtaining 23409, as the denominator; multiply the upper position by it, obtaining □□, as the area of the radius carrying a fraction (dai fen); [Note: Lodge it on the left.] afterwards put the yellow-breadth hypotenuse and multiply it by the celestial unknown, obtaining □□; again it should be divided by the bright hypotenuse; not dividing, lodge [it] as denominator, and simply take this as the whole diameter. Again halve it, obtaining □□, as the radius; multiply it by itself, obtaining □□□, as the equal number. Eliminate it against the left, obtaining the form below 丨□□□□; extract the root of the third-multiplication power, obtaining 72 bu, which is the bright base. The rest are each found according to the rules. This agrees with the question.\end{uniten}

\begin{unitzh}\unitlabel{ceyuan-009.20}又法：邊股内減二明弦，以邊股乘之，復以明弦冪乘之，為三乘方實；亷、從並同前。\end{unitzh}
\begin{uniten}Another rule: From the edge altitude subtract 2 bright hypotenuses, multiply by the edge altitude, and again multiply by the area of the bright hypotenuse; this is the dividend of the third-multiplication power. The edges and the follower are all the same as before.\end{uniten}

\begin{unitzh}\unitlabel{ceyuan-009.21}草曰：識别得二數相減，餘為髙股、虚弦共，又為髙弦、明勾共。此餘數内又去半徑，即明和也。明和、明弦相併，即股圓差；相減，則明黄方也。又倍明弦，加明黄，亦得股圓差也。邊股内減明勾，餘即大差弦也。立天元一為明勾，減於云數相減數，得□□，即髙弦也。以髙弦減邊股，得□□，即髙股也。以髙股減於云數相減數，得□□，即虚弦也。以天元又減虚弦，得□□，即□股也。乃置髙弦，以天元乘之，得□□，合明弦除之，不受除，便以此為髙勾也。{\zhsmall 即半徑。}髙勾自之，得丨□□□，為半徑冪，{\zhsmall 内帶明弦冪分母。}寄左。然後置邊股，以□股乘之，得□□，為半徑冪，又以明弦冪二萬三千四百零九分母通之，得□□，為同數。與左相消，得實、從、亷、隅五層，如前式。\end{unitzh}
\begin{uniten}Working: By recognition (shibie) one obtains: the two numbers subtracted from each other leave the high altitude and the empty hypotenuse (xu xian) combined; this is also the high hypotenuse and the bright base combined. From this remainder further remove the radius, and it is the bright sum (ming he). The bright sum and the bright hypotenuse added together are the difference of the altitude and the circle; subtracted from each other, they are the bright yellow square (ming huang fang). Again, doubling the bright hypotenuse and adding the bright yellow also gives the difference of the altitude and the circle. From the edge altitude subtract the bright base; the remainder is the great-difference hypotenuse (da cha xian). Set up the celestial unknown as the bright base; subtract it from the difference of the stated numbers, obtaining □□, which is the high hypotenuse. Subtract the high hypotenuse from the edge altitude, obtaining □□, which is the high altitude. Subtract the high altitude from the difference of the stated numbers, obtaining □□, which is the empty hypotenuse. Again subtract the celestial unknown from the empty hypotenuse, obtaining □□, which is the □ altitude. Then put the high hypotenuse and multiply it by the celestial unknown, obtaining □□; it should be divided by the bright hypotenuse, but it does not admit division; simply take this as the high base. [Note: That is, the radius.] Multiply the high base by itself, obtaining 丨□□□, as the area of the radius; [Note: Within it is carried the area of the bright hypotenuse as denominator.] lodge it on the left. Afterwards put the edge altitude and multiply it by the □ altitude, obtaining □□, as the area of the radius; again convert it with the denominator, the area of the bright hypotenuse, 23409, obtaining □□, as the equal number. Eliminate it against the left, obtaining dividend, follower, edges and corner in five layers, as in the form before.\end{uniten}

\begin{unitzh}\unitlabel{ceyuan-009.22}或問：邊股四百八十步，髙弦二百五十五步。問答同前。\end{unitzh}
\begin{uniten}Someone asks: The edge altitude is 480 bu and the high hypotenuse 255 bu. The question and answer are the same as before.\end{uniten}

\begin{unitzh}\unitlabel{ceyuan-009.23}法曰：以邊股減於二之髙弦，復以邊股乘之，開平方，得半徑。\end{unitzh}
\begin{uniten}Rule: Subtract the edge altitude from twice the high hypotenuse, and again multiply by the edge altitude; extract the square root, obtaining the radius.\end{uniten}

\begin{unitzh}\unitlabel{ceyuan-009.24}草曰：立天元一為半徑。先倍髙弦，内減邊股，餘□，復以邊股乘之，得□□，寄左。以天元冪與左相消，得丨□□，開平方，得數，倍之，即城徑也。合問。\end{unitzh}
\begin{uniten}Working: Set up the celestial unknown as the radius. First double the high hypotenuse and subtract the edge altitude from it, leaving □; again multiply it by the edge altitude, obtaining □□; lodge it on the left. Eliminate the area of the celestial unknown against the left, obtaining 丨□□; extract the square root, and double the number obtained: this is the city diameter. This agrees with the question.\end{uniten}

\begin{unitzh}\unitlabel{ceyuan-009.25}或問：邊股四百八十步，平弦一百三十六步。答問同前。\end{unitzh}
\begin{uniten}Someone asks: The edge altitude is 480 bu and the level hypotenuse (ping xian) 136 bu. The answer and question are the same as before.\end{uniten}

\begin{unitzh}\unitlabel{ceyuan-009.26}法曰：置平弦，以邊股再乘之，為實；以邊股自之，為益從；平弦為益亷；一虚隅。開立方，得半徑。\end{unitzh}
\begin{uniten}Rule: Put the level hypotenuse and multiply it twice by the edge altitude; this is the dividend. Multiply the edge altitude by itself; this is the negative follower. The level hypotenuse is the negative edge (yi lian). 1 is the empty corner (xu yu). Extract the cube root, obtaining the radius.\end{uniten}

\begin{unitzh}\unitlabel{ceyuan-009.27}草曰：别得平弦即皇極勾也。立天元一為半徑，副之。上位加平弦，得□□，即邊勾也；下位減於平弦，得□□，即□勾也。置□勾，以邊股乘之，得□□，合邊勾除，今不受除，寄為母，便以此為□股。乃以此邊股乘之，得□□，為半徑冪，{\zhsmall 内𢃄邊勾分母。}寄左。然後以天元為冪，以分母邊勾乘之，得丨□□，為同數。與左相消，得丨□□□，開立方，得一百二十步，倍之，即城徑也。合問。\end{unitzh}
\begin{uniten}Working: By distinction one obtains that the level hypotenuse is the imperial-pole base (huangji gou). Set up the celestial unknown as the radius and duplicate it. To the upper position add the level hypotenuse, obtaining □□, which is the edge base (bian gou); the lower position is subtracted from the level hypotenuse, obtaining □□, which is the □ base. Put the □ base and multiply it by the edge altitude, obtaining □□; it should be divided by the edge base, but now it does not admit division; lodge [it] as denominator, and simply take this as the □ altitude. Then multiply it by this edge altitude, obtaining □□, as the area of the radius; [Note: Within it is carried the edge base as denominator.] lodge it on the left. Afterwards make the celestial unknown into an area and multiply it by the denominator, the edge base, obtaining 丨□□, as the equal number. Eliminate it against the left, obtaining 丨□□□; extract the cube root, obtaining 120 bu; double it: this is the city diameter. This agrees with the question.\end{uniten}

\begin{unitzh}\unitlabel{ceyuan-009.28}或問：邊股四百八十步，明股明弦和二百八十八步。問答同前。\end{unitzh}
\begin{uniten}Someone asks: The edge altitude is 480 bu, and the sum of the bright altitude (ming gu) and the bright hypotenuse is 288 bu. The question and answer are the same as before.\end{uniten}

\begin{unitzh}\unitlabel{ceyuan-009.29}法曰：以云之云數相減，餘加邊股，復以減餘乘之，訖，又折半，於上；又以減餘自之，減上位，為實。併云數，半之，為法。得明勾。\end{unitzh}
\begin{uniten}Rule: Subtract the stated numbers (yun zhi yun shu) from each other [Translator's note: the text reads 云之云數, which appears garbled.]; add the edge altitude to the remainder, and again multiply by the remainder of the subtraction; when this is done, further halve it; [place it] above. Again multiply the remainder of the subtraction by itself and subtract it from the upper position; this is the dividend. Add the stated numbers together and halve [the sum]; this is the divisor. One obtains the bright base.\end{uniten}

\begin{unitzh}\unitlabel{ceyuan-009.30}草曰：别得二數相減，餘為大差勾。立天元一為明勾，減於大差勾，得□□，即半徑也。又以天元減半徑，得□□，為虚勾，於上。又以半徑加邊股，得□□，為通股，於下。上下相乘，得□□□，折半，得丨□□，為半徑冪，{\zhsmall 寄左。}然後以半徑冪丨□□為同數。與左相消，得□□，上法下實，得七十二步，即明勾也。合問。\end{unitzh}
\begin{uniten}Working: By distinction one obtains that the two numbers subtracted from each other leave the great-difference base (da cha gou). Set up the celestial unknown as the bright base; subtract it from the great-difference base, obtaining □□, which is the radius. Again subtract the celestial unknown from the radius, obtaining □□, as the empty base (xu gou); [place it] above. Again add the radius to the edge altitude, obtaining □□, as the through altitude (tong gu); [place it] below. Multiply the upper and the lower together, obtaining □□□; halve it, obtaining 丨□□, as the area of the radius; [Note: Lodge it on the left.] afterwards take the area of the radius, 丨□□, as the equal number. Eliminate it against the left, obtaining □□: the upper is the divisor, the lower the dividend; one obtains 72 bu, which is the bright base. This agrees with the question.\end{uniten}

\begin{unitzh}\unitlabel{ceyuan-009.31}或問：見邊股四百八十步，□勾□弦和五十步。問答同前。\end{unitzh}
\begin{uniten}Someone asks: Given the edge altitude, 480 bu, and the sum of the □ base and the □ hypotenuse, 50 bu. The question and answer are the same as before.\end{uniten}

\begin{unitzh}\unitlabel{ceyuan-010.1}法曰：半邊股、半和步相併，得為泛率。以汎半減邊股，以自之，又二之於上；以和步乘泛率，減上位，為實。以汎率減邊股，六之於上，内又加半个邊股、三个和步，為益從。三步常法。得□股。\end{unitzh}
\begin{uniten}Rule: Combine half the bian altitude (bian gu) and half the sum-bu (hebu), obtaining the floating rate (fanlü). Subtract half the floating rate from the bian altitude and multiply it by itself; again multiply it by 2, placing it above. Multiply the floating rate by the sum-bu and subtract it from the upper position, as the dividend (shi). Subtract the floating rate from the bian altitude and multiply it by 6, placing it above; into this further add half a bian altitude and 3 sum-bu, as the negative linear coefficient (yi cong). 3 bu is the constant divisor (changfa). This gives the □ altitude. [Translator's note: here and below, □ stands for a character or counting-rod diagram not reproduced in the transcription. As transmitted, this rule does not give 30, the value obtained in the working.]\end{uniten}

\begin{unitzh}\unitlabel{ceyuan-010.2}草曰：别得和步得□股，即小差也；小差、邊股共，即二中差{\zhsmall 按此句誤}。立天元一為□股，加和步，得□□，即小差也。以小差加邊股而半之，得□□，即中差也。中、小差相併，得□□，即大差也。以小差乘之，得□□□，為半段徑冪{\zhsmall 寄左}。然後置邊股，内減大差，得□□，為半徑。以自之，得□□□；又倍之，得下式□□□，與左相消，得下式□□□。開平方，得三十步，即□股也。合問。\end{unitzh}
\begin{uniten}Working: Separately we obtain: the sum-bu together with the □ altitude is the small difference (xiaocha); the small difference and the bian altitude together are twice the middle difference (zhongcha) [Siku editors: This sentence is wrong.]. Set up the celestial unknown (tianyuan) as the □ altitude; add the sum-bu to obtain □□, which is the small difference. Add the small difference to the bian altitude and halve it to obtain □□, which is the middle difference. Combine the middle and small differences to obtain □□, which is the large difference (dacha). Multiply it by the small difference to obtain □□□ as half a portion of the diameter area (mi) [Note: lodge it on the left.]. Then place the bian altitude, subtract the large difference from within it, and obtain □□ as the half-diameter (banjing). Multiply it by itself to obtain □□□; again double it to obtain the arrangement (shi) below, □□□; mutually eliminate it with the left to obtain the arrangement below, □□□. Extract the square root to obtain 30 bu, which is the □ altitude. This fits the question.\end{uniten}

\begin{unitzh}\unitlabel{ceyuan-010.3}按：草云「以小差、邊股共，即二中差」，有誤。蓋中差即勾股較，小差即股弦較，邊股即勾弦較與容圓半徑和。若設勾二十、股二十一、弦二十九，則勾弦較九，容圓半徑六，併之得十五，為邊股；股弦較八，為小差；小差、邊股共得二十三；勾股較一，為中差，倍之僅得二，則相差二十一矣。是知細草乃因題數之偶合而誤，非正法也。今依其術，另設法、草於後，以補其闕。\end{unitzh}
\begin{uniten}Note: The working says, "Taking the small difference and the bian altitude together is twice the middle difference"; this is wrong. For the middle difference is the base-altitude difference (gougu jiao), the small difference is the altitude-hypotenuse difference (guxian jiao), and the bian altitude is the sum of the base-hypotenuse difference (gouxian jiao) and the half-diameter of the contained circle (rongyuan). If one supposes a base (gou) of 20, an altitude (gu) of 21 and a hypotenuse (xian) of 29, then the base-hypotenuse difference is 9 and the half-diameter of the contained circle is 6; combining them gives 15, which is the bian altitude; the altitude-hypotenuse difference, 8, is the small difference; the small difference and the bian altitude together give 23; the base-altitude difference, 1, is the middle difference, and doubling it gives only 2; so they differ by 21. From this one knows that the working erred owing to a chance agreement of the numbers in the problem, and is not a correct method. Now, following its method, we separately set out a rule and a working below, to make up for what is lacking.\end{uniten}

\begin{unitzh}\unitlabel{ceyuan-010.4}法曰：以□勾弦和自之，邊股再乘，為實。倍邊股，加□勾弦和，再以□勾弦和乘之，為從。又倍□勾弦和，減邊股，餘為益亷。一為隅。𢃄縱立方開之，得□股。\end{unitzh}
\begin{uniten}Rule: Multiply the □ base-hypotenuse sum (gouxian he) by itself and multiply again by the bian altitude, as the dividend (shi). Double the bian altitude, add the □ base-hypotenuse sum, and again multiply it by the □ base-hypotenuse sum, as the linear coefficient (cong). Again double the □ base-hypotenuse sum and subtract it from the bian altitude; the remainder is the negative edge (yi lian). 1 is the corner (yu). Extract it as a cube root with accompanying extension (daizong) to obtain the □ altitude.\end{uniten}

\begin{unitzh}\unitlabel{ceyuan-010.5}草曰：别得邊股即髙股弦和，□股即髙股弦差，□股弦和即平勾也。立天天一為□股，自之，得丨□，應以□勾弦和除之；不除，便以為□勾弦較{\zhsmall 内寄□勾弦和分母}。轉以□勾弦和自之，得□，為□勾弦和；加□勾弦較，得丨○□，為倍□弦。又以□勾弦和分母乘倍□股，得□，為倍□股；與倍□弦相加，得丨□□，為倍□股弦和，即倍平勾。又於邊股内減□股，得□□，為倍髙股。倍髙股、倍平勾相乘，得□□□□，為圓徑冪，寄左。又以邊股、□股相乘，得□，為半徑冪；四因之，得□，為圓徑冪；又以□勾弦和分母乘之，得□，為同數。與左相消，得丨□□□，開帶縱立方，得□股三十步。合問。\end{unitzh}
\begin{uniten}Working: Separately we obtain: the bian altitude is the gao altitude-hypotenuse sum; the □ altitude is the gao altitude-hypotenuse difference (cha); the □ altitude-hypotenuse sum is the ping base. Set up the celestial unknown [Translator's note: the text writes 天天一 for 天元一] as the □ altitude and multiply it by itself to obtain 丨□. This ought to be divided by the □ base-hypotenuse sum; not dividing, simply take it as the □ base-hypotenuse difference [Note: lodged within it is the □ base-hypotenuse sum as denominator (fenmu).]. In turn multiply the □ base-hypotenuse sum by itself to obtain □ as the □ base-hypotenuse sum; add the □ base-hypotenuse difference to obtain 丨○□ as twice the □ hypotenuse. Again multiply twice the □ altitude by the denominator, the □ base-hypotenuse sum, to obtain □ as twice the □ altitude; add it to twice the □ hypotenuse to obtain 丨□□ as twice the □ altitude-hypotenuse sum, which is twice the ping base. Again subtract the □ altitude from within the bian altitude to obtain □□ as twice the gao altitude. Multiply twice the gao altitude and twice the ping base together to obtain □□□□ as the circle-diameter area (mi); lodge it on the left. Again multiply the bian altitude and the □ altitude together to obtain □ as the half-diameter area; multiply it by 4 to obtain □ as the circle-diameter area; again multiply it by the denominator, the □ base-hypotenuse sum, to obtain □ as the same number (tongshu). Mutually eliminate it with the left to obtain 丨□□□; extract the cube root with accompanying extension to obtain the □ altitude, 30 bu. This fits the question.\end{uniten}

\begin{unitzh}\unitlabel{ceyuan-010.6}《測圓海鏡》卷三\end{unitzh}
\begin{uniten}Sea Mirror of Circle Measurements, juan 3\end{uniten}

\begin{unitzh}\unitlabel{ceyuan-010.7}<子部,天文算法類,算書之屬,測圓海鏡>\end{unitzh}
\begin{uniten}<Masters section (zibu), Astronomy and Mathematics category (tianwen suanfa lei), Mathematical Books subcategory (suanshu zhi shu), Sea Mirror of Circle Measurements>\end{uniten}

\section*{卷04}\addcontentsline{toc}{section}{卷04}

\begin{unitzh}\unitlabel{ceyuan-011.1}欽定四庫全書\end{unitzh}
\begin{uniten}Imperially Commissioned Complete Library of the Four Treasuries (Qinding Siku quanshu)\end{uniten}

\begin{unitzh}\unitlabel{ceyuan-011.2}測圓海鏡卷四\end{unitzh}
\begin{uniten}Sea Mirror of Circle Measurements (Ceyuan haijing), juan 4\end{uniten}

\begin{unitzh}\unitlabel{ceyuan-011.3}元 李冶 撰\end{unitzh}
\begin{uniten}Composed by Li Ye of the Yuan\end{uniten}

\begin{unitzh}\unitlabel{ceyuan-011.4}底勾一十七問\end{unitzh}
\begin{uniten}Bottom base (di gou): 17 problems\end{uniten}

\begin{unitzh}\unitlabel{ceyuan-011.5}或問：乙出南門東行，不知步數而立；甲出北門東行二百步見之，就乙斜行二百七十二步，與乙相㑹。問答同前。\end{unitzh}
\begin{uniten}Someone asks: B goes out of the south gate and walks east an unknown number of bu (步), and stands; A goes out of the north gate, walks east 200 bu, and sees him; walking obliquely toward B 272 bu, he meets B. The question and answer are as before.\end{uniten}

\begin{unitzh}\unitlabel{ceyuan-011.6}法曰：二行差數乘甲東行，又四之，為平方實，得全徑。\end{unitzh}
\begin{uniten}Rule: Multiply A's eastward walk by the difference of the two walks, and further quadruple it, as the dividend (shi) for square-root extraction (pingfang); one obtains the full diameter (quan jing).\end{uniten}

\begin{unitzh}\unitlabel{ceyuan-011.7}草曰：識别得二行相減餘，即乙出南門東行數也。以甲東行減於就乙斜行，餘七十二步，以乘甲東行歩，得一萬四千四百步，又四之，得五萬七千六百步，為實。以平方開之，得二百四十步，即城徑也。合問。\end{unitzh}
\begin{uniten}Working: By identification (shibie) one finds that the remainder when the two walks are subtracted from each other is the number of B's eastward walk out of the south gate. Subtract A's eastward walk from the oblique walk toward B; the remainder is 72 bu. Multiply A's eastward walk in bu by it, obtaining 14400 bu; further quadruple it, obtaining 57600 bu, as the dividend. Extract its square root, obtaining 240 bu, which is the city diameter (cheng jing). This agrees with the question.\end{uniten}

\begin{unitzh}\unitlabel{ceyuan-011.8}或問：乙從坤隅南行三百六十步，甲出北門東行二百步見之。問答同前。\end{unitzh}
\begin{uniten}Someone asks: B walks south from the Kun corner (kun yu) 360 bu; A goes out of the north gate, walks east 200 bu, and sees him. The question and answer are as before.\end{uniten}

\begin{unitzh}\unitlabel{ceyuan-011.9}法曰：二行步相乘，倍之為實；乙南行為從；一步常法；得城徑。\end{unitzh}
\begin{uniten}Rule: Multiply the bu of the two walks together and double it, as the dividend; B's southward walk is the joining term (cong); 1 bu is the constant divisor (chang fa); one obtains the city diameter.\end{uniten}

\begin{unitzh}\unitlabel{ceyuan-011.10}草曰：立天元一為城徑，以減於二之甲東行步，得{\zhsmall □}□，為兩个小差。以乙南行步乘之，得□□，為城徑冪。{\zhsmall 寄左}然後以天元冪丨□與左相消，得丨□□。以平方開之，得二百四十步，即城徑也。合問。\end{unitzh}
\begin{uniten}Working: Set up the celestial unknown (tianyuan) as the city diameter; subtract it from twice A's eastward walk in bu, obtaining [Note: □]□ as two small differences (xiao cha) [Translator's note: here and below, each □ stands for an entry of a counting-rod array not reproduced in the transcription; 丨 is the rod numeral 1.]. Multiply it by B's southward walk in bu, obtaining □□ as the area (mi) of the city diameter [Note: lodge it on the left (ji zuo)]. Then make the area of the celestial unknown, 丨□, cancel (xiang xiao) with the left, obtaining 丨□□. Extract its square root, obtaining 240 bu, which is the city diameter. This agrees with the question.\end{uniten}

\begin{unitzh}\unitlabel{ceyuan-011.11}又法：半之乙南行步乘甲東行為實；半乙南行為從；一步常法；得半徑。\end{unitzh}
\begin{uniten}Another rule: Half of B's southward walk in bu multiplied by A's eastward walk is the dividend; half of B's southward walk is the joining term; 1 bu is the constant divisor; one obtains the half-diameter (ban jing).\end{uniten}

\begin{unitzh}\unitlabel{ceyuan-011.12}草曰：立天元一為半城徑，減甲東行，得□□，為小差。半乙南行步，得一百八十步，以乘小差，得□□，為半徑冪。{\zhsmall 寄左}然後以天元冪丨□與左相消，得下式丨□□。以平方開之，得一百二十步，倍之，即城徑也。合問。\end{unitzh}
\begin{uniten}Working: Set up the celestial unknown as half the city diameter; reduce A's eastward walk by it, obtaining □□ as the small difference. Halve B's southward walk in bu, obtaining 180 bu; multiply the small difference by it, obtaining □□ as the area of the half-diameter [Note: lodge it on the left]. Then make the area of the celestial unknown, 丨□, cancel with the left, obtaining the following array (xia shi), 丨□□. Extract its square root, obtaining 120 bu; doubling it, this is the city diameter. This agrees with the question.\end{uniten}

\begin{unitzh}\unitlabel{ceyuan-011.13}或問：乙從坤隅東行一百九十二歩而止，甲出北門東行二百步見乙。問答同前。\end{unitzh}
\begin{uniten}Someone asks: B walks east from the Kun corner 192 bu and stops; A goes out of the north gate, walks east 200 bu, and sees B. The question and answer are as before.\end{uniten}

\begin{unitzh}\unitlabel{ceyuan-011.14}法曰：兩行步相乘為實；甲東行為從；乙為隅；得半徑。\end{unitzh}
\begin{uniten}Rule: Multiply the bu of the two walks together as the dividend; A's eastward walk is the joining term; B is the corner (yu) [Translator's note: 乙 "B" is probably an error for 一 "one".]; one obtains the half-diameter.\end{uniten}

\begin{unitzh}\unitlabel{ceyuan-011.15}草曰：立天元一為半徑，減於乙東行，得□□。以甲行步乘之，得□□，為半徑冪。{\zhsmall 寄左}然後以天元冪丨□與左相消，得丨□□。以平方開之，得一百二十步，倍之，即城徑也。合問。\end{unitzh}
\begin{uniten}Working: Set up the celestial unknown as the half-diameter; subtract it from B's eastward walk, obtaining □□. Multiply it by A's walk in bu, obtaining □□ as the area of the half-diameter [Note: lodge it on the left]. Then make the area of the celestial unknown, 丨□, cancel with the left, obtaining 丨□□. Extract its square root, obtaining 120 bu; doubling it, this is the city diameter. This agrees with the question.\end{uniten}

\begin{unitzh}\unitlabel{ceyuan-011.16}或問：乙出南門直行一百三十五步，甲出北門東行二百步見乙。問答同前。\end{unitzh}
\begin{uniten}Someone asks: B goes out of the south gate and walks straight 135 bu; A goes out of the north gate, walks east 200 bu, and sees B. The question and answer are as before.\end{uniten}

\begin{unitzh}\unitlabel{ceyuan-011.17}法曰：以乙南行步乘甲東行冪，又四之，為實；從空；乙行為亷；一步常法；得城徑。\end{unitzh}
\begin{uniten}Rule: Multiply the area of A's eastward walk by B's southward walk in bu, and further quadruple it, as the dividend; the joining term is empty; B's walk is the edge (lian); 1 bu is the constant divisor; one obtains the city diameter.\end{uniten}

\begin{unitzh}\unitlabel{ceyuan-011.18}草曰：立天元一為城徑，加乙南行，得□□，為股率；其甲東行即勾率也；其乙南行□為小股。以勾率乘之，得□，合以股率除，今不除受，便以此為小勾。{\zhsmall 寄股率為母}乃以甲東行步乘之，得□，又四之，得□，為一段城徑冪。{\zhsmall 寄左}然後以天元城徑自之，又以股率分母通之，得丨□□，為同數，與左相消，得下式丨□□□。以立方開之，得二百四十步，即城徑也。合問。\end{unitzh}
\begin{uniten}Working: Set up the celestial unknown as the city diameter; add B's southward walk, obtaining □□ as the altitude rate (gu lü); A's eastward walk is the base rate (gou lü); B's southward walk, □, is the small altitude (xiao gu). Multiply it by the base rate, obtaining □; it ought to be divided by the altitude rate; now, not dividing, accept it [Translator's note: 不除受 is probably an inversion of 不受除, "does not accept division", as elsewhere.], and take this directly as the small base (xiao gou) [Note: lodge the altitude rate as denominator (mu)]. Then multiply it by A's eastward walk in bu, obtaining □; further quadruple it, obtaining □ as one piece (yi duan) of the area of the city diameter [Note: lodge it on the left]. Then multiply the celestial-unknown city diameter by itself, and further make it uniform (tong) with the altitude-rate denominator (fen mu), obtaining 丨□□ as the equal quantity (tong shu); make it cancel with the left, obtaining the following array, 丨□□□. Extract its cube root, obtaining 240 bu, which is the city diameter. This agrees with the question.\end{uniten}

\begin{unitzh}\unitlabel{ceyuan-011.19}又法：二行相乘，又以自乘，為實；以二行相乘倍之，為益方；南行冪為亷；八步益隅；立方開，得小勾七十二。\end{unitzh}
\begin{uniten}Another rule: Multiply the two walks together, and further multiply this by itself, as the dividend; double the product of the two walks, as the augmenting square (yi fang); the area of the southward walk is the edge; 8 bu is the augmenting corner (yi yu); extracting the cube root, one obtains the small base, 72.\end{uniten}

\begin{unitzh}\unitlabel{ceyuan-011.20}草曰：立天元一為小勾，以南行為小股，以東行二百步為大勾也。置大勾，内減天元，得□□，為中勾也。以小股乗之，得□□；以天元小勾除之，得□□，為中股，即城徑也。以自之，得□□□，為城徑冪也。{\zhsmall 寄左}又以天元小勾乘通勾二百步，得□；又四之，得□，為同數，與左相消，得□□□□。開立方，得七十二步，即小勾也。以乘通勾二百步為實，平方開，得一百二十步，倍之，即城徑也。合問。\end{unitzh}
\begin{uniten}Working: Set up the celestial unknown as the small base; take the southward walk as the small altitude, and the eastward walk of 200 bu as the large base (da gou). Set out the large base and subtract the celestial unknown within it, obtaining □□ as the middle base (zhong gou). Multiply it by the small altitude, obtaining □□; divide it by the celestial-unknown small base, obtaining □□ as the middle altitude (zhong gu), which is the city diameter. Multiply it by itself, obtaining □□□ as the area of the city diameter [Note: lodge it on the left]. Further multiply the through base (tong gou), 200 bu, by the celestial-unknown small base, obtaining □; further quadruple it, obtaining □ as the equal quantity; make it cancel with the left, obtaining □□□□. Extract the cube root, obtaining 72 bu, which is the small base. Multiply the through base, 200 bu, by it, as the dividend; extracting the square root, one obtains 120 bu; doubling it, this is the city diameter. This agrees with the question.\end{uniten}

\begin{unitzh}\unitlabel{ceyuan-011.21}又法求半徑：以南行步乘東行冪為實；從空；東行步為亷；二步常法；得半徑。\end{unitzh}
\begin{uniten}Another rule, to find the half-diameter: Multiply the area of the eastward walk by the southward walk in bu, as the dividend; the joining term is empty; the eastward walk in bu is the edge; 2 bu is the constant divisor; one obtains the half-diameter.\end{uniten}

\begin{unitzh}\unitlabel{ceyuan-011.22}草曰：立天元一為半徑，以二之，加南行步，得□□，為股率；以東行□為勾率；以南行為小股也。置小股，以勾率乘之，得□；以股率除之，不受除，只寄股率分母，便以此為小勾也。又以勾率乘之，得下式□，為半徑冪。{\zhsmall 寄左}再立天元半徑，以自之，又以分母股率乘之，得□□□，為同數，與左相消，得□□□□。開立方，得一百二十步，倍之，即城徑也。合問。\end{unitzh}
\begin{uniten}Working: Set up the celestial unknown as the half-diameter; double it and add the southward walk in bu, obtaining □□ as the altitude rate; take the eastward walk, □, as the base rate; take the southward walk as the small altitude. Set out the small altitude and multiply it by the base rate, obtaining □; divide it by the altitude rate: it does not accept division, so merely lodge the altitude rate as denominator and take this directly as the small base. Further multiply it by the base rate, obtaining the following array, □, as the area of the half-diameter [Note: lodge it on the left]. Again set up the celestial-unknown half-diameter, multiply it by itself, and further multiply it by the denominator, the altitude rate, obtaining □□□ as the equal quantity; make it cancel with the left, obtaining □□□□. Extract the cube root, obtaining 120 bu; doubling it, this is the city diameter. This agrees with the question.\end{uniten}

\begin{unitzh}\unitlabel{ceyuan-011.23}或問：乙出東門南行三十步而止，甲出北門東行二百步，望乙與城㕘相直。問答同前。\end{unitzh}
\begin{uniten}Someone asks: B goes out of the east gate, walks south 30 bu, and stops; A goes out of the north gate, walks east 200 bu, and looks toward B so that B and the city are in alignment (can xiang zhi). The question and answer are as before.\end{uniten}

\begin{unitzh}\unitlabel{ceyuan-011.24}法曰：以甲東行步乗乙南行冪為實；以乙南行冪為從；甲東行内減二之乙南行，為益亷；一步隅；得半徑。\end{unitzh}
\begin{uniten}Rule: Multiply the area of B's southward walk by A's eastward walk in bu, as the dividend; take the area of B's southward walk as the joining term; from A's eastward walk subtract twice B's southward walk, as the augmenting edge (yi lian); 1 bu is the corner; one obtains the half-diameter.\end{uniten}

\begin{unitzh}\unitlabel{ceyuan-011.25}草曰：立天元一為半城徑，減於甲東行步，得□□，為小勾。以天元加於乙南行步，得□□，為小股。乃以天元加東行步，得□□，為大勾。置大勾，以小股乗之，得丨□□，合以小勾除之，今不受除，便以此為大股。{\zhsmall 内帶小勾分母}又置天元半徑，以分母小勾乘之，得□□，減於大股，餘□□□。以乙南行步乗之，得□□□，為半徑冪，{\zhsmall 内有小勾分母}寄左。然後以天元為冪，又以小勾通之，得□□□，為同數，與左相消，得下式□□□□。以立方開之，得一百二十步，倍之，即城徑也。合問。{\zhsmall 翻法在記}\end{unitzh}
\begin{uniten}Working: Set up the celestial unknown as half the city diameter; subtract it from A's eastward walk in bu, obtaining □□ as the small base. Add the celestial unknown to B's southward walk in bu, obtaining □□ as the small altitude. Then add the celestial unknown to the eastward walk in bu, obtaining □□ as the large base. Set out the large base and multiply it by the small altitude, obtaining 丨□□; it ought to be divided by the small base, but now it does not accept division, so take this directly as the large altitude (da gu) [Note: it contains the small base as denominator]. Further set out the celestial-unknown half-diameter and multiply it by the denominator, the small base, obtaining □□; subtract it from the large altitude; the remainder is □□□. Multiply it by B's southward walk in bu, obtaining □□□ as the area of the half-diameter [Note: it contains the small base as denominator]; lodge it on the left. Then make the celestial unknown into an area, and further make it uniform with the small base, obtaining □□□ as the equal quantity; make it cancel with the left, obtaining the following array, □□□□. Extract its cube root, obtaining 120 bu; doubling it, this is the city diameter. This agrees with the question. [Note: the reversal method (fan fa) is in the Record (ji).]\end{uniten}

\begin{unitzh}\unitlabel{ceyuan-011.26}又法：乙南行乘甲東行為平實；二數相減為法；一隅；翻開，得半徑。\end{unitzh}
\begin{uniten}Another rule: B's southward walk multiplied by A's eastward walk is the level dividend (ping shi); the two numbers subtracted from each other are the divisor (fa); 1 is the corner; extracting by reversal (fan kai), one obtains the half-diameter.\end{uniten}

\begin{unitzh}\unitlabel{ceyuan-011.27}草曰：别得二數相併，為大勾内少一虚股；其二數相減，為小差弦也。立天元一為半徑，副置之：上位減於二百步，得□□，為勾圓差；{\zhsmall 即小差勾也}下位加三十步，得□□，為小差股。勾股相乘，得□□□，為一段小差積。{\zhsmall 寄左}再以小差勾減小差股，餘□□，為一較也。又以此較減於小差弦，得下式□□，為一个弦較較。以天元一乘之，得下式□□，為同數，與左相消，得□□□。開平方，得一百二十步，即半城徑也。合問。{\zhsmall 翻法在記}再立此法者，蓋從簡也。\end{unitzh}
\begin{uniten}Working: By distinguishing (bie) one finds that the two numbers added together are the large base less one empty altitude (xu gu), and that the two numbers subtracted from each other are the small-difference hypotenuse (xiao cha xian). Set up the celestial unknown as the half-diameter and set it out in duplicate (fu zhi): subtract the upper position from 200 bu, obtaining □□ as the base–circle difference (gou yuan cha) [Note: that is, the small-difference base (xiao cha gou)]; add 30 bu to the lower position, obtaining □□ as the small-difference altitude (xiao cha gu). Multiply base and altitude together, obtaining □□□ as one piece of small-difference area (xiao cha ji) [Note: lodge it on the left]. Again subtract the small-difference base from the small-difference altitude; the remainder, □□, is one difference (jiao). Further subtract this difference from the small-difference hypotenuse, obtaining the following array, □□, as one hypotenuse–difference difference (xian jiao jiao). Multiply it by the celestial unknown, obtaining the following array, □□, as the equal quantity; make it cancel with the left, obtaining □□□. Extract the square root, obtaining 120 bu, which is half the city diameter. This agrees with the question. [Note: the reversal method is in the Record.] The reason this method is set up again is presumably for the sake of simplicity.\end{uniten}

\begin{unitzh}\unitlabel{ceyuan-011.28}按：此乃以小差勾為平弦上弦較較，半徑為平股，故以小差弦上弦較較與半徑相乘，等於平弦上弦較較與小差股相乘，為一段小差積也。\end{unitzh}
\begin{uniten}Note: Here the small-difference base is taken as the hypotenuse–difference difference on the level hypotenuse (ping xian), and the half-diameter as the level altitude (ping gu); therefore the hypotenuse–difference difference on the small-difference hypotenuse multiplied by the half-diameter is equal to the hypotenuse–difference difference on the level hypotenuse multiplied by the small-difference altitude, making one piece of small-difference area.\end{uniten}

\begin{unitzh}\unitlabel{ceyuan-011.29}或問：乙出東門南行，不知步數而立；甲出北門東行二百步，望乙與城㕘相直，復就乙斜行一百七十步，與乙相㑹。問答同前。\end{unitzh}
\begin{uniten}Someone asks: B goes out of the east gate and walks south an unknown number of bu, and stands; A goes out of the north gate, walks east 200 bu, looks toward B so that B and the city are in alignment, and further walks obliquely toward B 170 bu and meets B. The question and answer are as before.\end{uniten}

\begin{unitzh}\unitlabel{ceyuan-011.30}法曰：以二行差乘甲東行為實；甲就乙斜行為方；一步常法；得半徑。\end{unitzh}
\begin{uniten}Rule: Multiply A's eastward walk by the difference of the two walks, as the dividend; A's oblique walk toward B is the square (fang); 1 bu is the constant divisor; one obtains the half-diameter.\end{uniten}

\begin{unitzh}\unitlabel{ceyuan-011.31}草曰：識别得二行相減，餘三十步，即乙出東門南行步也。{\zhsmall 更不湏用弦}立天元一以為半城徑，加乙南行，得□□，為小股。副置甲東行步：上位減天元，得下式□□，為小勾；下位加天元，得□□，為大勾也。乃置大勾，以小股乘之，得下式丨□□，合以小勾除，不受除，便以此為大股。{\zhsmall 内𢃄小勾分母}又倍天元，以小勾乘之，得□□，以減於大股，得□□□；又倍之，得下□□□，為兩个股圓差。合以勾圓差乘之，縁為其中已帶小勾分母，更不須乘，便以此為黄方冪，{\zhsmall 更無分母}寄左。然後倍天元以自之，得□□，為同數，與左相消，得□□□。上下俱半之，{\zhsmall 俱半之者蓋從簡也}得□□□。以平方開之，得一百二十步，倍之，即半徑也。合問。\end{unitzh}
\begin{uniten}Working: By identification one finds that the two walks subtracted from each other leave 30 bu, which is B's southward walk in bu out of the east gate [Note: there is no further need to use the hypotenuse]. Set up the celestial unknown as half the city diameter; add B's southward walk, obtaining □□ as the small altitude. Set out A's eastward walk in bu in duplicate: from the upper position subtract the celestial unknown, obtaining the following array, □□, as the small base; to the lower position add the celestial unknown, obtaining □□ as the large base. Then set out the large base and multiply it by the small altitude, obtaining the following array, 丨□□; it ought to be divided by the small base, but it does not accept division, so take this directly as the large altitude [Note: it contains the small base as denominator]. Further double the celestial unknown and multiply it by the small base, obtaining □□; subtract it from the large altitude, obtaining □□□; further double it, obtaining the following □□□ as two altitude–circle differences (gu yuan cha). It ought to be multiplied by the base–circle difference; because it already contains within it the small base as denominator, there is no further need to multiply, and this is taken directly as the area of the yellow square (huang fang) [Note: there is no longer any denominator]; lodge it on the left. Then double the celestial unknown and multiply it by itself, obtaining □□ as the equal quantity; make it cancel with the left, obtaining □□□. Halve both the upper and the lower [Note: halving both is presumably for the sake of simplicity], obtaining □□□. Extract its square root, obtaining 120 bu; doubling it, this is the half-diameter [Translator's note: so in the text.]. This agrees with the question.\end{uniten}

\begin{unitzh}\unitlabel{ceyuan-011.32}或問：乙出南門直行，不知步數而止；甲出北門東行二百步見之，復就乙斜行四百二十五步，與乙相㑹。問答同前。\end{unitzh}
\begin{uniten}Someone asks: B goes out of the south gate and walks straight an unknown number of bu, and stops; A goes out of the north gate, walks east 200 bu, and sees him, and further walks obliquely toward B 425 bu and meets B. The question and answer are as before.\end{uniten}

\begin{unitzh}\unitlabel{ceyuan-011.33}法曰：倍兩行差，以乘二之甲東行，為實；從空；四之甲東行於上，倍兩行差，加上位，為隅；得半徑。\end{unitzh}
\begin{uniten}Rule: Double the difference of the two walks and with it multiply twice A's eastward walk, as the dividend; the joining term is empty; put four times A's eastward walk above, double the difference of the two walks and add it to the upper position, as the corner; one obtains the half-diameter.\end{uniten}

\begin{unitzh}\unitlabel{ceyuan-011.34}草曰：識别得二行差二百二十五步，即半徑為勾之股也。立天元一以為半徑，便是小勾；其二行差，便是小股。乃置甲東行步，加天元，得□□，為大勾。以小股乗之，得下式□□；又以小勾除之，得□□，為大股。又倍天元以減之，得□□□，為股圓差。又倍之，得□□□，為兩个股圓差，於上。乃以天元減甲東行，得□□，為勾圓差，以乘上位，得下式□□○□，為城徑冪。{\zhsmall 寄左}然後倍天元一以自之，得□□，為同數，與左相消，得□□□。開平方，得一百二十步，倍之，即城徑也。合問。{\zhsmall 按此係得數各升一位然後開平方}\end{unitzh}
\begin{uniten}Working: By identification one finds that the difference of the two walks, 225 bu, is the altitude for which the half-diameter is the base. Set up the celestial unknown as the half-diameter; this is the small base, and the difference of the two walks is the small altitude. Then set out A's eastward walk in bu and add the celestial unknown, obtaining □□ as the large base. Multiply it by the small altitude, obtaining the following array, □□; further divide it by the small base, obtaining □□ as the large altitude. Further double the celestial unknown and subtract it from this, obtaining □□□ as the altitude–circle difference. Further double it, obtaining □□□ as two altitude–circle differences, placed above. Then reduce A's eastward walk by the celestial unknown, obtaining □□ as the base–circle difference; multiply the upper position by it, obtaining the following array, □□○□, as the area of the city diameter [Note: lodge it on the left]. Then double the celestial unknown and multiply it by itself, obtaining □□ as the equal quantity; make it cancel with the left, obtaining □□□. Extract the square root, obtaining 120 bu; doubling it, this is the city diameter. This agrees with the question. [Siku editors: here the numbers obtained are each raised by one place, and then the square root is extracted.]\end{uniten}

\begin{unitzh}\unitlabel{ceyuan-011.35}又法：併二數，以二數差乗之，開方，得底股。復以甲東行二百步乘之，為實；併二數而半之，以為法；如法，得二百四十步，即城徑也。合問。{\zhsmall 此用股上容圓求之，比前法極為簡易}\end{unitzh}
\begin{uniten}Another rule: Add the two numbers together and multiply by the difference of the two numbers; extract the square root, obtaining the bottom altitude (di gu). Then multiply it by A's eastward walk of 200 bu, as the dividend; add the two numbers together and halve, taking this as the divisor; dividing by the divisor, one obtains 240 bu, which is the city diameter. This agrees with the question. [Note: this finds it by using the circle contained on the altitude (gu shang rong yuan); compared with the previous method it is extremely simple and easy.]\end{uniten}

\begin{unitzh}\unitlabel{ceyuan-011.36}或問：乙從乾隅南行，不知步數而止；甲出北門東行二百步，望見之，復就乙斜行六百八十步，與乙相會。問答同前。\end{unitzh}
\begin{uniten}Someone asks: B walks south from the Qian corner (qian yu) an unknown number of bu, and stops; A goes out of the north gate, walks east 200 bu, looks and sees him, and further walks obliquely toward B 680 bu and meets B. The question and answer are as before.\end{uniten}

\begin{unitzh}\unitlabel{ceyuan-011.37}法曰：併二行，以二行差乘之，内減二行差冪，為實；併二行步及二行相減數{\zhsmall 按即倍乙斜行}，為從；二步常法；得半徑。\end{unitzh}
\begin{uniten}Rule: Add the two walks together, multiply by the difference of the two walks, and subtract from this the area of the difference of the two walks, as the dividend; add together the bu of the two walks and the number from subtracting the two walks from each other [Siku editors: that is, twice the oblique walk to B], as the joining term; 2 bu is the constant divisor; one obtains the half-diameter.\end{uniten}

\begin{unitzh}\unitlabel{ceyuan-011.38}草曰：識别得斜行六百八十步，即大弦也；其二行相減，餘四百八十步，即乙南行步内減半徑也。立天元一為半城徑，副置之：上位加二行相減數，得□□，為大股也；下位加甲東行步，得□□，為大勾也。乃以大股自增乘，得丨□□，為大股冪。{\zhsmall 寄左}乃併大勾、大弦，得□□，於上；又以大勾減大弦，得□□，為大差，以乘上位，得□□□，為同數，與左相消，得□□□。開平方，得一百二十步，倍之，即城徑也。合問。\end{unitzh}
\begin{uniten}Working: By identification one finds that the oblique walk, 680 bu, is the large hypotenuse (da xian); the two walks subtracted from each other leave 480 bu, which is B's southward walk in bu less the half-diameter. Set up the celestial unknown as half the city diameter and set it out in duplicate: to the upper position add the number from subtracting the two walks from each other, obtaining □□ as the large altitude; to the lower position add A's eastward walk in bu, obtaining □□ as the large base. Then multiply the large altitude by itself by increase-multiplication (zi zeng cheng), obtaining 丨□□ as the area of the large altitude [Note: lodge it on the left]. Then add the large base and the large hypotenuse, obtaining □□, placed above; further subtract the large base from the large hypotenuse, obtaining □□ as the large difference (da cha); multiply the upper position by it, obtaining □□□ as the equal quantity; make it cancel with the left, obtaining □□□. Extract the square root, obtaining 120 bu; doubling it, this is the city diameter. This agrees with the question.\end{uniten}

\begin{unitzh}\unitlabel{ceyuan-011.39}又法求大差。法曰：二行差自乘為實；置二之二行差於上，乃以甲東行步減二行差，又半之，以減於上，為益方；{\zhsmall 按三因斜行步二因東行步相減折半亦同}半步常法\end{unitzh}
\begin{uniten}Another rule, to find the large difference. Rule: The difference of the two walks multiplied by itself is the dividend; set out twice the difference of the two walks above; then subtract A's eastward walk in bu from the difference of the two walks, further halve it, and subtract it from the upper one, as the augmenting square [Siku editors: multiplying the oblique walk in bu by 3 and the eastward walk in bu by 2, subtracting them from each other and halving, is also the same]; 1/2 bu is the constant divisor [continues]\end{uniten}

\begin{unitzh}\unitlabel{ceyuan-012.1}草曰：立天元一為大差，減於二行差，得□□為半城徑。以自之，得丨□□為半徑冪{\zhsmall 寄左}。乃以半城徑減於甲東行，得下式丨□為小差。又以天元乘之，得丨□。又以半之，得□□為同數，與左相消，得下式□□□。以平方開之，得三百六十步，即大差也。合問。\end{unitzh}
\begin{uniten}Working: Set up the celestial unknown (tianyuan) as the large difference (dacha); subtract it from the difference of the two walks, obtaining □□ as half the city diameter. [Translator's note: here and below, □ stands for entries of the counting-rod diagrams that the transcription does not reproduce; they are given as the text lists them.] Multiply it by itself, obtaining 丨□□ as the area (mi) of the radius [Note: set aside on the left (ji zuo)]. Then subtract half the city diameter from A's eastward walk, obtaining the figure below (xiashi), 丨□, as the small difference (xiaocha). Further multiply it by the celestial unknown, obtaining 丨□. Further halve it, obtaining □□ as the equal quantity (tongshu). Cancel it against the left (xiangxiao), obtaining the figure below, □□□. Extract the square root, obtaining 360 bu (步), which is the large difference. This fits the question.\end{uniten}

\begin{unitzh}\unitlabel{ceyuan-012.2}或問：乙出東門，不知步數而立；甲出北門東行二百步，望乙與城叅相直；復就乙斜行一百三十六步，與乙相㑹。問答同前。\end{unitzh}
\begin{uniten}Suppose it is asked: B goes out of the east gate and stands still, the number of bu not being known. A goes out of the north gate, walks east 200 bu and looks toward B, B and the city being in line. He then walks obliquely toward B for 136 bu and meets B. Question and answer as before.\end{uniten}

\begin{unitzh}\unitlabel{ceyuan-012.3}法曰：甲東行步内減二之二行差{\zhsmall 按：倍斜行步内減東行步，亦同。}，餘以乘甲東行為實，一步常法，得半徑。\end{unitzh}
\begin{uniten}Rule: From A's eastward walk in bu subtract twice the difference of the two walks [Siku editors: doubling the oblique walk in bu and subtracting the eastward walk in bu from it is also the same]. Multiply A's eastward walk by the remainder, as the dividend (shi); 1 bu as the constant divisor (changfa); one obtains the radius.\end{uniten}

\begin{unitzh}\unitlabel{ceyuan-012.4}草曰：别得二行相減餘六十四步，即半徑為股之勾。立天元一為半城徑，就以為股率，其二行差即勾率也。乃置甲東行步，加天元，得□□為大勾；以天元股率乘之，得丨□，合以勾率除之，不受，便以此為大股{\zhsmall 内𢃄勾率分母}。乃倍天元，以勾率乘之，得□，以減大股，得丨□為一个大差，於上{\zhsmall 内𢃄勾率分母}。乃以天元減甲東行，得□□為小差，以乘上位，□□□為半段黄方冪{\zhsmall 内寄勾率為母}，寄左。然後以天元自之，又以勾率乘之，又倍之，得□□為同數，與左相消，得下式丨□□。以平方開之，得一百二十步，倍之，即城徑也。合問。\end{unitzh}
\begin{uniten}Working: By identification (bie) one obtains that the remainder of subtracting the two walks from each other, 64 bu, is the base (gou) of which the radius is the altitude (gu). Set up the celestial unknown as half the city diameter, and take it directly as the altitude rate (gu lü); the difference of the two walks is then the base rate (gou lü). Then put down A's eastward walk in bu and add the celestial unknown, obtaining □□ as the large base. Multiply it by the celestial unknown, the altitude rate, obtaining 丨□. It should be divided by the base rate, but it does not admit this, so take this as the large altitude [Note: carrying the base rate within as denominator (fenmu)]. Then double the celestial unknown and multiply it by the base rate, obtaining □. Subtract it from the large altitude, obtaining 丨□ as one large difference, in the upper position [Note: carrying the base rate within as denominator]. Then subtract the celestial unknown from A's eastward walk, obtaining □□ as the small difference, and multiply the upper position by it: □□□ is a half portion (banduan) of the area of the yellow square (huangfang) [Note: with the base rate lodged within as denominator]; set it aside on the left. After that, multiply the celestial unknown by itself, further multiply it by the base rate, and further double it, obtaining □□ as the equal quantity. Cancel it against the left, obtaining the figure below, 丨□□. Extract the square root, obtaining 120 bu; double it, and this is the city diameter. This fits the question.\end{uniten}

\begin{unitzh}\unitlabel{ceyuan-012.5}或問：乙出東門直行一十六步而止，甲出北門東行二百步，望見乙與城叅相直。問答同前。\end{unitzh}
\begin{uniten}Suppose it is asked: B goes out of the east gate, walks straight on 16 bu and stops. A goes out of the north gate, walks east 200 bu, and looks and sees B in line with the city. Question and answer as before.\end{uniten}

\begin{unitzh}\unitlabel{ceyuan-012.6}法曰：二行步相減，餘以自乘，内減乙東行冪，為實；二之甲東行，為益從；一步隅法；得半徑。\end{unitzh}
\begin{uniten}Rule: Subtract the two walks in bu from each other; multiply the remainder by itself, and subtract from it the area of B's eastward walk, as the dividend. Twice A's eastward walk is the subtractive joining term (yi cong); 1 bu is the corner divisor (yufa); one obtains the radius.\end{uniten}

\begin{unitzh}\unitlabel{ceyuan-012.7}草曰：立天元一以為半城徑，加乙行步，併以減於甲行步，得□□為平勾率，其天元半徑即平股率也。乃置乙東行一十六步為小勾，以股率乘之，得□，合以勾率除之，今不受除，便以此為小股{\zhsmall 内帶勾率分母}。又置乙東行，加二天元，得□□為大勾，以股率乘之，得□□，合以勾率除之，今不受除，便以此為大股{\zhsmall 内寄勾率為母}。以此小股、大股相乘，得□□□為半徑冪{\zhsmall 内寄勾率冪為母}，寄左。然後以勾率冪乘天元冪，得丨□□□為相同數，相消，得□□□□。開平方，得一百二十步，倍之，即城徑也。合問。{\zhsmall 按：此係得數各降二位，然後開平方。}\end{unitzh}
\begin{uniten}Working: Set up the celestial unknown and take it as half the city diameter. Add B's walk in bu, and subtract the sum from A's walk in bu, obtaining □□ as the level base rate (ping gou lü); the celestial unknown, the radius, is then the level altitude rate (ping gu lü). Then put down B's eastward walk, 16 bu, as the small base, and multiply it by the altitude rate, obtaining □. It should be divided by the base rate, but now it does not admit division, so take this as the small altitude [Note: carrying the base rate within as denominator]. Further put down B's eastward walk and add twice the celestial unknown, obtaining □□ as the large base. Multiply it by the altitude rate, obtaining □□. It should be divided by the base rate, but now it does not admit division, so take this as the large altitude [Note: with the base rate lodged within as denominator]. Multiply this small altitude and large altitude together, obtaining □□□ as the area of the radius [Note: with the area of the base rate lodged within as denominator]; set it aside on the left. After that, multiply the area of the celestial unknown by the area of the base rate, obtaining 丨□□□ as the mutually equal quantity (xiangtong shu). Cancel them against each other, obtaining □□□□. Extract the square root, obtaining 120 bu; double it, and this is the city diameter. This fits the question. [Siku editors: here the numbers obtained are each lowered by two places, and then the square root is extracted.]\end{uniten}

\begin{unitzh}\unitlabel{ceyuan-012.8}或問：甲乙二人同出北門，向東行，至東北十字道口分路，乙折南行一百五十步而立，甲又向東行，甲前後通行了二百步，廻望乙，恰與城相直。問答同前。\end{unitzh}
\begin{uniten}Suppose it is asked: two persons, A and B, go out of the north gate together and walk east. At the crossroads to the north-east they part ways: B turns, walks south 150 bu and stands still, while A walks on further east. A has walked 200 bu in all, before and after; looking back at B, B is exactly in line with the city. Question and answer as before.\end{uniten}

\begin{unitzh}\unitlabel{ceyuan-012.9}法曰：以二行步相乘於上，又以南行步乗之，為實；二行步相乘於上，又以乙南行減於甲東行，得數，復以乙南行乘之，加上位，共為法；得半徑。\end{unitzh}
\begin{uniten}Rule: Multiply the two walks in bu by each other, in the upper position; further multiply it by the southward walk in bu, as the dividend. Multiply the two walks in bu by each other, in the upper position; further subtract B's southward walk from A's eastward walk, and multiply the number obtained again by B's southward walk. Add it to the upper position, together making the divisor (fa); one obtains the radius.\end{uniten}

\begin{unitzh}\unitlabel{ceyuan-012.10}草曰：立天元一為半城徑，副之：上位加甲行步，得□□為大勾也；下位減於甲行步，餘□□為小勾也；其乙折南行即小股也。置大勾，以小股乘之，得□□，内寄小勾□□為母，便以為大股也。再置天元，以母乘之，得□□，減於大股，餘丨□□為半个矮梯底，於上{\zhsmall 内寄小勾為母}。再置乙折行步，内減天元，得□□為半个矮梯頭，以乘上位，得□□□□為半徑冪{\zhsmall 寄左}。乃以小勾分母乘天元冪，得下式□□□為同數，與左相消，得□□，上法下實，如法而一，得一百二十步，即城之半徑也。合問。\end{unitzh}
\begin{uniten}Working: Set up the celestial unknown as half the city diameter, and duplicate it (fu). To the upper position add A's walk in bu, obtaining □□ as the large base. The lower position is subtracted from A's walk in bu, the remainder □□ being the small base. B's turned southward walk is then the small altitude. Put down the large base and multiply it by the small altitude, obtaining □□, with the small base, □□, lodged within as denominator; take this as the large altitude. Again put down the celestial unknown and multiply it by the denominator, obtaining □□. Subtract it from the large altitude, the remainder 丨□□ being half of the low trapezoid's bottom, in the upper position [Note: with the small base lodged within as denominator]. Again put down B's turned walk in bu and subtract the celestial unknown from it, obtaining □□ as half of the low trapezoid's (aiti) head. Multiply the upper position by it, obtaining □□□□ as the area of the radius [Note: set aside on the left]. Then multiply the area of the celestial unknown by the small-base denominator, obtaining the figure below, □□□, as the equal quantity. Cancel it against the left, obtaining □□; the upper is the divisor and the lower the dividend. Divide by the divisor, obtaining 120 bu, which is the radius of the city. This fits the question.\end{uniten}

\begin{unitzh}\unitlabel{ceyuan-012.11}又法\end{unitzh}
\begin{uniten}Another method\end{uniten}

\begin{unitzh}\unitlabel{ceyuan-012.12}法曰：二行步相乘為實；倍甲東行，内減乙南行，為法。\end{unitzh}
\begin{uniten}Rule: Multiply the two walks in bu by each other as the dividend; double A's eastward walk and subtract B's southward walk from it, as the divisor.\end{uniten}

\begin{unitzh}\unitlabel{ceyuan-012.13}草曰：立天元一為半圓徑，副之：上位加甲東行，得□□為大勾；下位減甲東行，得□□為小勾，此小勾便是勾圓差也；其乙南行即小股也。置大勾，以小股乘之，得下式□□，内寄小勾□□為母，便以為大股也。再置天元，以二之，又以分母乘之，得□□為全徑，以減於大股，餘得□□□為股圓差也。合以勾圓差乘之，縁内已有小勾分母，故不湏再乘，便以此為兩段之半徑冪也，更無分母{\zhsmall 寄左}。再置天元，以自之，又二之，得□□為同數，與左相消，得□□，上法下實，一百二十步，即半城徑也。合問。\end{unitzh}
\begin{uniten}Working: Set up the celestial unknown as half the circle diameter, and duplicate it. To the upper position add A's eastward walk, obtaining □□ as the large base. The lower position is subtracted from A's eastward walk, obtaining □□ as the small base; this small base is precisely the base-circle difference (gouyuan cha). B's southward walk is then the small altitude. Put down the large base and multiply it by the small altitude, obtaining the figure below, □□, with the small base, □□, lodged within as denominator; take this as the large altitude. Again put down the celestial unknown, double it, and further multiply it by the denominator, obtaining □□ as the whole diameter. Subtract it from the large altitude, the remainder obtained, □□□, being the altitude-circle difference (guyuan cha). It should be multiplied by the base-circle difference; because the small-base denominator is already within, there is no need to multiply again, and this is taken as two portions of the area of the radius, with no further denominator [Note: set aside on the left]. Again put down the celestial unknown, multiply it by itself, and further double it, obtaining □□ as the equal quantity. Cancel it against the left, obtaining □□; the upper is the divisor and the lower the dividend: 120 bu, which is half the city diameter. This fits the question.\end{uniten}

\begin{unitzh}\unitlabel{ceyuan-012.14}或問：見底勾二百步，明弦一百五十三步。問答同前。\end{unitzh}
\begin{uniten}Suppose it is asked: the bottom base (di gou) is known to be 200 bu, and the bright hypotenuse (ming xian) 153 bu. Question and answer as before.\end{uniten}

\begin{unitzh}\unitlabel{ceyuan-012.15}法曰：半底勾乘明弦，為平實；併二云數而半之，為從；五分常法；得明勾。\end{unitzh}
\begin{uniten}Rule: Multiply half the bottom base by the bright hypotenuse, as the plane dividend (ping shi). Add the two stated numbers (yunshu) and halve the sum, as the joining term (cong); 5 fen (分) is the constant divisor; one obtains the bright base.\end{uniten}

\begin{unitzh}\unitlabel{ceyuan-012.16}草曰：立天元一為明勾，加明弦，得□□為髙股也。又以天元減底勾而半之，得下式□□為平勾也。勾股相乘，得□□□為半徑冪{\zhsmall 寄左}。然後以天元乘底勾，得下式□為同數，與左相消，得□□□。開平方，得七十二步，即明弦也。以明弦乗底勾，為平方實，如法開之，得一百二十步，倍之，即城徑也。合問。\end{unitzh}
\begin{uniten}Working: Set up the celestial unknown as the bright base; add the bright hypotenuse, obtaining □□ as the high altitude (gao gu). Further subtract the celestial unknown from the bottom base and halve it, obtaining the figure below, □□, as the level base (ping gou). Multiply base and altitude together, obtaining □□□ as the area of the radius [Note: set aside on the left]. After that, multiply the bottom base by the celestial unknown, obtaining the figure below, □, as the equal quantity. Cancel it against the left, obtaining □□□. Extract the square root, obtaining 72 bu, which is the bright hypotenuse. [Translator's note: so the text, here and in the next sentence; the unknown was set up as the bright base, the bright hypotenuse being given as 153 bu.] Multiply the bottom base by the bright hypotenuse, as the dividend for the square (pingfang shi); extract it according to the method, obtaining 120 bu; double it, and this is the city diameter. This fits the question.\end{uniten}

\begin{unitzh}\unitlabel{ceyuan-012.17}或問：見底勾二百步，□弦三十四步。問答同前。\end{unitzh}
\begin{uniten}Suppose it is asked: the bottom base is known to be 200 bu, and the □ hypotenuse 34 bu. Question and answer as before. [Translator's note: here and below, the □ in “□ hypotenuse”, “□ altitude”, “□ base”, “□ sum”, “□ yellow square” and “□ yellow” stands for a character not reproduced in the transcription.]\end{uniten}

\begin{unitzh}\unitlabel{ceyuan-012.18}法曰：底勾、□弦相減，餘倍之，内減去底勾{\zhsmall 按：倍□弦減底勾，亦同。}，復以底勾乗之，於上；又以□弦冪乘上位，為三乗方實；倍底勾，以□弦冪乗之，為從；二云數相減，餘以自之，為第一亷；二云數相減，餘又倍之，為第二益亷；一步隅法；得□股。\end{unitzh}
\begin{uniten}Rule: Subtract the bottom base and the □ hypotenuse from each other; double the remainder and subtract the bottom base from it [Siku editors: doubling the □ hypotenuse and subtracting it from the bottom base is also the same]. Again multiply it by the bottom base, in the upper position; further multiply the upper position by the area of the □ hypotenuse, as the dividend of the triply multiplied square (sancheng fang). Double the bottom base and multiply it by the area of the □ hypotenuse, as the joining term. Subtract the two stated numbers from each other and multiply the remainder by itself, as the first side term (lian). Subtract the two stated numbers from each other and further double the remainder, as the second subtractive side term (yi lian). 1 bu is the corner divisor; one obtains the □ altitude.\end{uniten}

\begin{unitzh}\unitlabel{ceyuan-012.19}草曰：立天元一為□股，加□弦，得□□為平勾。以平勾減底勾，餘□□為平弦。以倍之，得□□為黄長弦也。此弦内却減底勾，餘得下式□□為明勾也。復以底勾乘之，得□□，於上。又□弦自乘，得一千一百五十六為分母，以乗上位，得□□為帶分半徑冪{\zhsmall 寄左}。然後置黄長弦，以天元乗之，得□□，合以□弦除之，不除，寄為母，便以此為全徑也。以半之，得□□為半徑{\zhsmall 内帶□弦分母}，以自之，得丨□□□為同數，與左相消，得丨□□□□。開三乗方，得三十步，即□股也。餘各依數求之。合問。\end{unitzh}
\begin{uniten}Working: Set up the celestial unknown as the □ altitude; add the □ hypotenuse, obtaining □□ as the level base. Subtract the level base from the bottom base, the remainder □□ being the level hypotenuse (ping xian). Double it, obtaining □□ as the yellow-long hypotenuse (huangchang xian). From this hypotenuse in turn subtract the bottom base, the remainder obtained being the figure below, □□, as the bright base. Again multiply it by the bottom base, obtaining □□, in the upper position. Further, multiply the □ hypotenuse by itself, obtaining 1156 as the denominator; multiply the upper position by it, obtaining □□ as the area of the radius carrying a fraction (daifen) [Note: set aside on the left]. After that, put down the yellow-long hypotenuse and multiply it by the celestial unknown, obtaining □□. It should be divided by the □ hypotenuse; it is not divided but lodged as denominator, and this is taken as the whole diameter. Halve it, obtaining □□ as the radius [Note: carrying the □ hypotenuse within as denominator]. Multiply it by itself, obtaining 丨□□□ as the equal quantity. Cancel it against the left, obtaining 丨□□□□. Extract the root of the triply multiplied square, obtaining 30 bu, which is the □ altitude. The rest are each found according to the numbers. This fits the question.\end{uniten}

\begin{unitzh}\unitlabel{ceyuan-012.20}又法：底勾内減二□弦，復以底勾乘之，復以□弦冪乘之，為三乗方實；餘亷、從並與前同。\end{unitzh}
\begin{uniten}Another method: From the bottom base subtract twice the □ hypotenuse. Again multiply it by the bottom base, and again multiply it by the area of the □ hypotenuse, as the dividend of the triply multiplied square. The remaining side terms and the joining term are all the same as before.\end{uniten}

\begin{unitzh}\unitlabel{ceyuan-012.21}草曰：識别得二數相減餘一百六十六，為平勾、虚弦共，又為平弦、□股共。於此餘數内又去半徑，即□和也。□和、□弦相併，即勾圓差也；相減，則□黄方也。又倍□弦，加□黄，亦得勾圓差也。底勾内減□股，餘即小差弦也。\end{unitzh}
\begin{uniten}Working: By identification (shibie) one obtains that the remainder of subtracting the two numbers from each other, 166, is the level base and the void hypotenuse (xu xian) together, and is also the level hypotenuse and the □ altitude together. From this remainder further remove the radius, and this is the □ sum (he). The □ sum and the □ hypotenuse added together are the base-circle difference; subtracted from each other, they are the □ yellow square. Further, doubling the □ hypotenuse and adding the □ yellow also gives the base-circle difference. Subtract the □ altitude from the bottom base, and the remainder is the small-difference hypotenuse (xiaocha xian).\end{uniten}

\begin{unitzh}\unitlabel{ceyuan-012.22}立天元一為□股，減於云數相減數，得□□為平弦。以平弦減底勾，得□□即平勾。以平勾減於云數相減數，得□□即虚弦。以天元又減虚弦，得□□即明勾也。乃置平弦，以天元乘之，得□□，合□弦除，不除，寄為母，便以此為平股也{\zhsmall 即半徑}。平股自之，得丨□□□○為半徑冪{\zhsmall 内帶□弦冪分母}，寄左。然後置底勾，以明勾乗之，得□□，又以□弦冪一千一百五十六通之，得下式□□為同數，與左相消，得丨□□□□。亷、從一一如上。\end{unitzh}
\begin{uniten}Set up the celestial unknown as the □ altitude; subtract it from the number obtained by subtracting the stated numbers from each other, obtaining □□ as the level hypotenuse. Subtract the level hypotenuse from the bottom base, obtaining □□, which is the level base. Subtract the level base from the number obtained by subtracting the stated numbers from each other, obtaining □□, which is the void hypotenuse. Further subtract the celestial unknown from the void hypotenuse, obtaining □□, which is the bright base. Then put down the level hypotenuse and multiply it by the celestial unknown, obtaining □□. It should be divided by the □ hypotenuse; it is not divided but lodged as denominator, and this is taken as the level altitude (ping gu) [Note: that is, the radius]. Multiply the level altitude by itself, obtaining 丨□□□○ as the area of the radius [Note: carrying the area of the □ hypotenuse within as denominator]; set it aside on the left. After that, put down the bottom base and multiply it by the bright base, obtaining □□; further multiply it through (tong) by the area of the □ hypotenuse, 1156, obtaining the figure below, □□, as the equal quantity. Cancel it against the left, obtaining 丨□□□□. The side terms and the joining term are, one by one, as above.\end{uniten}

\begin{unitzh}\unitlabel{ceyuan-012.23}或問：見底勾二百步，平弦一百三十六步。問答同前。\end{unitzh}
\begin{uniten}Suppose it is asked: the bottom base is known to be 200 bu, and the level hypotenuse 136 bu. Question and answer as before.\end{uniten}

\begin{unitzh}\unitlabel{ceyuan-012.24}法曰：倍平弦，内減底勾，復以底勾乗之，開平方，得半徑。\end{unitzh}
\begin{uniten}Rule: Double the level hypotenuse and subtract the bottom base from it; again multiply it by the bottom base; extract the square root, and one obtains the radius.\end{uniten}

\begin{unitzh}\unitlabel{ceyuan-012.25}草曰：立天元為半徑。先倍平弦，内減底勾，餘□為明勾。復以底勾乗之，得□為半徑冪{\zhsmall 寄左}。然後以天元冪為同數，與左相消，得丨□□。開平方，得一百二十步，又倍之，即城徑也。合問。\end{unitzh}
\begin{uniten}Working: Set up the celestial unknown as the radius. First double the level hypotenuse and subtract the bottom base from it, the remainder □ being the bright base. Again multiply it by the bottom base, obtaining □ as the area of the radius [Note: set aside on the left]. After that, take the area of the celestial unknown as the equal quantity. Cancel it against the left, obtaining 丨□□. Extract the square root, obtaining 120 bu; further double it, and this is the city diameter. This fits the question.\end{uniten}

\begin{unitzh}\unitlabel{ceyuan-012.26}或問：底勾二百步，髙弦二百五十五步。問答同前。\end{unitzh}
\begin{uniten}Suppose it is asked: the bottom base is 200 bu, and the high hypotenuse (gao xian) 255 bu. Question and answer as before.\end{uniten}

\begin{unitzh}\unitlabel{ceyuan-012.27}法曰：底勾冪乗髙弦，為立實；底勾冪為從；髙弦為亷；一為隅；得半徑。\end{unitzh}
\begin{uniten}Rule: Multiply the high hypotenuse by the area of the bottom base, as the solid dividend (li shi). The area of the bottom base is the joining term; the high hypotenuse is the side term; 1 is the corner (yu); one obtains the radius.\end{uniten}

\begin{unitzh}\unitlabel{ceyuan-012.28}草曰：識别得髙弦即皇極股也。立天元一為半徑，副之：上位加髙弦，得□□即底股也；下位減於髙弦，得□□即明股也。置明股，以底勾乗之，得□□，合以底股除，不除，寄為母，便以此為明勾。又以底勾乗之，得□□為半徑冪{\zhsmall 内帶底股分母}，寄左。然後以天元冪乗底股，得丨□□，與左相消，得丨□□□。開立方，得一百二十步，倍之，即城徑也。合問。\end{unitzh}
\begin{uniten}Working: By identification one obtains that the high hypotenuse is the imperial-pole altitude (huangji gu). Set up the celestial unknown as the radius, and duplicate it. To the upper position add the high hypotenuse, obtaining □□, which is the bottom altitude (di gu). The lower position is subtracted from the high hypotenuse, obtaining □□, which is the bright altitude (ming gu). Put down the bright altitude and multiply it by the bottom base, obtaining □□. It should be divided by the bottom altitude; it is not divided but lodged as denominator, and this is taken as the bright base. Further multiply it by the bottom base, obtaining □□ as the area of the radius [Note: carrying the bottom altitude within as denominator]; set it aside on the left. After that, multiply the bottom altitude by the area of the celestial unknown, obtaining 丨□□. Cancel it against the left, obtaining 丨□□□. Extract the cube root, obtaining 120 bu; double it, and this is the city diameter. This fits the question.\end{uniten}

\begin{unitzh}\unitlabel{ceyuan-012.29}或問：底勾二百步，□勾、□弦和五十步。問答同前。\end{unitzh}
\begin{uniten}Suppose it is asked: the bottom base is 200 bu, and the sum of the □ base and the □ hypotenuse is 50 bu. Question and answer as before.\end{uniten}

\begin{unitzh}\unitlabel{ceyuan-012.30}法曰：以二云數相減，餘加底勾，復以減餘乗之，半之，於上；以減餘自之，減上位，為實；併云數，半之，為法；得□股。\end{unitzh}
\begin{uniten}Rule: Subtract the two stated numbers from each other and add the bottom base to the remainder. Again multiply it by the remainder of the subtraction and halve it, in the upper position. Multiply the remainder of the subtraction by itself and subtract it from the upper position, as the dividend. Add the stated numbers and halve the sum, as the divisor; one obtains the □ altitude.\end{uniten}

\begin{unitzh}\unitlabel{ceyuan-012.31}草曰：别得二數相減餘為小差股。立天元一為□股，減於小差股，得□□即半徑也。又以天元減半徑，得□□為虚股，於上。又以半徑加底勾，得□□為通勾，於下。上下相乗，得□□□，折半，得丨□□為半徑冪{\zhsmall 寄左}。然後以半徑自之，得下式丨□□為同數，與左相消，得□□，上法下實，得三十步，即□股也。合問。\end{unitzh}
\begin{uniten}Working: By identification one obtains that the remainder of subtracting the two numbers from each other is the small-difference altitude (xiaocha gu). Set up the celestial unknown as the □ altitude; subtract it from the small-difference altitude, obtaining □□, which is the radius. Further subtract the celestial unknown from the radius, obtaining □□ as the void altitude (xu gu), in the upper position. Further add the radius to the bottom base, obtaining □□ as the overall base (tong gou), in the lower position. Multiply upper and lower together, obtaining □□□; halve it, obtaining 丨□□ as the area of the radius [Note: set aside on the left]. After that, multiply the radius by itself, obtaining the figure below, 丨□□, as the equal quantity. Cancel it against the left, obtaining □□; the upper is the divisor and the lower the dividend, obtaining 30 bu, which is the □ altitude. This fits the question.\end{uniten}

\begin{unitzh}\unitlabel{ceyuan-012.32}或問：見底勾二百步，明股、明弦和二百八十八步。問答同前。\end{unitzh}
\begin{uniten}Suppose it is asked: the bottom base is known to be 200 bu, and the sum of the bright altitude and the bright hypotenuse 288 bu. Question and answer as before.\end{uniten}

\begin{unitzh}\unitlabel{ceyuan-012.33}法曰：二數相減，又半之，得數又減於底勾，餘為泛率。以泛率自之，又倍之，於上位；又二數相減而半之，以乗和步，所得減於上倍，為實。倍泛率於上位，又半底勾減和步，加上位，為法。得明勾。\end{unitzh}
\begin{uniten}Rule: Subtract the two numbers from each other and further halve the result; the number obtained is further subtracted from the bottom base, the remainder being the provisional rate (fanlü). Multiply the provisional rate by itself and further double it, in the upper position. Further subtract the two numbers from each other, halve, and multiply the sum in bu by it; subtract what is obtained from the doubled amount above, as the dividend. Double the provisional rate, in the upper position; further subtract half the bottom base from the sum in bu, and add it to the upper position, as the divisor. One obtains the bright base.\end{uniten}

\begin{unitzh}\unitlabel{ceyuan-013.1}草曰：别得和步得明勾為大差也，大差得底勾為二中差。立天元一為明勾，加和步，得□□，為股圓差也{\zhsmall 即大差}；内又加底勾，得□，折半，得□□，即通勾通股差也{\zhsmall 此即中差}。置大差，減中差，得下□□，即小差也。大小差相乘，得□□□，為半段圓徑冪{\zhsmall 寄左}。乃置底勾，内減小差，得□□，為半徑，以自之，得□□□，倍之，得下式□□□，為同數，與左相消，得□□，上法下實，得七十二步，即明勾也。合問。\end{unitzh}
\begin{uniten}Working: Separately we obtain: the sum in bu (he bu) together with the bright base (ming gou) makes the great difference (da cha); the great difference together with the bottom base (di gou) makes twice the middle difference (zhong cha). Set up the celestial unknown as the bright base; add the sum in bu, obtaining [array: □□], which is the altitude–circle difference (gu yuan cha) [Note: that is, the great difference]; to it further add the bottom base, obtaining [array: □]; halve it, obtaining [array: □□], which is the difference of the overall base (tong gou) and the overall altitude (tong gu) [Note: this is the middle difference]. Set out the great difference and subtract the middle difference, obtaining below [array: □□], which is the small difference (xiao cha). Multiply the great and small differences together, obtaining [array: □□□], which is a half portion (ban duan) of the area (mi) of the circle diameter [Note: lodge it on the left]. Then set out the bottom base and subtract the small difference from it, obtaining [array: □□], which is the radius; multiply it by itself, obtaining [array: □□□]; double it, obtaining the expression below, [array: □□□], as the equal number (tong shu); cancel it against the left (xiang xiao), obtaining [array: □□]; the upper is the divisor (fa) and the lower the dividend (shi); this gives 72 bu, which is the bright base. This agrees with the question.\end{uniten}

\begin{unitzh}\unitlabel{ceyuan-013.2}按：此條法草與三卷末以小差邊股共為二中差者同誤，依問另設於後。\end{unitzh}
\begin{uniten}Note: The rule and working of this item are mistaken in the same way as the one at the end of juan 3 that takes the small difference and the side altitude (bian gu) together as twice the middle difference. Following the question, another is set out separately below.\end{uniten}

\begin{unitzh}\unitlabel{ceyuan-013.3}法曰：以底勾乘明股弦和冪為實。倍底勾，以明股和乗之，加入明股弦和冪，為從。倍明股弦和，内減底勾，為亷。一為隅。開帶縱立方，得明勾。\end{unitzh}
\begin{uniten}Rule: Multiply the area (mi) of the sum of the bright altitude and hypotenuse (ming gu xian he) by the bottom base to make the dividend (shi). Double the bottom base, multiply it by the bright altitude sum (ming gu he) [Translator's note: evidently for the sum of the bright altitude and hypotenuse], and add in the area (mi) of the sum of the bright altitude and hypotenuse, to make the follower (cong). Double the sum of the bright altitude and hypotenuse and subtract the bottom base from it, to make the edge (lian). 1 is the corner (yu). Extract the cube root with extension (dai zong) to obtain the bright base.\end{uniten}

\begin{unitzh}\unitlabel{ceyuan-013.4}草曰：别得明弦得明勾為髙股，髙勾即半徑也；底勾為平勾弦和，明勾為平勾弦較，平股即半徑也。立天元一為明勾，自之，得丨□，應以明股弦和除之，不除，便以為明股弦較{\zhsmall 内寄明股弦和分母}。明股弦和自之，得□，為股弦和，以加股弦較，得丨□□，為倍明弦。以分母乗倍天元，得□，為倍明勾，與倍明弦相加，得丨□□，為倍髙股。置底勾，減天元，得□□，為倍平勾，與倍髙股相乘，得□□□□，為城徑冪{\zhsmall 内寄明股弦和分母}，寄左。又倍天元與倍底勾相乘，得□，以寄分母乘之，得□，為相同數，與左相消，得丨□□□，開立方，得明勾。合問。\end{unitzh}
\begin{uniten}Working: Separately we obtain: the bright hypotenuse (ming xian) together with the bright base makes the high altitude (gao gu); the high base (gao gou) is just the radius. The bottom base is the sum of the level base and hypotenuse (ping gou xian he); the bright base is the difference of the level base and hypotenuse (ping gou xian jiao); the level altitude (ping gu) is just the radius. Set up the celestial unknown as the bright base; multiply it by itself, obtaining [array: 丨□]; it ought to be divided by the sum of the bright altitude and hypotenuse; not dividing it, simply take it as the difference of the bright altitude and hypotenuse (ming gu xian jiao) [Note: within it the sum of the bright altitude and hypotenuse is lodged as denominator (fen mu)]. Multiply the sum of the bright altitude and hypotenuse by itself, obtaining [array: □], as the sum of altitude and hypotenuse (gu xian he); add it to the difference of altitude and hypotenuse (gu xian jiao), obtaining [array: 丨□□], as twice the bright hypotenuse. Multiply twice the celestial unknown by the denominator, obtaining [array: □], as twice the bright base; add it to twice the bright hypotenuse, obtaining [array: 丨□□], as twice the high altitude. Set out the bottom base and subtract the celestial unknown, obtaining [array: □□], as twice the level base; multiply it with twice the high altitude, obtaining [array: □□□□], as the area (mi) of the city diameter [Note: within it the sum of the bright altitude and hypotenuse is lodged as denominator]; lodge it on the left. Further, multiply twice the celestial unknown with twice the bottom base, obtaining [array: □]; multiply it by the lodged denominator, obtaining [array: □], as the mutually equal number (xiang tong shu); cancel it against the left, obtaining [array: 丨□□□]; extract the cube root to obtain the bright base. This agrees with the question.\end{uniten}

\begin{unitzh}\unitlabel{ceyuan-013.5}測圓海鏡卷四\end{unitzh}
\begin{uniten}Sea Mirror of Circle Measurements, juan 4\end{uniten}

\section*{卷05}\addcontentsline{toc}{section}{卷05}

\begin{unitzh}\unitlabel{ceyuan-014.1}欽定四庫全書\end{unitzh}
\begin{uniten}Imperially Approved Complete Library of the Four Treasuries (Qinding Siku quanshu)\end{uniten}

\begin{unitzh}\unitlabel{ceyuan-014.2}測圓海鏡卷五\end{unitzh}
\begin{uniten}Sea Mirror of Circle Measurements, juan (卷) 5\end{uniten}

\begin{unitzh}\unitlabel{ceyuan-014.3}元 李冶 撰\end{unitzh}
\begin{uniten}Composed by Li Ye of the Yuan\end{uniten}

\begin{unitzh}\unitlabel{ceyuan-014.4}大股一十八問\end{unitzh}
\begin{uniten}The great altitude (dagu): 18 problems\end{uniten}

\begin{unitzh}\unitlabel{ceyuan-014.5}或問：乙出南門直行一百三十五步而立，甲從乾隅南行六百步，望乙與城叅相直。問答同前。\end{unitzh}
\begin{uniten}Someone asks: B (yi 乙) goes out of the south gate, walks straight on 135 bu (步) and stands; A (jia 甲) walks south from the qian corner (乾隅) 600 bu, looks toward B, and B is aligned in a straight line with the city. Question and answer are the same as before. [Translator's note: the qian corner is the northwest corner of the square enclosing the round city.]\end{uniten}

\begin{unitzh}\unitlabel{ceyuan-014.6}法曰：倍二行差，内減甲南行步，復以乗甲南行步為實。{\zhsmall 倍二行差減甲南行步，即是甲南行步内減二之乙南行也。}四之甲南行步内減二之乙南行為從方，四為益隅，開平方得半徑。\end{unitzh}
\begin{uniten}Rule: Double the difference of the two walks, subtract A's southward walk in bu from it, and further multiply by A's southward walk in bu to make the dividend (shi). [Note: doubling the difference of the two walks and subtracting A's southward walk in bu is just A's southward walk in bu minus 2 times B's southward walk.] 4 times A's southward walk in bu minus 2 times B's southward walk is the joining square (congfang); 4 is the augmenting corner (yiyu). Extract the square root to obtain the half-diameter.\end{uniten}

\begin{unitzh}\unitlabel{ceyuan-014.7}草曰：立天元一為半徑，以二之加乙南行步，得□□為中股。以中股又減於甲南行步，得□□為股率。其天元半徑即勾率也。置甲南行為大股，以勾率乗之，得□，合以股率除之；不受除，便以此為大勾。{\zhsmall 内𢃄股率分母。}再置天元，以二之，以股率乘之，得□□，減於大勾，餘□□為勾圓差，於上。{\zhsmall 内有股率分母。}又以二之天元減甲南行，得□□為大差，以乘上位，得□□□為半段黄方冪。{\zhsmall 内寄股率分母。}然後以天元自之，又以股率乘之，又倍之，得□□□為同數，與左相消，得下式□□□。開平方得一百二十步，倍之即城徑也。合問。\end{unitzh}
\begin{uniten}Working: Set up the celestial unknown (tianyuan) as the half-diameter. Add 2 times it to B's southward walk in bu, obtaining □□ as the middle altitude (zhonggu). [Translator's note: □ here and below marks counting-rod arrays not reproduced in the transcription.] Further subtract the middle altitude from A's southward walk in bu, obtaining □□ as the altitude rate (gulü). The celestial-unknown half-diameter is the base rate (goulü). Set up A's southward walk as the great altitude and multiply it by the base rate, obtaining □, which should be divided by the altitude rate; as it does not admit division, simply take this as the great base (dagou). [Note: it contains the altitude rate as denominator (fenmu).] Again set up the celestial unknown, take 2 times it and multiply by the altitude rate, obtaining □□; subtract it from the great base; the remainder □□ is the base–circle difference (gouyuan cha), placed above. [Note: it has the altitude rate as denominator.] Further subtract 2 times the celestial unknown from A's southward walk, obtaining □□ as the great difference (dacha); with it multiply the upper position, obtaining □□□ as half a portion (duan) of the area (mi) of the yellow square (huangfang). [Note: the altitude-rate denominator is lodged in it.] Then multiply the celestial unknown by itself, further multiply by the altitude rate, and further double it, obtaining □□□ as the equal quantity (tongshu); eliminate (xiangxiao) it against the left, obtaining the following form □□□. Extract the square root, obtaining 120 bu; doubling it gives the city diameter. This agrees with the question.\end{uniten}

\begin{unitzh}\unitlabel{ceyuan-014.8}或問：乙出南門東行七十二步而止，甲從乾隅南行六百步，望乙與城叅相直。問答同前。\end{unitzh}
\begin{uniten}Someone asks: B goes out of the south gate, walks east 72 bu and stops; A walks south from the qian corner 600 bu, looks toward B, and B is aligned in a straight line with the city. Question and answer are the same as before.\end{uniten}

\begin{unitzh}\unitlabel{ceyuan-014.9}法曰：云數相乘為平實，甲南行為從，二益隅，得半徑。\end{unitzh}
\begin{uniten}Rule: Multiply the stated numbers together as the level dividend (pingshi); A's southward walk is the joining (cong); 2 is the augmenting corner; one obtains the half-diameter.\end{uniten}

\begin{unitzh}\unitlabel{ceyuan-014.10}草曰：别得虛勾乗通股，得半段圓徑冪，此與虚股乗通勾同。立天元一為半徑，内減乙東行，得□□為虚勾，以乘甲南行，得□□為半段徑冪。{\zhsmall 寄左。}再以天元為冪，又倍之，得□□為同數，與左相消，得□□□。開平方得一百二十步，即城徑也。合問。\end{unitzh}
\begin{uniten}Working: Separately one obtains that the empty base (xugou) multiplied by the comprehensive altitude (tonggu) gives half a portion of the area of the circle diameter; this is the same as the empty altitude (xugu) multiplied by the comprehensive base (tonggou). Set up the celestial unknown as the half-diameter; subtract B's eastward walk from it, obtaining □□ as the empty base; with it multiply A's southward walk, obtaining □□ as half a portion of the area of the diameter. [Note: lodge it on the left (jizuo).] Again make the celestial unknown into an area, and further double it, obtaining □□ as the equal quantity; eliminate it against the left, obtaining □□□. Extract the square root, obtaining 120 bu, which is the city diameter. This agrees with the question.\end{uniten}

\begin{unitzh}\unitlabel{ceyuan-014.11}或問：乙出東門直行一十六步，甲從乾隅南行六百步，望見乙。問答同前。\end{unitzh}
\begin{uniten}Someone asks: B goes out of the east gate and walks straight on 16 bu; A walks south from the qian corner 600 bu and looks and sees B. Question and answer are the same as before.\end{uniten}

\begin{unitzh}\unitlabel{ceyuan-014.12}法曰：以乙東行乘甲南行冪為實，二之乙東行乘甲行為從方，亷空，二步隅法，得半徑。\end{unitzh}
\begin{uniten}Rule: With B's eastward walk multiply the area of A's southward walk to make the dividend; 2 times B's eastward walk multiplied by A's walk is the joining square; the edge (lian) is empty; 2 bu is the corner divisor (yufa); one obtains the half-diameter.\end{uniten}

\begin{unitzh}\unitlabel{ceyuan-014.13}草曰：立天元一以為半城徑，以二之加於乙東行，得□□為勾率。又以天元減甲南行，得□□為股率。乃置乙東行，以股率乗之，得□□，合以勾率除；不除，便以此為小股。此小股即半梯之頭也。{\zhsmall 内帶勾率分母。}又以股率乗之，{\zhsmall 此股率即半梯之底也。}得□□□為半徑冪，{\zhsmall 内帶勾率分母。}寄左。然後置天元冪，以勾率通之，得□□□為同數，與左相消，得□○□□。開立方得一百二十步，倍之即城徑也。合問。\end{unitzh}
\begin{uniten}Working: Set up the celestial unknown to be the half city diameter; add 2 times it to B's eastward walk, obtaining □□ as the base rate. Further subtract the celestial unknown from A's southward walk, obtaining □□ as the altitude rate. Then set up B's eastward walk and multiply it by the altitude rate, obtaining □□, which should be divided by the base rate; without dividing, simply take this as the small altitude (xiaogu). This small altitude is the head (tou) of the half-trapezoid (banti). [Note: it contains the base rate as denominator.] Further multiply it by the altitude rate [Note: this altitude rate is the bottom (di) of the half-trapezoid], obtaining □□□ as the area of the half-diameter [Note: it contains the base rate as denominator]; lodge it on the left. Then set up the area of the celestial unknown and bring it to the common denominator (tong) with the base rate, obtaining □□□ as the equal quantity; eliminate it against the left, obtaining □○□□. Extract the cube root, obtaining 120 bu; doubling it gives the city diameter. This agrees with the question.\end{uniten}

\begin{unitzh}\unitlabel{ceyuan-014.14}或問：乙出東門南行三十步而立，甲從乾隅南行六百步，望見乙。問答同前。\end{unitzh}
\begin{uniten}Someone asks: B goes out of the east gate, walks south 30 bu and stands; A walks south from the qian corner 600 bu and looks and sees B. Question and answer are the same as before.\end{uniten}

\begin{unitzh}\unitlabel{ceyuan-014.15}法曰：二行步相乗為寳，以南行為從，一步常法，得半徑。\end{unitzh}
\begin{uniten}Rule: Multiply the two walks in bu together as the dividend; take the southward walk as the joining; 1 bu is the constant divisor (changfa); one obtains the half-diameter.\end{uniten}

\begin{unitzh}\unitlabel{ceyuan-014.16}草曰：立天元一為半徑，以減於甲南行，得□□為半梯底。以乙南行三十步為半梯頭，以乗之，得□□為半徑冪。{\zhsmall 寄左。}乃以天元冪與左相消，得丨□□。開平方得一百二十步，即半城徑也。合問。\end{unitzh}
\begin{uniten}Working: Set up the celestial unknown as the half-diameter; subtract it from A's southward walk, obtaining □□ as the bottom of the half-trapezoid. Take B's southward walk, 30 bu, as the head of the half-trapezoid and multiply by it, obtaining □□ as the area of the half-diameter. [Note: lodge it on the left.] Then eliminate the area of the celestial unknown against the left, obtaining 丨□□. Extract the square root, obtaining 120 bu, which is the half city diameter. This agrees with the question.\end{uniten}

\begin{unitzh}\unitlabel{ceyuan-014.17}或問：乙從艮隅南行一百五十步而立，甲從乾隅南行六百步，望見乙。問答同前。\end{unitzh}
\begin{uniten}Someone asks: B walks south from the gen corner (艮隅) 150 bu and stands; A walks south from the qian corner 600 bu and looks and sees B. Question and answer are the same as before. [Translator's note: the gen corner is the northeast corner of the square enclosing the city.]\end{uniten}

\begin{unitzh}\unitlabel{ceyuan-014.18}法曰：二行步相乗為實，并二行步為法，得半徑。\end{unitzh}
\begin{uniten}Rule: Multiply the two walks in bu together as the dividend; add the two walks in bu together as the divisor (fa); one obtains the half-diameter.\end{uniten}

\begin{unitzh}\unitlabel{ceyuan-014.19}草曰：立天元一為半徑，副置之：上以減於乙南行，得□□為半梯頭；下以減於甲南行，得□□為半梯底。上下相乗，得丨□□為半徑冪。{\zhsmall 寄左。}乃以天元冪與左相消，得下式□□，上法下實，如法而一，得一百二十步，倍之即城徑也。合問。\end{unitzh}
\begin{uniten}Working: Set up the celestial unknown as the half-diameter and set it out in duplicate (fuzhi): the upper one is subtracted from B's southward walk, obtaining □□ as the head of the half-trapezoid; the lower one is subtracted from A's southward walk, obtaining □□ as the bottom of the half-trapezoid. Multiply upper and lower together, obtaining 丨□□ as the area of the half-diameter. [Note: lodge it on the left.] Then eliminate the area of the celestial unknown against the left, obtaining the following form □□: the upper is the divisor, the lower is the dividend; dividing by the divisor gives 120 bu; doubling it gives the city diameter. This agrees with the question.\end{uniten}

\begin{unitzh}\unitlabel{ceyuan-014.20}或問：乙從艮隅東行八十步而立，甲從乾隅南行六百步，望見乙。問答同前。\end{unitzh}
\begin{uniten}Someone asks: B walks east from the gen corner 80 bu and stands; A walks south from the qian corner 600 bu and looks and sees B. Question and answer are the same as before.\end{uniten}

\begin{unitzh}\unitlabel{ceyuan-014.21}法曰：二行步相乘，又倍之，為實；二之乙東行為從；一步常法；得全徑。\end{unitzh}
\begin{uniten}Rule: Multiply the two walks in bu together and further double it, as the dividend; 2 times B's eastward walk is the joining; 1 bu is the constant divisor; one obtains the full diameter.\end{uniten}

\begin{unitzh}\unitlabel{ceyuan-014.22}草曰：别得乙東行八十步即小差也。立天元一為城徑，減於甲南行步，得□□為大差。以乙東行步乘之，得□□，又倍之，得□□為城徑冪。{\zhsmall 寄左。}然後以天元冪與左相消，得丨□□。開平方得二百四十步，即城徑也。合問。\end{unitzh}
\begin{uniten}Working: Separately one obtains that B's eastward walk, 80 bu, is the small difference (xiaocha). Set up the celestial unknown as the city diameter; subtract it from A's southward walk in bu, obtaining □□ as the great difference. Multiply it by B's eastward walk in bu, obtaining □□; further double it, obtaining □□ as the area of the city diameter. [Note: lodge it on the left.] Then eliminate the area of the celestial unknown against the left, obtaining 丨□□. Extract the square root, obtaining 240 bu, which is the city diameter. This agrees with the question.\end{uniten}

\begin{unitzh}\unitlabel{ceyuan-014.23}或問：南門東不知逺近有樹，甲從乾隅南行六百步，望樹與城㕘相直，復就樹斜行四百八步至樹。問答同前。\end{unitzh}
\begin{uniten}Someone asks: East of the south gate, at a distance not known, there is a tree. A walks south from the qian corner 600 bu, looks toward the tree, and the tree is aligned in a straight line with the city; then he walks obliquely toward the tree 408 bu and reaches the tree. Question and answer are the same as before.\end{uniten}

\begin{unitzh}\unitlabel{ceyuan-014.24}法曰：南行步冪内減兩段兩行相乘數為實，二之南行步為從，一步益隅，得城徑。\end{unitzh}
\begin{uniten}Rule: From the area of the southward walk in bu subtract 2 portions of the product of the two walks, as the dividend; 2 times the southward walk in bu is the joining; 1 bu is the augmenting corner; one obtains the city diameter.\end{uniten}

\begin{unitzh}\unitlabel{ceyuan-014.25}草曰：别得南行步内減城徑，即小股也；其斜行步，即小弦也；又二行相減，即大差為股之勾也。立天元一為圓徑，以減南行步，得□□為股圓差也。{\zhsmall 合為小股。}置南行步，以斜行步乘之，得□，合以小股除之；不受除，便以此為大弦。{\zhsmall 内帶小股分母。}再置南行步，以小股乗之，得□□為大股。{\zhsmall 亦帶小股分母。}以大股減大弦，得□□為小差也。合以大差乘之，縁於内帶大差分母，更不湏乘，便以為半段黄方冪。{\zhsmall 更無分母。}又二之，得□□為一段黄方冪。{\zhsmall 寄左。}然後以天元冪為同數，與左相消，得□□□。開平方得二百四十步，即城徑也。合問。\end{unitzh}
\begin{uniten}Working: Separately one obtains that the southward walk in bu minus the city diameter is the small altitude; the oblique walk in bu is the small hypotenuse (xiaoxian); and the two walks subtracted from each other give the base for which the great difference is the altitude. Set up the celestial unknown as the circle diameter; subtract it from the southward walk in bu, obtaining □□ as the altitude–circle difference (guyuan cha). [Note: it coincides with the small altitude.] Set up the southward walk in bu and multiply it by the oblique walk in bu, obtaining □, which should be divided by the small altitude; as it does not admit division, simply take this as the great hypotenuse (daxian). [Note: it contains the small altitude as denominator.] Again set up the southward walk in bu and multiply it by the small altitude, obtaining □□ as the great altitude. [Note: it also contains the small altitude as denominator.] Subtract the great altitude from the great hypotenuse, obtaining □□ as the small difference. It should be multiplied by the great difference; because it contains the great difference as denominator, there is no further need to multiply, and one simply takes it as half a portion of the area of the yellow square. [Note: there is no longer any denominator.] Further take 2 times it, obtaining □□ as 1 portion of the area of the yellow square. [Note: lodge it on the left.] Then take the area of the celestial unknown as the equal quantity and eliminate it against the left, obtaining □□□. Extract the square root, obtaining 240 bu, which is the city diameter. This agrees with the question.\end{uniten}

\begin{unitzh}\unitlabel{ceyuan-014.26}依前問，假令乙出南門東行，不知步數而立，甲從乾南行六百步，望乙與城相直，復就乙斜行四百八步。{\zhsmall 按：此即前問，以明又法。}\end{unitzh}
\begin{uniten}Following the preceding question, suppose B goes out of the south gate and walks east, an unknown number of bu, and stands; A walks south from qian 600 bu, looks toward B, and B is in line with the city; then he walks obliquely toward B 408 bu. [Note: Observe: this is just the preceding question, to make clear another method.]\end{uniten}

\begin{unitzh}\unitlabel{ceyuan-014.27}法曰：二行差冪乗甲南行為實，二之二行差以乗南行步為益方，二之二行差為隅，得半徑。\end{unitzh}
\begin{uniten}Rule: The area of the difference of the two walks multiplied by A's southward walk is the dividend; 2 times the difference of the two walks, multiplied by the southward walk in bu, is the augmenting square (yifang); 2 times the difference of the two walks is the corner (yu); one obtains the half-diameter.\end{uniten}

\begin{unitzh}\unitlabel{ceyuan-014.28}草曰：識别得二行相減，即半城徑與乙東行共也。得此數，更不須用斜。立天元為半徑，減於二行差一百九十二，得□□，即半梯頭也。又以二天元減甲南行步，得□□為股率；又以一百九十二為勾率。乃置甲南行，以勾率乘之，得□，合股率除；不除，便以此為大勾。{\zhsmall 内寄股率分母。}再置天元，以股率乘之，得□□，以減於大勾，得□□□為半梯底也。頭底相乘，得下□□□□為半城徑冪，{\zhsmall 内寄股率分母。}寄左。然後以股率乘天元冪，得□□□為同數，與左相消，得□□□。開平方得一百二十步，倍之即城徑也。合問。\end{unitzh}
\begin{uniten}Working: By discernment (shibie) one obtains that the two walks subtracted from each other give the half city diameter and B's eastward walk together. Having obtained this number, there is no further need to use the oblique walk. Set up the celestial unknown as the half-diameter and subtract it from the difference of the two walks, 192, obtaining □□, which is the head of the half-trapezoid. Further subtract 2 times the celestial unknown from A's southward walk in bu, obtaining □□ as the altitude rate; and take 192 as the base rate. Then set up A's southward walk and multiply it by the base rate, obtaining □, which should be divided by the altitude rate; without dividing, simply take this as the great base. [Note: the altitude-rate denominator is lodged in it.] Again set up the celestial unknown and multiply it by the altitude rate, obtaining □□; subtract it from the great base, obtaining □□□ as the bottom of the half-trapezoid. Multiply head and bottom together, obtaining the following □□□□ as the area of the half city diameter [Note: the altitude-rate denominator is lodged in it]; lodge it on the left. Then multiply the area of the celestial unknown by the altitude rate, obtaining □□□ as the equal quantity; eliminate it against the left, obtaining □□□. Extract the square root, obtaining 120 bu; doubling it gives the city diameter. This agrees with the question.\end{uniten}

\begin{unitzh}\unitlabel{ceyuan-014.29}或問：東門南不知逺近有樹，甲從乾隅南行六百步見樹，復向樹斜行五百一十步至樹。問答同前。\end{unitzh}
\begin{uniten}Someone asks: South of the east gate, at a distance not known, there is a tree. A walks south from the qian corner 600 bu and sees the tree; then he walks obliquely toward the tree 510 bu and reaches the tree. Question and answer are the same as before.\end{uniten}

\begin{unitzh}\unitlabel{ceyuan-014.30}法曰：二行差步乘甲南行步為實，二行之差步併甲南行步為從，二益隅。{\zhsmall 若欲從簡，上下俱折半。}\end{unitzh}
\begin{uniten}Rule: The difference of the two walks in bu multiplied by A's southward walk in bu is the dividend; the difference of the two walks in bu added to A's southward walk in bu is the joining; 2 is the augmenting corner. [Note: if one wishes for simplicity, halve both upper and lower.]\end{uniten}

\begin{unitzh}\unitlabel{ceyuan-014.31}草曰：别得二行相減數，即虚積之股也。立天元一為半徑，内減二行之差步，得□□為梯頭，於上。又以天元減於甲之南行步，得□□為梯底。上下相乗，得□□□為圓徑冪。{\zhsmall 寄左。}然後以天元冪與左相消，得□□□。開平方得一百二十步，即城徑也。合問。\end{unitzh}
\begin{uniten}Working: Separately one obtains that the number of the two walks subtracted from each other is the altitude of the empty area (xuji). Set up the celestial unknown as the half-diameter; subtract from it the difference of the two walks in bu, obtaining □□ as the trapezoid head, placed above. Further subtract the celestial unknown from A's southward walk in bu, obtaining □□ as the trapezoid bottom. Multiply upper and lower together, obtaining □□□ as the area of the circle diameter. [Note: lodge it on the left.] Then eliminate the area of the celestial unknown against the left, obtaining □□□. Extract the square root, obtaining 120 bu, which is the city diameter. This agrees with the question.\end{uniten}

\begin{unitzh}\unitlabel{ceyuan-014.32}或問：乙出東門直行，不知步數而立，甲從乾隅南行六百步，望見乙，復就乙斜行五百四十四步，與乙相㑹。問答同前。\end{unitzh}
\begin{uniten}Someone asks: B goes out of the east gate and walks straight on, an unknown number of bu, and stands; A walks south from the qian corner 600 bu and looks and sees B; then he walks obliquely toward B 544 bu and meets B. Question and answer are the same as before.\end{uniten}

\begin{unitzh}\unitlabel{ceyuan-014.33}法曰：以二行步相減，乘甲南行步得數，又半之南行步以乘之，為實；以二行差乘南行步於上，又以半之南行步乘南行步，加於上，為從方；二之南行步為益亷；一步常法；得半徑。\end{unitzh}
\begin{uniten}Rule: Subtract the two walks in bu from each other and multiply by A's southward walk in bu; the number obtained is further multiplied by half the southward walk in bu, as the dividend. Multiply the southward walk in bu by the difference of the two walks, placed above; further multiply the southward walk in bu by half the southward walk in bu and add it to the above, as the joining square; 2 times the southward walk in bu is the augmenting edge (yilian); 1 bu is the constant divisor; one obtains the half-diameter.\end{uniten}

\begin{unitzh}\unitlabel{ceyuan-014.34}草曰：别得二行相減，即半徑上勾股較。{\zhsmall 此股即半徑也。}又别得是大勾圓差不及平弦數。立天元一以為半城徑，以減南行步，得□□為中股，其斜行步即中弦也。乃立半城徑，以斜步乘之，得□，合以中股除；今不受除，便以此為平弦。{\zhsmall 内帶中股分母。}又以二行步相減，餘五十六步，為勾圓差不及平弦數。置此數，以中股乗之，得□□，復以減平弦，餘得□□為小差。{\zhsmall 内帶中股分母。}乃以二天元減甲南行步，得□□為大差，又半之，得□□，以乘小差，得□□□為半徑冪。{\zhsmall 寄左。}然後以天元自乗，又以中股通之，得□□□為同數，與左相消，得丨□□□。開立方得一百二十歩，倍之即城徑也。合問。{\zhsmall 翻法在記。}\end{unitzh}
\begin{uniten}Working: Separately one obtains that the two walks subtracted from each other give the base–altitude difference (gougu jiao) on the half-diameter. [Note: this altitude is the half-diameter.] Separately one further obtains that this is the amount by which the great base–circle difference falls short of the level hypotenuse (pingxian). Set up the celestial unknown to be the half city diameter; subtract it from the southward walk in bu, obtaining □□ as the middle altitude; the oblique walk in bu is the middle hypotenuse (zhongxian). Then set up the half city diameter and multiply it by the oblique bu, obtaining □, which should be divided by the middle altitude; now, as it does not admit division, simply take this as the level hypotenuse. [Note: it contains the middle altitude as denominator.] Further subtract the two walks in bu from each other; the remainder, 56 bu, is the amount by which the base–circle difference falls short of the level hypotenuse. Set up this number and multiply it by the middle altitude, obtaining □□; then subtract it from the level hypotenuse; the remainder obtained, □□, is the small difference. [Note: it contains the middle altitude as denominator.] Then subtract 2 times the celestial unknown from A's southward walk in bu, obtaining □□ as the great difference; further halve it, obtaining □□; with it multiply the small difference, obtaining □□□ as the area of the half-diameter. [Note: lodge it on the left.] Then multiply the celestial unknown by itself and further bring it to the common denominator with the middle altitude, obtaining □□□ as the equal quantity; eliminate it against the left, obtaining 丨□□□. Extract the cube root, obtaining 120 bu; doubling it gives the city diameter. This agrees with the question. [Note: the reversing method (fanfa) is in the record (ji).]\end{uniten}

\begin{unitzh}\unitlabel{ceyuan-014.35}或問：甲乙二人俱在乾隅，乙東行，不知步數而立，甲南行六百步，望見乙，復就乙斜行六百八十步，與乙相會。問答同前。\end{unitzh}
\begin{uniten}Someone asks: The two persons A and B are both at the qian corner. B walks east, an unknown number of bu, and stands; A walks south 600 bu and looks and sees B; then he walks obliquely toward B 680 bu and meets B. Question and answer are the same as before.\end{uniten}

\begin{unitzh}\unitlabel{ceyuan-014.36}法曰：以二行差乘二行併，開平方，得數内復減二行差，得全徑。\end{unitzh}
\begin{uniten}Rule: With the difference of the two walks multiply the sum of the two walks; extract the square root; from the number obtained further subtract the difference of the two walks; one obtains the full diameter.\end{uniten}

\begin{unitzh}\unitlabel{ceyuan-014.37}草曰：别得二行相減，即勾圓差也。先求大勾：立天元一為大勾，以二行相減餘八十步，以乘二行相併數一千二百八十步，得□為勾冪。開平方得三百二十步，即大勾也。大勾内減去勾圓差，餘二百四十步，即城徑也。合問。\end{unitzh}
\begin{uniten}Working: Separately one obtains that the two walks subtracted from each other give the base–circle difference. First seek the great base: set up the celestial unknown as the great base; with the remainder of the two walks subtracted from each other, 80 bu, multiply the number of the two walks added together, 1280 bu, obtaining □ as the area of the base. Extract the square root, obtaining 320 bu, which is the great base. From the great base subtract the base–circle difference; the remainder, 240 bu, is the city diameter. This agrees with the question.\end{uniten}

\begin{unitzh}\unitlabel{ceyuan-014.38}或問：南門外不知逺近有樹，甲從乾隅南行六百步，望樹與城㕘相直，復就樹斜行二百五十五步至樹。問答同前。\end{unitzh}
\begin{uniten}Someone asks: Outside the south gate, at a distance not known, there is a tree. A walks south from the qian corner 600 bu, looks toward the tree, and the tree is aligned in a straight line with the city; then he walks obliquely toward the tree 255 bu and reaches the tree. Question and answer are the same as before.\end{uniten}

\begin{unitzh}\unitlabel{ceyuan-015.1}法曰：倍二行相減數，内減甲南行，得數，復以乘甲南行，為實；倍二行相減數，為從；二步，益隅。得半徑。\end{unitzh}
\begin{uniten}Rule: Double the difference of the two walks; from it subtract Jia's southward walk; then multiply Jia's southward walk by the number obtained, as the dividend (shi). Double the difference of the two walks, as the following term (cong). 2 bu (步) as the augmenting corner (yi yu). One obtains the radius.\end{uniten}

\begin{unitzh}\unitlabel{ceyuan-015.2}草曰：識别得斜行步乃是樹至城心之數也。立天元一為半徑，加斜行步，得□□，為樹至城北門之步也。乃以減於甲南行，得□□，為小股率。其天元半徑即小勾率，其斜步即小弦數也。再置甲南行步，内減天元，得□□，為梯底，於上。又置梯底，内減二之小股率，得□□，即梯頭也。復以乘上位，得□□□，為半徑冪{\zhsmall 寄左}。然後以天元冪與左相消，得下式□□□。開平方，得一百二十歩，倍之，即城徑也。合問。\end{unitzh}
\begin{uniten}Working: Through discernment (shibie) one finds that the oblique-walk bu are the distance from the tree to the city center. Set up the celestial unknown (tianyuan) as the radius; add the oblique-walk bu, obtaining □□ [Translator's note: □ marks counting-rod notation not reproduced in the transcription] as the bu from the tree to the north gate of the city. Then subtract it from Jia's southward walk, obtaining □□ as the rate (lü) of the small altitude (gu). The celestial-unknown radius is the rate of the small base (gou), and the oblique bu are the number of the small hypotenuse (xian). Again set up Jia's southward-walk bu and subtract the celestial unknown from it, obtaining □□ as the trapezoid bottom (ti di), placed above. Again set up the trapezoid bottom and subtract from it 2 times the rate of the small altitude, obtaining □□, which is the trapezoid head (ti tou). Then multiply the upper position by it, obtaining □□□ as the area (mi) of the radius [Note: lodge it on the left (ji zuo)]. Afterwards eliminate the area of the celestial unknown against the left (xiang xiao), obtaining the following expression (xia shi) □□□. Extract the square root, obtaining 120 bu; double it: this is the city diameter. This agrees with the question.\end{uniten}

\begin{unitzh}\unitlabel{ceyuan-015.3}或問：東門外不知步數有槐樹一株。甲從乾隅南行，至柳樹下，望見槐樹，復斜行至槐樹下。甲自云：「我共行了一千一百四十四步。」乙從艮隅東行，望見槐樹與城相直，復斜行至槐樹下。乙自云：「我東行少，不及斜行五十六步。」問答同前。\end{unitzh}
\begin{uniten}Someone asks: Outside the east gate, at an unknown number of bu, there is one locust tree. Jia walks south from the Qian corner (qian yu) to the foot of a willow tree, looks and sees the locust tree, and then walks obliquely to the foot of the locust tree. Jia says of himself: "Altogether I walked 1144 bu." Yi walks east from the Gen corner (gen yu), looks and sees the locust tree in line with the city, and then walks obliquely to the foot of the locust tree. Yi says of himself: "My eastward walk is less: it falls short of the oblique walk by 56 bu." The question and answer are the same as before.\end{uniten}

\begin{unitzh}\unitlabel{ceyuan-015.4}法曰：甲斜行減於甲南行，以乘甲南行，得數，復以乘二之甲南行，為實。半之甲南行，以乘二之甲南行，於上；甲斜行減於甲南行，餘，復以乘甲南行，又倍之，加上位，為從方。二之甲南行，為益亷。五分，隅法{\zhsmall 按：五分隅法即半箇立方}。\end{unitzh}
\begin{uniten}Rule: Subtract Jia's oblique walk from Jia's southward walk and multiply Jia's southward walk by it; then multiply 2 times Jia's southward walk by the number obtained, as the dividend. Multiply 2 times Jia's southward walk by half of Jia's southward walk, placed above; subtract Jia's oblique walk from Jia's southward walk, then multiply Jia's southward walk by the remainder, further double it, and add it to the upper position, as the following square (cong fang). 2 times Jia's southward walk as the augmenting edge (yi lian). 5 fen (分) as the corner divisor (yu fa) [Siku editors: Note: a corner divisor of 5 fen is just half a cube].\end{uniten}

\begin{unitzh}\unitlabel{ceyuan-015.5}草曰：識别得五十六步是小差不及平弦數{\zhsmall 此小差即勾圓差也}，又為平弦上勾股差，又為甲斜行不及大股。乃副置甲共行在地，其上位加五十六步而半之，得六百步，即大股也；其下位減五十六步而半之，得五百四十四步，即今弦也。\end{unitzh}
\begin{uniten}Working: Through discernment one finds that the 56 bu are the amount by which the small difference (xiao cha) falls short of the level hypotenuse (ping xian) [Note: this small difference is the base–circle difference (gou yuan cha)]; they are also the base–altitude difference (gou gu cha) on the level hypotenuse, and also the amount by which Jia's oblique walk falls short of the large altitude (da gu). Then set up Jia's total walk in duplicate (fu zhi) on the ground: to its upper position add 56 bu and halve it, obtaining 600 bu, which is the large altitude; from its lower position subtract 56 bu and halve it, obtaining 544 bu, which is the present hypotenuse.\end{uniten}

\begin{unitzh}\unitlabel{ceyuan-015.6}立天元一為圓徑，以半之減於甲南行步，得□□，為中股；其斜行五百四十四步即中弦也。乃立半天元，以斜步乘之，得□，合以中股除之，今不受除，便以此為平弦{\zhsmall 内寄中股分母}。又置勾圓差不及平弦數，以中股乘之，得□□，復以減於平弦□□，為小差{\zhsmall 内帶小股分母}。又以天元減甲南行，倍之，得□□，為兩个大差，以乘小差，得□□□，為圓徑冪{\zhsmall 寄左}。然後以中股乘天元冪，得下式□□□，為同數，與左相消，得□□□□。開立方，得二百四十步，即城徑也。合問{\zhsmall 翻法在記}。\end{unitzh}
\begin{uniten}Set up the celestial unknown as the circle diameter; halve it and subtract it from Jia's southward-walk bu, obtaining □□ as the middle altitude (zhong gu); the oblique walk of 544 bu is the middle hypotenuse (zhong xian). Then set up half the celestial unknown and multiply it by the oblique bu, obtaining □; this should be divided by the middle altitude, but now it does not accept division, so simply take this as the level hypotenuse [Note: within it the middle altitude is lodged as denominator (fen mu)]. Again set up the amount by which the base–circle difference falls short of the level hypotenuse and multiply it by the middle altitude, obtaining □□; then subtract it from the level hypotenuse □□, as the small difference [Note: within it the small altitude is carried as denominator]. Again subtract the celestial unknown from Jia's southward walk and double it, obtaining □□ as two large differences (da cha); multiply the small difference by it, obtaining □□□ as the area of the circle diameter [Note: lodge it on the left]. Afterwards multiply the area of the celestial unknown by the middle altitude, obtaining the following expression □□□ as the equal quantity (tong shu); eliminate it against the left, obtaining □□□□. Extract the cube root, obtaining 240 bu, which is the city diameter. This agrees with the question [Note: the reversal method (fan fa) is in the Record (ji)].\end{uniten}

\begin{unitzh}\unitlabel{ceyuan-015.7}或問：出東門向南行，不知步數，有柳樹一株。甲從乾隅南行六百步，望見柳樹而止。乙出東門直行，不知步數，望柳樹與甲相直，却斜行三十四步，至柳樹下。問答同前。\end{unitzh}
\begin{uniten}Someone asks: Going out of the east gate and walking toward the south, at an unknown number of bu, there is one willow tree. Jia walks 600 bu south from the Qian corner, looks and sees the willow tree, and stops. Yi goes out of the east gate and walks straight ahead an unknown number of bu; looking at the willow tree, it is in line with Jia; then he walks obliquely 34 bu to the foot of the willow tree. The question and answer are the same as before.\end{uniten}

\begin{unitzh}\unitlabel{ceyuan-015.8}法曰：乙斜行乘甲南行數，以乗甲南行冪，為實。斜行乗甲南行冪，又三之，為從方。甲行冪内減兩段斜行南行相乘數{\zhsmall 按：甲南行内減二之乙斜行，以甲南行乘之}，為第一亷。二之南行步，為第二益亷。二步，常法。得半徑。\end{unitzh}
\begin{uniten}Rule: Multiply the number of Jia's southward walk by Yi's oblique walk, and multiply the area of Jia's southward walk by it, as the dividend. Multiply the area of Jia's southward walk by the oblique walk and further multiply it by 3, as the following square. From the area of Jia's walk subtract two portions (duan) of the product of the oblique walk and the southward walk [Siku editors: Note: from Jia's southward walk subtract 2 times Yi's oblique walk, and multiply it by Jia's southward walk], as the first edge. 2 times the southward-walk bu as the second augmenting edge. 2 bu as the constant divisor (chang fa). One obtains the radius.\end{uniten}

\begin{unitzh}\unitlabel{ceyuan-015.9}草曰：立天元一為半徑，以二之減甲南行，得□□，為大差。以自之，得□□□，為大差冪。加於南行冪，得□□□，又半之，得□□□，為大弦也，内帶大差□□分母，别寄。又置乙斜行，以大股六百步乘之，得□，合大弦除，不除，便以此為小股也{\zhsmall 内帶大弦分母}。\end{unitzh}
\begin{uniten}Working: Set up the celestial unknown as the radius; take it 2 times and subtract it from Jia's southward walk, obtaining □□ as the large difference. Multiply it by itself, obtaining □□□ as the area of the large difference. Add it to the area of the southward walk, obtaining □□□; further halve it, obtaining □□□ as the large hypotenuse (da xian); within it the large difference □□ is carried as denominator, lodged separately. Again set up Yi's oblique walk and multiply it by the large altitude, 600 bu, obtaining □; this should be divided by the large hypotenuse; not dividing, simply take this as the small altitude [Note: within it the large hypotenuse is carried as denominator].\end{uniten}

\begin{unitzh}\unitlabel{ceyuan-015.10}乃以天元減甲南行，得□□，即半梯底也。以乗小股半梯頭，得□□，為半徑冪，於上。此半徑冪内有大弦分母，縁别寄大弦分母元帶大差分母，故又用大差分母□□乘上半徑冪，得□□□，為帶分半徑冪也。所帶之分，謂只帶大弦分母也{\zhsmall 寄左}。然後以大弦乘天元冪，得□□□□，為同數，與左相消，得□□□□□。開三乘方，得一百二十步，即半城徑也。合問。\end{unitzh}
\begin{uniten}Then subtract the celestial unknown from Jia's southward walk, obtaining □□, which is half the trapezoid bottom. Multiply the small altitude, the half trapezoid head, by it, obtaining □□ as the area of the radius, placed above. Within this area of the radius there is the large-hypotenuse denominator; because the separately lodged large-hypotenuse denominator originally carries the large-difference denominator, therefore one further uses the large-difference denominator □□ to multiply the area of the radius above, obtaining □□□ as the area of the radius carrying parts (dai fen). The parts carried means that it carries only the large-hypotenuse denominator [Note: lodge it on the left]. Afterwards multiply the area of the celestial unknown by the large hypotenuse, obtaining □□□□ as the equal quantity; eliminate it against the left, obtaining □□□□□. Extract the root of the thrice-multiplied power (san cheng fang), obtaining 120 bu, which is half the city diameter. This agrees with the question.\end{uniten}

\begin{unitzh}\unitlabel{ceyuan-015.11}按：此條寄分内又帶寄分，則以所帶之分乘本條，仍以寄分乘次條者，蓋寄分為應除本條之數，而寄分内所帶之分又為應除寄分之數，今不除寄分而乘本條，則猶是寄分乘次條之理也。乗除之變，至斯而極矣。\end{unitzh}
\begin{uniten}Note: In this item the lodged parts (ji fen) again carry lodged parts within them; that one then multiplies the present item by the parts carried, and still multiplies the next item by the lodged parts, is because the lodged parts are the number by which the present item ought to be divided, while the parts carried within the lodged parts are in turn the number by which the lodged parts ought to be divided. Now, not dividing the lodged parts but multiplying the present item is still the same principle as multiplying the next item by the lodged parts. The transformations of multiplication and division reach their utmost here.\end{uniten}

\begin{unitzh}\unitlabel{ceyuan-015.12}又法：置甲南行冪於上，又置甲行冪，半之，以乗上位，為實。以斜行乗甲行冪，倍之，於上位；又以甲行再自乗，加上位，為益方。置甲行冪於上，以斜行乗甲南行，倍之，以減上位，為第一亷。甲南行步，為第二亷。半步，常法。得股圓差。\end{unitzh}
\begin{uniten}Another rule: Place the area of Jia's southward walk above; again set up the area of Jia's walk, halve it, and multiply the upper position by it, as the dividend. Multiply the area of Jia's walk by the oblique walk and double it, in the upper position; further multiply Jia's walk by itself twice (zai zi cheng) and add it to the upper position, as the augmenting square (yi fang). Place the area of Jia's walk above; multiply Jia's southward walk by the oblique walk, double it, and subtract it from the upper position, as the first edge. Jia's southward-walk bu as the second edge. 1/2 bu as the constant divisor. One obtains the altitude–circle difference (gu yuan cha).\end{uniten}

\begin{unitzh}\unitlabel{ceyuan-015.13}草曰：立天元一為股圓差{\zhsmall 即大差}，以自之為冪，以加甲南行冪，得丨□□，半之，又以天元除之，得□□□，為大弦。其甲南行即大股也。别置乙斜行三十四步，以大股乗之，得□，合大弦除，不除，便以為小股{\zhsmall 内寄大弦分母}。乃以天元加甲南行步，得□□，為全梯底也。以乗小股半梯頭，得□□，又倍之，得□□，為城徑冪{\zhsmall 内寄大弦為母}，寄左。置天元大差，減甲南行，餘□□，為圓徑，以自之，得丨□□，又以大弦分母乗之，得□□□□□，為同數，與左相消，得□□□□□。開三乘方，得三百六十步，即股圓差也。以股圓差減甲南行，餘二百四十步，即城徑也。合問。\end{unitzh}
\begin{uniten}Working: Set up the celestial unknown as the altitude–circle difference [Note: that is, the large difference]; multiply it by itself to make an area; add it to the area of Jia's southward walk, obtaining 丨□□; halve it, and further divide it by the celestial unknown, obtaining □□□ as the large hypotenuse. Jia's southward walk is the large altitude. Separately set up Yi's oblique walk, 34 bu, and multiply it by the large altitude, obtaining □; this should be divided by the large hypotenuse; not dividing, simply take it as the small altitude [Note: within it the large hypotenuse is lodged as denominator]. Then add the celestial unknown to Jia's southward-walk bu, obtaining □□ as the whole trapezoid bottom. Multiply the small altitude, the half trapezoid head, by it, obtaining □□; further double it, obtaining □□ as the area of the city diameter [Note: within it the large hypotenuse is lodged as denominator]; lodge it on the left. Set up the celestial unknown, the large difference, and subtract it from Jia's southward walk; the remainder □□ is the circle diameter; multiply it by itself, obtaining 丨□□; further multiply it by the large-hypotenuse denominator, obtaining □□□□□ as the equal quantity; eliminate it against the left, obtaining □□□□□. Extract the root of the thrice-multiplied power, obtaining 360 bu, which is the altitude–circle difference. Subtract the altitude–circle difference from Jia's southward walk; the remainder, 240 bu, is the city diameter. This agrees with the question.\end{uniten}

\begin{unitzh}\unitlabel{ceyuan-015.14}或問：甲從乾隅南行六百步而止。丙出南門直行，乙出南門東行，各不知步數而立。甲望乙、丙，悉與城㕘相直。既而乙就丙斜行一百五十三步相㑹。問答同前。\end{unitzh}
\begin{uniten}Someone asks: Jia walks 600 bu south from the Qian corner and stops. Bing goes out of the south gate and walks straight ahead; Yi goes out of the south gate and walks east; each walks an unknown number of bu and stands. Jia looks at Yi and Bing, and all are jointly in line with the city. Then Yi walks obliquely 153 bu toward Bing, and they meet. The question and answer are the same as before.\end{uniten}

\begin{unitzh}\unitlabel{ceyuan-015.15}法曰：以甲南行步再自之，於上；以斜行步乘甲南行冪，又倍之，減上位，為立方實。南行步自之，又四之，於上；以斜步乗甲南行，又倍之，減上位，為益從。六之甲行步，為從亷。四步，虚常法。得半徑。\end{unitzh}
\begin{uniten}Rule: Multiply Jia's southward-walk bu by itself twice, placed above; multiply the area of Jia's southward walk by the oblique-walk bu, further double it, and subtract it from the upper position, as the cube dividend (lifang shi). Multiply the southward-walk bu by itself, further multiply it by 4, placed above; multiply Jia's southward walk by the oblique bu, further double it, and subtract it from the upper position, as the augmenting following term (yi cong). 6 times Jia's walk bu as the following edge (cong lian). 4 bu as the void constant divisor (xu chang fa). One obtains the radius.\end{uniten}

\begin{unitzh}\unitlabel{ceyuan-015.16}草曰：立天元一為半徑，以二之減於甲南行，得□□，為大差也。以自之，得□□□，為大差冪也。乃置甲南行冪，内加大差冪而半之，得□□□，為大弦也{\zhsmall 内寄大差分母}。又置甲南行冪，内減大差冪而半之，得□□，為大勾也{\zhsmall 亦帶大差分母}。\end{unitzh}
\begin{uniten}Working: Set up the celestial unknown as the radius; take it 2 times and subtract it from Jia's southward walk, obtaining □□ as the large difference. Multiply it by itself, obtaining □□□ as the area of the large difference. Then set up the area of Jia's southward walk, add to it the area of the large difference, and halve it, obtaining □□□ as the large hypotenuse [Note: within it the large difference is lodged as denominator]. Again set up the area of Jia's southward walk, subtract from it the area of the large difference, and halve it, obtaining □□ as the large base (da gou) [Note: it also carries the large difference as denominator].\end{uniten}

\begin{unitzh}\unitlabel{ceyuan-015.17}乃置斜行步在地，以大勾乗之，得□□，合以大弦除，不除，便以此為小勾，内帶大弦為母{\zhsmall 其大勾内元有大差分母不用}，即半梯頭也{\zhsmall 寄上位}。再寄天元半徑，以大差乘之，得□□，以減於大勾，得□，為半梯底也。以乘上位，得□□□，為半徑冪也{\zhsmall 内帶大差及大弦為母}，寄左。然後置天元冪，以大差通之，又以大弦通之，得□□□□□，為同數，與左相消，得□□□□。開立方，得一百二十步，即半城徑也。合問。\end{unitzh}
\begin{uniten}Then set up the oblique-walk bu on the ground and multiply it by the large base, obtaining □□; this should be divided by the large hypotenuse; not dividing, simply take this as the small base, within which the large hypotenuse is carried as denominator [Note: the large-difference denominator originally within the large base is not used]; this is the half trapezoid head [Note: lodged in the upper position]. Again lodge the celestial-unknown radius and multiply it by the large difference, obtaining □□; subtract it from the large base, obtaining □ as the half trapezoid bottom. Multiply the upper position by it, obtaining □□□ as the area of the radius [Note: within it the large difference and the large hypotenuse are carried as denominators]; lodge it on the left. Afterwards set up the area of the celestial unknown, bring it to common terms (tong) with the large difference, and further bring it to common terms with the large hypotenuse, obtaining □□□□□ as the equal quantity; eliminate it against the left, obtaining □□□□. Extract the cube root, obtaining 120 bu, which is half the city diameter. This agrees with the question.\end{uniten}

\begin{unitzh}\unitlabel{ceyuan-015.18}依前問，假令南門外有樹，乙出南門東行，不知步數而立{\zhsmall 只云乙東行步少於樹去城步}。甲從乾隅向南行六百步，望樹與乙，悉與城㕘相直。乙就樹斜行一百五十三步至。問答同前。\end{unitzh}
\begin{uniten}Following the previous question: suppose that outside the south gate there is a tree. Yi goes out of the south gate, walks east an unknown number of bu, and stands [Note: it only says that Yi's eastward-walk bu are fewer than the bu by which the tree is distant from the city]. Jia walks 600 bu toward the south from the Qian corner and looks at the tree and Yi; all are jointly in line with the city. Yi walks obliquely 153 bu toward the tree and arrives. The question and answer are the same as before.\end{uniten}

\begin{unitzh}\unitlabel{ceyuan-015.19}法曰：以斜行步乗甲行冪，為立方實。以甲行冪半之，於上；以斜行步乘甲行步，減上位，為益從。亷空。五分，隅。得大勾大弦差。\end{unitzh}
\begin{uniten}Rule: Multiply the area of Jia's walk by the oblique-walk bu, as the cube dividend. Halve the area of Jia's walk, placed above; multiply Jia's walk bu by the oblique-walk bu and subtract it from the upper position, as the augmenting following term. The edge is empty (kong). 5 fen as the corner. One obtains the difference of the large base and the large hypotenuse.\end{uniten}

\begin{unitzh}\unitlabel{ceyuan-015.20}草曰：别得斜步即小弦，小弦得小和，即勾弦差也。立天元一為股圓差，以自之為冪，副之。上以加甲南行冪而半之，得□□□，為大弦也{\zhsmall 寄大差分母}；下以減於甲南行冪而半之，得下式□□□，為大勾也{\zhsmall 寄大差分母}。\end{unitzh}
\begin{uniten}Working: Distinguishing (bie), one finds that the oblique bu are the small hypotenuse; the small hypotenuse gives the small sum (xiao he), which is the base–hypotenuse difference (gou xian cha). Set up the celestial unknown as the altitude–circle difference; multiply it by itself to make an area, and duplicate it. The upper one: add it to the area of Jia's southward walk and halve it, obtaining □□□ as the large hypotenuse [Note: the large difference is lodged as denominator]. The lower one: subtract it from the area of Jia's southward walk and halve it, obtaining the following expression □□□ as the large base [Note: the large difference is lodged as denominator].\end{uniten}

\begin{unitzh}\unitlabel{ceyuan-015.21}乃置斜步，以大勾乗之，得下□□□，合以大弦除，不除，便以此為小勾{\zhsmall 寄大弦分母}。又置斜步，以甲南行乗之，得□，合以大弦除為小股，不除，而又以同母分通之，得□○，為同分小股也{\zhsmall 内只寄大弦分母}。注：{\zhsmall 大股乘時無大差分母，故今通之，以齊大勾上所有大差分母也。}又置斜步，以大弦通之，得□□□，為通分小弦也。三位相併，得□□，為股圓差{\zhsmall 寄左}。然後置天元大差，以大弦分母通之，得□○□，為同數，與左相消，得□○□□。開立方，得三百六十步，即股圓差也。以股圓差減於甲南行步，即城徑也。合問。\end{unitzh}
\begin{uniten}Then set up the oblique bu and multiply it by the large base, obtaining the following □□□; this should be divided by the large hypotenuse; not dividing, simply take this as the small base [Note: the large hypotenuse is lodged as denominator]. Again set up the oblique bu and multiply it by Jia's southward walk, obtaining □; this should be divided by the large hypotenuse to make the small altitude; not dividing, one further brings it to common terms with the parts of the same denominator, obtaining □○ as the small altitude in the same parts [Note: within it only the large hypotenuse is lodged as denominator]. Comment: [Note: when the multiplication by the large altitude was done there was no large-difference denominator; therefore it is now brought to common terms, so as to make it equal with the large-difference denominator present on the large base.] Again set up the oblique bu and bring it to common terms with the large hypotenuse, obtaining □□□ as the small hypotenuse brought to common parts. Add the three positions together, obtaining □□ as the altitude–circle difference [Note: lodge it on the left]. Afterwards set up the celestial unknown, the large difference, and bring it to common terms with the large-hypotenuse denominator, obtaining □○□ as the equal quantity; eliminate it against the left, obtaining □○□□. Extract the cube root, obtaining 360 bu, which is the altitude–circle difference. Subtract the altitude–circle difference from Jia's southward-walk bu: this is the city diameter. This agrees with the question.\end{uniten}

\begin{unitzh}\unitlabel{ceyuan-015.22}或問：東門外不知步數有樹。甲從乾南行六百步而止。乙出北門東行，斜望樹及甲，與城㕘相直，却就樹斜行一百三十六步。問答同前。\end{unitzh}
\begin{uniten}Someone asks: Outside the east gate, at an unknown number of bu, there is a tree. Jia walks 600 bu south from Qian and stops. Yi goes out of the north gate and walks east; looking obliquely, the tree and Jia are jointly in line with the city; then he walks obliquely 136 bu toward the tree. The question and answer are the same as before.\end{uniten}

\begin{unitzh}\unitlabel{ceyuan-015.23}法曰：二行步相乘，於上；又半甲南行乘之，為實。二行相乗，於上；又半甲南行，以乘甲南行，加上位，為益從。甲南行，為從亷。一步，益隅。開立方，得半徑。\end{unitzh}
\begin{uniten}Rule: Multiply the bu of the two walks together, placed above; further multiply it by half of Jia's southward walk, as the dividend. Multiply the two walks together, placed above; further multiply Jia's southward walk by half of Jia's southward walk and add it to the upper position, as the augmenting following term. Jia's southward walk as the following edge. 1 bu as the augmenting corner. Extract the cube root, obtaining the radius.\end{uniten}

\begin{unitzh}\unitlabel{ceyuan-015.24}草曰：立天元一為半徑，便以為小股，其斜行步即小弦也。乃以甲南行為大股，以小弦乘之，復以天元除之，得□□，即大弦也。又倍天元，減甲南行，餘□□，為大差；以減大弦，餘□□□，為大勾也。又倍天元以減勾，得□□□，為小差也。却以半大差□□乘之，得□□□，為半徑冪{\zhsmall 寄左}。乃以天元冪相消，得下式丨□□□。開立方，得一百二十步，即半徑。合問。\end{unitzh}
\begin{uniten}Working: Set up the celestial unknown as the radius, and simply take it as the small altitude; the oblique-walk bu are the small hypotenuse. Then take Jia's southward walk as the large altitude, multiply it by the small hypotenuse, and then divide by the celestial unknown, obtaining □□, which is the large hypotenuse. Again double the celestial unknown and subtract it from Jia's southward walk; the remainder □□ is the large difference; subtract it from the large hypotenuse; the remainder □□□ is the large base. Again double the celestial unknown and subtract it from the base, obtaining □□□ as the small difference. Then multiply it by half the large difference □□, obtaining □□□ as the area of the radius [Note: lodge it on the left]. Then eliminate with the area of the celestial unknown, obtaining the following expression 丨□□□. Extract the cube root, obtaining 120 bu, which is the radius. This agrees with the question.\end{uniten}

\begin{unitzh}\unitlabel{ceyuan-015.25}或問：南門外不知步數有槐樹一株，東門外不知步數有柳樹一株，槐、柳二樹相去二百八十九步。有人從乾南行六百步而止，斜望槐、柳，與城㕘相直。問答同前。\end{unitzh}
\begin{uniten}Someone asks: Outside the south gate, at an unknown number of bu, there is one locust tree; outside the east gate, at an unknown number of bu, there is one willow tree; the two trees, the locust and the willow, are 289 bu apart. A person walks 600 bu south from Qian and stops; looking obliquely, the locust and the willow are jointly in line with the city. The question and answer are the same as before.\end{uniten}

\begin{unitzh}\unitlabel{ceyuan-015.26}法曰：云數相乘得，又自增乗，為三乗方實。斜步冪乘南行步，又云之，為益從。二云數相乘，又倍之{\zhsmall 按：此下脫「内減斜步冪」五字}，為益亷。二之斜步，為第二從亷。二法，常法。得槐至城心步。\end{unitzh}
\begin{uniten}Rule: Multiply the stated numbers (yun shu) together; what is obtained is further multiplied by itself (zi zeng cheng), as the dividend of the thrice-multiplied power. Multiply the southward-walk bu by the area of the oblique bu, and further 云 (yun) it [Translator's note: 云 "say" is evidently corrupt here, perhaps for 二, "double"], as the augmenting following term. Multiply the two stated numbers together and further double it [Siku editors: Note: below this the five characters 内減斜步冪, "subtract from it the area of the oblique bu", are missing], as the augmenting edge. 2 times the oblique bu as the second following edge. 2 fa as the constant divisor [Translator's note: 二法 is perhaps an error for 二步, "2 bu"]. One obtains the bu from the locust tree to the city center.\end{uniten}

\begin{unitzh}\unitlabel{ceyuan-016.1}草曰：别得槐樹至城心步，即人所止至槐樹步也。乃立天元一為槐樹至城心步{\zhsmall 即人至槐處}，加於斜步，得□□，為邊弦也。以天元乗之，得丨□，合斜步除，不除，便以此為邊股{\zhsmall 寄斜步分母}。又以斜步乗南行步，得□，為大股；以邊股減之，餘□□□，為半城徑{\zhsmall 寄斜步分母}。以自之，得丨□□□□，為半徑冪{\zhsmall 内帶斜步為母}，寄左。又以天元減斜歩，得□□，為□弦；以天元乘之，得□□○，合斜步除，不除，寄為母，便以此為半梯頭；以邊股半梯底乗之，得□○□□，為同數。與左相消，得□□□□□，開三乘方，得二百五十五步，即槐樹至城心之步也，亦為皇極正股。又自之，得數以減斜冪，餘如平方而一，得城心至柳樹步，又為皇極正勾也。勾股相乘，倍之為實，如斜步而一，即城徑也。合問。\end{unitzh}
\begin{uniten}Working: By discernment (biede) one finds that the distance from the locust tree to the center of the city is the distance from where the person stopped to the locust tree. Then set up the celestial unknown as the distance from the locust tree to the center of the city [Note: that is, from the person to where the locust tree is]. Add it to the oblique distance (xie bu), obtaining □□ as the border hypotenuse (bian xian). [Translator's note: here and below, □ marks blanks in the transcription, mostly array entries; 丨 and ○ are rod numerals.] Multiply it by the celestial unknown, obtaining 丨□. It should be divided by the oblique distance, but is not divided; take this as the border altitude (bian gu) [Note: the oblique distance is kept as denominator (fenmu)]. Again multiply the southward walking distance by the oblique distance, obtaining □ as the great altitude (da gu). Subtract the border altitude from it; the remainder, □□□, is half the city diameter [Note: the oblique distance is kept as denominator]. Multiply it by itself, obtaining 丨□□□□ as the area (mi) of the radius [Note: within it the oblique distance is carried as denominator]; set it aside on the left. Again subtract the celestial unknown from the oblique distance, obtaining □□ as the □ hypotenuse. Multiply it by the celestial unknown, obtaining □□○. It should be divided by the oblique distance, but is not divided; keep it as denominator, and take this as the half-trapezoid head (ban ti tou). Multiply it by the border altitude, which is the half-trapezoid bottom (ban ti di), obtaining □○□□ as the equal number (tong shu). Eliminate it against the left (xiangxiao), obtaining □□□□□. Extract the root of the thrice-multiplied square (san cheng fang), obtaining 255 bu (步), which is the distance from the locust tree to the center of the city; it is also the proper altitude of the Huangji (huangji zheng gu). Again multiply it by itself, and subtract the number obtained from the oblique area (xie mi). Divide the remainder by the plane square (pingfang) [Translator's note: i.e., extract its square root], obtaining the distance from the center of the city to the willow tree, which is also the proper base of the Huangji (huangji zheng gou). Multiply the base and the altitude together and double it as the dividend (shi). Divide by the oblique distance: this is the city diameter. This matches the question.\end{uniten}

\begin{unitzh}\unitlabel{ceyuan-016.2}或問：甲從乾南行六百步而立，乙出南門直行，丙出東門直行，三人相望，俱與城相直，而乙丙共行了一百五十一步。問答同前。\end{unitzh}
\begin{uniten}Someone asks: Jia walks south from Qian (乾) 600 bu and stands; Yi goes out of the south gate and walks straight on; Bing goes out of the east gate and walks straight on. The three persons look at one another, and all are aligned with the city; and Yi and Bing together walked 151 bu. The question and the answer are the same as before.\end{uniten}

\begin{unitzh}\unitlabel{ceyuan-016.3}法曰：甲南行為冪，折半，又以自之，為實；倍共步，加甲南行，以乘半段甲行冪，為從方；甲行乘共數，為從亷；一个半甲南行，為第二益亷；二分五釐，為三乘方隅。\end{unitzh}
\begin{uniten}Rule: Make Jia's southward walk into an area (mi), halve it, and again multiply it by itself, as the dividend (shi). Double the combined distance (gong bu), add Jia's southward walk, and with this multiply the half portion of the area of Jia's walk, as the following square (congfang). Jia's walk multiplied by the combined number, as the following edge (conglian). 1 1/2 times Jia's southward walk, as the second increasing edge (yilian). 2 fen (分) 5 li (釐), as the corner (yu) of the thrice-multiplied square (san cheng fang).\end{uniten}

\begin{unitzh}\unitlabel{ceyuan-016.4}草曰：識别得共步加城徑，即皇極和也，又是半徑為勾之弦與半徑為股之弦相和步也。二之此數，内減去大弦，即皇極勾股内黄方面也，亦為太虚弦。乃立天元一為大差，以自之，副置二位：上位減於甲行冪，以天元除之，又折半，得□□□，為大勾也；下位加甲南行，以天元除之，又折半，得□□□，為大弦也。其甲南行即大股也。併勾大股，得下式□□□，即大和也。再以天元減甲南行，得□□，即圓徑也；加共步，得□□，即皇極和，又是半徑為勾之弦及半徑為股之弦共數也；又倍之，得□□，即全徑為勾之弦及全徑為股之弦共數也；内減大弦，得□□□，即小和内黄方面也。\end{unitzh}
\begin{uniten}Working: By discernment (shibie) one finds that the combined distance plus the city diameter is the Huangji sum (huangji he); it is also the summed distance of the hypotenuse whose base is the radius and the hypotenuse whose altitude is the radius. Multiply this number by 2 and subtract the great hypotenuse (da xian) from it: this is the side of the yellow square (huangfang mian) within the Huangji base and altitude (gougu); it is also the Taixu hypotenuse (taixu xian). Then set up the celestial unknown as the great difference (da cha). Multiply it by itself and set it up in duplicate in two positions. The upper position: subtract it from the area of Jia's walk, divide by the celestial unknown, and again halve, obtaining □□□ as the great base (da gou). The lower position: add Jia's southward walk, divide by the celestial unknown, and again halve, obtaining □□□ as the great hypotenuse. Jia's southward walk is the great altitude. Add together the base and the great altitude, obtaining the following expression (xia shi) □□□, which is the great sum (da he). Again subtract the celestial unknown from Jia's southward walk, obtaining □□, which is the circle diameter. Add the combined distance, obtaining □□, which is the Huangji sum, and is also the combined number of the hypotenuse whose base is the radius and the hypotenuse whose altitude is the radius. Again double it, obtaining □□, which is the combined number of the hypotenuse whose base is the full diameter and the hypotenuse whose altitude is the full diameter. Subtract the great hypotenuse from it, obtaining □□□, which is the side of the yellow square within the small sum (xiao he).\end{uniten}

\begin{unitzh}\unitlabel{ceyuan-016.5}乃置大和□□□，以小黄方面乘之，得□□□□□，合以小和除之，不除，便以此為大黄方也{\zhsmall 内寄小和為母}，寄左。然後以天元減甲南行，得□□，為大黄方；以小和乗之，得丨□□，為同數；與左相消，得□□□□□，開三乗方，得三百六十步，即股圓差也。以股圓差減於甲南行，餘二百四十步，即城徑也。合問。\end{unitzh}
\begin{uniten}Then set up the great sum □□□ and multiply it by the side of the small yellow square, obtaining □□□□□. It should be divided by the small sum, but is not divided; take this as the great yellow square (da huangfang) [Note: within it the small sum is kept as denominator]; set it aside on the left. Afterwards subtract the celestial unknown from Jia's southward walk, obtaining □□ as the great yellow square. Multiply it by the small sum, obtaining 丨□□ as the equal number. Eliminate it against the left, obtaining □□□□□. Extract the root of the thrice-multiplied square, obtaining 360 bu, which is the altitude–circle difference (gu yuan cha). Subtract the altitude–circle difference from Jia's southward walk; the remainder, 240 bu, is the city diameter. This matches the question.\end{uniten}

\begin{unitzh}\unitlabel{ceyuan-016.6}或問：丙出南門東行，乙出東門南行，各不知步數而立。甲從乾隅南行六百步，斜望乙丙，悉與城叅相直。乙就丙斜行一百二步相㑹。問答同前。\end{unitzh}
\begin{uniten}Someone asks: Bing goes out of the south gate and walks east, and Yi goes out of the east gate and walks south; the number of bu of each is unknown, and they stand. Jia walks south 600 bu from the Qian corner; looking obliquely, Yi and Bing are all aligned together with the city. Yi walks obliquely towards Bing 102 bu and they meet. The question and the answer are the same as before.\end{uniten}

\begin{unitzh}\unitlabel{ceyuan-016.7}法曰：以斜步乘甲南行冪，又倍之，為實；倍甲行冪於上，又以斜步乘二之甲南行，加於上，為從方；四之甲南行，為益亷；四步，常法。開立方，得半徑。\end{unitzh}
\begin{uniten}Rule: Multiply the area of Jia's southward walk by the oblique distance and again double it, as the dividend (shi). Double the area of Jia's walk at the top; again multiply 2 times Jia's southward walk by the oblique distance and add it to the top, as the following square (congfang). 4 times Jia's southward walk, as the increasing edge (yilian). 4 bu, the constant divisor (changfa). Extract the cube root, obtaining the radius.\end{uniten}

\begin{unitzh}\unitlabel{ceyuan-016.8}草曰：别得斜步為小弦也，以斜步減圓徑，餘為小和也。乃立天元為半徑，以二之減於甲南行，得□□，為大差也；以自之，得□□□，為大差冪也。置甲南行冪□，内加大差冪而半之，得□□□，為大弦也{\zhsmall 内帶大差為分母}；又置甲南行冪，内減大差冪而半之，得□□○，為大勾也{\zhsmall 帶大差分母}。又以大差乘股六百步，得□□，併入大勾，得□□□，為大和也{\zhsmall 帶大差分母}。乃先以小弦乘大和，得下式□□□，寄左。又以倍天元減斜步，得□□，為小和；以乘大弦，得□□□□，為同數；與左相消，得□□□□，開立方，得一百二十步，即半徑也。合問。\end{unitzh}
\begin{uniten}Working: By discernment one finds that the oblique distance is the small hypotenuse (xiao xian); subtract the oblique distance from the circle diameter, and the remainder is the small sum. Then set up the celestial unknown as the radius. Multiply it by 2 and subtract it from Jia's southward walk, obtaining □□ as the great difference. Multiply it by itself, obtaining □□□ as the area of the great difference. Set up the area of Jia's southward walk □, add the area of the great difference to it and halve, obtaining □□□ as the great hypotenuse [Note: within it the great difference is carried as denominator]. Again set up the area of Jia's southward walk, subtract the area of the great difference from it and halve, obtaining □□○ as the great base [Note: carrying the great-difference denominator]. Again multiply the altitude, 600 bu, by the great difference, obtaining □□; add it into the great base, obtaining □□□ as the great sum [Note: carrying the great-difference denominator]. Then first multiply the great sum by the small hypotenuse, obtaining the following expression □□□; set it aside on the left. Again subtract double the celestial unknown from the oblique distance, obtaining □□ as the small sum. With it multiply the great hypotenuse, obtaining □□□□ as the equal number. Eliminate it against the left, obtaining □□□□. Extract the cube root, obtaining 120 bu, which is the radius. This matches the question.\end{uniten}

\begin{unitzh}\unitlabel{ceyuan-016.9}依前問，假令乙出東門南行，丙出南門東行，各不知步數而立{\zhsmall 只云丙行步多於乙行步}。甲從乾隅南行六百步，望乙丙與城叅相直。乙復斜行就丙，行了一百二步，與丙相㑹。問答同前。\end{unitzh}
\begin{uniten}Following the previous question, suppose that Yi goes out of the east gate and walks south, and Bing goes out of the south gate and walks east; the number of bu of each is unknown, and they stand [Note: it only says that Bing walks more bu than Yi]. Jia walks south 600 bu from the Qian corner; looking, Yi and Bing are aligned together with the city. Yi then walks obliquely towards Bing; he walks 102 bu and meets Bing. The question and the answer are the same as before.\end{uniten}

\begin{unitzh}\unitlabel{ceyuan-016.10}法曰：以斜步乘甲行冪，又倍之，為立方實；甲行冪内加斜行南行相乗數，為從方；甲南行，為益亷；半步，為隅。得全徑。\end{unitzh}
\begin{uniten}Rule: Multiply the area of Jia's walk by the oblique distance and again double it, as the cube dividend (lifang shi). To the area of Jia's walk add the product of the oblique walk and the southward walk, as the following square (congfang). Jia's southward walk, as the increasing edge (yilian). 1/2 bu, as the corner (yu). One obtains the full diameter.\end{uniten}

\begin{unitzh}\unitlabel{ceyuan-016.11}草曰：别得相就步即小弦也，小弦得小和為直徑也。立天元一為城徑，以減於甲南行步，得□□，為大差；以自之，得丨□□，為太差冪也。置甲南行步，以自之為冪，副之：上以加大差冪而半之，得□□□，為大弦也{\zhsmall 内寄大差分母}；下以減大差冪而半之，得□□○，為大勾也{\zhsmall 内寄大差分母}。\end{unitzh}
\begin{uniten}Working: By discernment one finds that the approach distance (xiangjiu bu) is the small hypotenuse; the small hypotenuse joined with the small sum makes the diameter. Set up the celestial unknown as the city diameter. Subtract it from Jia's southward walking distance, obtaining □□ as the great difference. Multiply it by itself, obtaining 丨□□ as the area of the tai difference (tai cha) [Translator's note: evidently for the great difference, da cha]. Set up Jia's southward walking distance, multiply it by itself as an area, and duplicate it. To the upper one add the area of the great difference and halve, obtaining □□□ as the great hypotenuse [Note: within it the great-difference denominator is kept]. From the lower one subtract the area of the great difference and halve, obtaining □□○ as the great base [Note: within it the great-difference denominator is kept].\end{uniten}

\begin{unitzh}\unitlabel{ceyuan-016.12}乃置相就步在地，以大勾乗之，得□□，合大弦除，不除，寄為母，便以此為小勾也，寄大弦母。又置斜步{\zhsmall 即相就步也}，以甲南行乘之，得□，合以大弦除之，不除，寄為母，便以此為小股；而又以元分母大差乗之，得□□，為同分小股也，只寄大弦為母{\zhsmall 其大勾内元有大差分母，其大股内却無分母，故今乘過，復以大差通之，齊分母也}。又置斜行步，以大弦通之，得□□□，為小弦也。上三位相併，得□□，為城徑也{\zhsmall 内寄大弦分母}，寄左。然後置天元，以大弦通之，得□□□，為同數；與左相消，得□□□□，開立方，得二百四十步，即城徑也。合問。\end{unitzh}
\begin{uniten}Then set up the approach distance on the ground (zai di) and multiply it by the great base, obtaining □□. It should be divided by the great hypotenuse, but is not divided; keep it as denominator, and take this as the small base (xiao gou), keeping the great-hypotenuse denominator. Again set up the oblique distance [Note: that is, the approach distance] and multiply it by Jia's southward walk, obtaining □. It should be divided by the great hypotenuse, but is not divided; keep it as denominator, and take this as the small altitude (xiao gu). Further multiply it by the original denominator, the great difference, obtaining □□ as the small altitude with the same denominator (tongfen); only the great hypotenuse is kept as denominator [Note: within the great base there was originally the great-difference denominator, whereas within the great altitude there is no denominator; therefore now, after the multiplication, it is again brought into common terms (tong) by the great difference, to equalize the denominators]. Again set up the oblique walking distance and bring it into common terms with the great hypotenuse, obtaining □□□ as the small hypotenuse. Add together the three positions above, obtaining □□ as the city diameter [Note: within it the great-hypotenuse denominator is kept]; set it aside on the left. Afterwards set up the celestial unknown and bring it into common terms with the great hypotenuse, obtaining □□□ as the equal number. Eliminate it against the left, obtaining □□□□. Extract the cube root, obtaining 240 bu, which is the city diameter. This matches the question.\end{uniten}

\begin{unitzh}\unitlabel{ceyuan-016.13}測圓海鏡卷五\end{unitzh}
\begin{uniten}Sea Mirror of Circle Measurements, juan 5\end{uniten}

\section*{卷06}\addcontentsline{toc}{section}{卷06}

\begin{unitzh}\unitlabel{ceyuan-017.1}欽定四庫全書\end{unitzh}
\begin{uniten}Imperially Commissioned Complete Library of the Four Treasuries (Qinding Siku quanshu)\end{uniten}

\begin{unitzh}\unitlabel{ceyuan-017.2}測圓海鏡卷六\end{unitzh}
\begin{uniten}Sea Mirror of Circle Measurements (Ceyuan haijing), juan 6\end{uniten}

\begin{unitzh}\unitlabel{ceyuan-017.3}元 李冶 撰\end{unitzh}
\begin{uniten}Composed by Li Ye of the Yuan\end{uniten}

\begin{unitzh}\unitlabel{ceyuan-017.4}大勾一十八問\end{unitzh}
\begin{uniten}Eighteen problems on the great base (da gou)\end{uniten}

\begin{unitzh}\unitlabel{ceyuan-017.5}或問：乙從東門直行一十六步，甲從乾隅東行三百二十步，望乙與城叅相直。問答同前。\end{unitzh}
\begin{uniten}Someone asks: Yi goes straight from the east gate 16 bu (步). Jia goes east from the Qian corner [Translator's note: Qian 乾 is the northwest; Kun 坤, below, is the southwest.] 320 bu and looks toward Yi in alignment with the city. The question and the answer are the same as before.\end{uniten}

\begin{unitzh}\unitlabel{ceyuan-017.6}法曰：甲東行内減二之乙南行，復以乘甲東行，為實；四之東行内減二之乙東行，為從；四益隅；得半徑。\end{unitzh}
\begin{uniten}Rule: From Jia's eastward walk subtract twice Yi's southward walk [Translator's note: Yi walks east in this problem, as in the working; 南 "south" appears to be an error for 東 "east".], and again multiply by Jia's eastward walk, as the dividend (shi). From 4 times the eastward walk subtract twice Yi's eastward walk, as the joining (cong). 4 as the augmenting corner (yi yu). One obtains the radius.\end{uniten}

\begin{unitzh}\unitlabel{ceyuan-017.7}草曰：立天元一為半徑，以二之加乙東行，得□□，為中勾。以中勾減於甲東行，得□□，為勾率也。其天元半徑即股率也。置甲東行為大勾，以股率乗之，得□，合以勾率除之，不受除，便以此為大股{\zhsmall 内帶勾率分母}。再置天元，以二之，以勾率乗之，得□□，減於大股，餘□□，為股圓差，於上{\zhsmall 内有勾率分母}。又以二之天元減甲東行，得□□，為小差，以乗上位，得□□□，為半段黄方冪{\zhsmall 内有勾率分母}，寄左。然後以天元自之，又以勾率乘之，又就分倍之，得□□□，為同數，與左相消，得□□□，開平方，得一百二十步，倍之，即城徑也。合問。\end{unitzh}
\begin{uniten}Working: Set up the celestial unknown (tianyuan) as the radius. Double it and add Yi's eastward walk, obtaining □□ [Translator's note: each □ stands for an entry of a counting-rod array not reproduced in the transcription; 丨 and ○ below are the rod numerals 1 and 0.], as the middle base (zhong gou). Subtract the middle base from Jia's eastward walk, obtaining □□, as the base rate (gou lü). The celestial-unknown radius is the altitude rate (gu lü). Set up Jia's eastward walk as the great base and multiply it by the altitude rate, obtaining □. It ought to be divided by the base rate; it does not accept division, so simply take this as the great altitude (da gu) [Note: it carries the base rate within as denominator (fenmu)]. Again set up the celestial unknown, double it, and multiply it by the base rate, obtaining □□. Subtract it from the great altitude; the remainder, □□, is the difference of altitude and circle (gu yuan cha), [placed] above [Note: within there is the base rate as denominator]. Further subtract twice the celestial unknown from Jia's eastward walk, obtaining □□, as the small difference (xiao cha). Multiply the upper position by it, obtaining □□□, as half a piece of the area (mi) of the yellow square (huang fang) [Note: within there is the base rate as denominator]; lodge it on the left. Then multiply the celestial unknown by itself, further multiply it by the base rate, and further, according to the parts (jiu fen), double it, obtaining □□□, as the equal quantity (tong shu). Cancel it against the left (xiang xiao), obtaining □□□. Extracting the square root gives 120 bu; double it: this is the city diameter. This agrees with the question.\end{uniten}

\begin{unitzh}\unitlabel{ceyuan-017.8}或問：乙出東門南行三十步而立，甲從乾隅東行三百二十步，望乙與城叅相直。問答同前。\end{unitzh}
\begin{uniten}Someone asks: Yi goes out of the east gate, walks south 30 bu and stands. Jia goes east from the Qian corner 320 bu and looks toward Yi in alignment with the city. The question and the answer are the same as before.\end{uniten}

\begin{unitzh}\unitlabel{ceyuan-017.9}法曰：甲乙相乘為實，甲東行為從，二虚法，得半徑。\end{unitzh}
\begin{uniten}Rule: Multiply Jia and Yi together as the dividend (shi); Jia's eastward walk as the joining (cong); 2 as the empty divisor (xu fa). One obtains the radius.\end{uniten}

\begin{unitzh}\unitlabel{ceyuan-017.10}草曰：識别具見大股第二問中。立天元為半徑，内減乙南行，得□□，為虚股，以乘通勾甲東行，得□□，為半段城徑冪{\zhsmall 寄左}。然後以天元自之，又就分二之，得□□，為同數，與左相消，得□□□，開平方，得一百二十步，倍之，即城徑也。合問。\end{unitzh}
\begin{uniten}Working: The identification (shibie) is fully shown in the second problem on the great altitude. Set up the celestial unknown as the radius. From it subtract Yi's southward walk, obtaining □□, as the empty altitude (xu gu). With it multiply the through base (tong gou), Jia's eastward walk, obtaining □□, as half a piece of the area of the city diameter [Note: lodge it on the left]. Then multiply the celestial unknown by itself and further, according to the parts, double it, obtaining □□, as the equal quantity. Cancel it against the left, obtaining □□□. Extracting the square root gives 120 bu; double it: this is the city diameter. This agrees with the question.\end{uniten}

\begin{unitzh}\unitlabel{ceyuan-017.11}或問：乙出南門直行一百三十五步而立，甲從乾隅東行三百二十步，望見乙。問答同前。\end{unitzh}
\begin{uniten}Someone asks: Yi goes out of the south gate, walks straight 135 bu and stands. Jia goes east from the Qian corner 320 bu and catches sight of Yi. The question and the answer are the same as before.\end{uniten}

\begin{unitzh}\unitlabel{ceyuan-017.12}法曰：以乙南行乘甲東行冪，為實；二之乙南行乘甲東行，為從方；亷空；二步常法；得半徑。\end{unitzh}
\begin{uniten}Rule: Multiply the area of Jia's eastward walk by Yi's southward walk, as the dividend (shi). Twice Yi's southward walk multiplied by Jia's eastward walk, as the joining square (cong fang). The edge (lian) is empty. 2 bu as the constant divisor (chang fa). One obtains the radius.\end{uniten}

\begin{unitzh}\unitlabel{ceyuan-017.13}草曰：立天元一為半城徑，以二之，加於乙南行，得□□，為股率；以天元減甲東行，得□□，為勾率。乃置乙南行，以勾率乘之，得□□，合股率除，不除，便以此為小勾，此即半梯之頭{\zhsmall 内帶股率分母}。又以勾率乘之，得□□□，為半徑冪{\zhsmall 内𢃄股率分母}，寄左。乃以股率乘天元冪，得□□□，為同數，與左相消，得□○□□，開立方，得一百二十步，倍之，即城徑也。合問。\end{unitzh}
\begin{uniten}Working: Set up the celestial unknown as half the city diameter. Double it and add it to Yi's southward walk, obtaining □□, as the altitude rate. Subtract the celestial unknown from Jia's eastward walk, obtaining □□, as the base rate. Then set up Yi's southward walk and multiply it by the base rate, obtaining □□. It ought to be divided by the altitude rate; not dividing, simply take this as the small base (xiao gou): this is the head of the half trapezoid (ti) [Note: it carries the altitude rate within as denominator]. Again multiply it by the base rate, obtaining □□□, as the area of the radius [Note: it carries the altitude rate within as denominator]; lodge it on the left. Then multiply the area of the celestial unknown by the altitude rate, obtaining □□□, as the equal quantity. Cancel it against the left, obtaining □○□□. Extracting the cube root gives 120 bu; double it: this is the city diameter. This agrees with the question.\end{uniten}

\begin{unitzh}\unitlabel{ceyuan-017.14}或問：乙出南門東行七十二步，甲從西北隅取直行三百二十步，見乙。問答同前。\end{unitzh}
\begin{uniten}Someone asks: Yi goes out of the south gate and walks east 72 bu. Jia, from the northwest corner, takes a straight walk of 320 bu and sees Yi. The question and the answer are the same as before.\end{uniten}

\begin{unitzh}\unitlabel{ceyuan-017.15}法曰：二行相乘為實，以東行為從，一步常法，得半徑。\end{unitzh}
\begin{uniten}Rule: Multiply the two walks together as the dividend (shi); take the eastward walk as the joining (cong); 1 bu as the constant divisor. One obtains the radius.\end{uniten}

\begin{unitzh}\unitlabel{ceyuan-017.16}草曰：立天元一為半城徑，以減甲東行步，得□□，為梯底；以乙東行七十二步為梯頭，以乘之，得□□，為半徑冪{\zhsmall 寄左}。然後以天元冪與左相消，得丨□□，以平方開之，得一百二十步，倍之，即城徑也。合問。\end{unitzh}
\begin{uniten}Working: Set up the celestial unknown as half the city diameter. Subtract it from Jia's eastward walk in bu, obtaining □□, as the bottom of the trapezoid (ti di). Take Yi's eastward walk of 72 bu as the head of the trapezoid (ti tou) and multiply by it, obtaining □□, as the area of the radius [Note: lodge it on the left]. Then cancel the area of the celestial unknown against the left, obtaining 丨□□. Extracting its square root gives 120 bu; double it: this is the city diameter. This agrees with the question.\end{uniten}

\begin{unitzh}\unitlabel{ceyuan-017.17}或問：乙從西南隅直東行一百九十二步，甲從西北隅直東行三百二十步，望見乙。問答同前。\end{unitzh}
\begin{uniten}Someone asks: Yi goes straight east from the southwest corner 192 bu. Jia goes straight east from the northwest corner 320 bu and catches sight of Yi. The question and the answer are the same as before.\end{uniten}

\begin{unitzh}\unitlabel{ceyuan-017.18}法曰：二行步相乘為實，二行相併為法，得半徑。\end{unitzh}
\begin{uniten}Rule: Multiply the bu of the two walks together as the dividend (shi); add the two walks together as the divisor (fa). One obtains the radius.\end{uniten}

\begin{unitzh}\unitlabel{ceyuan-017.19}草曰：立天元一為半徑，副置之。上以減於乙東行，得□□，為梯頭，於上；下位減於甲東行，得□□，為梯底；以乘上位，得丨□□，為半徑冪{\zhsmall 寄左}。然後以天元冪與左相消，得□□，上法下實，即半徑也。合問。\end{unitzh}
\begin{uniten}Working: Set up the celestial unknown as the radius and set it out in duplicate (fu zhi). The upper one: subtract it from Yi's eastward walk, obtaining □□, as the head of the trapezoid, [placed] above. The lower position: subtract it from Jia's eastward walk, obtaining □□, as the bottom of the trapezoid. Multiply the upper position by it, obtaining 丨□□, as the area of the radius [Note: lodge it on the left]. Then cancel the area of the celestial unknown against the left, obtaining □□: the upper is the divisor (fa), the lower is the dividend (shi); this is the radius. This agrees with the question.\end{uniten}

\begin{unitzh}\unitlabel{ceyuan-017.20}或問：乙從坤隅直南行三百六十步而止，甲從乾隅直東行三百二十步，望見乙。問答同前。\end{unitzh}
\begin{uniten}Someone asks: Yi goes straight south from the Kun corner 360 bu and stops. Jia goes straight east from the Qian corner 320 bu and catches sight of Yi. The question and the answer are the same as before.\end{uniten}

\begin{unitzh}\unitlabel{ceyuan-017.21}法曰：二行步相乗，倍之，為實；二之甲東行，為從；一步常法；得城徑。\end{unitzh}
\begin{uniten}Rule: Multiply the bu of the two walks together and double it, as the dividend (shi). Twice Jia's eastward walk, as the joining (cong). 1 bu as the constant divisor. One obtains the city diameter.\end{uniten}

\begin{unitzh}\unitlabel{ceyuan-017.22}草曰：立天元一以為城徑，加一南行，得□□，為股；二行步相併，得六百八十步，為弦；甲東行為勾。勾股相乘，得□□，又倍之，得□□，為二直積{\zhsmall 寄左}。然後以勾股弦相併，得□□，為三事和，以天元乘之，得丨□，為同數，與左相消，得丨□□，開平方，得二百四十步，即城徑也。合問。\end{unitzh}
\begin{uniten}Working: Set up the celestial unknown to be the city diameter. Add the one southward walk, obtaining □□, as the altitude (gu). Add the bu of the two walks together, obtaining 680 bu, as the hypotenuse (xian). Jia's eastward walk is the base (gou). Multiply base and altitude together, obtaining □□; further double it, obtaining □□, as two rectangular areas (zhi ji) [Note: lodge it on the left]. Then add base, altitude and hypotenuse together, obtaining □□, as the sum of the three things (san shi he). Multiply it by the celestial unknown, obtaining 丨□, as the equal quantity. Cancel it against the left, obtaining 丨□□. Extracting the square root gives 240 bu: this is the city diameter. This agrees with the question.\end{uniten}

\begin{unitzh}\unitlabel{ceyuan-017.23}或問：東門南不知逺近有樹，甲從乾隅東行三百二十步，望樹與城叅相直，復就樹斜行一百七十步至樹。問答同前。\end{unitzh}
\begin{uniten}Someone asks: South of the east gate, at a distance not known, there is a tree. Jia goes east from the Qian corner 320 bu and looks toward the tree in alignment with the city. Then, going toward the tree, he walks obliquely 170 bu and reaches the tree. The question and the answer are the same as before.\end{uniten}

\begin{unitzh}\unitlabel{ceyuan-017.24}法曰：兩段東行步冪内減兩段東行斜行相乗數，為實{\zhsmall 按：或云倍東行步以二行差東之，亦同}；二之東行，為從；一益隅；得城徑。\end{unitzh}
\begin{uniten}Rule: From two pieces of the area of the eastward walk in bu subtract two pieces of the product of the eastward walk and the oblique walk, as the dividend (shi) [Note: Remark: some say: double the eastward walk in bu and "east" (東) it by the difference of the two walks; this is also the same.] [Translator's note: 東 "east" here is evidently an error for 乘 "multiply".]. Twice the eastward walk, as the joining (cong). 1 as the augmenting corner. One obtains the city diameter.\end{uniten}

\begin{unitzh}\unitlabel{ceyuan-017.25}草曰：别得東行步即大勾，斜行步即小弦也。乃立天元一為城徑，減東行步，得□□，為勾圓差也{\zhsmall 今為小勾}。置東行步，以斜步乘之，得□，合以小勾除之，今不受除，便以此為大弦{\zhsmall 内帶小勾分母}。再置東行步，以小勾乘之，得□□，為大勾；以減大弦，得□□，為大差，合以小差乗之{\zhsmall 縁内帶小差分母}，更不湏乗，便以此為半段黄方冪{\zhsmall 更無分母}。又二之，得□□，為一段黄方冪{\zhsmall 寄左}。然後以天元冪與左相消，得□□□，開平方，得二百四十步，即城徑也。合問。\end{unitzh}
\begin{uniten}Working: By discernment one obtains that the eastward walk in bu is the great base and the oblique walk in bu is the small hypotenuse (xiao xian). Then set up the celestial unknown as the city diameter. Subtract it from the eastward walk in bu, obtaining □□, as the difference of base and circle (gou yuan cha) [Note: now it is the small base]. Set up the eastward walk in bu and multiply it by the oblique walk in bu, obtaining □. It ought to be divided by the small base; now it does not accept division, so simply take this as the great hypotenuse (da xian) [Note: it carries the small base within as denominator]. Again set up the eastward walk in bu and multiply it by the small base, obtaining □□, as the great base. Subtract it from the great hypotenuse, obtaining □□, as the great difference (da cha). It ought to be multiplied by the small difference [Note: because it carries the small difference within as denominator]; there is no further need to multiply: simply take this as half a piece of the area of the yellow square [Note: there is no longer any denominator]. Further double it, obtaining □□, as one piece of the area of the yellow square [Note: lodge it on the left]. Then cancel the area of the celestial unknown against the left, obtaining □□□. Extracting the square root gives 240 bu: this is the city diameter. This agrees with the question.\end{uniten}

\begin{unitzh}\unitlabel{ceyuan-017.26}依前問，假令乙出東門南行，不知步數而止，甲從乾東行三百二十步，望乙與城相直，復就乙斜行一百七十步。\end{unitzh}
\begin{uniten}Following the previous question: suppose Yi goes out of the east gate and walks south, the number of bu not known, and stops. Jia goes east from Qian 320 bu and looks toward Yi in alignment with the city. Then, going toward Yi, he walks obliquely 170 bu.\end{uniten}

\begin{unitzh}\unitlabel{ceyuan-017.27}法曰：以甲東行乘二行差冪，為實；以甲東行乘二之二行差，為從方；二之二行差，為隅法；得半徑。\end{unitzh}
\begin{uniten}Rule: Multiply the area of the difference of the two walks by Jia's eastward walk, as the dividend (shi). Multiply twice the difference of the two walks by Jia's eastward walk, as the joining square. Twice the difference of the two walks, as the corner divisor (yu fa). One obtains the radius.\end{uniten}

\begin{unitzh}\unitlabel{ceyuan-017.28}草曰：識别得二行相減餘一百五十，即半城徑與乙南行共數也。得此數，更不湏用斜。立天元一為半徑，減於二行差，得□□，即半梯頭也。又以二天元減甲東行步，得□□，為勾率；又以一百五十為股率。乃置甲東行，以股率乘之，得□，合勾率除，不除，便以此為大股{\zhsmall 内寄勾率分母}。再置天元，以勾率乘之，得□□，以減於大股，得□□□，為半梯底也。頭底相乘，得下□□□□，為半徑冪也{\zhsmall 内帶勾率分母}，寄左。然後以勾率乘天元冪，得□□□，為同數，與左相消，得□□□，開平方，得一百二十步，倍之，即城徑也。合問。\end{unitzh}
\begin{uniten}Working: By identification one obtains that the remainder of subtracting the two walks from each other, 150, is the combined number of half the city diameter and Yi's southward walk. Having obtained this number, there is no further need to use the oblique [walk]. Set up the celestial unknown as the radius. Subtract it from the difference of the two walks, obtaining □□: this is the head of the half trapezoid. Further subtract 2 celestial unknowns from Jia's eastward walk in bu, obtaining □□, as the base rate; further take 150 as the altitude rate. Then set up Jia's eastward walk and multiply it by the altitude rate, obtaining □. It ought to be divided by the base rate; not dividing, simply take this as the great altitude [Note: the base rate is lodged within as denominator]. Again set up the celestial unknown and multiply it by the base rate, obtaining □□. Subtract it from the great altitude, obtaining □□□, as the bottom of the half trapezoid. Multiply head and bottom together, obtaining, below, □□□□, as the area of the radius [Note: it carries the base rate within as denominator]; lodge it on the left. Then multiply the area of the celestial unknown by the base rate, obtaining □□□, as the equal quantity. Cancel it against the left, obtaining □□□. Extracting the square root gives 120 bu; double it: this is the city diameter. This agrees with the question.\end{uniten}

\begin{unitzh}\unitlabel{ceyuan-017.29}或問：南門東不知逺近有樹，甲從乾隅東行三百二十步見樹，復向樹斜行二百七十二步至樹。問答同前。\end{unitzh}
\begin{uniten}Someone asks: East of the south gate, at a distance not known, there is a tree. Jia goes east from the Qian corner 320 bu and sees the tree. Then, facing the tree, he walks obliquely 272 bu and reaches the tree. The question and the answer are the same as before.\end{uniten}

\begin{unitzh}\unitlabel{ceyuan-017.30}法曰：二之二行差乘二之甲東行，為實；併二之二行差及二之甲東行，為從；二步益隅；得城徑。\end{unitzh}
\begin{uniten}Rule: Multiply twice the difference of the two walks by twice Jia's eastward walk, as the dividend (shi). Add together twice the difference of the two walks and twice Jia's eastward walk, as the joining (cong). 2 bu as the augmenting corner. One obtains the city diameter.\end{uniten}

\begin{unitzh}\unitlabel{ceyuan-017.31}草曰：别得二行相減餘四十八步，即虚積之勾也。立天元一為城徑，内減二之二行差，得□□，為梯頭，於上。置甲東行步，以二之，内減天元，得□□，為梯底，以乘上位，得□□□，為城徑冪{\zhsmall 寄左}。然後以天元冪與左相消，得□□□，開平方，得二百四十步，即城徑也。合問。\end{unitzh}
\begin{uniten}Working: By discernment one obtains that the remainder of subtracting the two walks from each other, 48 bu, is the base of the empty area (xu ji). Set up the celestial unknown as the city diameter. From it subtract twice the difference of the two walks, obtaining □□, as the head of the trapezoid, [placed] above. Set up Jia's eastward walk in bu, double it, and from it subtract the celestial unknown, obtaining □□, as the bottom of the trapezoid. Multiply the upper position by it, obtaining □□□, as the area of the city diameter [Note: lodge it on the left]. Then cancel the area of the celestial unknown against the left, obtaining □□□. Extracting the square root gives 240 bu: this is the city diameter. This agrees with the question.\end{uniten}

\begin{unitzh}\unitlabel{ceyuan-017.32}或問：甲從乾隅東行三百二十步而止，乙出南門直行，不知步數，望見甲，復就甲斜行四百二十五步，與甲相㑹。問答同前。\end{unitzh}
\begin{uniten}Someone asks: Jia goes east from the Qian corner 320 bu and stops. Yi goes out of the south gate and walks straight, the number of bu not known, and catches sight of Jia. Then, going toward Jia, he walks obliquely 425 bu and meets Jia. The question and the answer are the same as before.\end{uniten}

\begin{unitzh}\unitlabel{ceyuan-017.33}法曰：二行步相減，以乘東行冪，得數半之，為實；以半之東行步乗東行步，於上；二行步相減餘乗東行步，減上位，為從；二之東行步，為益亷；一步常法；得半徑。\end{unitzh}
\begin{uniten}Rule: Subtract the bu of the two walks from each other, multiply the area of the eastward walk by this, and halve the number obtained, as the dividend (shi). Multiply the eastward walk in bu by half the eastward walk in bu, [placed] above. Multiply the eastward walk in bu by the remainder of subtracting the bu of the two walks from each other and subtract it from the upper position, as the joining (cong). Twice the eastward walk in bu, as the augmenting edge (yi lian). 1 bu as the constant divisor. One obtains the radius.\end{uniten}

\begin{unitzh}\unitlabel{ceyuan-017.34}草曰：識别得二行相減，是髙積上勾股較{\zhsmall 此勾即半徑也}；又别得是髙弦不及股圓差數。乃立天元為半城徑，以減東行步，得□□，為中勾；其斜行步即中弦也。又置半城徑，以斜步乗之，得□，合以中勾除之，不受除，便以此為髙弦{\zhsmall 内寄中勾為母}。又以二行步相減餘一百五步，為髙弦不及股圓差數，置此數，以中勾乘之，得□□，加入髙弦，得□□，為大差，於上{\zhsmall 内帶中勾分母}。又倍天元減東行步，得□□，為小差，又半之，得□□，以乘上位，得□□□，為半徑冪{\zhsmall 内有中勾分母}，寄左。乃以天元自乗，又以中勾乘之，得□□□，為同數，與左相消，得□□□□，以立方開，得一百二十步，倍之，即城徑也。合問。\end{unitzh}
\begin{uniten}Working: By identification one obtains that the subtraction of the two walks from each other is the difference of base and altitude (gou gu jiao) on the height area (gao ji) [Note: this base is the radius]; further one discerns that it is the amount by which the height hypotenuse (gao xian) falls short of the difference of altitude and circle. Then set up the celestial unknown as half the city diameter. Subtract it from the eastward walk in bu, obtaining □□, as the middle base; the oblique walk in bu is the middle hypotenuse (zhong xian). Further set up half the city diameter and multiply it by the oblique walk in bu, obtaining □. It ought to be divided by the middle base; it does not accept division, so simply take this as the height hypotenuse [Note: the middle base is lodged within as denominator]. Further take the remainder of subtracting the bu of the two walks from each other, 105 bu, as the amount by which the height hypotenuse falls short of the difference of altitude and circle. Set up this number and multiply it by the middle base, obtaining □□. Add it into the height hypotenuse, obtaining □□, as the great difference, [placed] above [Note: it carries the middle base within as denominator]. Further subtract twice the celestial unknown from the eastward walk in bu, obtaining □□, as the small difference; further halve it, obtaining □□. Multiply the upper position by it, obtaining □□□, as the area of the radius [Note: within there is the middle base as denominator]; lodge it on the left. Then multiply the celestial unknown by itself and further multiply it by the middle base, obtaining □□□, as the equal quantity. Cancel it against the left, obtaining □□□□. Extracting the cube root gives 120 bu; double it: this is the city diameter. This agrees with the question.\end{uniten}

\begin{unitzh}\unitlabel{ceyuan-017.35}或問：甲乙二人俱在乾隅，乙直南行，不知步數而立；甲直東行三百二十步，望見乙，復就乙斜行六百八十步，與乙相㑹。問答同前。\end{unitzh}
\begin{uniten}Someone asks: Two persons, Jia and Yi, are both at the Qian corner. Yi goes straight south, the number of bu not known, and stands. Jia goes straight east 320 bu and catches sight of Yi. Then, going toward Yi, he walks obliquely 680 bu and meets Yi. The question and the answer are the same as before.\end{uniten}

\begin{unitzh}\unitlabel{ceyuan-017.36}法曰：以二行差乘甲東行步，又二之，為實；以二之二行差，為從；一步常法；得城徑。\end{unitzh}
\begin{uniten}Rule: Multiply Jia's eastward walk in bu by the difference of the two walks, and further double it, as the dividend (shi). Take twice the difference of the two walks as the joining (cong). 1 bu as the constant divisor. One obtains the city diameter.\end{uniten}

\begin{unitzh}\unitlabel{ceyuan-017.37}草曰：别得二行步相減餘三百六十步，即股圓差也。乃立天元一為圓徑，以減於甲東行步，得□□，為小差；以東行斜行差三百六十步乘之，得□□；倍之，得□□，為一段城徑冪{\zhsmall 寄左}。乃以天元冪與左相消，得丨□□，開平方，得二百四十步，即城徑也。合問。\end{unitzh}
\begin{uniten}Working: By discernment one obtains that the remainder of subtracting the bu of the two walks from each other, 360 bu, is the difference of altitude and circle. Then set up the celestial unknown as the circle diameter. Subtract it from Jia's eastward walk in bu, obtaining □□, as the small difference. Multiply it by the difference of the eastward walk and the oblique walk, 360 bu, obtaining □□. Double it, obtaining □□, as one piece of the area of the city diameter [Note: lodge it on the left]. Then cancel the area of the celestial unknown against the left, obtaining 丨□□. Extracting the square root gives 240 bu: this is the city diameter. This agrees with the question.\end{uniten}

\begin{unitzh}\unitlabel{ceyuan-017.38}或問：東門外不知逺近有樹，甲從乾隅東行三百二十步，望樹與城叅相直，復就樹斜行一百三十六步至樹。問答同前。\end{unitzh}
\begin{uniten}Someone asks: Outside the east gate, at a distance not known, there is a tree. Jia goes east from the Qian corner 320 bu and looks toward the tree in alignment with the city. Then, going toward the tree, he walks obliquely 136 bu and reaches the tree. The question and the answer are the same as before.\end{uniten}

\begin{unitzh}\unitlabel{ceyuan-017.39}法曰：倍二行相減數，内減甲東行，得數復以乘甲東行，為實{\zhsmall 按：或云倍斜步以減甲東行，餘以甲東行乗之，亦同}；倍二行差，為從；二步虚常法；得半徑。\end{unitzh}
\begin{uniten}Rule: Double the number from subtracting the two walks from each other, subtract Jia's eastward walk from it, and again multiply the number obtained by Jia's eastward walk, as the dividend (shi) [Note: Remark: some say: double the oblique walk in bu and subtract it from Jia's eastward walk, then multiply the remainder by Jia's eastward walk; this is also the same.]. Double the difference of the two walks, as the joining (cong). 2 bu as the empty constant divisor (xu chang fa). One obtains the radius.\end{uniten}

\begin{unitzh}\unitlabel{ceyuan-018.1}草曰：識别得斜行步乃樹至城心步也。立天元一為半徑，加斜行步，得□□，即樹至城西門之步也。乃以減於甲東行，得下□□，為小勾率；其天元半徑即小股率；其斜步即小弦數也。再置甲東行步，内減天元，得□□，為梯底，於上。又置梯底，内減二之小勾率，得□□{\zhsmall 按：倍小勾得三百六十八步少二元，以少二元減梯底之少一元，反為多一元；以三百六十八步減梯底之三百二十步，反為少四十八步也。}，以乘上位，得□□□，為半徑冪。乃以天元冪與左相消，得下式□□□，以平方開之，得一百二十步，倍之即城徑也。合問。\end{unitzh}
\begin{uniten}Working: By the distinctions (shibie) one obtains that the oblique-walk bu (步) are the bu from the tree to the city's center. Set up the celestial unknown as the radius; add the oblique-walk bu, obtaining □□ [Translator's note: □ marks blanks in the transcription, standing for counting-rod numerals or characters; 丨 is the rod numeral 1, ○ is zero.], which is the bu from the tree to the city's west gate. Then subtract this from A's eastward walk, obtaining the following □□, as the small base rate (xiao gou lü); the celestial-unknown radius is the small altitude rate (xiao gu lü); the oblique bu are the small hypotenuse number (xiao xian shu). Again set A's eastward-walk bu, subtract the celestial unknown from it, obtaining □□, as the trapezoid bottom (ti di), placed above. Further set the trapezoid bottom, subtract from it twice the small base rate, obtaining □□ [Siku editors: Doubling the small base gives 368 bu less 2 unknowns (yuan); subtracting the "less 2 unknowns" from the "less 1 unknown" of the trapezoid bottom turns into "more 1 unknown" instead; subtracting 368 bu from the 320 bu of the trapezoid bottom turns into "less 48 bu" instead.]; multiply the upper position by it, obtaining □□□, as the area (mi) of the radius. Then eliminate (xiang xiao) the area of the celestial unknown against the left, obtaining the following form □□□; extract its square root, obtaining 120 bu; doubling it gives the city's diameter. It fits the question.\end{uniten}

\begin{unitzh}\unitlabel{ceyuan-018.2}或問：南門外不知步數有槐一株。甲從乾隅直東行至柳樹下，望見槐樹，復斜行至槐樹下。甲自云：「我共行了七百四十五步。」乙從坤隅南行，望見槐柳與城叅相直，復斜行至槐樹下。乙自云：「我南行步多於斜行步一百五步。」\end{unitzh}
\begin{uniten}Someone asks: Outside the south gate, at an unknown number of bu, there is 1 pagoda tree (huai). A goes from the northwest (qian) corner straight east to below a willow tree, looks and sees the pagoda tree, and then walks obliquely to below the pagoda tree. A says of himself: "I have walked 745 bu altogether." B goes south from the southwest (kun) corner, looks and sees the pagoda tree and the willow in line (can xiang zhi) with the city, and then walks obliquely to below the pagoda tree. B says of himself: "My southward-walk bu exceed my oblique-walk bu by 105 bu."\end{uniten}

\begin{unitzh}\unitlabel{ceyuan-018.3}按：此問下有草無法，今依細草補之。\end{unitzh}
\begin{uniten}Note: Under this question there is a working but no rule; now it is supplied according to the working (xi cao).\end{uniten}

\begin{unitzh}\unitlabel{ceyuan-018.4}法曰：置甲共步，内減乙較步，餘數折半，自之，再倍乙較步乗之，為立方實。置上減餘折半數，又減二之乙較步，復以減餘折半數乗之，為從。甲共步内減乙較步，為亷。五分為負隅。開立方，得城徑。\end{unitzh}
\begin{uniten}Rule: Set A's total bu, subtract B's difference (jiao) bu from it, halve the remaining number, multiply it by itself; again double B's difference bu and multiply by it, as the dividend (shi) for the cube root. Set the above number obtained by halving the remainder of the subtraction, further subtract twice B's difference bu, and again multiply by the number obtained by halving the remainder of the subtraction, as the joining (cong). From A's total bu subtract B's difference bu, as the edge (lian). 5 fen as the negative corner (fu yu) [Translator's note: fen and li here are decimal fractions, 1/10 and 1/100.]. Extract the cube root, obtaining the city's diameter.\end{uniten}

\begin{unitzh}\unitlabel{ceyuan-018.5}草曰：識别得一百五步是大差多於髙弦數，又為髙弦上勾股差數。又别得是甲斜行多於東行數也。乃副置甲共行七百四十五步在地：其上位加一百五步而半之，得四百二十五步，即甲斜行也；其下位減一百五步而半之，得三百二十步，即甲東行也。乃立天元一為圓徑，以半之，減於甲東行步，得□□，為中勾；其甲斜行四百二十五步即中弦也。再置天元，以半之，為小勾，以中弦乘之，得□，合以中勾除，不除，便以為髙弦，於上{\zhsmall 内𢃄中勾分母}。别置乙多步一百五步，以中勾乘之，得□□，為大差多於髙弦數也；以加入上位，得下式□□，為一个大差也。置甲東行，以天元減之，又倍之，得□□，為二个小差；以乗大差，得下□□□，為一段黄方冪{\zhsmall 内帶中股分母}，寄左。然後置天元冪丨□，以中勾通之，得□□□，與左相消，得□□□□，開立方，得二百四十步，即城徑也。合問。\end{unitzh}
\begin{uniten}Working: By the distinctions one obtains that 105 bu is the amount by which the large difference (da cha) exceeds the high hypotenuse (gao xian), and is also the base–altitude difference on the high hypotenuse. Further, by distinction (bie) one obtains that it is the amount by which A's oblique walk exceeds his eastward walk. Then set A's total walk, 745 bu, in duplicate (fu zhi) on the ground: to its upper position add 105 bu and halve, obtaining 425 bu, which is A's oblique walk; from its lower position subtract 105 bu and halve, obtaining 320 bu, which is A's eastward walk. Then set up the celestial unknown as the circle's diameter; halve it and subtract it from A's eastward-walk bu, obtaining □□, as the middle base (zhong gou); A's oblique walk, 425 bu, is the middle hypotenuse (zhong xian). Again set the celestial unknown, halve it, as the small base; multiply it by the middle hypotenuse, obtaining □; it should be divided by the middle base; not dividing, take it directly as the high hypotenuse, placed above [Note: it carries within it the middle-base denominator]. Separately set B's excess bu, 105 bu, multiply it by the middle base, obtaining □□, as the amount by which the large difference exceeds the high hypotenuse; add it into the upper position, obtaining the following form □□, as 1 large difference. Set A's eastward walk, subtract the celestial unknown from it, and further double it, obtaining □□, as 2 small differences (xiao cha); multiply the large difference by it, obtaining the following □□□, as 1 portion of the yellow-square area (huang fang mi) [Note: it carries within it the middle-altitude denominator]; lodge it on the left (ji zuo). Then set the area of the celestial unknown, 丨□, make it commensurable (tong) by the middle base, obtaining □□□; eliminate it against the left, obtaining □□□□; extract the cube root, obtaining 240 bu, which is the city's diameter. It fits the question.\end{uniten}

\begin{unitzh}\unitlabel{ceyuan-018.6}或問：出東門直行，不知步數，有槐樹一株；出南門東行，不知步數，有柳樹一株。槐栁斜相距一百五十三步。甲從乾東行三百二十步，望槐柳與城㕘相直。問答同前。\end{unitzh}
\begin{uniten}Someone asks: Going out of the east gate and walking straight on, at an unknown number of bu, there is 1 pagoda tree; going out of the south gate and walking east, at an unknown number of bu, there is 1 willow tree. The pagoda tree and the willow are 153 bu apart obliquely. A goes east from the northwest 320 bu and looks at the pagoda tree and the willow: they are in line with the city. The question and answer are the same as before.\end{uniten}

\begin{unitzh}\unitlabel{ceyuan-018.7}法曰：二行相乘訖，又以乗甲東行冪，為實。斜行乗甲東行冪，又三之，為從方。甲東行冪内減兩段二行相乘數，為第一亷。二之甲東行，為益二亷。二步常法。開三乗方，得半徑。\end{unitzh}
\begin{uniten}Rule: Having multiplied the two walks together, further multiply the area of A's eastward walk by it, as the dividend. Multiply the area of A's eastward walk by the oblique walk, and further triple it, as the joining square (cong fang). From the area of A's eastward walk subtract 2 portions of the product of the two walks, as the first edge. Twice A's eastward walk as the augmenting (yi) second edge. 2 bu as the constant divisor (chang fa). Extract the root of the thrice-multiplied power (san cheng fang), obtaining the radius.\end{uniten}

\begin{unitzh}\unitlabel{ceyuan-018.8}草曰：立天元一為半徑，以二之，減於甲東行，得□□，為小差。以自之，得□□□，加於甲東行冪，復半之，得□□□，為大弦{\zhsmall 内寄小差分母}。又置斜相距步，以大勾乘之，得□，合大弦除，不除，便以此為小勾{\zhsmall 内𢃄大弦分母}。乃以天元減甲東行數，得□□，為半梯底；以乘小勾半梯頭，得□□，為半徑冪，於上。此半徑冪内有大弦分母，此大弦分母元𢃄小差分母，故先用小差分母以乗上半徑冪，得□□□，為半徑冪也，内𢃄本大弦分母{\zhsmall 寄左}。然後以大弦乘天元冪，得□□□□，為同數，與左相消，得□□□□□，開三乗方，得一百二十步，即半城徑也。合問。\end{unitzh}
\begin{uniten}Working: Set up the celestial unknown as the radius; double it and subtract it from A's eastward walk, obtaining □□, as the small difference. Multiply it by itself, obtaining □□□; add it to the area of A's eastward walk and then halve, obtaining □□□, as the large hypotenuse (da xian) [Note: within it is lodged the small-difference denominator]. Further set the bu of the oblique distance between them, multiply by the large base (da gou), obtaining □; it should be divided by the large hypotenuse; not dividing, take this directly as the small base [Note: it carries within it the large-hypotenuse denominator]. Then subtract the celestial unknown from the number of A's eastward walk, obtaining □□, as the half trapezoid bottom; multiply by it the small base, the half trapezoid top (ban ti tou), obtaining □□, as the area of the radius, placed above. This area of the radius contains the large-hypotenuse denominator; this large-hypotenuse denominator originally carries the small-difference denominator; therefore first use the small-difference denominator to multiply the upper area of the radius, obtaining □□□, as the area of the radius, which carries within it the original large-hypotenuse denominator [Note: lodge it on the left]. Then multiply the area of the celestial unknown by the large hypotenuse, obtaining □□□□, as the equal number (tong shu); eliminate it against the left, obtaining □□□□□; extract the root of the thrice-multiplied power, obtaining 120 bu, which is half the city's diameter. It fits the question.\end{uniten}

\begin{unitzh}\unitlabel{ceyuan-018.9}或問：甲從乾隅東行三百二十步而止；丙出東門南行，乙出東門直行，各不知步數而立。甲廻望乙丙，悉與城叅相直。既而乙就丙斜行三十四步相㑹。問答同前。\end{unitzh}
\begin{uniten}Someone asks: A goes east from the northwest corner 320 bu and stops; C goes out of the east gate and walks south, B goes out of the east gate and walks straight on, each an unknown number of bu, and they stand. A turns and looks at B and C: all are in line with the city. After that, B goes obliquely toward C, 34 bu, and they meet. The question and answer are the same as before.\end{uniten}

\begin{unitzh}\unitlabel{ceyuan-018.10}法曰：甲東行再自之，於上；以二之斜行步乘甲東行冪，減上位，為立方實。兩段南行冪内減東行斜行相乘數，為益從。以甲東行加五{\zhsmall 按：加五即加半}，為從亷。五分虚隅。得全徑。\end{unitzh}
\begin{uniten}Rule: Multiply A's eastward walk by itself twice, placing it above; multiply the area of A's eastward walk by twice the oblique-walk bu, subtract it from the upper position, as the dividend for the cube root. From 2 portions of the area of the southward walk subtract the product of the eastward walk and the oblique walk, as the augmenting joining [Translator's note: C's southward walk is not given, so 南行 here appears corrupt.]. Take A's eastward walk with five added (jia wu) [Siku editors: adding five is adding a half], as the joining edge (cong lian). 5 fen as the empty corner (xu yu). One obtains the whole diameter.\end{uniten}

\begin{unitzh}\unitlabel{ceyuan-018.11}草曰：立天元一為城徑，以減於甲東行步，得□□，為小差。以自之，得丨□□，為小差冪也。乃置甲東行冪，内加小差冪而半之，得□□□，為大弦也{\zhsmall 内帶小差分母}。又置甲東行冪，乃減小差冪而半之，得□□○，為大股也{\zhsmall 内帶小差分母}。乃置斜行步在地，以大股乘之，得□□，合以大弦除之，不除，而又倍之，得□□，為梯頭也{\zhsmall 即兩个小股，内寄大弦為母，權寄}。乃置天元圓徑，以半之，以小差分母通之，得□□，以減於大股，餘得□，又倍之，得□，為梯底也{\zhsmall 即兩个邊股，内亦有小差分母}。以乘權寄，得□□□，為城徑冪也{\zhsmall 内寄大弦及小差分母}，寄左。然後以天元自之為冪，以大弦通之，又以小差通之，得□□□□□，為同數，與左相消，得□□□□，開立方，得二百四十步，即城徑也。合問。\end{unitzh}
\begin{uniten}Working: Set up the celestial unknown as the city's diameter; subtract it from A's eastward-walk bu, obtaining □□, as the small difference. Multiply it by itself, obtaining 丨□□, as the area of the small difference. Then set the area of A's eastward walk, add to it the area of the small difference and halve, obtaining □□□, as the large hypotenuse [Note: it carries within it the small-difference denominator]. Further set the area of A's eastward walk, then subtract the area of the small difference and halve, obtaining □□○, as the large altitude (da gu) [Note: it carries within it the small-difference denominator]. Then set the oblique-walk bu on the ground, multiply by the large altitude, obtaining □□; it should be divided by the large hypotenuse; not dividing, but further doubling it, obtaining □□, as the trapezoid top [Note: that is, 2 small altitudes, within which the large hypotenuse is lodged as denominator; provisionally lodge it (quan ji)]. Then set the celestial-unknown circle diameter, halve it, make it commensurable by the small-difference denominator, obtaining □□; subtract it from the large altitude, the remainder obtained being □; further double it, obtaining □, as the trapezoid bottom [Note: that is, 2 side altitudes (bian gu), within which there is also the small-difference denominator]. Multiply the provisionally lodged quantity by it, obtaining □□□, as the area of the city's diameter [Note: within it are lodged the large-hypotenuse and small-difference denominators]; lodge it on the left. Then multiply the celestial unknown by itself to make the area, make it commensurable by the large hypotenuse, and further by the small difference, obtaining □□□□□, as the equal number; eliminate it against the left, obtaining □□□□; extract the cube root, obtaining 240 bu, which is the city's diameter. It fits the question.\end{uniten}

\begin{unitzh}\unitlabel{ceyuan-018.12}依前問，假令東門外有樹，乙出東門南行，不知步數而立{\zhsmall 只云樹去城步少於乙南行步}。甲從乾隅向東行三百二十步，望乙與樹，悉與城叅相直。乙復就樹斜行三十四步到樹。問答同前。\end{unitzh}
\begin{uniten}Following the previous question, suppose that outside the east gate there is a tree; B goes out of the east gate, walks south an unknown number of bu, and stands [Note: it is only said that the bu from the tree to the city are fewer than B's southward-walk bu]. A goes from the northwest corner toward the east 320 bu and looks at B and the tree: all are in line with the city. B then goes obliquely to the tree, 34 bu, reaching the tree. The question and answer are the same as before.\end{uniten}

\begin{unitzh}\unitlabel{ceyuan-018.13}法曰：甲東行自之，又以斜步乘之，為立方實。置半段甲東行冪於上，以斜步乗甲東行，減上位，為從。亷空。半步常法。得勾圓差。\end{unitzh}
\begin{uniten}Rule: Multiply A's eastward walk by itself, and further multiply it by the oblique bu, as the dividend for the cube root. Set half a portion of the area of A's eastward walk above; multiply A's eastward walk by the oblique bu, subtract it from the upper position, as the joining. The edge is empty (kong). 1/2 bu as the constant divisor. One obtains the base–circle difference (gou yuan cha).\end{uniten}

\begin{unitzh}\unitlabel{ceyuan-018.14}草曰：别得乙斜行即□弦也，□弦得小勾股即大股弦較也。乃立天元一為勾圓差，以自之為冪，副之：上以加於甲東行冪而半之，得□□□，為大弦也{\zhsmall 寄小差分母}；下以減於甲東行冪而半之，得□□□，為大股也{\zhsmall 寄小差分母}。乃置斜步，以大股乘之，得□□□，合大弦除，不除，便以此為小股{\zhsmall 寄大弦分母}。又置斜步，以甲東行乗之，得□，合大弦除，不除，便以此為小勾，而又以通母分通之，得□，為同分小勾也{\zhsmall 寄大弦分母}。注：{\zhsmall 大股乘時有小差分母，今大勾無母，故又以齊同之。}又置斜步，以大弦通之，得□□□，為同分小弦也。三位相併，得□□，為勾圓差也{\zhsmall 寄左}。然後置天元，以大弦通之，得□○□，為同數，與左相消，得□○□□，開立方，得八十步，即勾圓差也。以勾圓差減於甲東行步，餘二百四十步，即城徑也。合問。\end{unitzh}
\begin{uniten}Working: By distinction one obtains that B's oblique walk is the □ hypotenuse; the □ hypotenuse, obtaining the small base and altitude, is the difference between the large altitude and hypotenuse. Then set up the celestial unknown as the base–circle difference; multiply it by itself to make the area, and set it in duplicate: the upper one add to the area of A's eastward walk and halve, obtaining □□□, as the large hypotenuse [Note: the small-difference denominator is lodged]; the lower one subtract from the area of A's eastward walk and halve, obtaining □□□, as the large altitude [Note: the small-difference denominator is lodged]. Then set the oblique bu, multiply by the large altitude, obtaining □□□; it should be divided by the large hypotenuse; not dividing, take this directly as the small altitude [Note: the large-hypotenuse denominator is lodged]. Further set the oblique bu, multiply by A's eastward walk, obtaining □; it should be divided by the large hypotenuse; not dividing, take this directly as the small base, and further make it commensurable by the common-denominator part (tong mu fen), obtaining □, as the small base in the same parts (tong fen) [Note: the large-hypotenuse denominator is lodged]. Annotation (zhu): [Note: When multiplying by the large altitude there was the small-difference denominator; now the large base has no denominator; therefore one further homogenizes and equalizes (qi tong) it.] Further set the oblique bu, make it commensurable by the large hypotenuse, obtaining □□□, as the small hypotenuse in the same parts. Add the 3 positions together, obtaining □□, as the base–circle difference [Note: lodge it on the left]. Then set the celestial unknown, make it commensurable by the large hypotenuse, obtaining □○□, as the equal number; eliminate it against the left, obtaining □○□□; extract the cube root, obtaining 80 bu, which is the base–circle difference. Subtract the base–circle difference from A's eastward-walk bu; the remainder, 240 bu, is the city's diameter. It fits the question.\end{uniten}

\begin{unitzh}\unitlabel{ceyuan-018.15}或問：南門外不知步數有樹。甲從乾東行三百二十步而立。乙出西門，便南行，望樹及甲與城叅相直，却就樹斜行二百五十五步至樹。問答同前。\end{unitzh}
\begin{uniten}Someone asks: Outside the south gate, at an unknown number of bu, there is a tree. A goes east from the northwest 320 bu and stands. B goes out of the west gate and then walks south; he looks at the tree and A, in line with the city, and then goes obliquely to the tree, 255 bu, reaching the tree. The question and answer are the same as before.\end{uniten}

\begin{unitzh}\unitlabel{ceyuan-018.16}法曰：二行相乘於上，以半之甲東行乗之，為實。二行相乘於上，又半之甲東行，以乘甲東行，加上位，為益從。甲東行為從亷。一步虚法。開立方，得半徑。\end{unitzh}
\begin{uniten}Rule: Multiply the two walks together, placing it above; multiply it by half A's eastward walk, as the dividend. Multiply the two walks together, placing it above; further, multiply A's eastward walk by half A's eastward walk, add it to the upper position, as the augmenting joining. A's eastward walk as the joining edge. 1 bu as the empty divisor (xu fa). Extract the cube root, obtaining the radius.\end{uniten}

\begin{unitzh}\unitlabel{ceyuan-018.17}草曰：立天元一為半徑，便以為小勾；其斜行即小弦也。乃以甲東行為大勾，以小弦乘之，復以天元除之，得□□，即大弦也。又倍天元，減東行，餘□□，為小差；以減大弦，餘□□□，為大股也。又倍天元，以減股，餘□□，為大差也。却以半小差□□乗之，得下式□□□，為半徑冪{\zhsmall 寄左}。乃以天元冪與左相消，得丨□□□，開立方，得一百二十步，倍之即城徑也。合問。\end{unitzh}
\begin{uniten}Working: Set up the celestial unknown as the radius, and take it directly as the small base; the oblique walk is the small hypotenuse. Then take A's eastward walk as the large base, multiply it by the small hypotenuse, and then divide by the celestial unknown, obtaining □□, which is the large hypotenuse. Further double the celestial unknown and subtract it from the eastward walk, the remainder □□ being the small difference; subtract it from the large hypotenuse, the remainder □□□ being the large altitude. Further double the celestial unknown and subtract it from the altitude, the remainder □□ being the large difference. Then multiply it by half the small difference, □□, obtaining the following form □□□, as the area of the radius [Note: lodge it on the left]. Then eliminate the area of the celestial unknown against the left, obtaining 丨□□□; extract the cube root, obtaining 120 bu; doubling it gives the city's diameter. It fits the question.\end{uniten}

\begin{unitzh}\unitlabel{ceyuan-018.18}或問：南門外不知步數有槐樹一株，東門外不知步數有柳樹一株，槐柳相距二百八十九步。甲從乾東行三百二十步，斜望槐柳與城叅相直。問答同前。\end{unitzh}
\begin{uniten}Someone asks: Outside the south gate, at an unknown number of bu, there is 1 pagoda tree; outside the east gate, at an unknown number of bu, there is 1 willow tree; the pagoda tree and the willow are 289 bu apart. A goes east from the northwest 320 bu and looks obliquely at the pagoda tree and the willow: they are in line with the city. The question and answer are the same as before.\end{uniten}

\begin{unitzh}\unitlabel{ceyuan-018.19}法曰：二行相乗得數，又自增乘，為實。斜行冪乘甲東行，又倍之，為益從。兩行相乘，又倍之，為益亷。二之斜步，為第二亷。二步常法。開三乘方，得栁至城心步。\end{unitzh}
\begin{uniten}Rule: Multiply the two walks together to get a number, and further multiply it by itself (zi zeng cheng), as the dividend. Multiply A's eastward walk by the area of the oblique walk, and further double it, as the augmenting joining. Multiply the two walks together and further double it, as the augmenting edge. Twice the oblique bu as the second edge. 2 bu as the constant divisor. Extract the root of the thrice-multiplied power, obtaining the bu from the willow to the city's center.\end{uniten}

\begin{unitzh}\unitlabel{ceyuan-018.20}草曰：别得柳至城心步即甲立處柳樹步也。立天元一為柳至城心步，加斜步，得□□，為底弦；以天元乘之，得丨□○，合斜步除，不除，便以此為底勾{\zhsmall 寄斜步分母}。乃再置通勾，以斜步乘之，得□，為帶母通勾；内減底勾，餘□□□，為半徑；以自之，得丨□□□□，為半徑冪，内帶斜步冪分母{\zhsmall 寄左}。乃以天元減斜步，得□□，為明弦；以天元乘之，得□□，合斜步除，不除，便以此為半梯頭{\zhsmall 寄斜步為母}。復以底勾半梯底乘之，得□□□□，為同數，與左相消，得□□□□□，開三乘方，得一百三十六步，即柳至城心步也。合問。\end{unitzh}
\begin{uniten}Working: By distinction one obtains that the bu from the willow to the city's center are the bu from the place where A stands to the willow tree. Set up the celestial unknown as the bu from the willow to the city's center; add the oblique bu, obtaining □□, as the bottom hypotenuse (di xian); multiply it by the celestial unknown, obtaining 丨□○; it should be divided by the oblique bu; not dividing, take this directly as the bottom base (di gou) [Note: the oblique-bu denominator is lodged]. Then again set the through base (tong gou), multiply it by the oblique bu, obtaining □, as the through base carrying the denominator; subtract the bottom base from it, the remainder □□□ being the radius; multiply it by itself, obtaining 丨□□□□, as the area of the radius, carrying within it the area of the oblique bu as denominator [Note: lodge it on the left]. Then subtract the celestial unknown from the oblique bu, obtaining □□, as the bright hypotenuse (ming xian); multiply it by the celestial unknown, obtaining □□; it should be divided by the oblique bu; not dividing, take this directly as the half trapezoid top [Note: the oblique bu is lodged as denominator]. Further multiply it by the bottom base, the half trapezoid bottom, obtaining □□□□, as the equal number; eliminate it against the left, obtaining □□□□□; extract the root of the thrice-multiplied power, obtaining 136 bu, which is the bu from the willow to the city's center. It fits the question.\end{uniten}

\begin{unitzh}\unitlabel{ceyuan-018.21}或問：甲從乾隅東行三百二十步而立，乙出城東行，丙出城南行，三人相望，俱與城相直。乙丙共行了一百五十一步。問答同前。\end{unitzh}
\begin{uniten}Someone asks: A goes east from the northwest corner 320 bu and stands; B goes out of the city and walks east; C goes out of the city and walks south; the 3 men look at one another, and all are in line with the city. B and C have walked 151 bu altogether. The question and answer are the same as before.\end{uniten}

\begin{unitzh}\unitlabel{ceyuan-018.22}法曰：以甲東行為冪，折半，又以自之，為三乘方實。倍共步，加甲東行，以乗半段甲行冪，為從方。甲行乗共數，為從亷。甲東行加五，為第二益亷。二分五釐常法。得小差。\end{unitzh}
\begin{uniten}Rule: Take A's eastward walk to make an area, halve it, and further multiply it by itself, as the dividend for the thrice-multiplied power. Double the total bu, add A's eastward walk, and multiply half a portion of the area of A's walk by it, as the joining square. A's walk multiplied by the total number, as the joining edge. A's eastward walk with five added, as the second augmenting edge. 2 fen 5 li as the constant divisor. One obtains the small difference.\end{uniten}

\begin{unitzh}\unitlabel{ceyuan-019.1}草曰：别得乙丙共行步，即明股□勾共也。立天元一為小差，以自之，副置二位：上位減於甲東行冪，以天元除之，又折半，得□○□，即大股也；下位加甲行冪，以天元除之，又折半，得□○□，為大弦也。其甲東行即大勾也。併大勾、大股，得□□□，即大和也。再立天元，以減甲東行步，得□□，即圓徑也。以圓徑加共行步，得□□，即皇極和也{\zhsmall 即小和，又為髙弦、平弦共數}。又倍之，得□□，即黄長弦、黄廣弦共也。内減大弦，得下式□□□，為皇極内小黄方也{\zhsmall 亦為虚弦}。再置大和□□□，以小黄方乘之，得下式□□□□□，合以小和除之，不除，便以為城徑，内寄小和為母{\zhsmall 寄左}。然後天元減甲東行，得□□，為大黄方，以小和乘之，得丨□□，為同數，與左相消，得□□□□□，開三乗方，得八十步，即小差也。以小差減甲東行，餘二百四十步，即城徑也。合問。\end{unitzh}
\begin{uniten}Working: Separately, one obtains that the combined walk in bu (步) of Yi and Bing is the bright altitude (ming gu) and the □ base (gou) combined. [Translator's note: □ here and below marks characters or rod numerals not reproduced in the transcription; ○ and 丨 are the rod numerals 0 and 1.] Set up the celestial unknown (tianyuan) as the small difference (xiao cha); multiply it by itself and set it down in duplicate in two positions. The upper position: subtract it from the area (mi) of Jia's eastward walk, divide by the celestial unknown, and also halve, obtaining □○□, which is the large altitude (da gu). The lower position: add the area of Jia's walk, divide by the celestial unknown, and also halve, obtaining □○□, which is the large hypotenuse (da xian). Jia's eastward walk is the large base (da gou). Add together the large base and the large altitude, obtaining □□□, which is the large sum (da he). Again set up the celestial unknown and subtract it from Jia's eastward walk in bu, obtaining □□, which is the circle diameter (yuan jing). Add the circle diameter to the combined walk in bu, obtaining □□, which is the Imperial Pole sum (huangji he) [Note: that is, the small sum (xiao he); it is also the combined number of the height hypotenuse (gao xian) and the level hypotenuse (ping xian)]. Further double it, obtaining □□, which is the yellow-long hypotenuse (huangchang xian) and the yellow-broad hypotenuse (huangguang xian) combined. From it subtract the large hypotenuse, obtaining the array below (xia shi), □□□, which is the small yellow square (xiao huangfang) within the Imperial Pole [Note: it is also the void hypotenuse (xu xian)]. Again set down the large sum, □□□, and multiply it by the small yellow square, obtaining the array below, □□□□□. It ought to be divided by the small sum; do not divide, but simply take it as the city diameter (cheng jing), with the small sum lodged within it as denominator (mu) [Note: lodge it on the left]. After that, subtract the celestial unknown from Jia's eastward walk, obtaining □□, which is the large yellow square (da huangfang); multiply it by the small sum, obtaining 丨□□, which is the same quantity (tong shu). Eliminate it against the left (xiang xiao), obtaining □□□□□; extract the root of the thrice-multiplied square (san cheng fang), obtaining 80 bu, which is the small difference. Subtract the small difference from Jia's eastward walk; the remainder, 240 bu, is the city diameter. This matches the question.\end{uniten}

\begin{unitzh}\unitlabel{ceyuan-019.2}或問：丙出南門東行，乙出東門南行，各不知步數而立。甲從乾隅東行三百二十步，望乙、丙，悉與城㕘相直。乙就丙斜行一百二步相㑹。問答同前。\end{unitzh}
\begin{uniten}Someone asks: Bing goes out of the south gate and walks east; Yi goes out of the east gate and walks south; each stands after an unknown number of bu. Jia walks east 320 bu from the Qian corner (qian yu); looking at Yi and Bing, all are aligned with the city (can xiang zhi). Yi goes toward Bing, walking obliquely 102 bu, and they meet. The question and the answer are the same as before.\end{uniten}

\begin{unitzh}\unitlabel{ceyuan-019.3}法曰：甲東行自之於上，倍斜行步乘之，為立方實；倍斜行步乘甲東行於上，加兩段甲東行冪，為從；四之甲東行，為益亷；四為隅法。得半城徑。\end{unitzh}
\begin{uniten}Rule: Multiply Jia's eastward walk by itself and set it above; multiply it by twice the oblique walk in bu; this is the dividend (shi) for the cube root (lifang). Multiply Jia's eastward walk by twice the oblique walk in bu and set it above; add 2 pieces (duan) of the area (mi) of Jia's eastward walk; this is the joining term (cong). 4 times Jia's eastward walk is the augmenting edge (yi lian); 4 is the corner divisor (yu fa). One obtains half the city diameter.\end{uniten}

\begin{unitzh}\unitlabel{ceyuan-019.4}草曰：别得斜步即虛弦，減於全徑，即小和也。乃立天元一為半徑，以二之，減於甲東行，得□□，為小差也。以自之，得□□□，為小差冪也。置甲東行冪，内加小差冪而半之，得下□□□，為大弦{\zhsmall 内帶小差分母}。置甲東行冪，内減小差冪而半之，得□□，為大股也，内亦帶小差為母。又以小差乘大勾，得□□，併入大股，得□□□，為大和也{\zhsmall 帶小差母}。乃先以小弦乗大和，得下□□□，寄左。次以斜步減於二天元，得□□，為小和，以乗大弦，得下式□□□□，為同數，與左相消，得□□□□，開立方，得一百二十步，即半城徑也。合問。\end{unitzh}
\begin{uniten}Working: Separately, one obtains that the oblique bu are the void hypotenuse; subtracted from the whole diameter (quan jing), this is the small sum. Then set up the celestial unknown as the half-diameter (ban jing); double it and subtract it from Jia's eastward walk, obtaining □□, which is the small difference. Multiply it by itself, obtaining □□□, which is the area (mi) of the small difference. Set down the area of Jia's eastward walk, add into it the area of the small difference, and halve, obtaining below □□□, which is the large hypotenuse [Note: it carries the small difference within it as denominator (fen mu)]. Set down the area of Jia's eastward walk, subtract from it the area of the small difference, and halve, obtaining □□, which is the large altitude; it, too, carries the small difference within it as denominator. Further multiply the large base by the small difference, obtaining □□; add it into the large altitude, obtaining □□□, which is the large sum [Note: carrying the small difference as denominator]. Then first multiply the large sum by the small hypotenuse (xiao xian), obtaining below □□□; lodge it on the left. Next subtract the oblique bu from twice the celestial unknown, obtaining □□, which is the small sum; multiply the large hypotenuse by it, obtaining the array below, □□□□, which is the same quantity. Eliminate it against the left, obtaining □□□□; extract the cube root, obtaining 120 bu, which is half the city diameter. This matches the question.\end{uniten}

\begin{unitzh}\unitlabel{ceyuan-019.5}依前問，假令乙出東門南行，丙出南門東行，各不知步數而立{\zhsmall 只云丙行多於乙行步}。甲從乾隅東行三百二十步，望乙、丙，與城㕘相直。其乙丙共行一百二步。問答同前。\end{unitzh}
\begin{uniten}Following the preceding question, suppose that Yi goes out of the east gate and walks south, and Bing goes out of the south gate and walks east; each stands after an unknown number of bu [Note: it only says that Bing's walk exceeds Yi's walk in bu]. Jia walks east 320 bu from the Qian corner; looking at Yi and Bing, they are aligned with the city. Yi and Bing together walk 102 bu. The question and the answer are the same as before.\end{uniten}

\begin{unitzh}\unitlabel{ceyuan-019.6}法曰：倍共步，以乗甲東行冪，為立方實；共步乗甲東行於上，又以甲東行自之，加上位，為益從；甲東行為從亷；五分隅常法。得城徑。\end{unitzh}
\begin{uniten}Rule: Double the combined bu and multiply the area (mi) of Jia's eastward walk by it; this is the dividend (shi) for the cube root. Multiply Jia's eastward walk by the combined bu and set it above; further multiply Jia's eastward walk by itself and add it to the upper position; this is the augmenting joining term (yi cong). Jia's eastward walk is the joining edge (cong lian); 5 fen is the corner constant divisor (yu chang fa) [Translator's note: 5 fen = 0.5]. One obtains the city diameter.\end{uniten}

\begin{unitzh}\unitlabel{ceyuan-019.7}草曰：别得共步便為小弦，得小勾、小股，即與圓徑同。立天元為城徑，以減乙東行，得□□，為小差，以自之，得□□□，為小差冪也。乃置甲東行，以自之為冪，副之：上以加小差冪而半之，得□□□，為大弦也{\zhsmall 内寄小差分母}；下以減小差冪而半之，得□□○，為大股也{\zhsmall 内寄小差分母}。乃置共步在地，以大股乘之，得□□，合大弦除，不除，便以此為小股也{\zhsmall 寄大弦分母}。又置共步，以甲東行乘之，得□，合以大弦除，不除，便以此為小勾，而又以元分母小差乘之，得□□，為同分小勾{\zhsmall 只寄大弦分母}。注{\zhsmall 其大弦内元帶小差分母，其大勾内却無分母，故母故今復以小差通之，齊同其分母也}。又置共步，以大弦通之，得□□□，同分小弦也。三位相併，得□□，為城徑也{\zhsmall 内有大弦分母}，寄左。然後置天元城徑□，以大弦分母通之，得□□□○，為同數，與左相消，得□□□□，開立方，得二百四十步，即城徑也。合問。\end{unitzh}
\begin{uniten}Working: Separately, one obtains that the combined bu are the small hypotenuse; obtaining the small base (xiao gou) and the small altitude (xiao gu), this is then the same as the circle diameter. Set up the celestial unknown as the city diameter and subtract it from Yi's eastward walk [Translator's note: 乙 (Yi) here appears to be an error for 甲 (Jia)], obtaining □□, which is the small difference; multiply it by itself, obtaining □□□, which is the area of the small difference. Then set down Jia's eastward walk and multiply it by itself to make the area (mi); set it down in duplicate. For the upper, add the area of the small difference and halve, obtaining □□□, which is the large hypotenuse [Note: the small difference is lodged within it as denominator]; for the lower, subtract the area of the small difference and halve, obtaining □□○, which is the large altitude [Note: the small difference is lodged within it as denominator]. Then set down the combined bu on the ground (zai di) and multiply it by the large altitude, obtaining □□. It ought to be divided by the large hypotenuse; do not divide, but simply take this as the small altitude [Note: the large hypotenuse is lodged as denominator]. Again set down the combined bu and multiply it by Jia's eastward walk, obtaining □. It ought to be divided by the large hypotenuse; do not divide, but simply take this as the small base; and further multiply it by the original denominator, the small difference, obtaining □□, which is the small base of the same denominator (tong fen) [Note: only the large hypotenuse is lodged as denominator]. Annotation (zhu): [Note: The large hypotenuse originally carries the small difference within it as denominator, while the large base has no denominator within it; therefore denominator therefore [Translator's note: 故母故 appears garbled] now we again bring it to a common denominator (tong) with the small difference, homogenizing and equalizing (qi tong) their denominators.] Again set down the combined bu and bring it to a common denominator with the large hypotenuse, obtaining □□□, the small hypotenuse of the same denominator. Add the three positions together, obtaining □□, which is the city diameter [Note: within it there is the large hypotenuse as denominator]; lodge it on the left. After that, set down the celestial unknown, the city diameter, □, and bring it to a common denominator with the large-hypotenuse denominator, obtaining □□□○, which is the same quantity. Eliminate it against the left, obtaining □□□□; extract the cube root, obtaining 240 bu, which is the city diameter. This matches the question.\end{uniten}

\begin{unitzh}\unitlabel{ceyuan-019.8}測圓海鏡卷六\end{unitzh}
\begin{uniten}Sea Mirror of Circle Measurements (Ceyuan haijing), juan 6\end{uniten}

\begin{unitzh}\unitlabel{ceyuan-019.9}<子部,天文算法類,算書之屬,測圓海鏡>\end{unitzh}
\begin{uniten}<Masters branch (zibu), Astronomy and Calculation category (tianwen suanfa lei), Mathematical Books subcategory (suanshu zhi shu), Sea Mirror of Circle Measurements (Ceyuan haijing)>\end{uniten}

\section*{卷07}\addcontentsline{toc}{section}{卷07}

\begin{unitzh}\unitlabel{ceyuan-020.1}欽定四庫全書\end{unitzh}
\begin{uniten}Imperially Authorized Complete Library of the Four Treasuries (Qinding Siku quanshu)\end{uniten}

\begin{unitzh}\unitlabel{ceyuan-020.2}測圓海鏡卷七\end{unitzh}
\begin{uniten}Sea Mirror of Circle Measurements (Ceyuan haijing), juan 7\end{uniten}

\begin{unitzh}\unitlabel{ceyuan-020.3}元·李冶撰\end{unitzh}
\begin{uniten}Composed by Li Ye of the Yuan\end{uniten}

\begin{unitzh}\unitlabel{ceyuan-020.4}明□前一十八問\end{unitzh}
\begin{uniten}18 Problems in Front of Ming [illegible]\end{uniten}

\begin{unitzh}\unitlabel{ceyuan-020.5}或問：出南門東行七十二步有樹，出東門南行三十步見之。問答同前。\end{unitzh}
\begin{uniten}Someone asks: Going out of the south gate and walking east 72 bu (步), there is a tree; going out of the east gate and walking south 30 bu, one sees it. The question and the answer are the same as before.\end{uniten}

\begin{unitzh}\unitlabel{ceyuan-020.6}法曰：倍南行以乘倍東行，為平實；併二行，又倍之，為從；一虚隅；得城徑。\end{unitzh}
\begin{uniten}Rule: Double the southward walk and with it multiply double the eastward walk, as the plane dividend (pingshi); add the two walks and double again, as the follower (cong); 1 as the empty corner (xuyu); one obtains the city diameter (chengjing).\end{uniten}

\begin{unitzh}\unitlabel{ceyuan-020.7}草曰：識别得此問名為弦外容圓，又為内率求虚積。其二行步相併為虚弦，若以相減即虚較也。又倍東行為弦較和，倍南行即弦較較，此二數相乘則兩虚積也。若直以二行相乘，則半个虚積也。又倍東行減於城徑，餘即二虚勾也；倍南行減於城徑，則二虚股也。虚積上三事和即城徑也。乃立天元一為圓徑，便以為三事和也。倍二行步減之，得□□，為黄方。一天元乘之，得□□，為二虚積{\zhsmall 寄左}。然後倍東行以乗倍南行，得八千六百四十，為同數。與左相消，得丨□□，益積開平方，得二百四十步，即城徑也。合問。\end{unitzh}
\begin{uniten}Working: One discerns that this question is named "the circle contained outside the hypotenuse (xian wai rongyuan)", and that it is also "seeking the void area (xuji) by the inner rate (neilü)". The two walking distances added together are the void hypotenuse (xuxian); if they are subtracted from each other, it is the void difference (xujiao). Also, double the eastward walk is the hypotenuse-difference sum (xianjiao he), and double the southward walk is the hypotenuse-difference difference (xianjiao jiao); these two numbers multiplied together give 2 void areas. If one directly multiplies the two walks together, it is half a void area. Also, subtract double the eastward walk from the city diameter: the remainder is 2 void bases (xugou); subtract double the southward walk from the city diameter: then it is 2 void altitudes (xugu). The sum of the three things (sanshi he) on the void area is the city diameter. Then set up the celestial unknown (tianyuan) as the circle diameter, and this is taken as the sum of the three things. Double the two walking distances and subtract them from it, obtaining □□ [Translator's note: □ here and below stands for counting-rod numbers or arrays not transcribed in the source; 丨 and ○ are rod signs as transcribed.] as the yellow square (huangfang). Multiply it by 1 celestial unknown, obtaining □□ as 2 void areas [Note: deposit on the left (jizuo)]. Afterwards double the eastward walk and with it multiply double the southward walk, obtaining 8640 as the equal quantity (tongshu). Eliminate it against the left, obtaining 丨□□; extract the square root with augmented area (yiji), obtaining 240 bu, which is the city diameter. This agrees with the question.\end{uniten}

\begin{unitzh}\unitlabel{ceyuan-020.8}又法：二行步相乘為實；二行步相併為從；一步虛法；得半徑。\end{unitzh}
\begin{uniten}Another rule: Multiply the two walking distances together as the dividend (shi); add the two walking distances together as the follower; 1 bu as the empty divisor (xufa); one obtains the half-diameter (banjing).\end{uniten}

\begin{unitzh}\unitlabel{ceyuan-020.9}草曰：立天元一為半徑，副置二位。上加東行步，得□□，為大差勾；下加□股，得□□，為小差股。此二數相乘，得下式丨□□，為半段黄方冪{\zhsmall 寄左}。然後立天元以自之，又二之，與左相消，得丨□□，益積開平方，得一百二十步，即半城徑也。\end{unitzh}
\begin{uniten}Working: Set up the celestial unknown as the half-diameter and place it in duplicate in 2 positions. To the upper add the eastward walking distance, obtaining □□ as the great-difference base (dacha gou); to the lower add the [illegible] altitude, obtaining □□ as the small-difference altitude (xiaocha gu). Multiply these two numbers together, obtaining the following form 丨□□ as half a portion (duan) of the area (mi) of the yellow square [Note: deposit on the left]. Afterwards set up the celestial unknown, multiply it by itself, and double it; eliminate it against the left, obtaining 丨□□; extract the square root with augmented area, obtaining 120 bu, which is half the city diameter.\end{uniten}

\begin{unitzh}\unitlabel{ceyuan-020.10}又法：二云數相乘，倍之於上，加云數差冪，權寄。併二云數，又自增乗，得數内減上位，為平實；併云數而倍之，為從；二步益隅；得半徑。\end{unitzh}
\begin{uniten}Another rule: Multiply the two stated numbers (yunshu) together and double it, placing it above; add the area of the difference of the stated numbers, and set it aside provisionally (quanji). Add the two stated numbers and also multiply the sum by itself (zizengcheng); from the resulting number subtract the upper position, as the plane dividend; add the stated numbers and double it, as the follower; 2 bu as the augmenting corner (yiyu); one obtains the half-diameter.\end{uniten}

\begin{unitzh}\unitlabel{ceyuan-020.11}草曰：立天元一為半徑，副之。上減明勾，得下□□，為虚勾；下減□股，得□□，為虚股。勾股相乘，得丨□□；又倍之，得□□□；又加二行差冪□，得□□□，為弦冪{\zhsmall 寄左}。然後併云步以自之，得□，為同數。與左相消，得□□□，益積開平方，得一百二十步，即半城徑也。\end{unitzh}
\begin{uniten}Working: Set up the celestial unknown as the half-diameter and duplicate it. From the upper subtract the Ming base (minggou), obtaining the following □□ as the void base; from the lower subtract the [illegible] altitude, obtaining □□ as the void altitude. Multiply base and altitude together, obtaining 丨□□; double it, obtaining □□□; add the area of the difference of the two walks, □, obtaining □□□ as the hypotenuse area (xianmi) [Note: deposit on the left]. Afterwards add the stated distances and multiply the sum by itself, obtaining □ as the equal quantity. Eliminate it against the left, obtaining □□□; extract the square root with augmented area, obtaining 120 bu, which is half the city diameter.\end{uniten}

\begin{unitzh}\unitlabel{ceyuan-020.12}又法：云數相乘，又倍之，為平實；云數相減，為從；一常法；得虚勾。\end{unitzh}
\begin{uniten}Another rule: Multiply the stated numbers together and double it, as the plane dividend; subtract the stated numbers from each other, as the follower; 1 as the regular divisor (changfa); one obtains the void base.\end{uniten}

\begin{unitzh}\unitlabel{ceyuan-020.13}草曰：立天元一為虚勾。以南行減東行，餘四十二步，為虚較也。以虚較加天元，得丨□，為虛股。以天元乘之，得下丨□，為直積{\zhsmall 寄左}。然後倍南行乘東行，得□，與左相消，得丨□□，開平方，得四十八步，即虚勾也。以勾除積，得九十步，即虚股也。併勾股，得□，為虚和也。内加入二行併□，得□，即圓徑也。\end{unitzh}
\begin{uniten}Working: Set up the celestial unknown as the void base. Subtract the southward walk from the eastward walk; the remainder, 42 bu, is the void difference. Add the void difference to the celestial unknown, obtaining 丨□ as the void altitude. Multiply it by the celestial unknown, obtaining the following 丨□ as the rectangular area (zhiji) [Note: deposit on the left]. Afterwards double the southward walk and multiply the eastward walk by it, obtaining □; eliminate it against the left, obtaining 丨□□; extract the square root, obtaining 48 bu, which is the void base. Divide the area by the base, obtaining 90 bu, which is the void altitude. Add base and altitude, obtaining □ as the void sum (xuhe). To it add the sum of the two walks, □, obtaining □, which is the circle diameter.\end{uniten}

\begin{unitzh}\unitlabel{ceyuan-020.14}又法：併兩行步以自乘於上；又倍南行乘倍東行，加上位，為平實；一隅法；得小和。\end{unitzh}
\begin{uniten}Another rule: Add the two walking distances and multiply the sum by itself, placing it above; also double the southward walk and multiply double the eastward walk by it, and add it to the upper position, as the plane dividend; 1 as the corner divisor (yufa); one obtains the small sum (xiaohe).\end{uniten}

\begin{unitzh}\unitlabel{ceyuan-020.15}草曰：立天元一為小和，併二行步加之，得□□，為三事和也。倍二行步而併之，得□，以減三事和，餘□□，為黄方。却以三事和乘之，得下丨□□，為二虚積也{\zhsmall 寄左}。乃倍南行以乘倍東行，得□，為同數。與左相消，得丨□□，開平方，得一百三十八步，即虚和也。加入二行步，得二百四十步，即城徑也。合問。\end{unitzh}
\begin{uniten}Working: Set up the celestial unknown as the small sum; add the two walking distances to it, obtaining □□ as the sum of the three things. Double the two walking distances and add them together, obtaining □; subtract it from the sum of the three things; the remainder, □□, is the yellow square. Then multiply it by the sum of the three things, obtaining the following 丨□□ as 2 void areas [Note: deposit on the left]. Then double the southward walk and with it multiply double the eastward walk, obtaining □ as the equal quantity. Eliminate it against the left, obtaining 丨□□; extract the square root, obtaining 138 bu, which is the void sum. Add the two walking distances, obtaining 240 bu, which is the city diameter. This agrees with the question.\end{uniten}

\begin{unitzh}\unitlabel{ceyuan-020.16}或問：丙出南門直行一百三十五步而立，甲出東門直行一十六步見之。問答同前。\end{unitzh}
\begin{uniten}Someone asks: Bing goes out of the south gate, walks straight ahead 135 bu and stands; Jia goes out of the east gate, walks straight ahead 16 bu and sees him. The question and the answer are the same as before.\end{uniten}

\begin{unitzh}\unitlabel{ceyuan-020.17}法曰：以丙行步一百三十五步再自之，得二百四十六萬零三百七十五，於上。又以甲行步一十六乘丙行冪一萬八千二百二十五，得二十九萬一千六百，以乘上位，得七千一百七十四億四千五百三十五萬，為三乘方實。以二行步相乘，又倍之，得四千三百二十，以乘丙行步再自之數，得一百六億二千八百八十二萬，為益從。第一亷空。以甲行乘丙行冪，得二十九萬一千六百，又倍之，得五十八萬三千二百，於上；四之甲行冪一千零二十四，以乘丙行步，得一十三萬八千二百四十，減上位，餘四十四萬四千九百六十，為第二亷。二行步相乘，得二千一百六十，為虚常法。得丙行步上勾弦差八十一。\end{unitzh}
\begin{uniten}Rule: Take Bing's walking distance, 135 bu, and multiply it by itself twice, obtaining 2460375, placed above. Also with Jia's walking distance, 16, multiply the area (mi) of Bing's walk, 18225, obtaining 291600; with it multiply the upper position, obtaining 717445350000 as the dividend of the thrice-multiplied power (sancheng fang). Multiply the two walking distances together and double it, obtaining 4320; with it multiply the number of Bing's walking distance multiplied by itself twice, obtaining 10628820000, as the augmenting follower (yicong). The first edge (lian) is empty. With Jia's walk multiply the area of Bing's walk, obtaining 291600; double it, obtaining 583200, placed above; with 4 times the area of Jia's walk, 1024, multiply Bing's walking distance, obtaining 138240; subtract it from the upper position; the remainder, 444960, is the second edge. Multiply the two walking distances together, obtaining 2160, as the empty regular divisor. One obtains the base-hypotenuse difference (gouxian cha) on Bing's walking distance, 81.\end{uniten}

\begin{unitzh}\unitlabel{ceyuan-020.18}按：法中載數自此始，亦擇其數繁者詳之，使人易曉也。\end{unitzh}
\begin{uniten}Note: The recording of numbers in the rules begins from here; also, those whose numbers are complicated are selected and set out in detail, so that people may understand them easily.\end{uniten}

\begin{unitzh}\unitlabel{ceyuan-020.19}草曰：識得二數相併，以減於皇極弦，餘即虚勾虛股併也。若以二數相減，餘為髙弦内減平弦，又為皇極弦内少个小差弦，又為大差弦内減个皇極弦也。立天元一為丙行大差數。置丙行步一百三十五，自乘，得□，用天元除之，得□□，為勾弦併也。上減天元，得□□□，為二丙勾也。復用丙南行乘之，得□□□，為二積也。又以天元除之，得□□○□，為丙勾外容圓徑{\zhsmall 泛寄}。\end{unitzh}
\begin{uniten}Working: One recognizes that when the two numbers are added together and subtracted from the Imperial Pole hypotenuse (huangji xian), the remainder is the sum of the void base and the void altitude. If the two numbers are subtracted from each other, the remainder is the high hypotenuse (gaoxian) minus the level hypotenuse (pingxian); it is also the Imperial Pole hypotenuse less a small-difference hypotenuse (xiaocha xian); it is also the great-difference hypotenuse (dacha xian) minus an Imperial Pole hypotenuse. Set up the celestial unknown as the great-difference number of Bing's walk. Set up Bing's walking distance, 135, and multiply it by itself, obtaining □; divide it by the celestial unknown, obtaining □□ as the sum of base and hypotenuse (gouxian bing). From the above subtract the celestial unknown, obtaining □□□ as 2 Bing bases. Again multiply it by Bing's southward walk, obtaining □□□ as 2 areas. Also divide it by the celestial unknown, obtaining □□○□ as the diameter of the circle contained outside Bing's base [Note: floating deposit (fanji)].\end{uniten}

\begin{unitzh}\unitlabel{ceyuan-020.20}别置丙南行，用二甲勾乘之，得□，合用二丙勾除之，不受除，便以此為甲股{\zhsmall 内寄二丙勾為分母}。復用二甲勾三十二乘之，得□，為二个甲直積也。又置丙南行，内減天元，得□□，為黄方。以自乘，得丨□□，為丙上勾弦差乘股弦差二段。以天元除之，得□□□，為两个丙小差也。乃用甲股乗之，得下式□□□；復用丙南行除之，得□□□；又折半，得□□□，為一个甲步股弦差也，内亦帶前二丙勾分母。\end{unitzh}
\begin{uniten}Separately set up Bing's southward walk and multiply it by 2 Jia bases, obtaining □; it ought to be divided by 2 Bing bases, but it does not admit division, so this is then taken as the Jia altitude [Note: within it, 2 Bing bases are deposited as the denominator (fenmu)]. Again multiply it by 2 Jia bases, 32, obtaining □ as 2 Jia rectangular areas. Also set up Bing's southward walk and subtract the celestial unknown from it, obtaining □□ as the yellow square. Multiply it by itself, obtaining 丨□□ as 2 portions of the base-hypotenuse difference on Bing multiplied by the altitude-hypotenuse difference (guxian cha). Divide it by the celestial unknown, obtaining □□□ as 2 Bing small differences (xiaocha). Then multiply it by the Jia altitude, obtaining the following form □□□; again divide it by Bing's southward walk, obtaining □□□; also halve it, obtaining □□□ as 1 altitude-hypotenuse difference of Jia's walk; within it is also carried the earlier denominator of 2 Bing bases.\end{uniten}

\begin{unitzh}\unitlabel{ceyuan-020.21}復置二个甲直積，内已寄此甲股弦差分母，便為甲步股外容圓徑{\zhsmall 寄左}。乃再置先求到泛寄{\zhsmall 按即前所寄□□○□之數}，用甲股弦差分母乘之，得□□○□□，為同數。與左相消，得下式□□○□□，開三乗方，得八十一步，即丙步上勾弦差也。\end{unitzh}
\begin{uniten}Again set up the 2 Jia rectangular areas; within them the denominator of this Jia altitude-hypotenuse difference has already been deposited, so they are then the diameter of the circle contained outside the altitude of Jia's walk [Note: deposit on the left]. Then again set up the floating deposit found earlier [Siku editors: that is, the number □□○□ deposited earlier]; multiply it by the denominator of the Jia altitude-hypotenuse difference, obtaining □□○□□ as the equal quantity. Eliminate it against the left, obtaining the following form □□○□□; extract the root of the thrice-multiplied power (kai sancheng fang), obtaining 81 bu, which is the base-hypotenuse difference on Bing's walk.\end{uniten}

\begin{unitzh}\unitlabel{ceyuan-020.22}《鈐經》載此法：以勾弦差率冪減丙行差冪，復以丙行乘之為實，以差率冪為法，如法得徑。此法只是以勾外求容圓半，合以大差除陪積，而今皆以大差冪為分母也。依法求之：勾弦差八十一，自之，得六千五百六十一，以減於丙行冪一萬八千二百二十五，餘一萬一千六百六十四；復以丙行一百三十五乘之，得一百五十七萬四千六百四十，為實；以大差冪六千五百六十一為法；如法得二百四十步，即城徑也。\end{unitzh}
\begin{uniten}The Qianjing records this method: subtract the area of the base-hypotenuse difference rate (lü) from the difference area of Bing's walk, and again multiply by Bing's walk, as the dividend; take the area of the difference rate as the divisor (fa); dividing by the divisor gives the diameter. This method is simply to seek the half of the circle contained outside the base: it ought to divide the accompanying area (peiji) by the great difference, but now in every case the area of the great difference is taken as the denominator. Seeking it according to the method: the base-hypotenuse difference, 81, multiplied by itself, gives 6561; subtract it from the area of Bing's walk, 18225; the remainder is 11664; again multiply by Bing's walk, 135, obtaining 1574640 as the dividend; take the area of the great difference, 6561, as the divisor; dividing by the divisor gives 240 bu, which is the city diameter.\end{uniten}

\begin{unitzh}\unitlabel{ceyuan-020.23}又法：二行相乘，得數又自之，為三乘方實；併二行步，以乗二行相乘數，又倍之，為從；二行相併數以自乘於上，又二行相減數自乗，減上位，為第一亷；第二亷空；一益隅；益積開之，得半徑{\zhsmall 其第一亷只是四段二行相乗數}。\end{unitzh}
\begin{uniten}Another rule: Multiply the two walks together, and multiply the resulting number by itself again, as the dividend of the thrice-multiplied power; add the two walking distances and with them multiply the product of the two walks, and double it, as the follower; multiply the sum of the two walks by itself, placing it above; also multiply the difference of the two walks by itself and subtract it from the upper position, as the first edge; the second edge is empty; 1 as the augmenting corner; extract it with augmented area, obtaining the half-diameter [Note: its first edge is simply 4 portions of the product of the two walks].\end{uniten}

\begin{unitzh}\unitlabel{ceyuan-020.24}草曰：立天元一為半城徑，副置之。上加南行步，得□□，為股；下位加東行步，得□□，為勾。勾股相乘，得丨□□，為直積一段。以天元除之，得丨□□，為弦。以自之，得丨□□□□，為弦冪{\zhsmall 寄左}。乃以勾自之，得丨□□；又以股自之，得丨□□；二位相併，得□□□，為同數。與左相消，得丨○□□□，益積開三乘方，得一百二十步，即半城徑也。\end{unitzh}
\begin{uniten}Working: Set up the celestial unknown as half the city diameter and place it in duplicate. To the upper add the southward walking distance, obtaining □□ as the altitude; to the lower position add the eastward walking distance, obtaining □□ as the base. Multiply base and altitude together, obtaining 丨□□ as 1 portion of the rectangular area. Divide it by the celestial unknown, obtaining 丨□□ as the hypotenuse. Multiply it by itself, obtaining 丨□□□□ as the hypotenuse area [Note: deposit on the left]. Then multiply the base by itself, obtaining 丨□□; also multiply the altitude by itself, obtaining 丨□□; add the 2 positions together, obtaining □□□ as the equal quantity. Eliminate it against the left, obtaining 丨○□□□; extract the root of the thrice-multiplied power with augmented area, obtaining 120 bu, which is half the city diameter.\end{uniten}

\begin{unitzh}\unitlabel{ceyuan-020.25}又法：條段同前。\end{unitzh}
\begin{uniten}Another rule: the sections (tiaoduan) are the same as before.\end{uniten}

\begin{unitzh}\unitlabel{ceyuan-020.26}草曰：以前求得勾股率，置出南門步為小股，以勾率乘之，得□□，合以股率除，不除，寄為母，便以此為半梯頭，於上。又置南行步，加二天元，得□□，為大股，以勾率乘之，得□□□，合以股率除，不除，寄為母，便以此為梯底。以乘上位，得□□□□，為半徑自乘數，内帶股率冪為母{\zhsmall 寄左}。然後置天元以自之，又以股率冪乘之，得下丨□□□，為同數。與左相消，得數一如前答。\end{unitzh}
\begin{uniten}Working: With the base and altitude rates found before, set up the distance walked out of the south gate as the small altitude; multiply it by the base rate, obtaining □□; it ought to be divided by the altitude rate; not dividing, deposit it as the denominator (mu), and then take this as half the trapezoid head (titou), placed above. Also set up the southward walking distance and add 2 celestial unknowns, obtaining □□ as the great altitude; multiply it by the base rate, obtaining □□□; it ought to be divided by the altitude rate; not dividing, deposit it as the denominator, and then take this as the trapezoid bottom (tidi). With it multiply the upper position, obtaining □□□□ as the number of the half-diameter multiplied by itself, which carries within it the area of the altitude rate as the denominator [Note: deposit on the left]. Afterwards set up the celestial unknown, multiply it by itself, and also multiply it by the area of the altitude rate, obtaining the following 丨□□□ as the equal quantity. Eliminate it against the left; the numbers obtained are entirely as in the previous answer.\end{uniten}

\begin{unitzh}\unitlabel{ceyuan-020.27}又法：以二行差冪數自乗，又倍之，為實；併二行步，以乘二行差冪，又四之，為益從；四段南行冪内減二段差冪，於上，又二段差冪内減四段東行冪，餘以減上位{\zhsmall 按併二行冪減二行差冪四因之亦同}，為第一亷；四之二行共，為第二亷；二步虚法；益積開之，得皇極弦二百八十九。\end{unitzh}
\begin{uniten}Another rule: Multiply the number of the area of the difference of the two walks by itself, and double it, as the dividend; add the two walking distances, with them multiply the area of the difference of the two walks, and quadruple it, as the augmenting follower; 4 portions of the area of the southward walk less 2 portions of the difference area, placed above; also 2 portions of the difference area less 4 portions of the area of the eastward walk; subtract the remainder from the upper position [Siku editors: adding the areas of the two walks, subtracting the area of the difference of the two walks, and multiplying it by 4 is also the same], as the first edge; 4 times the two walks combined, as the second edge; 2 bu as the empty divisor; extract it with augmented area, obtaining the Imperial Pole hypotenuse, 289.\end{uniten}

\begin{unitzh}\unitlabel{ceyuan-020.28}草曰：立天元一為皇極弦，以自之為弦冪，於上。以二行步相減，餘□，以自之，得□，為較冪，以減上，得丨□□，為二直積。復以天元除之，得□○□，為一个城徑也。副置之：上位加二之東行步，得□□□，為二勾也，以自增乘，得丨□□□□，為四段勾冪，於上；下位加二之南行，得□□□，為二股也，以自增乘，得丨□□□□，為四段股冪也。併入上位，得下式□□□□□，為四段弦冪{\zhsmall 寄左}。然後以天元為冪，四之，為同數。與左相消，得下式□□□□□，益積開三乘方，得二百八十九步，即皇極弦也。\end{unitzh}
\begin{uniten}Working: Set up the celestial unknown as the Imperial Pole hypotenuse; multiply it by itself as the hypotenuse area, placed above. Subtract the two walking distances from each other; the remainder, □, multiplied by itself, gives □ as the difference area (jiaomi); subtract it from the upper, obtaining 丨□□ as 2 rectangular areas. Again divide by the celestial unknown, obtaining □○□ as 1 city diameter. Place it in duplicate: to the upper position add 2 times the eastward walking distance, obtaining □□□ as 2 bases; multiply it by itself, obtaining 丨□□□□ as 4 portions of the base area, placed above; to the lower position add 2 times the southward walk, obtaining □□□ as 2 altitudes; multiply it by itself, obtaining 丨□□□□ as 4 portions of the altitude area. Add it into the upper position, obtaining the following form □□□□□ as 4 portions of the hypotenuse area [Note: deposit on the left]. Afterwards make the celestial unknown into an area and quadruple it, as the equal quantity. Eliminate it against the left, obtaining the following form □□□□□; extract the root of the thrice-multiplied power with augmented area, obtaining 289 bu, which is the Imperial Pole hypotenuse.\end{uniten}

\begin{unitzh}\unitlabel{ceyuan-020.29}欲見城徑者：别立天元半徑，副之，加東行為勾，加南行為股，勾股各為冪，併之，與弦冪相消，開方，得城徑也。\end{unitzh}
\begin{uniten}If one wishes to see the city diameter: separately set up the celestial unknown, the half-diameter, and duplicate it; add the eastward walk, as the base; add the southward walk, as the altitude; make the base and the altitude each into an area, add them, eliminate against the hypotenuse area, extract the root, and one obtains the city diameter.\end{uniten}

\begin{unitzh}\unitlabel{ceyuan-021.1}又法：以二行差一百一十九自乘，得一萬四千一百六十一，為差冪。以東行步乘之，得二十二萬六千五百七十六，為汎率。又自增乗，得五百一十三億三千六百六十八萬三千七百七十六，為五乘方實。倍東行步，得三十二，以二行差一百一十九乘之，得三千八百八，為小汎。以乘泛率，又倍之，得一十七億二千五百六十○萬二千八百一十六，為從方。併兩行而倍之，得三百二，以乘泛率，得六千八百四十二萬五千九百五十二於上位；以小泛冪一千四百五十萬○○八百六十四加入上位，共得八千二百九十二萬六千八百一十六，為第一亷。併兩行而倍之，得三百二，以乗小泛，得一百一十五萬○○一十六，為寄數。倍二行差以乘差冪，得三百三十七萬零三百一十八，内減寄數，餘二百二十二萬零三百零二，為第二亷。六段二行差冪八萬四千九百六十六，内減二行併數冪二萬二千八百一，餘六萬二千一百六十五，為第三益亷。六之二行差七百一十四，為第四益亷。二步虚法，得□弦三十四步。\end{unitzh}
\begin{uniten}Another rule: Multiply the difference of the two walks, 119, by itself, obtaining 14161 as the area (mi) of the difference. Multiply it by the eastward bu (步), obtaining 226576 as the general rate (fan lü). Again multiply it by itself, obtaining 51336683776 as the dividend (shi) of the fifth-multiplied power (wucheng fang). Double the eastward bu, obtaining 32; multiply it by the difference of the two walks, 119, obtaining 3808 as the small general (xiaofan). With it multiply the general rate, and again double, obtaining 1725602816 as the joined side (congfang). Combine the two walks and double, obtaining 302; with it multiply the general rate, obtaining 68425952, in the upper position; add the area of the small general, 14500864, into the upper position, altogether obtaining 82926816 as the first edge (lian). Combine the two walks and double, obtaining 302; with it multiply the small general, obtaining 1150016 as the stored number (jishu). Double the difference of the two walks and with it multiply the area of the difference, obtaining 3370318; subtract the stored number from it; the remainder, 2220302, is the second edge. 6 portions (duan) of the area of the difference of the two walks, 84966; subtract from it the area of the combined number of the two walks, 22801; the remainder, 62165, is the third augmenting edge (yilian). 6 times the difference of the two walks, 714, is the fourth augmenting edge. 2 bu as the empty divisor (xufa). Obtain the □ hypotenuse: 34 bu. [Translator's note: □ in "□ hypotenuse", "□ base", "□ altitude" and "□ slant" stands for a character not reproduced in the transcription.]\end{uniten}

\begin{unitzh}\unitlabel{ceyuan-021.2}草曰：立天元一為皇極弦上股弦差，{\zhsmall 即東行步上斜也，亦謂□斜。}以元加二行差，得□□，即明弦也。{\zhsmall 此即皇極弦上勾弦差也。}以天元乗之，又倍之，得□□□，即皇極内黄方冪也。{\zhsmall 泛寄。}置皇極弦上勾弦差，以東行步乘之，得□□，以天元除之，得□□，為明勾也。又置天元，以南行乘之，得□□，合用明弦除，不除，寄為母，便以此為□股於上。{\zhsmall 寄明弦母。}乃再置明勾，以明弦乘之，得□□□，亦為帶分明勾，加入上位，得□□□，即是一个虚弦也。以自增乘，得下式□□□□□，為一段虚弦冪也，内帶明弦冪分母。{\zhsmall 寄左。}然後置明弦，以自之，得丨□□，為明弦冪，以乘泛寄，得□□□□，為同數。與左相消，得下式□□□□□□□。開五乗方，得三十四步，為東行步上斜步也。{\zhsmall 即□弦。}其東行十六步即□勾也。勾弦各自為冪，以相減，餘九百步，開方得三十步，即□股也。既各得此數，乃以股外容圓半法求圓徑，得二百四十步，即城徑也。合問。\end{unitzh}
\begin{uniten}Working: Set up the celestial unknown (tianyuan) as the altitude–hypotenuse difference on the huangji hypotenuse [Note: that is, the slant on the eastward bu; it is also called the □ slant]. Add the difference of the two walks to the unknown (yuan), obtaining □ □, which is the ming hypotenuse [Note: this is the base–hypotenuse difference on the huangji hypotenuse]. [Translator's note: in the counting-rod arrays each □ stands for one row not reproduced in the transcription, and 1 renders the rod numeral 丨.] Multiply it by the celestial unknown and again double it, obtaining □ □ □, which is the area of the yellow square (huangfang) within the huangji [Note: general store (fanji)]. Set up the base–hypotenuse difference on the huangji hypotenuse and multiply it by the eastward bu, obtaining □ □; divide it by the celestial unknown, obtaining □ □ as the ming base. Again set up the celestial unknown and multiply it by the southward walk, obtaining □ □; it ought to be divided by the ming hypotenuse; not dividing, store it as the denominator, and simply take this as the □ altitude, on top [Note: storing the ming hypotenuse as denominator]. Then again set up the ming base and multiply it by the ming hypotenuse, obtaining □ □ □, which is likewise the ming base carrying a fraction (daifen); add it into the upper position, obtaining □ □ □, which is precisely one xu hypotenuse. Multiply it by itself, obtaining the following array □ □ □ □ □ as one portion of the area of the xu hypotenuse; within it the area of the ming hypotenuse is carried as denominator [Note: set aside on the left]. Then set up the ming hypotenuse and multiply it by itself, obtaining 1 □ □ as the area of the ming hypotenuse; with it multiply the general store, obtaining □ □ □ □ as the equal quantity (tongshu). Eliminate it against the left (xiangxiao), obtaining the following array □ □ □ □ □ □ □. Extract the root of the fifth-multiplied power, obtaining 34 bu, which is the slant bu on the eastward bu [Note: that is, the □ hypotenuse]. The eastward walk of 16 bu is the □ base. Make base and hypotenuse each into an area by itself and subtract one from the other; the remainder is 900 bu; extracting the square root gives 30 bu, which is the □ altitude. Having obtained each of these numbers, then by the half-method of the circle contained outside the altitude (gu wai rong yuan) find the circle diameter, obtaining 240 bu, which is the city diameter. This fits the question.\end{uniten}

\begin{unitzh}\unitlabel{ceyuan-021.3}按：此草又法求□弦，至開帶縱五乘方，法愈繁，數愈賾，而天元一之用愈見其妙。苐所得帶縱五乘方亷隅積數雖具，而未習其法者不能信其數之必然。今姑取已得之□弦數，按亷隅數推其積數，以明其數之無可疑焉。\end{unitzh}
\begin{uniten}Note: In this working of the other rule, finding the □ hypotenuse goes as far as extracting a fifth-multiplied power with attached extension (daizong); the method grows ever more complicated and the numbers ever more abstruse, and the marvel of the use of the celestial unknown is seen all the more. However, although the edge, corner (yu) and area numbers of the fifth-multiplied power with attached extension so obtained are all given, those who have not studied its method cannot believe that its numbers are necessarily so. Now let us for the moment take the number of the □ hypotenuse already obtained and, according to the edge and corner numbers, work out its area number, so as to make clear that its numbers leave no room for doubt.\end{uniten}

\begin{unitzh}\unitlabel{ceyuan-021.4}置五乘方數二，以□弦三十四乘之，得六十八，與四乘方數七百一十四相加，得七百八十二；又以□弦乘之，得二萬六千五百八十八，與三乘方數六萬二千一百六十五相加，得八萬八千七百五十三；又以□弦乘之，得三百零一萬七千六百零二，與立方數二百二十二萬零三百零二相加，得五百二十三萬七千九百零四；又以□弦乘之，得一億七千八百零八萬八千七百三十六，内減所少平方數八千二百九十二萬六千八百一十六，餘九千五百一十六萬一千九百二十；又以□弦乘之，得三十二億三千五百五十萬零五千二百八十，内減所少元數十七億二千五百六十萬零二千八百一十六，餘十五億零九百九十萬零二千四百六十四；又以□弦乗之，得五百一十三億三千六百六十八萬三千七百七十六，為積數，與草中積數合。{\zhsmall 此即無次商帶縱五乘方法。}\end{unitzh}
\begin{uniten}Set up the fifth-multiplied power number, 2; multiply it by the □ hypotenuse, 34, obtaining 68; add it to the fourth-multiplied power number, 714, obtaining 782. Again multiply by the □ hypotenuse, obtaining 26588; add it to the third-multiplied power number, 62165, obtaining 88753. Again multiply by the □ hypotenuse, obtaining 3017602; add it to the cube number, 2220302, obtaining 5237904. Again multiply by the □ hypotenuse, obtaining 178088736; subtract from it the deficient square number, 82926816; the remainder is 95161920. Again multiply by the □ hypotenuse, obtaining 3235505280; subtract from it the deficient yuan number, 1725602816; the remainder is 1509902464. Again multiply by the □ hypotenuse, obtaining 51336683776 as the area number; it agrees with the area number in the working. [Note: this is the method of the fifth-multiplied power with attached extension without a second quotient (cishang).]\end{uniten}

\begin{unitzh}\unitlabel{ceyuan-021.5}或問：出東門一十六步有樹，出南門東行七十二步見之。問答同前。\end{unitzh}
\begin{uniten}Someone asks: Going out of the east gate 16 bu there is a tree; going out of the south gate and walking east 72 bu, one sees it. Question and answer as before.\end{uniten}

\begin{unitzh}\unitlabel{ceyuan-021.6}法曰：二行步相減，得數以自之於上；又以出東門步自之，減上位，為平方實；二之出南門東行步，為益從；一步常法；翻開，得半徑。\end{unitzh}
\begin{uniten}Rule: Subtract the bu of the two walks from each other; multiply the resulting number by itself, on top. Again multiply the bu out of the east gate by itself and subtract it from the upper position, as the dividend of the square (pingfang shi). Twice the bu walked east out of the south gate is the augmenting joining term (yicong); 1 bu is the constant divisor (changfa); extract with reversal (fankai), obtaining the radius.\end{uniten}

\begin{unitzh}\unitlabel{ceyuan-021.7}草曰：别得人到樹即平弦也，半圓徑即平股也，其東行七十二步則平勾平弦差也。乃立天元一為半徑，加一十六，減七十二，得□□，為勾也。以自之，得丨□□，為勾冪；又加入天元股冪，得□□□，為弦冪。{\zhsmall 寄左。}再立天元一為半徑，加出東門步，得□□，即弦也。以自之，得丨□□，為同數。與左相消，得□□□，翻法開之，得一百二十步，即半城徑也。合問。\end{unitzh}
\begin{uniten}Working: Separately one obtains: the stretch from the person to the tree is the ping hypotenuse; half the circle diameter is the ping altitude; the eastward walk of 72 bu is then the difference of the ping base and the ping hypotenuse. Then set up the celestial unknown as the radius; add 16 and subtract 72, obtaining □ □ as the base. Multiply it by itself, obtaining 1 □ □ as the area of the base; again add in the celestial unknown's area of the altitude, obtaining □ □ □ as the area of the hypotenuse [Note: set aside on the left]. Again set up the celestial unknown as the radius; add the bu out of the east gate, obtaining □ □, which is the hypotenuse. Multiply it by itself, obtaining 1 □ □ as the equal quantity. Eliminate it against the left, obtaining □ □ □; extract it by the reversal method (fanfa), obtaining 120 bu, which is half the city diameter. This fits the question.\end{uniten}

\begin{unitzh}\unitlabel{ceyuan-021.8}或問：出南門一百三十五步有樹，出東門南行三十步見之。問答同前。\end{unitzh}
\begin{uniten}Someone asks: Going out of the south gate 135 bu there is a tree; going out of the east gate and walking south 30 bu, one sees it. Question and answer as before.\end{uniten}

\begin{unitzh}\unitlabel{ceyuan-021.9}法曰：樹去城步内減南行步，餘以為冪於上；又以樹去城步為冪，内減上位，為平實；倍樹去城步，為從；一虚隅；翻法，得半城徑。\end{unitzh}
\begin{uniten}Rule: From the bu of the tree's distance from the city subtract the southward bu; make the remainder into an area, on top. Again make the bu of the tree's distance from the city into an area, and subtract the upper position from it, as the square dividend (pingshi). Double the bu of the tree's distance from the city as the joining term (cong); 1 as the empty corner (xuyu); by the reversal method, obtain half the city diameter.\end{uniten}

\begin{unitzh}\unitlabel{ceyuan-021.10}草曰：别得人距樹即髙弦也，半圓徑即髙勾也，其南行三十步即髙弦上小差也。乃立天元一為半徑，加樹去城步為弦，内減小差□，得□□，即股也。以自之，得丨□□，為股冪；内加入天元冪，得□□□，為弦冪。{\zhsmall 寄左。}再置弦□□，自之，得丨□□，為同數。與左相消，得丨□□，翻開，得一百二十步，即半城徑也。合問。\end{unitzh}
\begin{uniten}Working: Separately one obtains: the person's distance from the tree is the gao hypotenuse; half the circle diameter is the gao base; the southward walk of 30 bu is the small difference (xiaocha) on the gao hypotenuse. Then set up the celestial unknown as the radius; add the bu of the tree's distance from the city, as the hypotenuse; subtract from it the small difference □, obtaining □ □, which is the altitude. Multiply it by itself, obtaining 1 □ □ as the area of the altitude; add into it the area of the celestial unknown, obtaining □ □ □ as the area of the hypotenuse [Note: set aside on the left]. Again set up the hypotenuse □ □ and multiply it by itself, obtaining 1 □ □ as the equal quantity. Eliminate it against the left, obtaining 1 □ □; extract with reversal, obtaining 120 bu, which is half the city diameter. This fits the question.\end{uniten}

\begin{unitzh}\unitlabel{ceyuan-021.11}或問：乙出東門，不知逺近而立；甲出南門東行七十二步，望見乙，就乙斜行一百三十六步，與乙相會。問答同前。\end{unitzh}
\begin{uniten}Someone asks: B (乙) goes out of the east gate, how far is not known, and stands; A (甲) goes out of the south gate and walks east 72 bu, sees B in the distance, goes toward B walking obliquely 136 bu, and meets B. Question and answer as before.\end{uniten}

\begin{unitzh}\unitlabel{ceyuan-021.12}法曰：以斜行步自之於上，以二行相減餘自為冪，減上位，為平實；從空；一步常法；得半徑。\end{unitzh}
\begin{uniten}Rule: Multiply the bu of the oblique walk by itself, on top; make the remainder from subtracting the two walks from each other into an area by itself, and subtract it from the upper position, as the square dividend; the joining term is empty; 1 bu as the constant divisor; obtain the radius.\end{uniten}

\begin{unitzh}\unitlabel{ceyuan-021.13}草曰：别得七十二步即大差也，斜行即弦，半徑即股也。立天元一為半徑，以自之，為股冪；又以二行差六十四以自之，得□，為勾冪；併二冪，得丨□□，為弦冪。{\zhsmall 寄左。}然後以斜行步自之，得□，為同數。與左相消，得丨□□，開平方，得一百二十步，倍之即城徑也。合問。\end{unitzh}
\begin{uniten}Working: Separately one obtains: the 72 bu is the large difference (dacha); the oblique walk is the hypotenuse; the radius is the altitude. Set up the celestial unknown as the radius and multiply it by itself, as the area of the altitude; again multiply the difference of the two walks, 64, by itself, obtaining □ as the area of the base; combine the two areas, obtaining 1 □ □ as the area of the hypotenuse [Note: set aside on the left]. Then multiply the bu of the oblique walk by itself, obtaining □ as the equal quantity. Eliminate it against the left, obtaining 1 □ □; extract the square root, obtaining 120 bu; doubling it gives the city diameter. This fits the question.\end{uniten}

\begin{unitzh}\unitlabel{ceyuan-021.14}或問：甲出南門，不知逺近而立；乙出東門南行三十步，望見甲，却就甲斜行二百五十五步，與甲相㑹。問答同前。\end{unitzh}
\begin{uniten}Someone asks: A goes out of the south gate, how far is not known, and stands; B goes out of the east gate and walks south 30 bu, sees A in the distance, then goes toward A walking obliquely 255 bu, and meets A. Question and answer as before.\end{uniten}

\begin{unitzh}\unitlabel{ceyuan-021.15}法曰：二行差自之為冪，以減於斜行冪，為平實；一虚隅；得半徑。\end{unitzh}
\begin{uniten}Rule: Multiply the difference of the two walks by itself as an area, and subtract it from the area of the oblique walk, as the square dividend; 1 as the empty corner; obtain the radius.\end{uniten}

\begin{unitzh}\unitlabel{ceyuan-021.16}草曰：别得南行步即股弦差也，斜步即弦也，半徑即勾也。乃立天元一為半城徑，以自之，為冪；以二行相減餘二百二十五，以自之，得□，為股冪；二冪相併，得丨□□，為弦冪。{\zhsmall 寄左。}然後以斜行自之，得□，為同數。與左相消，得下丨□□，開平方，得一百二十步，即半徑也。合問。\end{unitzh}
\begin{uniten}Working: Separately one obtains: the southward bu is the difference of altitude and hypotenuse; the oblique bu is the hypotenuse; the radius is the base. Then set up the celestial unknown as half the city diameter and multiply it by itself, as an area; subtract the two walks from each other, the remainder being 225, and multiply it by itself, obtaining □ as the area of the altitude; combine the two areas, obtaining 1 □ □ as the area of the hypotenuse [Note: set aside on the left]. Then multiply the oblique walk by itself, obtaining □ as the equal quantity. Eliminate it against the left, obtaining below 1 □ □; extract the square root, obtaining 120 bu, which is the radius. This fits the question.\end{uniten}

\begin{unitzh}\unitlabel{ceyuan-021.17}或問：甲出南門東行，不知步數而立；乙出東門南行三十步，望見甲，斜行一百二步相㑹。問答同前。\end{unitzh}
\begin{uniten}Someone asks: A goes out of the south gate and walks east, the number of bu not known, and stands; B goes out of the east gate and walks south 30 bu, sees A in the distance, walks obliquely 102 bu, and they meet. Question and answer as before.\end{uniten}

\begin{unitzh}\unitlabel{ceyuan-021.18}法曰：二行相乘，四之於上，又加入斜行冪，為平實，得虚和一百三十八。\end{unitzh}
\begin{uniten}Rule: Multiply the two walks together and multiply by 4, on top; again add in the area of the oblique walk, as the square dividend; obtain the xu sum (xuhe), 138.\end{uniten}

\begin{unitzh}\unitlabel{ceyuan-021.19}草曰：别得斜步内減南行，為甲東行步也。此問以弦外容圓入之。以二行相減數乘乙南行三十步，得□，又四之，得□，為二直積也；又加入斜步冪□，共得□，即和冪也。平方而一，得一百三十八步，即虚和也。又加斜步，得二百四十步，即城徑也。合問。\end{unitzh}
\begin{uniten}Working: Separately one obtains: the oblique bu minus the southward walk is A's eastward walk in bu. This question is entered by way of the circle contained outside the hypotenuse (xian wai rong yuan). With the number from subtracting the two walks multiply B's southward walk of 30 bu, obtaining □; again multiply by 4, obtaining □, as two rectangle areas (zhiji); again add in the area of the oblique bu, □, altogether obtaining □, which is the area of the sum. Taking its square root (pingfang er yi) gives 138 bu, which is the xu sum. Again add the oblique bu, obtaining 240 bu, which is the city diameter. This fits the question.\end{uniten}

\begin{unitzh}\unitlabel{ceyuan-021.20}或問：乙出東門南行，不知步數而立；甲出南門東行七十二步，望見乙，斜行一百二步，與乙相㑹。問答同前。\end{unitzh}
\begin{uniten}Someone asks: B goes out of the east gate and walks south, the number of bu not known, and stands; A goes out of the south gate and walks east 72 bu, sees B in the distance, walks obliquely 102 bu, and meets B. Question and answer as before.\end{uniten}

\begin{unitzh}\unitlabel{ceyuan-021.21}法曰：倍相減步，以乘倍東行，得數復以減於斜步冪，餘為實，平方而一，得較也。又以二行相減數乘倍東行，為平實；以較為從方；得勾。勾較共為長。又以斜步併入勾股共，即城徑。\end{unitzh}
\begin{uniten}Rule: Double the bu from the subtraction and with it multiply twice the eastward walk; the resulting number, subtract it in turn from the area of the oblique bu; the remainder is the dividend; taking its square root gives the difference (jiao). Again with the number from subtracting the two walks multiply twice the eastward walk, as the square dividend; take the difference as the joined side; obtain the base. The base and the difference together make the length (chang). Again combine the oblique bu into the sum of base and altitude: this is the city diameter.\end{uniten}

\begin{unitzh}\unitlabel{ceyuan-021.22}草曰：别得二行相減餘□，為乙南行步也。以此數又減於甲東行，餘四十二步，即較也。乃以二行相減數□乘倍東行，得□，為平實；以較為從；平方開，得四十八，即勾也。勾内加較，得九十步，即股也。勾股共得一百三十八，又加入斜步，共得二百四十步，即城徑也。合問。\end{unitzh}
\begin{uniten}Working: Separately one obtains: the remainder from subtracting the two walks, □, is B's southward walk in bu. Again subtract this number from A's eastward walk; the remainder, 42 bu, is the difference. Then with the number from subtracting the two walks, □, multiply twice the eastward walk, obtaining □ as the square dividend; take the difference as the joining term; extracting the square root gives 48, which is the base. Add the difference to the base, obtaining 90 bu, which is the altitude. Base and altitude together give 138; again add in the oblique bu, altogether obtaining 240 bu, which is the city diameter. This fits the question.\end{uniten}

\begin{unitzh}\unitlabel{ceyuan-021.23}或問：乙出南門東行，甲出東門南行，兩相望見。既而乙云：「我東行不及城徑一百六十八步。」甲云：「我南行不及城徑二百一十。」問答同前。\end{unitzh}
\begin{uniten}Someone asks: B goes out of the south gate and walks east; A goes out of the east gate and walks south; the two see each other. Thereupon B says: "My eastward walk falls short of the city diameter by 168 bu." A says: "My southward walk falls short of the city diameter by 210." Question and answer as before.\end{uniten}

\begin{unitzh}\unitlabel{ceyuan-021.24}法曰：半甲不及步，以自之為冪；半甲不及步内減云數差，以自之為冪；二冪相併，内却減差冪，為平實；二之乙不及，為益從；三步半虚法；得甲南行。\end{unitzh}
\begin{uniten}Rule: Halve A's shortfall in bu and multiply it by itself as an area; from half of A's shortfall in bu subtract the difference of the stated numbers, and multiply it by itself as an area; combine the two areas, and from them in turn subtract the area of the difference, as the square dividend; twice B's shortfall as the augmenting joining term; 3 1/2 bu as the empty divisor; obtain A's southward walk.\end{uniten}

\begin{unitzh}\unitlabel{ceyuan-021.25}草曰：别得乙不及為虚勾半徑共，又為徑内減明勾也；甲不及為虛股半徑共，又為徑内減□股也；又二云數相併，為虚和圓徑共也；云數相減，即虚較也。乃立天元一為甲南行，以減於甲不及步，又半之，得□□，為虚股也。虚股内減虚較，得□□，為虚勾。勾自之，得□□□，為勾冪也；又股自之，得下式□□□，為股冪也；二冪相併，得□□□，為弦冪。{\zhsmall 寄左。}然後以天元加虚較，得□□，為乙東行；又加入天元甲南行，得□□，為虚弦。以自之，得□□□，為同數。與左相消，得□□□，開平方，得三十步，即甲南行也。内加少步，即城徑也。合問。\end{unitzh}
\begin{uniten}Working: Separately one obtains: B's shortfall is the xu base and the radius together, and is also the diameter minus the ming base; A's shortfall is the xu altitude and the radius together, and is also the diameter minus the □ altitude; again, the two stated numbers combined are the xu sum and the circle diameter together; the stated numbers subtracted from each other give the xu difference (xujiao). Then set up the celestial unknown as A's southward walk; subtract it from A's shortfall in bu, and again halve it, obtaining □ □ as the xu altitude. From the xu altitude subtract the xu difference, obtaining □ □ as the xu base. Multiply the base by itself, obtaining □ □ □ as the area of the base; again multiply the altitude by itself, obtaining the following array □ □ □ as the area of the altitude; combine the two areas, obtaining □ □ □ as the area of the hypotenuse [Note: set aside on the left]. Then add the xu difference to the celestial unknown, obtaining □ □ as B's eastward walk; again add in the celestial unknown, A's southward walk, obtaining □ □ as the xu hypotenuse. Multiply it by itself, obtaining □ □ □ as the equal quantity. Eliminate it against the left, obtaining □ □ □; extract the square root, obtaining 30 bu, which is A's southward walk. Add to it the lacking bu; this is the city diameter. This fits the question.\end{uniten}

\begin{unitzh}\unitlabel{ceyuan-022.1}或問：丙出南門直行，甲出東門直行，兩相望見。既而丙云：「我行少於城徑一百五步。」甲云：「我行少於城徑二百二十四步。」問答同前。\end{unitzh}
\begin{uniten}Someone asks: Bing goes out of the south gate and walks straight on; Jia goes out of the east gate and walks straight on; the two see each other. Then Bing says: "My walk is less than the city diameter by 105 bu (步)." Jia says: "My walk is less than the city diameter by 224 bu." The question and answer are the same as before.\end{uniten}

\begin{unitzh}\unitlabel{ceyuan-022.2}法曰：二少歩相乘訖，又自乗為實；六之共步，乘云數相乘數，為益從；十八之云數相乘數，於上；又三之共步自乘，加上位，内復減丙少步冪、甲少步冪，為從亷；四十八之共步，為益二亷；六十三步，常法；翻法開三乗方，得一百二十步，即半徑。\end{unitzh}
\begin{uniten}Rule: Multiply the two deficiency bu together; when this is done, multiply the result by itself again; this is the dividend (shi). 6 times the combined bu, multiplied by the product of the stated numbers, is the subtractive following term (yi cong). 18 times the product of the stated numbers is placed at the top. Further, multiply 3 times the combined bu by itself and add it to the upper position; from this again subtract the area (mi) of Bing's deficiency bu and the area of Jia's deficiency bu; this is the following edge (cong lian). 48 times the combined bu is the subtractive second edge (yi er lian). 63 bu is the constant divisor (chang fa). By the reversing method (fan fa), extract the root of the thrice-multiplied power (kai san cheng fang), obtaining 120 bu, which is the half diameter.\end{uniten}

\begin{unitzh}\unitlabel{ceyuan-022.3}草曰：别得云數共，減於倍城徑，為甲丙共數；又云數相減，即皇極差，亦為甲行不及丙行數。立天元一為半城徑，以三之，副置二位：上位減丙少步，得□□，為皇極股也；下位減甲少步，得□□，為皇極勾也。勾股相乘，得□□□，以天元除之，得□□□，為弦也。弦自之，得□□□□□，為弦冪{\zhsmall 寄左}。然後以股自之，得下□□□，為股冪，於上；又以勾自之，得□□□，為勾冪，併以加入上位，得□□□，為同數，與左相消，得□□□□□，翻法開三乘方，得一百二十步，即半城徑也。合問。\end{unitzh}
\begin{uniten}Working: Separately we obtain: the stated numbers combined, subtracted from twice the city diameter, give the combined number of Jia and Bing. Further, the stated numbers subtracted from each other give the Supreme Pole difference (huangji cha), which is also the number by which Jia's walk falls short of Bing's walk. Set up the celestial unknown (tianyuan) as half the city diameter; multiply it by 3 and set it down in duplicate in two positions. From the upper position subtract Bing's deficiency bu, obtaining □□ [Translator's note: here and below, □ marks an array or a character not reproduced in the transcription], which is the Supreme Pole altitude (gu). From the lower position subtract Jia's deficiency bu, obtaining □□, which is the Supreme Pole base (gou). Multiply the base and the altitude together, obtaining □□□; divide it by the celestial unknown, obtaining □□□, which is the hypotenuse (xian). Multiply the hypotenuse by itself, obtaining □□□□□, which is the area of the hypotenuse [Note: set aside on the left (ji zuo)]. Then multiply the altitude by itself, obtaining, below, □□□, which is the area of the altitude, placed at the top. Further multiply the base by itself, obtaining □□□, which is the area of the base, and add it into the upper position, obtaining □□□, which is the equal quantity (tong shu). Cancel it against the left (xiang xiao), obtaining □□□□□. By the reversing method extract the root of the thrice-multiplied power, obtaining 120 bu, which is half the city diameter. This matches the question.\end{uniten}

\begin{unitzh}\unitlabel{ceyuan-022.4}或問：甲出東門直行，丙出南門直行，各不知步數而立。乙望見甲，就甲斜行了二百八十九步，與甲相㑹。其二直行共一百五十一步。問答同前。\end{unitzh}
\begin{uniten}Someone asks: Jia goes out of the east gate and walks straight on; Bing goes out of the south gate and walks straight on; each stops, the number of bu not being known. Yi sees Jia and walks obliquely toward Jia 289 bu, meeting Jia. Their two straight walks together are 151 bu. The question and answer are the same as before.\end{uniten}

\begin{unitzh}\unitlabel{ceyuan-022.5}法曰：斜冪内減共步冪，為平實；倍共步，内減斜步，為從；一，常法。得徑。\end{unitzh}
\begin{uniten}Rule: From the area of the oblique subtract the area of the combined bu; this is the square dividend (ping shi). Double the combined bu and subtract the oblique bu from it; this is the following term (cong). 1 is the constant divisor. This gives the diameter.\end{uniten}

\begin{unitzh}\unitlabel{ceyuan-022.6}草曰：别得共數、城徑併，即皇極和也。立天元一為圓徑，加共步，得□□，為皇極和。以自之，得丨□□，於上。以斜行冪□減上位，餘丨□□，為二直積{\zhsmall 寄左}。然後以天元乘斜步，得□□，與左相消，得丨□□，開平方，得二百四十步，即城徑也。合問。\end{unitzh}
\begin{uniten}Working: Separately we obtain: the combined number and the city diameter added together are the Supreme Pole sum (huangji he). Set up the celestial unknown as the diameter of the circle; add the combined bu, obtaining □□, which is the Supreme Pole sum. Multiply it by itself, obtaining 丨□□, placed at the top. Subtract the area of the oblique walk, □, from the upper position, leaving 丨□□, which is two rectangle areas (er zhi ji) [Note: set aside on the left]. Then multiply the oblique bu by the celestial unknown, obtaining □□. Cancel it against the left, obtaining 丨□□. Extract the square root, obtaining 240 bu, which is the city diameter. This matches the question.\end{uniten}

\begin{unitzh}\unitlabel{ceyuan-022.7}或問：甲出東門直行，乙出東門南行，丙出南門直行，丁出南門東行，各不知步數而立。四人遥相望，悉與城叅相直。只云甲丙共行了一百五十一步，乙丁立處相距一百二步。又云丙直行步多於甲直行步。問答同前。\end{unitzh}
\begin{uniten}Someone asks: Jia goes out of the east gate and walks straight on; Yi goes out of the east gate and walks south; Bing goes out of the south gate and walks straight on; Ding goes out of the south gate and walks east; each stops, the number of bu not being known. The four people look at one another from afar, all being in alignment with the city. It is only said that Jia and Bing together walked 151 bu, and that the places where Yi and Ding stand are 102 bu apart. It is further said that Bing's straight walk in bu is more than Jia's straight walk in bu. The question and answer are the same as before.\end{uniten}

\begin{unitzh}\unitlabel{ceyuan-022.8}法曰：共步、距步相減，得數，自之，於上；以共步為冪，内減上，為平實；二之距步，内減共步距步差，為從；一步，虚法。得城徑。\end{unitzh}
\begin{uniten}Rule: Subtract the combined bu and the distance bu from each other; multiply the number obtained by itself and place it at the top. Make the combined bu into an area and subtract the top from it; this is the square dividend. Take 2 times the distance bu and subtract from it the difference between the combined bu and the distance bu; this is the following term. 1 bu is the empty divisor (xu fa). This gives the city diameter.\end{uniten}

\begin{unitzh}\unitlabel{ceyuan-022.9}草曰：别得共步得城徑，即皇極和也；相距步，即虚弦也；皇極和内減虚弦，即皇極弦也；又共步距步差□，即皇極弦内減城徑也{\zhsmall 此名旁差}。乃立天元一為城徑，加共步，得□□，為皇極和也。以自之，得丨□□，於上。以共步距步差□加天元，得□□，為皇極弦也。以自之，得下式丨□□，減上位，餘得□□，為二直積{\zhsmall 寄左}。然後以天元徑乘皇極弦，得丨□，為同數，與左相消，得丨□□，開平方，得二百四十步，即城徑也。合問。\end{unitzh}
\begin{uniten}Working: Separately we obtain: the combined bu, gaining the city diameter, is the Supreme Pole sum. The distance bu is the empty hypotenuse (xu xian). The Supreme Pole sum minus the empty hypotenuse is the Supreme Pole hypotenuse (huangji xian). Further, the difference between the combined bu and the distance bu, □, is the Supreme Pole hypotenuse minus the city diameter [Note: this is called the side difference (pang cha)]. Then set up the celestial unknown as the city diameter; add the combined bu, obtaining □□, which is the Supreme Pole sum. Multiply it by itself, obtaining 丨□□, placed at the top. Add the difference between the combined bu and the distance bu, □, to the celestial unknown, obtaining □□, which is the Supreme Pole hypotenuse. Multiply it by itself, obtaining the array below, 丨□□; subtract it from the upper position, the remainder obtained being □□, which is two rectangle areas [Note: set aside on the left]. Then multiply the Supreme Pole hypotenuse by the celestial-unknown diameter, obtaining 丨□, which is the equal quantity. Cancel it against the left, obtaining 丨□□. Extract the square root, obtaining 240 bu, which is the city diameter. This matches the question.\end{uniten}

\begin{unitzh}\unitlabel{ceyuan-022.10}或問：甲出南門東行，不知步數而立。乙出東門南行，望見甲，復就甲斜行，與甲相㑹。乙通計行了一百三十二步，其乙南行步不及斜行七十二步，其甲東行多於乙南行。問答同前。\end{unitzh}
\begin{uniten}Someone asks: Jia goes out of the south gate, walks east, and stops, the number of bu not being known. Yi goes out of the east gate and walks south, sees Jia, and then walks obliquely toward Jia, meeting Jia. Yi walked 132 bu in all; Yi's southward walk in bu falls short of the oblique walk by 72 bu; and Jia's eastward walk is more than Yi's southward walk. The question and answer are the same as before.\end{uniten}

\begin{unitzh}\unitlabel{ceyuan-022.11}法曰：倍不及步，在地；以不及步減通步，以乗之，為實；以四之不及步為法。得乙南行三十步。\end{unitzh}
\begin{uniten}Rule: Double the shortfall bu and set it on the ground (zai di). Subtract the shortfall bu from the total bu, and with this multiply the former; this is the dividend. Take 4 times the shortfall bu as the divisor (fa). This gives Yi's southward walk, 30 bu.\end{uniten}

\begin{unitzh}\unitlabel{ceyuan-022.12}草曰：别得乙南行，即□股也；以減通步，即虚弦也；以減不及步，即虚較也；其不及步，即甲東行也。立天元一為乙南行，置不及步，以天元乘之，又四之，得□，為二直積{\zhsmall 寄左}。然後倍不及步，以為弦較和，於上□；以不及步減通步，得□，為弦較較；以乗上位，得□，為同數，與左相消，得□□，上法下實，得三十步，為乙南行也。餘各以數求之。\end{unitzh}
\begin{uniten}Working: Separately we obtain: Yi's southward walk is the □ altitude. Subtracting it from the total bu gives the empty hypotenuse. Subtracting it from the shortfall bu gives the empty difference (xu jiao). The shortfall bu is Jia's eastward walk. Set up the celestial unknown as Yi's southward walk. Set down the shortfall bu, multiply it by the celestial unknown and further by 4, obtaining □, which is two rectangle areas [Note: set aside on the left]. Then double the shortfall bu to make the hypotenuse-difference sum (xian jiao he), placed at the top, □. Subtract the shortfall bu from the total bu, obtaining □, which is the hypotenuse-difference difference (xian jiao jiao). Multiply the upper position by it, obtaining □, which is the equal quantity. Cancel it against the left, obtaining □□. The upper is the divisor and the lower the dividend; this gives 30 bu, which is Yi's southward walk. The rest are each found from the numbers.\end{uniten}

\begin{unitzh}\unitlabel{ceyuan-022.13}又法：别得通行步，為兩个乙南行、一个甲東行共也；其不及步，即東行步也。云步相併，即兩个虚弦；相減，即兩个乙南行也。\end{unitzh}
\begin{uniten}Another rule: Separately we obtain: the total walked bu is two of Yi's southward walk and one of Jia's eastward walk together; the shortfall bu is the eastward-walk bu. The stated bu added together are two empty hypotenuses; subtracted from each other, they are two of Yi's southward walk.\end{uniten}

\begin{unitzh}\unitlabel{ceyuan-022.14}或問：甲出南門東行，不知步數而立。乙出東門南行，望見甲，復斜行，與甲相㑹。二人共行了二百四步。又云：甲行不及乙一百三十二。{\zhsmall 按：甲不及乙六十步，非一百三十二步，當云甲行不及共步，方合。}問答同前。\end{unitzh}
\begin{uniten}Someone asks: Jia goes out of the south gate, walks east, and stops, the number of bu not being known. Yi goes out of the east gate and walks south, sees Jia, and then walks obliquely, meeting Jia. The two people together walked 204 bu. It is further said: Jia's walk falls short of Yi's by 132. [Siku editors: Jia falls short of Yi by 60 bu, not by 132 bu; it ought to say "Jia's walk falls short of the combined bu", and then it would agree.] The question and answer are the same as before.\end{uniten}

\begin{unitzh}\unitlabel{ceyuan-022.15}法曰：别得二行共，即兩个虛弦也；其不及步，即乙南行與一虚弦共也。置不及步，内減一弦，餘三十步，即乙南行也。以乙南行反以減虚弦，餘七十二步，即甲東行也。以乙南行減甲東行，餘即虚較也。此問無草。\end{unitzh}
\begin{uniten}Rule: Separately we obtain: the two walks combined are two empty hypotenuses; the shortfall bu is Yi's southward walk together with one empty hypotenuse. Set down the shortfall bu and subtract one hypotenuse from it, leaving 30 bu, which is Yi's southward walk. In turn, subtract Yi's southward walk from the empty hypotenuse, leaving 72 bu, which is Jia's eastward walk. Subtract Yi's southward walk from Jia's eastward walk; the remainder is the empty difference. This question has no working.\end{uniten}

\begin{unitzh}\unitlabel{ceyuan-022.16}按：右二問語若淺近，然以發明加減乘除相通之義，最為深切。集中倣此者，可類推之。\end{unitzh}
\begin{uniten}Note: The wording of the two questions on the right may seem shallow and simple, but for bringing to light the principle that addition, subtraction, multiplication and division pass into one another, they are most profound and incisive. Those in the collection that imitate these can be inferred by analogy.\end{uniten}

\begin{unitzh}\unitlabel{ceyuan-022.17}或問：乙出東門南行，甲出西門南行，甲望見乙，斜行五百一十步，相㑹。乙云：「我南行少於城徑二百一十步。」問答同前。\end{unitzh}
\begin{uniten}Someone asks: Yi goes out of the east gate and walks south; Jia goes out of the west gate and walks south; Jia sees Yi and walks obliquely 510 bu, and they meet. Yi says: "My southward walk is less than the city diameter by 210 bu." The question and answer are the same as before.\end{uniten}

\begin{unitzh}\unitlabel{ceyuan-022.18}法曰：少步冪為平實；四斜步，内減二少步，為益從；五步，常法。得乙南行。\end{unitzh}
\begin{uniten}Rule: The area of the deficiency bu is the square dividend. Take 4 times the oblique bu and subtract 2 times the deficiency bu from it; this is the subtractive following term. 5 bu is the constant divisor. This gives Yi's southward walk.\end{uniten}

\begin{unitzh}\unitlabel{ceyuan-022.19}草曰：别得少步為徑内減叀股。立天元一為乙南行，以二之，減於倍斜行步，得□□，為梯底也。以二之天元乘之，得□□，為徑冪{\zhsmall 寄左}。再置天元，加少步，得下式□□，為城徑，以自之，得丨□□，與左相消，得□□□，開平方，得三十步，即乙南行也。加少步，即城徑也。合問。\end{unitzh}
\begin{uniten}Working: Separately we obtain: the deficiency bu is the diameter minus the zhuan altitude (zhuan gu). Set up the celestial unknown as Yi's southward walk; multiply it by 2 and subtract it from twice the oblique walk in bu, obtaining □□, which is the trapezoid base (ti di). Multiply it by 2 times the celestial unknown, obtaining □□, which is the area of the diameter [Note: set aside on the left]. Again set down the celestial unknown and add the deficiency bu, obtaining the array below, □□, which is the city diameter. Multiply it by itself, obtaining 丨□□. Cancel it against the left, obtaining □□□. Extract the square root, obtaining 30 bu, which is Yi's southward walk. Adding the deficiency bu gives the city diameter. This matches the question.\end{uniten}

\begin{unitzh}\unitlabel{ceyuan-022.20}或問：乙出南門東行，甲出北門東行，甲望見乙，斜行二百七十二步，與乙相㑹。乙云：「我東行不及城徑一百六十八步。」問答同前。\end{unitzh}
\begin{uniten}Someone asks: Yi goes out of the south gate and walks east; Jia goes out of the north gate and walks east; Jia sees Yi, walks obliquely 272 bu, and meets Yi. Yi says: "My eastward walk falls short of the city diameter by 168 bu." The question and answer are the same as before.\end{uniten}

\begin{unitzh}\unitlabel{ceyuan-022.21}法曰：以不及步冪之為實；四斜，内減二之不及步，為虚從；五，常法；平實開，得乙東行七十二。\end{unitzh}
\begin{uniten}Rule: Take the shortfall bu and form its area as the dividend. Take 4 [times] the oblique and subtract 2 times the shortfall bu from it; this is the empty following term (xu cong). 5 is the constant divisor. Extracting the root with the square dividend gives Yi's eastward walk, 72.\end{uniten}

\begin{unitzh}\unitlabel{ceyuan-022.22}草曰：别得不及步，為城徑減明勾也。立天元一為乙東行，以倍之，減於二之斜行步，得下□□，為梯底也。倍天元乘之，得□□，為徑冪{\zhsmall 寄左}。再置天元，加不及步，得□□，為城徑，以自之，得丨□□，為同數，與左相消，得□□□，開平方，得七十二步，即乙東行也。加入少步，即城徑也。合問。\end{unitzh}
\begin{uniten}Working: Separately we obtain: the shortfall bu is the city diameter minus the bright base (ming gou). Set up the celestial unknown as Yi's eastward walk; double it and subtract it from 2 times the oblique walk in bu, obtaining, below, □□, which is the trapezoid base. Multiply it by twice the celestial unknown, obtaining □□, which is the area of the diameter [Note: set aside on the left]. Again set down the celestial unknown and add the shortfall bu, obtaining □□, which is the city diameter. Multiply it by itself, obtaining 丨□□, which is the equal quantity. Cancel it against the left, obtaining □□□. Extract the square root, obtaining 72 bu, which is Yi's eastward walk. Adding in the deficiency bu gives the city diameter. This matches the question.\end{uniten}

\begin{unitzh}\unitlabel{ceyuan-022.23}或問：乙出南門東行，丁出東門南行，却有甲、丙二人共在西北隅，甲向東行，丙向南行。四人遥相望見，俱與城叅相直。既而相㑹。甲云：「我多乙二百四十八步。」丙云：「我多於丁五百七十步。」問答同前。\end{unitzh}
\begin{uniten}Someone asks: Yi goes out of the south gate and walks east; Ding goes out of the east gate and walks south; and there are, moreover, two people, Jia and Bing, together at the northwest corner; Jia walks toward the east and Bing walks toward the south. The four people see one another from afar, all being in alignment with the city. Then they meet. Jia says: "I [walked] more than Yi by 248 bu." Bing says: "I [walked] more than Ding by 570 bu." The question and answer are the same as before.\end{uniten}

\begin{unitzh}\unitlabel{ceyuan-022.24}法曰：二多步相乗，為平實；併二多步而半之，為從；七分半，常法。得城徑。\end{unitzh}
\begin{uniten}Rule: Multiply the two excess bu together; this is the square dividend. Add the two excess bu and halve; this is the following term. 7 1/2 fen (分) is the constant divisor. This gives the city diameter.\end{uniten}

\begin{unitzh}\unitlabel{ceyuan-022.25}草曰：别得甲多步，為大勾内減明勾也；丙多步，為大股内少叀股也。又乙東行得一虚勾為半徑，丁南行得一虚股為半徑。又二多數相併，得□，為大和内少虚弦也。又二多數相減，餘□，為兩个角差。又甲多步内減半徑，即勾方差也；丙多步内減半徑，即股方差也。立天元一為城徑，以半之，減於甲多步，得□□，為勾方差；又以半徑減於丙多步，得□□，為股方差。二差相乘，得□□□，為徑冪{\zhsmall 寄左}。然後以天元冪與左相消，得下式□□□，開平方，得二百四十步，即城徑也。合問。\end{unitzh}
\begin{uniten}Working: Separately we obtain: Jia's excess bu is the great base (da gou) minus the bright base; Bing's excess bu is the great altitude (da gu) less the zhuan altitude. Further, Yi's eastward walk, gaining one empty base (xu gou), is the half diameter; Ding's southward walk, gaining one empty altitude (xu gu), is the half diameter. Further, the two excess numbers added together give □, which is the great sum (da he) less the empty hypotenuse. Further, the two excess numbers subtracted from each other leave □, which is two corner differences (jiao cha). Further, Jia's excess bu minus the half diameter is the base-square difference (gou fang cha); Bing's excess bu minus the half diameter is the altitude-square difference (gu fang cha). Set up the celestial unknown as the city diameter; halve it and subtract it from Jia's excess bu, obtaining □□, which is the base-square difference. Further, subtract the half diameter from Bing's excess bu, obtaining □□, which is the altitude-square difference. Multiply the two differences together, obtaining □□□, which is the area of the diameter [Note: set aside on the left]. Then cancel the area of the celestial unknown against the left, obtaining the array below, □□□. Extract the square root, obtaining 240 bu, which is the city diameter. This matches the question.\end{uniten}

\begin{unitzh}\unitlabel{ceyuan-022.26}或問：甲、丙二人俱在西北隅，甲向東行，丙向南行；又乙出南門東行，丁出東門南行，各不知步數而立。四人遥相望見，悉與城叅相直。既而相㑹。甲云：「我與乙共行了三百九十二步。」丙云：「我與丁共行六百三十步。」問答同前。\end{unitzh}
\begin{uniten}Someone asks: Two people, Jia and Bing, are both at the northwest corner; Jia walks toward the east and Bing walks toward the south. Further, Yi goes out of the south gate and walks east, and Ding goes out of the east gate and walks south; each stops, the number of bu not being known. The four people see one another from afar, all being in alignment with the city. Then they meet. Jia says: "Yi and I together walked 392 bu." Bing says: "Ding and I together walked 630 bu." The question and answer are the same as before.\end{uniten}

\begin{unitzh}\unitlabel{ceyuan-022.27}法曰：甲乙共自之為冪，丙丁共自之為冪，二冪又相乘，為三乘方實。甲乙共自之為冪，以丙丁共乘之，於上；又以丙丁共自之為冪，以甲乙共乘之，加上位，為益從。甲乙共自之為冪，丙丁共自之為冪，併，以七分半乘之，於上；又以甲乙共乘丙丁共，得數，減上位，為第一益亷。併二共數，以七分半乘之，為第二亷。以七分半自之，得五分六釐二毫五絲，於上位；以一步内減上位，餘四分三釐七毫五絲，為虚隅。得城徑。\end{unitzh}
\begin{uniten}Rule: Multiply Jia and Yi combined by itself to make an area; multiply Bing and Ding combined by itself to make an area; multiply the two areas together again; this is the dividend of the thrice-multiplied power. Multiply Jia and Yi combined by itself to make an area and multiply it by Bing and Ding combined, placing it at the top. Further, multiply Bing and Ding combined by itself to make an area, multiply it by Jia and Yi combined, and add it to the upper position; this is the subtractive following term. Multiply Jia and Yi combined by itself to make an area, and multiply Bing and Ding combined by itself to make an area; add them, multiply by 7 1/2 fen, and place it at the top. Further, multiply Bing and Ding combined by Jia and Yi combined, and subtract the number obtained from the upper position; this is the first subtractive edge (di yi yi lian). Add the two combined numbers and multiply by 7 1/2 fen; this is the second edge (di er lian). Multiply 7 1/2 fen by itself, obtaining 5 fen 6 li (釐) 2 hao (毫) 5 si (絲), placed in the upper position. From 1 bu subtract the upper position, leaving 4 fen 3 li 7 hao 5 si, which is the empty corner (xu yu). This gives the city diameter.\end{uniten}

\begin{unitzh}\unitlabel{ceyuan-022.28}草曰：别得甲為大勾，乙為明勾，丙為大股，丁為叀股也。甲乙共内減半徑，即是黄長弦也；丙丁共内減半徑，即黄廣弦也。黄長弦、黄廣弦二數相減，餘為兩个皇極差也。乃立天元為城徑，半之，副置二位：上以減於甲乙共數，得□□，即黄長弦也；以自之，得□□□，為黄長弦冪也；内減天元一冪，餘得下式□□□，為勾方差冪也。下位以減於丙丁共，得下式□□，即黄廣弦也；以自之，得□□□，為黄廣弦冪也；内減天元一冪，餘得□□□，為殷方差冪也。再以勾方差冪、股方差冪相乘，得□□□□□，為徑冪{\zhsmall 寄左}。然後以天元為冪，又以冪自之，與左相消，得下式□□□□□，開三乘方，得二百四十步，即城徑也。合問。\end{unitzh}
\begin{uniten}Working: Separately we obtain: Jia is the great base, Yi is the bright base, Bing is the great altitude, and Ding is the zhuan altitude. Jia and Yi combined minus the half diameter is the yellow-long hypotenuse (huang chang xian); Bing and Ding combined minus the half diameter is the yellow-broad hypotenuse (huang guang xian). The two numbers, the yellow-long hypotenuse and the yellow-broad hypotenuse, subtracted from each other, leave two Supreme Pole differences. Then set up the celestial unknown as the city diameter; halve it and set it down in duplicate in two positions. Subtract the upper from the combined number of Jia and Yi, obtaining □□, which is the yellow-long hypotenuse. Multiply it by itself, obtaining □□□, which is the area of the yellow-long hypotenuse. Subtract from it the area of the celestial unknown, the remainder obtained being the array below, □□□, which is the area of the base-square difference. Subtract the lower position from Bing and Ding combined, obtaining the array below, □□, which is the yellow-broad hypotenuse. Multiply it by itself, obtaining □□□, which is the area of the yellow-broad hypotenuse. Subtract from it the area of the celestial unknown, the remainder obtained being □□□, which is the area of the yin-square difference [Translator's note: 殷 "yin" is evidently an error for 股 "altitude"]. Again multiply the area of the base-square difference and the area of the altitude-square difference together, obtaining □□□□□, which is the area of the diameter [Note: set aside on the left]. Then make the celestial unknown into an area, and further multiply the area by itself. Cancel it against the left, obtaining the array below, □□□□□. Extract the root of the thrice-multiplied power, obtaining 240 bu, which is the city diameter. This matches the question.\end{uniten}

\begin{unitzh}\unitlabel{ceyuan-022.29}測圓海鏡卷七\end{unitzh}
\begin{uniten}Sea Mirror of Circle Measurements, juan 7\end{uniten}

\section*{卷08}\addcontentsline{toc}{section}{卷08}

\begin{unitzh}\unitlabel{ceyuan-023.1}《欽定四庫全書》 《測圓海鏡》卷八 元·李冶撰\end{unitzh}
\begin{uniten}Imperially Commissioned Complete Library of the Four Treasuries (Qinding Siku quanshu). Sea Mirror of Circle Measurements (Ceyuan haijing), juan 8. Compiled by Li Ye of the Yuan.\end{uniten}

\begin{unitzh}\unitlabel{ceyuan-023.2}明叀後一十六問\end{unitzh}
\begin{uniten}Ming and zhuan, latter part: 16 problems\end{uniten}

\begin{unitzh}\unitlabel{ceyuan-023.3}或問：出南門向東有槐樹一株，出東門向南有柳樹一株。丙、丁俱出南門，丙直行，丁往至槐樹下；甲、乙俱出東門，甲直行，乙往至柳樹下。四人遥相望見，各不知所行步數，只云丙、丁共行了二百七步，甲、乙共行四十六步，又云甲、丙立處相距二百八十九步。問答同前。\end{unitzh}
\begin{uniten}Someone asks: Going out of the south gate, toward the east there is one scholar tree (huai); going out of the east gate, toward the south there is one willow tree. Bing and Ding both go out of the south gate: Bing walks straight on, Ding goes to the foot of the scholar tree. Jia and Yi both go out of the east gate: Jia walks straight on, Yi goes to the foot of the willow. The four persons see one another from afar. The number of bu (步) each has walked is not known; it is only said that Bing and Ding together walked 207 bu and Jia and Yi together walked 46 bu; it is further said that the places where Jia and Bing stand are 289 bu apart. The question and the answer are the same as before.\end{uniten}

\begin{unitzh}\unitlabel{ceyuan-023.4}法曰：以二共相減數，又以減距數為實，二為法，得平勾。\end{unitzh}
\begin{uniten}Rule: Take the number obtained by subtracting the two totals from each other, and with it further subtract from the number of the distance to make the dividend (shi); 2 is the divisor (fa); one obtains the ping base (gou).\end{uniten}

\begin{unitzh}\unitlabel{ceyuan-023.5}草曰：識别得丙丁共即明和也，甲乙共即叀和也，相距步即極弦也。二共相併，即極弦内少个虚黄也，又為極和内少个虚和也；二共相減，餘為平勾髙股差也，又為虚差極差共也，又為通差内減極差也。立天元為平勾，加入二共相減數，得□□，為髙弦；又加天元，得□□，為極弦{\zhsmall 寄左}。以相距步二百八十九與左相消，得□□，上法下實，如法得六十四，即平勾也。以二共相減數加平匀，得二百二十五，為髙股；復以平勾乘之，得一萬四千四百步，開平方，得一百二十步，即城半徑也。合問。\end{unitzh}
\begin{uniten}Working: By identification (shibie) one finds that Bing and Ding's total is the ming sum (he), Jia and Yi's total is the zhuan sum, and the distance in bu is the ji hypotenuse (xian). The two totals added together are the ji hypotenuse lacking one xu yellow (huang); they are also the ji sum lacking one xu sum. The two totals subtracted from each other leave the difference (cha) of the ping base and the gao altitude (gu); this is also the xu difference and the ji difference together; it is also the tong difference less the ji difference. Set up the celestial unknown (tianyuan) as the ping base; add in the difference of the two totals, obtaining □□ [Translator's note: □ marks counting-rod entries or characters not transcribed in the source.] as the gao hypotenuse [Translator's note: the same quantity is called the gao altitude below.]; add the celestial unknown again, obtaining □□ as the ji hypotenuse [Note: lodged on the left (jizuo)]. Cancel the distance, 289 bu, against what is on the left (xiangxiao), obtaining □□; the upper is the divisor, the lower the dividend; dividing by the divisor gives 64, which is the ping base. Add the difference of the two totals to the ping yun [Translator's note: so the text (匀), evidently for 勾 "base".], obtaining 225 as the gao altitude; multiply it again by the ping base, obtaining 14400 bu; extracting the square root gives 120 bu, which is the radius (banjing) of the city. This agrees with the question.\end{uniten}

\begin{unitzh}\unitlabel{ceyuan-023.6}又法：二共數併，以減相距數，餘者半為泛率；以泛率加丙丁共為長，以泛率加甲乙共為闊，長闊相乘為平方實，得半徑。\end{unitzh}
\begin{uniten}Another rule: Add the two totals together and subtract them from the number of the distance; half of the remainder is the provisional rate (fanlü). Add the provisional rate to Bing and Ding's total to make the length; add the provisional rate to Jia and Yi's total to make the width; multiply length and width together to make the dividend for the square root (pingfang shi); one obtains the radius.\end{uniten}

\begin{unitzh}\unitlabel{ceyuan-023.7}草曰：置極弦，内減二共併數，餘三十六步，即虚黄也。半之，副置二位：上以加明和，得二百二十五步，為髙股也；下以加叀和，得六十四步，為平勾也。二位相乘，得一萬四千四百步，開平方，得一百二十步，即半徑也。合問。\end{unitzh}
\begin{uniten}Working: Set out the ji hypotenuse and subtract from it the sum of the two totals; the remainder, 36 bu, is the xu yellow. Halve it and set it out in two auxiliary positions: to the upper add the ming sum, obtaining 225 bu, which is the gao altitude; to the lower add the zhuan sum, obtaining 64 bu, which is the ping base. Multiply the two positions together, obtaining 14400 bu; extracting the square root gives 120 bu, which is the radius. This agrees with the question.\end{uniten}

\begin{unitzh}\unitlabel{ceyuan-023.8}或問：依前見丙丁共二百七步，甲乙共四十六步，又云二樹相去一百二步。問答同前。\end{unitzh}
\begin{uniten}Someone asks: As before, Bing and Ding's total is seen to be 207 bu and Jia and Yi's total 46 bu; it is further said that the two trees are 102 bu apart. The question and the answer are the same as before.\end{uniten}

\begin{unitzh}\unitlabel{ceyuan-023.9}法曰：以甲乙共乘樹相去步，得數，又以自之，為平實；從空；併二共數為冪於上，内減甲乙共自之數、丙丁共自之數{\zhsmall 按：或云二共數相乘倍之，亦同。}，為益隅；得叀弦。\end{unitzh}
\begin{uniten}Rule: Multiply the bu of the distance between the trees by Jia and Yi's total; multiply the number obtained further by itself to make the square dividend (pingshi); the joining term (cong) is empty. Add the two totals together and make an area (mi) of it, placed above; subtract from it the number of Jia and Yi's total multiplied by itself and the number of Bing and Ding's total multiplied by itself [Siku editors: Note: some say that multiplying the two totals together and doubling gives the same as well.] to make the subtractive corner (yiyu); one obtains the zhuan hypotenuse.\end{uniten}

\begin{unitzh}\unitlabel{ceyuan-023.10}草曰：識别得兩樹相去步即虚弦也，餘數具前。立天元一為叀弦，置明和，以天元乘之，合叀和除，不除，便以□為明弦也{\zhsmall 内帶□和分母}。乃置虚弦，以分母叀和乘之，得□，加入明弦，得□□，為極股也，内帶叀和分母。以自之，得下式□□□，為極股冪{\zhsmall 内寄叀和羃為分母}。又以天元加虚弦，得□□，為極勾；以自之，得丨□□；又以叀和冪□乘之，得□□□，為勾冪也。勾股相併，得□□□，為兩積一較冪也，内有叀和冪分母{\zhsmall 寄左}。然後置明弦□於上，以叀和乘天元，得□，加上位，得□，為二弦併；又置虚弦，以叀和乘之，得□，併入上位，得下式□□，為極弦；以自之，得□□□，為同數。與左相消，得□□□，開平方，得三十四步，即叀弦也。\end{unitzh}
\begin{uniten}Working: By identification one finds that the distance in bu between the two trees is the xu hypotenuse; the other numbers are as set out before. Set up the celestial unknown as the zhuan hypotenuse. Set out the ming sum and multiply it by the celestial unknown; it ought to be divided by the zhuan sum; without dividing, simply take □ as the ming hypotenuse [Note: it contains the □ sum as denominator (fenmu)]. Then set out the xu hypotenuse and multiply it by the denominator, the zhuan sum, obtaining □; add it into the ming hypotenuse, obtaining □□ as the ji altitude; it contains the zhuan sum as denominator. Multiply it by itself, obtaining the form below (xiashi), □□□, as the area of the ji altitude [Note: within it the area of the zhuan sum is lodged as denominator]. Further add the celestial unknown to the xu hypotenuse, obtaining □□ as the ji base; multiply it by itself, obtaining 丨□□ [Translator's note: 丨 is the rod numeral 1.]; further multiply it by the area of the zhuan sum, □, obtaining □□□ as the base area. Combining base and altitude gives □□□ as two areas (ji) and one difference area (jiao mi); it contains the area of the zhuan sum as denominator [Note: lodged on the left]. Then set out the ming hypotenuse □ above; multiply the celestial unknown by the zhuan sum, obtaining □; add it to the upper position, obtaining □ as the two hypotenuses combined; further set out the xu hypotenuse and multiply it by the zhuan sum, obtaining □; add it into the upper position, obtaining the form below, □□, as the ji hypotenuse; multiply it by itself, obtaining □□□ as the equal number (tongshu). Cancel it against what is on the left, obtaining □□□; extracting the square root gives 34 bu, which is the zhuan hypotenuse.\end{uniten}

\begin{unitzh}\unitlabel{ceyuan-023.11}又法：以樹相去步自之，又以甲乙共乘之，為平實；從空；倍丙丁共，為虚隅；得叀弦。\end{unitzh}
\begin{uniten}Another rule: Multiply the bu of the distance between the trees by itself, and further multiply by Jia and Yi's total, to make the square dividend; the joining term is empty; double Bing and Ding's total to make the void corner (xuyu); one obtains the zhuan hypotenuse.\end{uniten}

\begin{unitzh}\unitlabel{ceyuan-023.12}草曰：立天元一為叀弦，依前術求得明弦□，便以為皇極勾弦差也{\zhsmall 内帶叀和分母}。以天元□弦便為皇極股弦差，以乘之，又倍之，得□□，為虚弦冪{\zhsmall 内有叀和分母，寄左}。然後以虚弦自之，又以分母□乘之，得四十七萬八千五百八十四，為同數。與左相消，得□○□，開平方，得三十四步，即叀弦也。合問。\end{unitzh}
\begin{uniten}Working: Set up the celestial unknown as the zhuan hypotenuse. Following the previous method (shu), find the ming hypotenuse □ and simply take it as the huangji base–hypotenuse difference [Note: it contains the zhuan sum as denominator]. Take the celestial unknown □ hypotenuse simply as the huangji altitude–hypotenuse difference; multiply by it and further double it, obtaining □□ as the area of the xu hypotenuse [Note: it contains the zhuan sum as denominator; lodged on the left]. Then multiply the xu hypotenuse by itself, and further multiply by the denominator □, obtaining 478584 as the equal number. Cancel it against what is on the left, obtaining □○□; extracting the square root gives 34 bu, which is the zhuan hypotenuse. This agrees with the question.\end{uniten}

\begin{unitzh}\unitlabel{ceyuan-023.13}或問：皇極大小差共一百八十七步，明黄、叀黄共六十六步。問答同前。\end{unitzh}
\begin{uniten}Someone asks: The huangji greater and lesser differences (dacha, xiaocha) together are 187 bu; the ming yellow and the zhuan yellow together are 66 bu. The question and the answer are the same as before.\end{uniten}

\begin{unitzh}\unitlabel{ceyuan-023.14}法曰：後數自乘為實，前後數相減，餘為法，得虚黄方。\end{unitzh}
\begin{uniten}Rule: Multiply the latter number by itself to make the dividend; subtract the former and latter numbers from each other, and the remainder is the divisor; one obtains the xu yellow square (huangfang).\end{uniten}

\begin{unitzh}\unitlabel{ceyuan-023.15}草曰：别得一百八十七即明叀二弦共也，其六十六即太虚大小差共也。又二數相併，得□，即明叀二和共；若以相減，餘□，即明叀四差共也。立天元一為太虚黄方面，加二黄共，得□□，即虚弦也；倍虚弦，又加天元，得□□，即城徑也；又以虚弦加皇極大小差，得□□，即極弦也。以極弦乘城徑，得□□□，為兩段皇極勾股積{\zhsmall 寄左}。再以極弦、虚弦相併，得□□，即皇極勾股共也；自之，得□□□，内減皇極弦冪丨□□，得□□□，為同數。與寄左相消，得□□，上法下實，如法得三十六步，即太虚黄方靣也。合問。\end{unitzh}
\begin{uniten}Working: By distinction one finds that the 187 is the two hypotenuses of ming and zhuan together, and the 66 is the taixu greater and lesser differences together. Further, the two numbers added together give □, which is the two sums of ming and zhuan together; if they are subtracted from each other, the remainder □ is the four differences of ming and zhuan together. Set up the celestial unknown as the side (mian) of the taixu yellow square; add the total of the two yellows, obtaining □□, which is the xu hypotenuse; double the xu hypotenuse and add the celestial unknown again, obtaining □□, which is the diameter of the city; further add the huangji greater and lesser differences to the xu hypotenuse, obtaining □□, which is the ji hypotenuse. Multiply the diameter of the city by the ji hypotenuse, obtaining □□□ as two portions (duan) of the huangji base-and-altitude area (ji) [Note: lodged on the left]. Again add the ji hypotenuse and the xu hypotenuse together, obtaining □□, which is the huangji base and altitude together; multiply it by itself, obtaining □□□; subtract from it the area of the huangji hypotenuse, 丨□□, obtaining □□□ as the equal number. Cancel it against what was lodged on the left, obtaining □□; the upper is the divisor, the lower the dividend; dividing by the divisor gives 36 bu, which is the side of the taixu yellow square. This agrees with the question.\end{uniten}

\begin{unitzh}\unitlabel{ceyuan-023.16}或問：東門南有柳一株，南門東有槐一株。甲出東門直行，丙出東門直行，甲、丙、槐、柳悉與城㕘相直。既而甲就柳樹斜行三十四步至柳樹下，丙就槐樹斜行一百五十三步至槐樹下。問答同前。\end{unitzh}
\begin{uniten}Someone asks: South of the east gate there is one willow; east of the south gate there is one scholar tree. Jia goes out of the east gate and walks straight on; Bing goes out of the east gate [Translator's note: so the text; elsewhere Bing leaves by the south gate.] and walks straight on; Jia, Bing, the scholar tree and the willow are all aligned together with the city. Then Jia, making for the willow, walks obliquely 34 bu to the foot of the willow; Bing, making for the scholar tree, walks obliquely 153 bu to the foot of the scholar tree. The question and the answer are the same as before.\end{uniten}

\begin{unitzh}\unitlabel{ceyuan-023.17}法曰：云數相乘，倍之，便為平方實，開方，得虚弦一百二步。以此弦加甲行步，即極勾；以此弦加丙行步，即極股。餘各依法求之。識别：甲斜行即叀弦也，丙斜行即明弦也。無草。\end{unitzh}
\begin{uniten}Rule: Multiply the stated numbers together and double; this is then the dividend for the square root; extracting the square root gives the xu hypotenuse, 102 bu. Add this hypotenuse to the bu of Jia's walk: this is the ji base; add this hypotenuse to the bu of Bing's walk: this is the ji altitude. The rest are each found according to the rules. Identification: Jia's oblique walk is the zhuan hypotenuse; Bing's oblique walk is the ming hypotenuse. There is no working.\end{uniten}

\begin{unitzh}\unitlabel{ceyuan-023.18}或問：東門南有柳一株，南門東有槐一株。甲出東門直行，丙出南門直行，二人遙相望，槐、栁與城邊悉相直。既而甲復斜行至柳樹下，丙復斜行至槐樹下，各不知步數，只云丙共行了二百八十八步，甲斜行與柳至東門步共得六十四步。問答同前。\end{unitzh}
\begin{uniten}Someone asks: South of the east gate there is one willow; east of the south gate there is one scholar tree. Jia goes out of the east gate and walks straight on; Bing goes out of the south gate and walks straight on; the two persons see each other from afar, and the scholar tree and the willow are all in line with the edge of the city. Then Jia again walks obliquely to the foot of the willow, and Bing again walks obliquely to the foot of the scholar tree. The numbers of bu are not known in either case; it is only said that Bing walked 288 bu in all, and that Jia's oblique walk and the bu from the willow to the east gate together come to 64 bu. The question and the answer are the same as before.\end{uniten}

\begin{unitzh}\unitlabel{ceyuan-023.19}法曰：二云數相乘於上；以六十四步自之，又二之，減上位，為平實。十四之六十四於上，倍丙行，減上位，為從{\zhsmall 按：倍丙行乃數偶合，當云九个半六十四内減丙行為從。}。二十，常法。得甲直行步。\end{unitzh}
\begin{uniten}Rule: Multiply the two stated numbers together, placing the result above; multiply 64 bu by itself, further double it, and subtract from the upper position to make the square dividend. Take 14 times 64, placed above; double Bing's walk and subtract from the upper position to make the joining term [Siku editors: Note: "double Bing's walk" is a chance coincidence of the numbers; it should say: from 9 1/2 times 64 subtract Bing's walk to make the joining term.]. 20 is the constant divisor (changfa). One obtains the bu of Jia's straight walk.\end{uniten}

\begin{unitzh}\unitlabel{ceyuan-023.20}草曰：别得丙共步即明股明弦和也，六十四即平勾也，内甲斜行即叀弦也，柳至東門步即叀股也。又云二數相併，即明差與極弦共也；二云數相減，即明差與平勾髙股差共也。又平勾内減叀勾，即虚勾也。立天元一為叀勾，置丙共步，以天元乘之，復以六十四除之，得□□呔，為明勾也；又以天元減於六十四，得□□，為虚勾也；併虚明二勾□□，為半徑也。以自之，得□□□□，倍之，得□□□□，為半段圓城徑冪{\zhsmall 寄左}。乃以天元加六十四，得□□，為勾圓差於上；又以明勾加丙共步，得□□□，為股圓差於下；上下相乘，得□□□□，為同數。與左相消，得□□□，開平方，得一十六步，即叀勾也。此叀勾乃甲出東門直行步也。餘皆依數求。合問。\end{unitzh}
\begin{uniten}Working: By distinction one finds that Bing's total in bu is the sum of the ming altitude and the ming hypotenuse; 64 is the ping base; within it, Jia's oblique walk is the zhuan hypotenuse, and the bu from the willow to the east gate are the zhuan altitude. Further, the stated two numbers added together are the ming difference and the ji hypotenuse together; the two stated numbers subtracted from each other are the ming difference and the difference of the ping base and the gao altitude together. Further, the ping base less the zhuan base is the xu base. Set up the celestial unknown as the zhuan base. Set out Bing's total in bu, multiply it by the celestial unknown, and then divide by 64, obtaining □□呔 [Translator's note: 呔 so in the text.] as the ming base; further subtract the celestial unknown from 64, obtaining □□ as the xu base; the two bases xu and ming added together, □□, are the radius. Multiply it by itself, obtaining □□□□; double it, obtaining □□□□ as half a portion of the area of the round city's diameter [Note: lodged on the left]. Then add the celestial unknown to 64, obtaining □□ as the base–circle difference (gouyuan cha), placed above; further add the ming base to Bing's total in bu, obtaining □□□ as the altitude–circle difference (guyuan cha), placed below; multiply upper and lower together, obtaining □□□□ as the equal number. Cancel it against what is on the left, obtaining □□□; extracting the square root gives 16 bu, which is the zhuan base. This zhuan base is the bu Jia walked straight on out of the east gate. The rest are all found according to the numbers. This agrees with the question.\end{uniten}

\begin{unitzh}\unitlabel{ceyuan-023.21}或問：東門南有柳樹一株，南門東有槐樹一株。甲出東門直行，丙出南門直行，二人遥相望，槐、柳與城邊悉相直。既而甲復斜行至柳樹下，丙復斜行至槐樹下，各不知步數，只云甲共行五十步，丙斜行與槐至南門步共得二百二十五步。問答同前。\end{unitzh}
\begin{uniten}Someone asks: South of the east gate there is one willow tree; east of the south gate there is one scholar tree. Jia goes out of the east gate and walks straight on; Bing goes out of the south gate and walks straight on; the two persons see each other from afar, and the scholar tree and the willow are all in line with the edge of the city. Then Jia again walks obliquely to the foot of the willow, and Bing again walks obliquely to the foot of the scholar tree. The numbers of bu are not known in either case; it is only said that Jia walked 50 bu in all, and that Bing's oblique walk and the bu from the scholar tree to the south gate together come to 225 bu. The question and the answer are the same as before.\end{uniten}

\begin{unitzh}\unitlabel{ceyuan-023.22}法曰：以二百二十五步自之為冪，又以此冪自為冪，於上；置甲共行，以二百二十五步三度乘之，得數，復折半，減上位，為平實。置二百二十五步自之數，以二云數相減數乘之，又倍之，於上；倍五十步在地，以二百二十五步自之數乘之，復折半，加上位，為益從。云數相減自乘，於上；以云數相乘，復折半，減上位，為常法。得明股。\end{unitzh}
\begin{uniten}Rule: Multiply 225 bu by itself to make an area (mi); further multiply this area by itself to make an area, placed above. Set out Jia's total walk and multiply it three times by 225 bu; halve the number obtained and subtract it from the upper position to make the square dividend. Set out the number of 225 bu multiplied by itself, multiply it by the difference of the two stated numbers, and further double it, placed above; double 50 bu, setting it in the lower (earth, di) position, multiply it by the number of 225 bu multiplied by itself, then halve it, and add it to the upper position to make the subtractive joining term (yicong). Multiply the difference of the stated numbers by itself, placed above; multiply the stated numbers together, then halve, and subtract from the upper position to make the constant divisor. One obtains the ming altitude.\end{uniten}

\begin{unitzh}\unitlabel{ceyuan-023.23}草曰：識别得甲共步即叀勾叀弦共也，二百二十五即髙股也，内丙斜行即明弦，槐至南門步即明勾也。又二云數相併，即極弦内減一个叀差也；云數相減，即叀差與髙股平勾差共也。又髙股内減明股，即虚股也。立天元一為明股，即丙出南門直行步也。置五十步，以天元乘之，得□，合髙股除，不除，便以此□為叀股也，内帶髙股□分母。再置髙股，内減天元，得□□，為虚股；以分母髙股乘之，得下式□□，加入叀股，得□□，即半徑也。以自增乘，得下□□□，為半徑冪也，内帶髙股冪為母{\zhsmall 寄左}。然後置甲共步，以分母髙股乘之，得□，加入叀股，得□□，為勾圓差於上{\zhsmall 内帶髙股分母}；又以天元加髙股，得□□，為股圓差於下；上下相乘，得□□□；又以分母髙股乘之，得□□□；復折半，得□□□，為同數。與左相消，得□□□，開平方，得一百三十五步，即明股也。合問。\end{unitzh}
\begin{uniten}Working: By identification one finds that Jia's total in bu is the zhuan base and the zhuan hypotenuse together; 225 is the gao altitude; within it, Bing's oblique walk is the ming hypotenuse, and the bu from the scholar tree to the south gate are the ming base. Further, the two stated numbers added together are the ji hypotenuse less one zhuan difference; the stated numbers subtracted from each other are the zhuan difference and the difference of the gao altitude and the ping base together. Further, the gao altitude less the ming altitude is the xu altitude. Set up the celestial unknown as the ming altitude, which is the bu Bing walked straight on out of the south gate. Set out 50 bu and multiply it by the celestial unknown, obtaining □; it ought to be divided by the gao altitude; without dividing, simply take this □ as the zhuan altitude, which contains the gao altitude □ as denominator. Again set out the gao altitude and subtract the celestial unknown from it, obtaining □□ as the xu altitude; multiply it by the denominator, the gao altitude, obtaining the form below, □□; add it into the zhuan altitude, obtaining □□, which is the radius. Multiply it by itself (zizeng cheng), obtaining, below, □□□ as the area of the radius, which contains the area of the gao altitude as denominator (mu) [Note: lodged on the left]. Then set out Jia's total in bu and multiply it by the denominator, the gao altitude, obtaining □; add it into the zhuan altitude, obtaining □□ as the base–circle difference, placed above [Note: it contains the gao altitude as denominator]; further add the celestial unknown to the gao altitude, obtaining □□ as the altitude–circle difference, placed below; multiply upper and lower together, obtaining □□□; further multiply by the denominator, the gao altitude, obtaining □□□; then halve, obtaining □□□ as the equal number. Cancel it against what is on the left, obtaining □□□; extracting the square root gives 135 bu, which is the ming altitude. This agrees with the question.\end{uniten}

\begin{unitzh}\unitlabel{ceyuan-023.24}或問：通勾通弦共一千步，叀勾叀弦共五十步。問答同前。\end{unitzh}
\begin{uniten}Someone asks: The tong base and the tong hypotenuse together are 1000 bu; the zhuan base and the zhuan hypotenuse together are 50 bu. The question and the answer are the same as before.\end{uniten}

\begin{unitzh}\unitlabel{ceyuan-023.25}法曰：置一千，減二之五十步，為汎率；以自乘，復半之，於上；又置泛率，復以五十乘之，加上位，為平實。二十二之泛率於上{\zhsmall 按：二十二乃此題叀和除通和所得通倍叀數加二數之數，易題則數不同矣，當直云通倍叀數加二數乘泛率。}；以四十二{\zhsmall 按：四十二乃此題倍通倍叀數加二數之數，當直云倍通倍叀數加二數。}乘五十，得數，内減泛率，加上位，為益從。二百{\zhsmall 按：二百乃此題通倍叀數加二數自乘折半於上，又倍通倍叀數併二數以減上位之數，當同上，不必載數。}為常法。得叀股。\end{unitzh}
\begin{uniten}Rule: Set out 1000 and subtract twice 50 bu to make the provisional rate; multiply it by itself and then halve it, placed above; further set out the provisional rate, multiply it by 50, and add to the upper position to make the square dividend. Take 22 times the provisional rate, placed above [Siku editors: Note: 22 is, in this problem, the multiple of tong over zhuan (tong bei zhuan shu), obtained by dividing the tong sum by the zhuan sum, plus 2; if the problem is changed, the number will differ; it ought to say directly: multiply the provisional rate by the multiple of tong over zhuan plus 2.]; with 42 [Siku editors: Note: 42 is, in this problem, twice the multiple of tong over zhuan, plus 2; it ought to say directly: twice the multiple of tong over zhuan plus 2.] multiply 50; from the number obtained subtract the provisional rate, and add to the upper position to make the subtractive joining term. 200 [Siku editors: Note: 200 is, in this problem, the number obtained when the multiple of tong over zhuan plus 2 is multiplied by itself, halved and placed above, and then twice the multiple of tong over zhuan with 2 added is subtracted from the upper position; it ought to be treated as above, and the number need not be recorded.] is the constant divisor. One obtains the zhuan altitude.\end{uniten}

\begin{unitzh}\unitlabel{ceyuan-024.1}草曰：立天元一為叀股。置一千，以天元乘之，以五十除之，得□，為通股也。又以天元加五十步，得□□，即小差也。通股加小差，得□□，即通弦也。以通弦減一千，得□□，即通勾也。以小差減通勾，得□□，即圓徑也。以圓徑減通股，得□□，即大差也。置大差，以小差乘之，得□□□{\zhsmall 寄左}。然後置圓徑，以自之，得□□□；折半，得□□□；與左相消，得□□□；開平方，得三十步，即叀股也。合問。\end{unitzh}
\begin{uniten}Working: Set up the celestial unknown as the zhuan (叀) altitude (gu). Set 1000, multiply it by the celestial unknown, and divide it by 50, obtaining □ [Translator's note: □ marks a counting-rod figure not reproduced in the transcription.], which is the tong (通) altitude. Again add the celestial unknown to 50 bu (步), obtaining □□, which is the small difference (xiaocha). Add the small difference to the tong altitude, obtaining □□, which is the tong hypotenuse (xian). Subtract the tong hypotenuse from 1000, obtaining □□, which is the tong base (gou). Subtract the small difference from the tong base, obtaining □□, which is the diameter of the circle. Subtract the diameter of the circle from the tong altitude, obtaining □□, which is the great difference (dacha). Set the great difference and multiply it by the small difference, obtaining □□□ [Note: lodge it on the left (jizuo)]. Then set the diameter of the circle and multiply it by itself, obtaining □□□; halve it, obtaining □□□; cancel it against the left (xiangxiao), obtaining □□□; extract the square root, obtaining 30 bu, which is the zhuan altitude. This agrees with the question.\end{uniten}

\begin{unitzh}\unitlabel{ceyuan-024.2}按：此題通勾弦和為叀勾弦和度盡之數，則不用寄分而用除法，以從省便。作者蓋舉一以例其餘也。\end{unitzh}
\begin{uniten}Note: In this problem the sum of the tong base and hypotenuse is a number exactly measured out by the sum of the zhuan base and hypotenuse; therefore lodged fractions are not used, but division is used, so as to follow economy and convenience. The author presumably cites one case to exemplify the rest.\end{uniten}

\begin{unitzh}\unitlabel{ceyuan-024.3}或問：通勾、通弦共一千步，明勾、明弦共二百二十五步。問答同前。\end{unitzh}
\begin{uniten}Someone asks: the tong base and the tong hypotenuse together are 1000 bu; the ming (明) base and the ming hypotenuse together are 225 bu. The question and the answer are the same as before.\end{uniten}

\begin{unitzh}\unitlabel{ceyuan-024.4}法曰：以後數再自乘，又以前數乘之，為平實；以後數為冪，又以前數乘之，為從；以前數冪為常法。得明股。\end{unitzh}
\begin{uniten}Rule: Multiply the latter number by itself a second time, and further multiply it by the former number; this is the level dividend (pingshi). Make the latter number into an area (mi), and further multiply it by the former number; this is the joint (cong). Take the area of the former number as the constant divisor (changfa). Obtain the ming altitude.\end{uniten}

\begin{unitzh}\unitlabel{ceyuan-024.5}草曰：别得二百二十五步，即髙股也。立天元一為明股。置一千，以天元乘之，合以髙股除，不除，便以此□為通股{\zhsmall 内帶髙股為母}。以天元加髙股，□□，即大差也。置大差，以髙股分母乘之，得□□，即帶分大差也。以此減於通股，餘□□，即圓徑也。以自增乘，得□□□，寄左{\zhsmall 内𢃄髙股冪分母}。然後置一千，以髙股分母通之，得□；内減帶分大差，得□□，為兩个通勾也；内減兩个圓徑，得□□，為兩个小差也。以帶分大差乘之，得下式□□□，為同數。與左相消，得□□，開平方，得一百三十五步，即明股也。合問。\end{unitzh}
\begin{uniten}Working: By distinction one obtains 225 bu, which is the gao (髙) altitude. Set up the celestial unknown as the ming altitude. Set 1000 and multiply it by the celestial unknown; it ought to be divided by the gao altitude; not dividing, simply take this □ as the tong altitude [Note: within, it carries the gao altitude as denominator]. Add the celestial unknown to the gao altitude: □□, which is the great difference. Set the great difference and multiply it by the gao-altitude denominator, obtaining □□, which is the great difference carrying the fraction. Subtract this from the tong altitude; the remainder □□ is the diameter of the circle. Multiply it by itself (zizengcheng), obtaining □□□; lodge it on the left [Note: within, it carries the area (mi) of the gao altitude as denominator]. Then set 1000 and convert it (tong) by the gao-altitude denominator, obtaining □; from it subtract the great difference carrying the fraction, obtaining □□, which is two tong bases; from it subtract two diameters of the circle, obtaining □□, which is two small differences. Multiply it by the great difference carrying the fraction, obtaining the form below, □□□, which is the same number (tongshu). Cancel it against the left, obtaining □□; extract the square root, obtaining 135 bu, which is the ming altitude. This agrees with the question.\end{uniten}

\begin{unitzh}\unitlabel{ceyuan-024.6}或問：通股、通弦共一千二百八十步，叀股、叀弦共六十四步。問答同前。\end{unitzh}
\begin{uniten}Someone asks: the tong altitude and the tong hypotenuse together are 1280 bu; the zhuan altitude and the zhuan hypotenuse together are 64 bu. The question and the answer are the same as before.\end{uniten}

\begin{unitzh}\unitlabel{ceyuan-024.7}法曰：云數相乘為平實，前數為益從。置前數，以後數除之，得二十，為泛率；泛率減一，以自乘，於上；又倍泛率減一，加上位，為常法。倒積開，得叀勾。\end{unitzh}
\begin{uniten}Rule: Multiply the said numbers together; this is the level dividend. The former number is the increasing joint (yicong). Set the former number and divide it by the latter number, obtaining 20; this is the general rate (fanlü). Subtract 1 from the general rate and multiply it by itself, putting it above; again double the general rate, subtract 1, and add it to the upper position; this is the constant divisor. Extract with the area inverted (daoji kai), obtaining the zhuan base.\end{uniten}

\begin{unitzh}\unitlabel{ceyuan-024.8}草曰：别得六十四步，即平勾也。立天元一為叀勾。置前數，以天元乘之，以後數除之，得□，即通勾也。又置天元，加後數，得□□，即小差也。以小差減通勾，餘□□，即圓徑也。以自之，得□□□{\zhsmall 寄左}。然後以小差減於前數，得□□，為二通股；内減兩个圓徑，得□□，為二大差也。以小差乘之，得下□□□，與左相消，得□□□，開平方，得一十六步，即叀勾也。合問。\end{unitzh}
\begin{uniten}Working: By distinction one obtains 64 bu, which is the ping (平) base. Set up the celestial unknown as the zhuan base. Set the former number, multiply it by the celestial unknown, and divide it by the latter number, obtaining □, which is the tong base. Again set the celestial unknown and add the latter number, obtaining □□, which is the small difference. Subtract the small difference from the tong base; the remainder □□ is the diameter of the circle. Multiply it by itself, obtaining □□□ [Note: lodge it on the left]. Then subtract the small difference from the former number, obtaining □□, which is two tong altitudes; from it subtract two diameters of the circle, obtaining □□, which is two great differences. Multiply it by the small difference, obtaining below □□□; cancel it against the left, obtaining □□□; extract the square root, obtaining 16 bu, which is the zhuan base. This agrees with the question.\end{uniten}

\begin{unitzh}\unitlabel{ceyuan-024.9}或問：通股、通弦共一千二百八十步，明股、明弦共二百八十八步。問答同前。\end{unitzh}
\begin{uniten}Someone asks: the tong altitude and the tong hypotenuse together are 1280 bu; the ming altitude and the ming hypotenuse together are 288 bu. The question and the answer are the same as before.\end{uniten}

\begin{unitzh}\unitlabel{ceyuan-024.10}法曰：二數相減，以後數乘之，内減後數冪，又半之，為泛率；以自乘，為平實{\zhsmall 按：或云前數内減二後數，餘以後數乘之，折半，自之，亦同}。置前數，加二之後數而半之，為次率；以乘泛率，於上；以後數乘泛率，減上位{\zhsmall 按：或云二數相加，以乘前折半數，亦同}，為益從。次率自乘之，於上；以前數加次率，復以後數乘之，減上位{\zhsmall 按：或云前數折半，内減後數，又以半前數乘之，亦同}，為隅法。得明勾。\end{unitzh}
\begin{uniten}Rule: Subtract the two numbers from each other, multiply by the latter number, from it subtract the area (mi) of the latter number, and further halve it; this is the general rate. Multiply it by itself; this is the level dividend [Siku editors: Some say: from the former number subtract twice the latter number, multiply the remainder by the latter number, halve it, and multiply it by itself; it is also the same]. Set the former number, add twice the latter number, and halve it; this is the secondary rate (cilü). Multiply the general rate by it, putting it above; multiply the general rate by the latter number and subtract it from the upper position [Siku editors: Some say: add the two numbers together and with it multiply the previously halved number; it is also the same]; this is the increasing joint. Multiply the secondary rate by itself, putting it above; add the former number to the secondary rate, again multiply it by the latter number, and subtract it from the upper position [Siku editors: Some say: halve the former number, from it subtract the latter number, and further multiply it by half the former number; it is also the same]; this is the corner divisor (yufa). Obtain the ming base.\end{uniten}

\begin{unitzh}\unitlabel{ceyuan-024.11}草曰：别得二數相減，餘□，為通勾、通股及明勾共也。立天元一為明勾。置前數，以天元乘之，合以後數除之，不除，便以此□為通勾也{\zhsmall 内寄後數分母}。又以二數相減，得數内又減天元，得□□，為通和也。乃以分母二百八十八乘之，得下式□□；内減通勾，餘□□，為通股也。又以天元加後數，又以分母{\zhsmall 即後數也}通之，得□□，為大差也。以此大差減於通股，得下式□□，為一个圓徑也。半之，得□□，以自得之□□□，為半徑冪{\zhsmall 寄左}。然後以半圓徑減通勾，得□□，為底勾；又以天元乘之，又以分母二百八十八乘之，得□□呔，為同數。與左相消，得□□□，開平方，得七十二步，即明勾也。合問。\end{unitzh}
\begin{uniten}Working: By distinction one obtains: subtracting the two numbers from each other, the remainder □ is the tong base, the tong altitude and the ming base together. Set up the celestial unknown as the ming base. Set the former number and multiply it by the celestial unknown; it ought to be divided by the latter number; not dividing, simply take this □ as the tong base [Note: within, the latter number is lodged as denominator]. Again subtract the two numbers from each other; from the number obtained further subtract the celestial unknown, obtaining □□, which is the tong sum. Then multiply it by the denominator 288, obtaining the form below, □□; from it subtract the tong base; the remainder □□ is the tong altitude. Again add the celestial unknown to the latter number, and further convert it (tong) by the denominator [Note: that is, the latter number], obtaining □□, which is the great difference. Subtract this great difference from the tong altitude, obtaining the form below, □□, which is one diameter of the circle. Halve it, obtaining □□; taking it by itself, one obtains it as □□□, which is the area (mi) of the radius [Note: lodge it on the left]. Then subtract half the diameter of the circle from the tong base, obtaining □□, which is the di (底) base; further multiply it by the celestial unknown, and further multiply it by the denominator 288, obtaining □□呔 [Translator's note: 呔 apparently stands for a counting-rod figure marked 太.], which is the same number. Cancel it against the left, obtaining □□□; extract the square root, obtaining 72 bu, which is the ming base. This agrees with the question.\end{uniten}

\begin{unitzh}\unitlabel{ceyuan-024.12}或問：明股、明弦併二百八十八步，叀勾、叀弦併五十步，又云明股、叀勾併多於虚弦四十九步。問答同前。\end{unitzh}
\begin{uniten}Someone asks: the ming altitude and the ming hypotenuse combined are 288 bu; the zhuan base and the zhuan hypotenuse combined are 50 bu; it is further said that the ming altitude and the zhuan base combined exceed the xu (虚) hypotenuse by 49 bu. The question and the answer are the same as before.\end{uniten}

\begin{unitzh}\unitlabel{ceyuan-024.13}法曰：前二數相併，内減二之多步，即圓徑。又只以前二數相乘，便是半徑冪。\end{unitzh}
\begin{uniten}Rule: Combine the former two numbers; from it subtract twice the excess bu; this is the diameter of the circle. Further, simply multiply the former two numbers together; this then is the area (mi) of the radius.\end{uniten}

\begin{unitzh}\unitlabel{ceyuan-024.14}草曰：識别得前二數相減而半之，即極差也。其多步名傍差，又圓徑不及極弦數。\end{unitzh}
\begin{uniten}Working: By discernment (shibie) one obtains: subtract the former two numbers from each other and halve; this is the ji (極) difference. The excess bu is called the side difference (pangcha); it is also the number by which the diameter of the circle falls short of the ji hypotenuse.\end{uniten}

\begin{unitzh}\unitlabel{ceyuan-024.15}或問：平差、髙差共一百六十一步，明股、叀勾併多於虚弦四十九步。問答同前。\end{unitzh}
\begin{uniten}Someone asks: the ping difference and the gao difference together are 161 bu; the ming altitude and the zhuan base combined exceed the xu hypotenuse by 49 bu. The question and the answer are the same as before.\end{uniten}

\begin{unitzh}\unitlabel{ceyuan-024.16}法曰：二數相減，又半之，以自乘，為實；後數為法。得平勾。\end{unitzh}
\begin{uniten}Rule: Subtract the two numbers from each other, further halve it, and multiply it by itself; this is the dividend (shi). The latter number is the divisor (fa). Obtain the ping base.\end{uniten}

\begin{unitzh}\unitlabel{ceyuan-024.17}草曰：立天元一為平勾，以加前數，得□□，為髙股也。又以天元加髙股，得□□，為極弦；内減後數，得□□；又半之，得□□，為半徑；以自之，得丨□□{\zhsmall 寄左}。然後以天元乘髙股，得丨□，為同數。與左相消，得□□，上法下實，得六十四步，即平勾也。合問。\end{unitzh}
\begin{uniten}Working: Set up the celestial unknown as the ping base; add it to the former number, obtaining □□, which is the gao altitude. Again add the celestial unknown to the gao altitude, obtaining □□, which is the ji hypotenuse; from it subtract the latter number, obtaining □□; further halve it, obtaining □□, which is the radius; multiply it by itself, obtaining 丨□□ [Translator's note: 丨 is the counting-rod numeral 1.] [Note: lodge it on the left]. Then multiply the gao altitude by the celestial unknown, obtaining 丨□, which is the same number. Cancel it against the left, obtaining □□; the upper is the divisor and the lower is the dividend; obtain 64 bu, which is the ping base. This agrees with the question.\end{uniten}

\begin{unitzh}\unitlabel{ceyuan-024.18}或問：平勾、髙股差一百六十一步，明差、叀差併七十七步，又云極弦多於城徑四十九步。問答同前。\end{unitzh}
\begin{uniten}Someone asks: the difference of the ping base and the gao altitude is 161 bu; the ming difference and the zhuan difference combined are 77 bu; it is further said that the ji hypotenuse exceeds the diameter of the city by 49 bu. The question and the answer are the same as before.\end{uniten}

\begin{unitzh}\unitlabel{ceyuan-024.19}法曰：併上二位而半之，為平率；其四十九，即旁率也。副置平率，上加旁率，下減旁率，以相乘，為實；倍旁差為法。得勾圓差{\zhsmall 按：求實數有誤，當云併上二位而半之，内減後數，於上；又置上前數，内減後數，以乘上位，為實，方合}。\end{unitzh}
\begin{uniten}Rule: Combine the upper two positions and halve; this is the level rate (pinglü); the 49 is the side rate (panglü). Set the level rate in duplicate (fuzhi): to the upper add the side rate, from the lower subtract the side rate, and multiply them together; this is the dividend. Double the side difference as the divisor. Obtain the base–circle difference (gouyuancha) [Siku editors: There is an error in finding the dividend; it should say: combine the upper two positions and halve, from it subtract the latter number, putting it above; again set the upper former number, from it subtract the latter number, and multiply the upper position by it; this is the dividend; only then does it agree].\end{uniten}

\begin{unitzh}\unitlabel{ceyuan-024.20}草曰：識别得平勾、髙股差名為角差。副置角差，上加七十七而半之，得□，即極差也；下減七十七而半之，得□，即虚差也。角差加極差，得□，即通差也。又極弦多於城徑步，名為旁差。副置角差，上加旁差，得□，為兩个髙段上勾股較；下減傍差，得□，為兩个平段上勾股較也。又副置極差，上加傍差，得□，為股圓差上勾股較；下減旁差，□，為勾圓差上勾股較也。立天元一為勾圓差。依法求得通差，加入天元，得□□，即大差也。以天元乘之，得丨□，為半段圓徑冪{\zhsmall 寄左}。乃置大差□□，内減股圓差上勾股較□，餘有□□，為股圓差之勾，於上；再置天元，内加勾圓差上勾股較□，得□□，為勾圓差之股；以乘上位，得丨□□，為同數。與左相消，得□□，上法下實，得八十歩，即勾圓差也。\end{unitzh}
\begin{uniten}Working: By discernment one obtains: the difference of the ping base and the gao altitude is called the corner difference (jiaocha). Set the corner difference in duplicate: to the upper add 77 and halve, obtaining □, which is the ji difference; from the lower subtract 77 and halve, obtaining □, which is the xu difference. Add the ji difference to the corner difference, obtaining □, which is the tong difference. Further, the bu by which the ji hypotenuse exceeds the diameter of the city are called the side difference. Set the corner difference in duplicate: to the upper add the side difference, obtaining □, which is two base–altitude differences (gougujiao) on the gao segment (duan); from the lower subtract the side difference, obtaining □, which is two base–altitude differences on the ping segment. Again set the ji difference in duplicate: to the upper add the side difference, obtaining □, which is the base–altitude difference on the altitude–circle difference (guyuancha); from the lower subtract the side difference: □, which is the base–altitude difference on the base–circle difference. Set up the celestial unknown as the base–circle difference. Following the method, find the tong difference and add the celestial unknown into it, obtaining □□, which is the great difference. Multiply it by the celestial unknown, obtaining 丨□, which is half a piece of the area (mi) of the diameter of the circle [Note: lodge it on the left]. Then set the great difference □□; from it subtract the base–altitude difference on the altitude–circle difference, □; the remainder is □□, which is the base of the altitude–circle difference, putting it above. Again set the celestial unknown; to it add the base–altitude difference on the base–circle difference, □, obtaining □□, which is the altitude of the base–circle difference; multiply the upper position by it, obtaining 丨□□, which is the same number. Cancel it against the left, obtaining □□; the upper is the divisor and the lower is the dividend; obtain 80 bu, which is the base–circle difference.\end{uniten}

\begin{unitzh}\unitlabel{ceyuan-024.21}又依前問：見角差一百六十一步，見明差、叀差併七十七步，又見太虚弦較較六十步。問答同前。\end{unitzh}
\begin{uniten}Again, following the previous question: the corner difference is given as 161 bu; the ming difference and the zhuan difference combined are given as 77 bu; further, the taixu (太虚) hypotenuse–difference difference (xianjiaojiao) is given as 60 bu. The question and the answer are the same as before.\end{uniten}

\begin{unitzh}\unitlabel{ceyuan-024.22}法曰：前二數相減而半之，得數加入半之太虚弦較較，為泛率；以自乘，為平實。置一百六十一，内減二之泛率，為從；一常法。得平勾。\end{unitzh}
\begin{uniten}Rule: Subtract the former two numbers from each other and halve; to the number obtained add half the taixu hypotenuse–difference difference; this is the general rate. Multiply it by itself; this is the level dividend. Set 161, from it subtract twice the general rate; this is the joint. 1 is the constant divisor. Obtain the ping base.\end{uniten}

\begin{unitzh}\unitlabel{ceyuan-024.23}草曰：别得□，即二叀股也。立天元一為平勾。先以前二數相減而半之，得□，為虛差；以虚差加叀股，得□，即明勾也。以明勾加天元，得丨□，為平弦；以自之，得丨□□；内減天元冪，得□□，為半徑冪{\zhsmall 寄左}。然後以天元加一百六十一，為髙股；以天元乘之，得丨□，為同數。與左相消，得丨□□，開平方，得六十四步，即平勾也。\end{unitzh}
\begin{uniten}Working: By distinction one obtains □, which is two zhuan altitudes. Set up the celestial unknown as the ping base. First subtract the former two numbers from each other and halve, obtaining □, which is the xu difference; add the xu difference to the zhuan altitude, obtaining □, which is the ming base. Add the ming base to the celestial unknown, obtaining 丨□, which is the ping hypotenuse; multiply it by itself, obtaining 丨□□; from it subtract the area (mi) of the celestial unknown, obtaining □□, which is the area of the radius [Note: lodge it on the left]. Then add the celestial unknown to 161; this is the gao altitude; multiply it by the celestial unknown, obtaining 丨□, which is the same number. Cancel it against the left, obtaining 丨□□; extract the square root, obtaining 64 bu, which is the ping base.\end{uniten}

\begin{unitzh}\unitlabel{ceyuan-024.24}又法曰：前數内加半之太虚弦較較，以自乘{\zhsmall 按：此語内有誤，當云倍角差加半太虚較，以半太虚較乘之}，為實；前數内減太虚弦較較，為從；一常法。開平方，得平勾。此更不用明差、叀差併也。\end{unitzh}
\begin{uniten}Another rule: To the former number add half the taixu hypotenuse–difference difference, and multiply it by itself [Siku editors: Within this statement there is an error; it should say: double the corner difference, add half the taixu difference, and multiply it by half the taixu difference]; this is the dividend. From the former number subtract the taixu hypotenuse–difference difference; this is the joint. 1 is the constant divisor. Extract the square root, obtaining the ping base. This moreover does not use the combined ming difference and zhuan difference.\end{uniten}

\begin{unitzh}\unitlabel{ceyuan-024.25}草曰：依前求平勾。前髙股内加叀股，得□□，為髙弦也；以自之，得丨□□，於上位；内減髙股冪丨□□，餘得□□，為半徑冪{\zhsmall 寄左}。然後以天元乘髙股，得丨□，為同數。與左相消，得下丨□□，開平方，得六十四步，即平勾也。合問。\end{unitzh}
\begin{uniten}Working: Seek the ping base as before. To the previous gao altitude add the zhuan altitude, obtaining □□, which is the gao hypotenuse; multiply it by itself, obtaining 丨□□, in the upper position; from it subtract the area (mi) of the gao altitude, 丨□□; the remainder obtained, □□, is the area of the radius [Note: lodge it on the left]. Then multiply the gao altitude by the celestial unknown, obtaining 丨□, which is the same number. Cancel it against the left, obtaining below 丨□□; extract the square root, obtaining 64 bu, which is the ping base. This agrees with the question.\end{uniten}

\begin{unitzh}\unitlabel{ceyuan-024.26}或問：髙差、平差併一百六十一步，明差、叀差併七十七步。問答同前。\end{unitzh}
\begin{uniten}Someone asks: the gao difference and the ping difference combined are 161 bu; the ming difference and the zhuan difference combined are 77 bu. The question and the answer are the same as before.\end{uniten}

\begin{unitzh}\unitlabel{ceyuan-024.27}法曰：以前數自乘，於上；二數相併而半之，以自乘，減上位，得數復自增乘，為平實。前數自之，於上；又以四之前數乘之，寄位。以前數自之，於上；併二數而半之，以自乘，減上位，得數又以四之前數乘之{\zhsmall 按：此下落「又倍之」三字}，減於寄位，為從。前數自之，又四之，於上；又以四之前數為冪，加上位，權寄。以前數為冪，於上；併二數而半之，以自乘，減上位，得數復八之，加上位；又以四之前數為冪，加入上位，併以減於權寄，為常法{\zhsmall 按：或云二和併而自之，又半之，以減髙平共差冪，又四之，為常法，亦同}。得平勾。\end{unitzh}
\begin{uniten}Rule: Multiply the former number by itself, putting it above; combine the two numbers and halve, multiply it by itself, and subtract it from the upper position; multiply the number obtained again by itself (zizengcheng); this is the level dividend. Multiply the former number by itself, putting it above; further multiply it by four times the former number; lodge it in a position. Multiply the former number by itself, putting it above; combine the two numbers and halve, multiply it by itself, and subtract it from the upper position; multiply the number obtained further by four times the former number [Siku editors: Below this the three characters 又倍之, "further double it", have dropped out]; subtract it from the lodged position; this is the joint. Multiply the former number by itself, further multiply it by 4, putting it above; further make four times the former number into an area (mi), and add it to the upper position; lodge it provisionally (quanji). Make the former number into an area, putting it above; combine the two numbers and halve, multiply it by itself, and subtract it from the upper position; multiply the number obtained again by 8, and add it to the upper position; further make four times the former number into an area, and add it into the upper position; take the combination and subtract it from what was provisionally lodged; this is the constant divisor [Siku editors: Some say: combine the two sums and multiply by itself, further halve it, subtract it from the area (mi) of the combined difference of gao and ping, and further multiply it by 4; this is the constant divisor; it is also the same]. Obtain the ping base.\end{uniten}

\begin{unitzh}\unitlabel{ceyuan-024.28}草曰：識别得二位相併而半之，得□，即極差也。立天元一為平勾，加一百六十一，得□□，為髙股；髙股内又加天元，得□□，為極弦；以自之，得□□□，於上；内減極差冪一萬四千一百六十一，餘□□□，為兩段極積。合以極弦除，不除，寄為母，便以此為城徑；以自增乘，得□□□□□，為圓徑冪{\zhsmall 内有極弦冪分母，寄左}。然後以天元乘髙股，又四之，得□□；又以分母極弦冪□□□通之，得□□□□呔，為同數。與左相消，得□□□，開平方，得六十四步，即平勾也。合問。\end{unitzh}
\begin{uniten}Working: By discernment one obtains: combine the two positions and halve, obtaining □, which is the ji difference. Set up the celestial unknown as the ping base; add 161, obtaining □□, which is the gao altitude; to the gao altitude further add the celestial unknown, obtaining □□, which is the ji hypotenuse; multiply it by itself, obtaining □□□, putting it above; from it subtract the area (mi) of the ji difference, 14161; the remainder □□□ is two pieces of the ji area. It ought to be divided by the ji hypotenuse; not dividing, lodge it as denominator, and simply take this as the diameter of the city; multiply it by itself (zizengcheng), obtaining □□□□□, which is the area of the diameter of the circle [Note: within, there is the area (mi) of the ji hypotenuse as denominator; lodge it on the left]. Then multiply the gao altitude by the celestial unknown, and further multiply it by 4, obtaining □□; further convert it (tong) by the denominator, the area of the ji hypotenuse, □□□, obtaining □□□□呔, which is the same number. Cancel it against the left, obtaining □□□; extract the square root, obtaining 64 bu, which is the ping base. This agrees with the question.\end{uniten}

\begin{unitzh}\unitlabel{ceyuan-025.1}或問：見明和二百七步，叀和四十六步。問答同前。\end{unitzh}
\begin{uniten}Someone asks: Given the Ming sum (minghe) 207 bu (步) and the Zhuan sum (zhuanhe) 46 bu. The question and the answer are the same as before.\end{uniten}

\begin{unitzh}\unitlabel{ceyuan-025.2}法曰：二和上下相減，數同則止，名為泛率。又以二和直相減，餘為泛實{\zhsmall 此則角差也}。乃以泛率除汎實，所得為差率也。以差率加減泛率，若半訖與勾股相應者，其泛率便為和率，其泛實便為較率乘和率也；若不相應，則直取差率以消息之，定為相管和率{\zhsmall 其勾股數少、得見弦黄而相為率者：勾三股四，則其和七而其較一也；勾五股十二，則其和一十七而其較七也；勾八股十五，則其和二十三而其較亦得七七；勾七股二十四，則其和三十一而其較一十七也；勾九股二十，則其和二十九而其較一十一也。此消息之大畧也，餘皆倣此}。乃以和率約二和，其明和所得為明壘率，其叀和所得為叀壘率也。又副置和率，上加差率而半之，則為股率也；下位減差率而半之，則為勾率也。既見勾股及差三率，各以壘率乘之，即各得勾股及差之真數也。\end{unitzh}
\begin{uniten}Rule: Subtract the two sums from each other, upper and lower, stopping when the numbers are equal; call this the rough rate (fanlü). Further, subtract the two sums from each other directly; the remainder is the rough dividend (fanshi) [Note: this is the corner difference (jiaocha)]. Then divide the rough dividend by the rough rate; what is obtained is the difference rate (chalü). Add the difference rate to the rough rate and subtract it from it; if, when the halving is done, these correspond to a base (gou) and an altitude (gu), then the rough rate is the sum rate (helü), and the rough dividend is the gap rate (jiaolü) multiplied by the sum rate. If they do not correspond, then take the difference rate directly and adjust it (xiaoxi), fixing the associated sum rate (xiangguan helü) [Note: Base and altitude numbers that are small, for which a hypotenuse (xian) and a yellow (huang) are obtained, and which serve as rates for each other: base 3 and altitude 4, then their sum is 7 and their gap 1; base 5 and altitude 12, then their sum is 17 and their gap 7; base 8 and altitude 15, then their sum is 23 and their gap also gives 7 7 [Translator's note: the doubled 七 "seven" appears to be corrupt.]; base 7 and altitude 24, then their sum is 31 and their gap 17; base 9 and altitude 20, then their sum is 29 and their gap 11. This is the general outline of the adjusting; all the rest imitate this]. Then reduce (yue) the two sums by the sum rate: what is obtained from the Ming sum is the Ming piling rate (leilü), and what is obtained from the Zhuan sum is the Zhuan piling rate. Further set out the sum rate in duplicate (fuzhi): to the upper add the difference rate and halve it, which gives the altitude rate; from the lower position subtract the difference rate and halve it, which gives the base rate. Having found the three rates of base, altitude and difference, multiply each by the piling rate; then one obtains in each case the true numbers (zhenshu) of the base, altitude and difference.\end{uniten}

\begin{unitzh}\unitlabel{ceyuan-025.3}按：此用約分，以勾股率數求之，甚為省便。然必兩數度盡而得數最小者方可用；若兩數不能度盡，或雖度盡而得數尚大者，轉屬繁難，故又設後法。\end{unitzh}
\begin{uniten}Note: This uses reduction of fractions (yuefen) and seeks it by the rate numbers of base and altitude; it is very economical and convenient. But it can be used only when the two numbers are measured out completely (dujin) and the numbers obtained are the smallest; if the two numbers cannot be measured out completely, or if, though they are measured out completely, the numbers obtained are still large, it turns into something cumbersome and difficult; therefore a later method is also set out.\end{uniten}

\begin{unitzh}\unitlabel{ceyuan-025.4}又法：二云數相併，以自乘於上，二之；云數相乘，又四之；以相併，以四分半乘之，又四之，以併入上位，為從方；以七十步零四分三釐七毫五絲為常法，得叀小差四步。\end{unitzh}
\begin{uniten}Another method: Add the two stated numbers (yunshu) together and multiply it by itself, in the upper position; multiply it by 2. Multiply the stated numbers together and further multiply by 4. Add them together, multiply it by 4 1/2 fen (分), further multiply by 4, and add it into the upper position: this is the joining square (congfang). Take 70 bu and 4 fen 3 li (釐) 7 hao (毫) 5 si (絲) as the constant divisor (changfa); one obtains the Zhuan minor difference (xiaocha), 4 bu.\end{uniten}

\begin{unitzh}\unitlabel{ceyuan-025.5}按：此法未求實數，其求從隅皆用本題數，不可通用。今依細草意另演一法於後，亦惟二和數可以度盡者用之；若不能度盡者，仍用寄分為便。\end{unitzh}
\begin{uniten}Note: This method has not found the number of the dividend (shi); its finding of the joining (cong) and the corner (yu) all uses the numbers of this problem, and it cannot be used generally. Now, following the idea of the detailed working (xicao), another method is separately worked out below; it too is to be used only where the numbers of the two sums can be measured out completely; where they cannot be measured out completely, it is still convenient to use lodged fractions (jifen).\end{uniten}

\begin{unitzh}\unitlabel{ceyuan-025.6}法曰：二和數相減，自之，為平方實。叀和除明和得數，自而倍之，内減四之除得數，再加二單數，以乘二和相併之數，為從。除得數自而四之於上；又以除得數自乘，内減四之除得數，外加一單數，自之，以減上位，為常法。得叀小差。\end{unitzh}
\begin{uniten}Rule: Subtract the numbers of the two sums from each other and multiply it by itself: this is the dividend (shi) for the square root (pingfang). Divide the Ming sum by the Zhuan sum; multiply the number obtained by itself and double it, from this subtract 4 times the number obtained by the division, and again add 2 units; with this multiply the number of the two sums added together: this is the joining (cong). Multiply the number obtained by the division by itself and multiply it by 4, in the upper position; further multiply the number obtained by the division by itself, from this subtract 4 times the number obtained by the division, and add on besides 1 unit; multiply this by itself and subtract it from the upper position: this is the constant divisor. One obtains the Zhuan minor difference.\end{uniten}

\begin{unitzh}\unitlabel{ceyuan-025.7}草曰：以二和相約，命得叀率一、明率四步半，其兩數大小差率並同。又别得明小差、叀大差俱為半虚黄也。立天元為叀小差，以四歩半乘之，得□元，為□大差也，又為明小差，又為半虚黄。置此□大差，又以四步半乘之，得□，為明大差也。其四差相併，得□，減於二和併，得□□，即兩段太虚大小差併也。内加三段虚黄方□，得□□，合成一个太虚三事和，即圓城徑也。以自增乘，得□□□，為徑冪{\zhsmall 寄左}。乃置叀和，加半虚黄，得□□，為平勾；又置明和，内加半虚黄，得□□，為髙股。勾相乘，得下式□□□；又四之，得□□□，為同數，與左相消，得下式□□□。開平方，得四步，即叀小差也。合問。\end{unitzh}
\begin{uniten}Working: Reduce the two sums by each other; naming the result gives the Zhuan rate 1 and the Ming rate 4 1/2 bu, and for the two numbers the rates of the major and minor differences are likewise the same. Further, it is separately obtained that the Ming minor difference and the Zhuan major difference are both half the Xu yellow (xuhuang). Set up the celestial unknown (tianyuan) as the Zhuan minor difference; multiply it by 4 1/2 bu, obtaining □, marked yuan (元) [Translator's note: □ marks a counting-rod number, diagram or character not reproduced in the transcription.], as the □ major difference (dacha); it is also the Ming minor difference, and also half the Xu yellow. Set out this □ major difference and further multiply it by 4 1/2 bu, obtaining □ as the Ming major difference. Add the four differences together, obtaining □; subtract this from the two sums added together, obtaining □□, which is 2 portions (duan) of the Taixu major and minor differences added together. To this add 3 portions of the Xu yellow square (xuhuangfang), □, obtaining □□, which together make one Taixu sum of the three items (sanshi he), that is, the diameter of the round city (yuancheng). Multiply it by itself (zizeng cheng), obtaining □□□ as the area (mi) of the diameter [Note: set aside on the left (jizuo)]. Then set out the Zhuan sum and add half the Xu yellow, obtaining □□ as the Ping base (pinggou); further set out the Ming sum and add to it half the Xu yellow, obtaining □□ as the Gao altitude (gaogu). Multiply it with the base, obtaining the following form (xiashi) □□□; further multiply it by 4, obtaining □□□ as the equal quantity (tongshu), and cancel it against the left (xiangxiao), obtaining the following form □□□. Extract the square root, obtaining 4 bu, which is the Zhuan minor difference. This agrees with the question.\end{uniten}

\begin{unitzh}\unitlabel{ceyuan-025.8}或問：明叀二勾共八十八步，明叀二股共一百六十五步。問答同前。\end{unitzh}
\begin{uniten}Someone asks: The two bases of Ming and Zhuan together are 88 bu, and the two altitudes of Ming and Zhuan together are 165 bu. The question and the answer are the same as before.\end{uniten}

\begin{unitzh}\unitlabel{ceyuan-025.9}法曰：先識别得二大差共、二小差共及四差共，乃以二大差、二小差相乘為實，以四差共為法，如法得半之虚黄方。\end{unitzh}
\begin{uniten}Rule: First, by discernment (shibie), obtain the two major differences together, the two minor differences together and the four differences together; then multiply the two major differences and the two minor differences together as the dividend (shi), and take the four differences together as the divisor (fa); dividing by the divisor gives half of the Xu yellow square.\end{uniten}

\begin{unitzh}\unitlabel{ceyuan-025.10}草曰：先置前後云數，以約法約之，得一十一，即壘率也。復各置前後數，如壘率而一，前得八，即勾率也；後得一十五，即股率也。再以勾股率求得較率七、和率二十三、弦率一十七、黄方率六、大差率九、小差率二，即見諸率。各以壘率乘之：其二和共得□，二較共得□，二弦共得□，二黄共得□，二大差共□，二小差共□，四差共□。已上皆為明叀所得之共數也。乃立天元一為半虚黄，便為明小差，又為叀大差也。以減於大差共，得□□，即明大差也；又以減於小差共，得□□，即叀小差也。以二數相增乘，得丨□□{\zhsmall 寄左}。以天元冪與寄左相消，得□□，上法下實，得一十八步，即半之虚黄方也。\end{unitzh}
\begin{uniten}Working: First set out the preceding and following stated numbers and reduce them by the method of reduction (yuefa), obtaining 11, which is the piling rate. Again set out each of the preceding and following numbers and divide by the piling rate: the preceding gives 8, which is the base rate; the following gives 15, which is the altitude rate. Further, from the base and altitude rates find the gap rate 7, the sum rate 23, the hypotenuse rate 17, the yellow square rate (huangfang lü) 6, the major difference rate 9 and the minor difference rate 2; then all the rates are found. Multiply each by the piling rate: the two sums together give □, the two gaps together give □, the two hypotenuses together give □, the two yellows together give □, the two major differences together □, the two minor differences together □, the four differences together □. All the above are the combined numbers obtained for Ming and Zhuan. Then set up the celestial unknown as half the Xu yellow; this is then the Ming minor difference, and also the Zhuan major difference. Subtract it from the major differences together, obtaining □□, which is the Ming major difference; further subtract it from the minor differences together, obtaining □□, which is the Zhuan minor difference. Multiply the two numbers by each other (xiangzeng cheng), obtaining 丨□□ [Translator's note: 丨 is the counting-rod numeral 1.] [Note: set aside on the left]. Cancel the area (mi) of the celestial unknown against what is set aside on the left, obtaining □□; the upper is the divisor (fa) and the lower the dividend (shi); one obtains 18 bu, which is half of the Xu yellow square.\end{uniten}

\begin{unitzh}\unitlabel{ceyuan-025.11}以倍之，得□，又加於二黄共六十六，共得一百二，即明勾叀股共也，又為極黄方，又為虚弦也。又以三十六減於一百八十七，餘一百五十一，即明股叀勾共也。此數内減虚弦，餘□，為明叀二差較也，此名傍差。以旁差減二弦共一百八十七，餘得□，即太虚和也。却加入虚弦一百二，併得□，為太虚三事和，即圓城徑也。合問。\end{unitzh}
\begin{uniten}Double it, obtaining □; further add it to the two yellows together, 66, obtaining together 102, which is the Ming base and the Zhuan altitude together; it is also the Ji yellow square (ji huangfang), and also the Xu hypotenuse (xuxian). Further subtract 36 from 187, leaving 151, which is the Ming altitude and the Zhuan base together. From this number subtract the Xu hypotenuse, leaving □, which is the gap of the two differences of Ming and Zhuan; this is called the side difference (pangcha). Subtract the side difference from the two hypotenuses together, 187; the remainder obtained, □, is the Taixu sum. Then add in the Xu hypotenuse, 102; combined this gives □, which is the Taixu sum of the three items, that is, the diameter of the round city. This agrees with the question.\end{uniten}

\begin{unitzh}\unitlabel{ceyuan-025.12}又或以虚黄方加於上和共二百五十三，得□，為極弦也；以旁差減極弦，餘二百四十步，亦同。\end{unitzh}
\begin{uniten}Or again: add the Xu yellow square to the sums together above, 253, obtaining □, which is the Ji hypotenuse (jixian); subtract the side difference from the Ji hypotenuse, leaving 240 bu; this is also the same.\end{uniten}

\begin{unitzh}\unitlabel{ceyuan-025.13}又或前後副置勾股較和弦黄六率在地，前以小差率二因之，則勾得□，股得□，較得□，和得□，弦得□，黄得□，即叀段各數也；後以大差率九因之，則勾得□，股得□，較得□，和得□，弦得□，黄得□，即明段各數也。既得明叀各數，餘可知{\zhsmall 按：此因明弦即皇極形勾弦差，叀弦即皇極形股弦差，故以小差率乘各率，即得叀段各數；以大差率乘各率，即得明段各數也}。\end{unitzh}
\begin{uniten}Or again: in front and behind, set out in duplicate the six rates of base, altitude, gap, sum, hypotenuse and yellow on the ground. Multiply (yin) the front ones by the minor difference rate 2: then the base gives □, the altitude gives □, the gap gives □, the sum gives □, the hypotenuse gives □, the yellow gives □; these are the several numbers of the Zhuan portion. Multiply the back ones by the major difference rate 9: then the base gives □, the altitude gives □, the gap gives □, the sum gives □, the hypotenuse gives □, the yellow gives □; these are the several numbers of the Ming portion. Once the several numbers of Ming and Zhuan are obtained, the rest can be known [Siku editors: This is because the Ming hypotenuse is the base–hypotenuse difference of the Huangji figure, and the Zhuan hypotenuse is the altitude–hypotenuse difference of the Huangji figure; therefore multiplying the rates by the minor difference rate gives the several numbers of the Zhuan portion, and multiplying the rates by the major difference rate gives the several numbers of the Ming portion].\end{uniten}

\begin{unitzh}\unitlabel{ceyuan-025.14}按：右二卷明叀前十八問、後十六問，在集中尤為神妙，惜其中有偶爾思省未至者，亦未暇修飾故耳。\end{unitzh}
\begin{uniten}Note: In the two juan on the right, the Ming–Zhuan former 18 problems and latter 16 problems are especially marvellous within the collection; it is a pity that among them there are some in which reflection occasionally fell short, which is simply because there was also no leisure to polish them.\end{uniten}

\begin{unitzh}\unitlabel{ceyuan-025.15}測圓海鏡卷八\end{unitzh}
\begin{uniten}Ceyuan haijing (Sea Mirror of Circle Measurements), juan 8\end{uniten}

\begin{unitzh}\unitlabel{ceyuan-025.16}子部，天文算法類，算書之屬，測圓海鏡\end{unitzh}
\begin{uniten}Masters section (zibu); category of astronomy and mathematics; subcategory of mathematical books; Ceyuan haijing.\end{uniten}

\section*{卷09}\addcontentsline{toc}{section}{卷09}

\begin{unitzh}\unitlabel{ceyuan-026.1}欽定四庫全書。《測圓海鏡》卷九。元·李冶撰。\end{unitzh}
\begin{uniten}Imperially Commissioned Complete Library of the Four Treasuries (Siku quanshu). Sea Mirror of Circle Measurements, juan 9. Composed by Li Ye of the Yuan.\end{uniten}

\begin{unitzh}\unitlabel{ceyuan-026.2}大斜四問\end{unitzh}
\begin{uniten}Four problems on the great oblique (da xie)\end{uniten}

\begin{unitzh}\unitlabel{ceyuan-026.3}或問：甲、丙俱在中心，丙望南門直行，不知步數而止；甲出東門直行，不知步數，望見丙，斜行與丙相㑹。二人共行了六百八十步，仍云甲直行少於丙直行一百一十九步。問答同前。\end{unitzh}
\begin{uniten}Someone asks: Jia and Bing are both at the center. Bing, facing the south gate, walks straight, not knowing the number of bu (步), and stops. Jia goes out of the east gate and walks straight, not knowing the number of bu, catches sight of Bing, and walks obliquely to meet Bing. The two people together walked 680 bu, and it is further said that Jia's straight walk is less than Bing's straight walk by 119 bu. The question and answer are the same as before.\end{uniten}

\begin{unitzh}\unitlabel{ceyuan-026.4}法曰：二數相減，餘以為冪，内却減差冪，為平實。二數相減，又四之於上，又加入二之差步，為益從。二步常法。得皇極勾。\end{unitzh}
\begin{uniten}Rule: Subtract the two numbers from each other and make the remainder into an area (mi); from it, in turn, subtract the area of the difference; this is the square dividend (ping shi). Subtract the two numbers from each other, again multiply by 4, [placing it] above; again add in twice the bu of the difference; this is the augmenting cong (yi cong) [Translator's note: cong is the coefficient of the first power; yi, "augmenting", marks it as negative]. 2 bu is the constant divisor (chang fa). One obtains the huangji ("august pole") base (gou).\end{uniten}

\begin{unitzh}\unitlabel{ceyuan-026.5}草曰：别得共步即皇極三事和，少步即勾股差也。立天元一為皇極勾，加少步得□□，為股也。又以天元加股得□□，為和也。以和減共步得□□，為弦也。弦自之得□□□，為一段弦冪{\zhsmall 寄左}。然後置股，以天元乘之，又倍之，得□□，為二直積；加入少步冪□，共得□□□，為同數。與左相消，得□□□，平方而一，得一百三十六，即勾也。勾加差為股，勾股相乘倍之為實，勾股和減共步為法，得城徑。\end{unitzh}
\begin{uniten}Working: By distinguishing, one obtains that the total bu is the huangji sum of the three elements (sanshi he), and the "less" bu (shao bu) is the difference of base and altitude (gu). Set up the celestial unknown as the huangji base. Add the "less" bu, obtaining □□ [Translator's note: □ marks a row of a counting-rod array not reproduced in the transcription] as the altitude. Again add the celestial unknown to the altitude, obtaining □□ as the sum. Subtract the sum from the total bu, obtaining □□ as the hypotenuse (xian). Multiply the hypotenuse by itself, obtaining □□□ as 1 portion of the hypotenuse area [Note: set aside on the left (ji zuo)]. Then set out the altitude, multiply it by the celestial unknown, and double it again, obtaining □□ as 2 rectangle areas (zhi ji). Add in the area of the "less" bu, □, obtaining altogether □□□ as the equal quantity (tong shu). Cancel it against the left (xiang xiao), obtaining □□□, and divide by square root extraction (pingfang er yi), obtaining 136, which is the base. Add the difference to the base to make the altitude. Multiply base and altitude together and double as the dividend (shi); subtract the sum of base and altitude from the total bu as the divisor (fa); one obtains the city diameter (cheng jing).\end{uniten}

\begin{unitzh}\unitlabel{ceyuan-026.6}又法：云數併與云數差相乘{\zhsmall 按：此句有誤，當云和數與倍差相加相減，二得數相乘}，為平實。云數併與二數差相併，得數以減於八之共步，為益從{\zhsmall 按：此只云六因和步為益從，亦同}。一步常法。得皇極黄方。\end{unitzh}
\begin{uniten}Another rule: Multiply the sum of the stated numbers (yun shu) by the difference of the stated numbers [Siku editors: We note that this sentence is in error; it should say: add and subtract the sum number and twice the difference, and multiply the two numbers obtained together] as the square dividend. Add the sum of the stated numbers and the difference of the two numbers, and subtract the number obtained from 8 times the total bu, as the augmenting cong [Siku editors: We note that here one may simply say "multiply the sum bu by 6 as the augmenting cong"; it is the same]. 1 bu is the constant divisor. One obtains the huangji yellow square (huangfang).\end{uniten}

\begin{unitzh}\unitlabel{ceyuan-026.7}草曰：立天元一為黄方{\zhsmall 即虚弦也}，副置之。上位加共步，得□□，為二和也；下位減共步，得□□，為二弦也。先以二和自乘，得丨□□，為四段和冪；又以二弦自乘，得丨□□，為四段弦冪。二數相減，餘得□，又倍之，得下式□，為十六段直積於天元位{\zhsmall 寄左}。然後副置二和，上位加二之少步，得□□，為四股；下位減二之少步，得□□，為四勾。勾股相乘，得丨□□，為同數。與左相消，得□□□，平方而一，得一百二步，即皇極黄方也。餘各依法求之，合問。\end{unitzh}
\begin{uniten}Working: Set up the celestial unknown as the yellow square [Note: that is, the void hypotenuse (xu xian)], and set it out in duplicate (fu zhi). To the upper position add the total bu, obtaining □□ as twice the sum; the lower position is subtracted from the total bu, obtaining □□ as twice the hypotenuse. First multiply twice the sum by itself, obtaining 丨□□ [Translator's note: 丨 is the rod numeral 1] as 4 portions of the sum area; again multiply twice the hypotenuse by itself, obtaining 丨□□ as 4 portions of the hypotenuse area. Subtract the two numbers from each other; the remainder obtained is □. Double it again, obtaining the array below (xia shi), □, as 16 portions of the rectangle area, in the position of the celestial unknown [Note: set aside on the left]. Then set out twice the sum in duplicate. To the upper position add twice the "less" bu, obtaining □□ as 4 altitudes; from the lower position subtract twice the "less" bu, obtaining □□ as 4 bases. Multiply base and altitude together, obtaining 丨□□ as the equal quantity. Cancel against the left, obtaining □□□, and divide by square root extraction, obtaining 102 bu, which is the huangji yellow square. For the rest, find each according to the rules. This agrees with the question.\end{uniten}

\begin{unitzh}\unitlabel{ceyuan-026.8}或問：甲、丙俱在西北隅起，丙向南行，不知步數而立；甲向東行，望見丙，就丙斜行六百八十步，與丙相㑹。丙云：「我南行步多於甲東行二百八十步。」問答同前。\end{unitzh}
\begin{uniten}Someone asks: Jia and Bing both set out from the northwest corner. Bing walks south, not knowing the number of bu, and stands. Jia walks east, catches sight of Bing, goes obliquely toward Bing 680 bu, and meets Bing. Bing says: "My southward bu exceed Jia's eastward walk by 280 bu." The question and answer are the same as before.\end{uniten}

\begin{unitzh}\unitlabel{ceyuan-026.9}法曰：以云數差乘云數併，為實；倍多步，為從；二為平隅。得大勾。\end{unitzh}
\begin{uniten}Rule: Multiply the sum of the stated numbers by the difference of the stated numbers as the dividend; double the "more" bu (duo bu) as the cong; 2 is the square corner (ping yu). One obtains the great base (da gou).\end{uniten}

\begin{unitzh}\unitlabel{ceyuan-026.10}草曰：立天元為大平{\zhsmall 按：大平即大勾}，加差得□□，為股。倍天元乘之，得□□，為二積{\zhsmall 寄左}。然後以斜步多步併□與斜步多步較□相乘，得□，為同數。與左相消，得□□□，開平方，得三百二十步，即大勾也。合問。\end{unitzh}
\begin{uniten}Working: Set up the celestial unknown as the great level (da ping) [Siku editors: We note that the great level is the great base]. Add the difference, obtaining □□ as the altitude. Multiply it by twice the celestial unknown, obtaining □□ as 2 areas [Note: set aside on the left]. Then multiply the sum of the oblique bu and the "more" bu, □, by the difference (jiao) of the oblique bu and the "more" bu, □, obtaining □ as the equal quantity. Cancel against the left, obtaining □□□, and extract the square root, obtaining 320 bu, which is the great base. This agrees with the question.\end{uniten}

\begin{unitzh}\unitlabel{ceyuan-026.11}或問：甲、乙二人共立於艮隅，乙南行，過城外而立；甲東行，望乙，與城叅相直而止。丙、丁二人共立於坤隅，丁向東行，過城門而立；丙向南行，望丁及甲、乙，悉與城俱相直。丙復就甲斜行六百八十步，與甲相㑹。乙、丁又云：「吾二人直行共得三百四十二。」問答同前。\end{unitzh}
\begin{uniten}Someone asks: Jia and Yi, two people, stand together at the Gen (northeast) corner. Yi walks south, passes outside the city, and stands. Jia walks east, looks at Yi in alignment with the city (can xiang zhi), and stops. Bing and Ding, two people, stand together at the Kun (southwest) corner. Ding walks east, passes the city gate, and stands. Bing walks south and looks at Ding as well as Jia and Yi, all aligned together with the city. Bing then goes obliquely toward Jia 680 bu and meets Jia. Yi and Ding further say: "The straight walks of us two together come to 342." The question and answer are the same as before.\end{uniten}

\begin{unitzh}\unitlabel{ceyuan-026.12}法曰：二云數相乘，倍之為實。倍斜行於上，以二云數相減，加上位，為從。一步常法。開平方，得城徑。\end{unitzh}
\begin{uniten}Rule: Multiply the two stated numbers together and double as the dividend. Double the oblique walk, [placing it] above; subtract the two stated numbers from each other and add [the remainder] to the upper position, as the cong. 1 bu is the constant divisor. Extract the square root, obtaining the city diameter.\end{uniten}

\begin{unitzh}\unitlabel{ceyuan-026.13}草曰：别得斜步即大弦也，其共步則一徑、一虚弦共也，其二數相併為一大和、一虚弦共數也。立天元為徑，減於共步，得□□，為虚弦也。以虚弦復減於天元，得□□，為虚和；以斜步乘之，得□□{\zhsmall 寄左}。乃以天元加斜步，得□□，為大和；以虚弦乘之，得□□□，為同數。與左相消，得丨□□，開平方，得二百四十步，即城徑也。合問。\end{unitzh}
\begin{uniten}Working: By distinguishing, one obtains that the oblique bu is the great hypotenuse (da xian); the total bu is one diameter and one void hypotenuse together; and the two numbers added together are the number of one great sum (da he) and one void hypotenuse together. Set up the celestial unknown as the diameter and subtract it from the total bu, obtaining □□ as the void hypotenuse. Again subtract the void hypotenuse from the celestial unknown, obtaining □□ as the void sum (xu he); multiply it by the oblique bu, obtaining □□ [Note: set aside on the left]. Then add the oblique bu to the celestial unknown, obtaining □□ as the great sum; multiply it by the void hypotenuse, obtaining □□□ as the equal quantity. Cancel against the left, obtaining 丨□□, and extract the square root, obtaining 240 bu, which is the city diameter. This agrees with the question.\end{uniten}

\begin{unitzh}\unitlabel{ceyuan-026.14}或問：甲從北門向東直行，庚從西門穿城東行，丙從西門向南直行，壬從北門穿城南行，四人遥相望，悉與城叅相直。只云丙相望處六百八十步，庚、壬穿城共行了六百三十一步。問答同前。\end{unitzh}
\begin{uniten}Someone asks: Jia, from the north gate, walks straight eastward; Geng, from the west gate, goes through the city walking east; Bing, from the west gate, walks straight southward; Ren, from the north gate, goes through the city walking south. The four people look at each other from afar, all aligned with the city. It is only said that where Bing [and Jia] look at each other it is 680 bu, and that Geng and Ren, going through the city, together walked 631 bu. The question and answer are the same as before.\end{uniten}

\begin{unitzh}\unitlabel{ceyuan-026.15}法曰：共步自之得數；以共步減斜，餘自乘，以減上，為實。二之斜步，加入共步減斜餘數，為從。一步常法。得城徑。\end{unitzh}
\begin{uniten}Rule: Multiply the total bu by itself to obtain a number. Subtract the total bu from the oblique, multiply the remainder by itself, and subtract it from the above, as the dividend. Take twice the oblique bu and add in the remainder of the total bu subtracted from the oblique, as the cong. 1 bu is the constant divisor. One obtains the city diameter.\end{uniten}

\begin{unitzh}\unitlabel{ceyuan-026.16}草曰：共行步為一徑與皇和共也，又為大和皇弦差也；甲、丙相望即大弦也。以共步減大弦，餘□，為皇極弦上減一徑也。立天元一為圓徑，減於共步，得□□，為皇極和也；以自之，得丨□□，於上。弦内減共步，餘□，又以天元加之，得□□，為皇極弦；以自之，得丨□□，減上位，餘得□□，為兩个皇直積{\zhsmall 寄左}。乃以天元乘皇弦，得下式丨□，為同數。與左相消，得丨□□，開平方，得二百四十步，即城徑也。合問。\end{unitzh}
\begin{uniten}Working: The bu walked together are one diameter and the huang sum (huang he) together; they are also the difference of the great sum and the huang hypotenuse (huang xian). [The distance at which] Jia and Bing look at each other is the great hypotenuse. Subtract the total bu from the great hypotenuse; the remainder, □, is the huangji hypotenuse with one diameter subtracted. Set up the celestial unknown as the circle diameter (yuan jing) and subtract it from the total bu, obtaining □□ as the huangji sum. Multiply it by itself, obtaining 丨□□, [placed] above. From the hypotenuse subtract the total bu, the remainder being □. Again add the celestial unknown to it, obtaining □□ as the huangji hypotenuse. Multiply it by itself, obtaining 丨□□, and subtract it from the upper position; the remainder obtained, □□, is 2 huang rectangle areas [Note: set aside on the left]. Then multiply the huang hypotenuse by the celestial unknown, obtaining the array below, 丨□, as the equal quantity. Cancel against the left, obtaining 丨□□, and extract the square root, obtaining 240 bu, which is the city diameter. This agrees with the question.\end{uniten}

\begin{unitzh}\unitlabel{ceyuan-026.17}大和八問\end{unitzh}
\begin{uniten}Eight problems on the great sum (da he)\end{uniten}

\begin{unitzh}\unitlabel{ceyuan-026.18}或問：庚從西門穿城東行二百五十六步而立，壬從北門穿城南行三百七十五步而立。又有甲、丙二人俱在乾隅，甲向東行，丙向南行，各不知步數而立。四人遙相望，只云甲、丙共行了九百二十步。問答同前。\end{unitzh}
\begin{uniten}Someone asks: Geng, from the west gate, goes through the city walking east 256 bu and stands; Ren, from the north gate, goes through the city walking south 375 bu and stands. There are also two people, Jia and Bing, both at the Qian (northwest) corner. Jia walks east and Bing walks south, each not knowing the number of bu, and they stand. The four people look at each other from afar. It is only said that Jia and Bing together walked 920 bu. The question and answer are the same as before.\end{uniten}

\begin{unitzh}\unitlabel{ceyuan-026.19}法曰：庚東行冪、壬南行冪相併於上。併庚壬步而倍之，内減大和，餘復減於庚壬共，得數{\zhsmall 按：或云併庚壬步以減大和，亦同}，以自乘，減上位，為平實。併庚壬步，為益從。半步為隅法。得城徑。\end{unitzh}
\begin{uniten}Rule: Add together the area of Geng's eastward walk and the area of Ren's southward walk, [placing it] above. Add Geng's and Ren's bu and double; from this subtract the great sum; subtract the remainder in turn from the total of Geng and Ren, obtaining a number [Siku editors: We note that one may also say: add Geng's and Ren's bu and subtract it from the great sum; it is the same]. Multiply it by itself and subtract it from the upper position, as the square dividend. Add Geng's and Ren's bu as the augmenting cong. 1/2 bu is the corner divisor (yu fa). One obtains the city diameter.\end{uniten}

\begin{unitzh}\unitlabel{ceyuan-026.20}草曰：立天元一為圓徑，以半之，副置二位：上以減於庚東行，得□□，為平弦也；下以減於壬南行，得□□，為髙弦也。二弦相併，得□□，為皇弦虚弦共也。倍此數，得□□，為大弦虚弦共也。以大弦虚弦共減於大和，餘□□，為虚勾虚股共也。天元内減虚勾虚股共，餘□□，即虚弦也。復置皇弦虚弦共，内減虚弦，餘□□，即皇極弦也。以自之，得□□{\zhsmall 寄左}。然後以平弦自之，得下式□□□，為勾冪也；又以髙弦自之，得□□□，為股冪也。二冪相併，得□□□，為同數。與左相消，得□□□，平方而一，得二百四十步，即城徑也。合問。\end{unitzh}
\begin{uniten}Working: Set up the celestial unknown as the circle diameter; halve it and set it out in two positions. The upper is subtracted from Geng's eastward walk, obtaining □□ as the level hypotenuse (ping xian); the lower is subtracted from Ren's southward walk, obtaining □□ as the height hypotenuse (gao xian). Add the two hypotenuses together, obtaining □□ as the huang hypotenuse and the void hypotenuse together. Double this number, obtaining □□ as the great hypotenuse and the void hypotenuse together. Subtract the great hypotenuse and void hypotenuse together from the great sum; the remainder, □□, is the void base (xu gou) and void altitude (xu gu) together. From the celestial unknown subtract the void base and void altitude together; the remainder, □□, is the void hypotenuse. Again set out the huang hypotenuse and void hypotenuse together and from it subtract the void hypotenuse; the remainder, □□, is the huangji hypotenuse. Multiply it by itself, obtaining □□ [Note: set aside on the left]. Then multiply the level hypotenuse by itself, obtaining the array below, □□□, as the base area; again multiply the height hypotenuse by itself, obtaining □□□ as the altitude area. Add the two areas together, obtaining □□□ as the equal quantity. Cancel against the left, obtaining □□□, and divide by square root extraction, obtaining 240 bu, which is the city diameter. This agrees with the question.\end{uniten}

\begin{unitzh}\unitlabel{ceyuan-026.21}或問：丙、甲俱在西北隅，甲向東行，不知步數而立；丙向南行，望見甲，就甲斜行，與甲相㑹。甲直行、丙直行共九百二十步{\zhsmall 甲步少於丙步}。又出東門南行有柳樹一株，出南門東行有槐樹一株。戊、己二人同在巽隅，戊就柳樹，已從槐樹，亦與甲、乙遥相望。只云已行少於戊行數與兩樹相距數相併，得一百四十四步；其二數相減，餘六十步。問答同前。\end{unitzh}
\begin{uniten}Someone asks: Bing and Jia are both at the northwest corner. Jia walks east, not knowing the number of bu, and stands. Bing walks south, catches sight of Jia, goes obliquely toward Jia, and meets Jia. Jia's straight walk and Bing's straight walk together are 920 bu [Note: Jia's bu are fewer than Bing's bu]. Further, going out of the east gate and walking south, there is one willow tree; going out of the south gate and walking east, there is one locust tree. Wu and Ji, two people, are together at the Xun (southeast) corner. Wu goes to the willow tree and Ji [Translator's note: here and below the text writes 已 for 己 Ji] makes for the locust tree, and they too look at Jia and Yi from afar. It is only said that the number by which Ji's walk is less than Wu's walk and the number of the distance between the two trees, added together, give 144 bu, and that the two numbers subtracted from each other leave 60 bu. The question and answer are the same as before.\end{uniten}

\begin{unitzh}\unitlabel{ceyuan-026.22}法曰：二云數相併而半之，為虚弦，以乘大和九百二十步於上；以一百四十四減大和，以虚較乘之，減上位，為平實。以一百四十四減大和，又二之於上，以二之虚較減上位{\zhsmall 按：或云倍甲丙直行共，加己戊較與兩樹距之較，減三之己戊較與兩樹距之和，亦同}，為從。四虚隅。得太虚勾。\end{unitzh}
\begin{uniten}Rule: Add the two stated numbers together and halve, as the void hypotenuse; multiply the great sum, 920 bu, by it, [placing it] above. Subtract 144 from the great sum, multiply by the void difference (xu jiao), and subtract from the upper position, as the square dividend. Subtract 144 from the great sum, again double it, [placing it] above; subtract twice the void difference from the upper position [Siku editors: We note that one may also say: double the total of Jia's and Bing's straight walks, add the difference between the Ji–Wu difference and the distance of the two trees, and subtract 3 times the sum of the Ji–Wu difference and the distance of the two trees; it is the same], as the cong. 4 is the empty corner (xu yu) [Translator's note: xu, "empty", here marks the corner coefficient as negative]. One obtains the base of the great void (taixu).\end{uniten}

\begin{unitzh}\unitlabel{ceyuan-026.23}草曰：别得甲丙直行共即大和也，戊就柳樹步即虚股也，已就槐樹步即虚勾也，其一百四十四步即二明勾，其六十步即二叀股也。立天元一為虚勾，加明勾，得□□，為半徑也；倍之，得□□，即城徑也{\zhsmall 又為虚弦上三事和}。二云數相併而半之，得□，即小弦也；相減而半之，得□，即小較也。以天元加較，得□□，即小股也。小勾股共，得□□，即小和也。以小三事減大和，得□□，即大弦也。乃先置小和，以大弦乘之，得下式□□□{\zhsmall 寄左}。次以小弦乘大和，得□□，與左相消，得下式□□□，開平方，得四十八步，即虚勾也。加明勾，又倍之，得二百四十步，即城徑也。合問。\end{unitzh}
\begin{uniten}Working: By distinguishing, one obtains that the total of Jia's and Bing's straight walks is the great sum; Wu's bu going to the willow tree are the void altitude; Ji's bu going to the locust tree are the void base; the 144 bu are 2 ming ("bright") bases (ming gou); and the 60 bu are 2 zhuan altitudes (zhuan gu). Set up the celestial unknown as the void base. Add the ming base, obtaining □□ as half the diameter. Double it, obtaining □□, which is the city diameter [Note: it is also the sum of the three elements on the void hypotenuse]. Add the two stated numbers together and halve, obtaining □, which is the small hypotenuse (xiao xian); subtract them from each other and halve, obtaining □, which is the small difference (xiao jiao). Add the difference to the celestial unknown, obtaining □□, which is the small altitude. The small base and altitude together give □□, which is the small sum. Subtract the small three elements from the great sum, obtaining □□, which is the great hypotenuse. Then first set out the small sum and multiply it by the great hypotenuse, obtaining the array below, □□□ [Note: set aside on the left]. Next multiply the great sum by the small hypotenuse, obtaining □□. Cancel against the left, obtaining the array below, □□□, and extract the square root, obtaining 48 bu, which is the void base. Add the ming base and again double, obtaining 240 bu, which is the city diameter. This agrees with the question.\end{uniten}

\begin{unitzh}\unitlabel{ceyuan-026.24}或問：甲從乾隅東行，乙從艮隅南行，丙從乾隅南行，丁從坤隅東行，四人遥相望見。既而甲還至艮隅就乙，丙還至坤隅就丁。甲、丙直行共九百二十步，甲還就乙共二百三十步，丙還就丁共五百五十二步。問答同前。\end{unitzh}
\begin{uniten}Someone asks: Jia walks east from the Qian corner, Yi walks south from the Gen corner, Bing walks south from the Qian corner, and Ding walks east from the Kun corner; the four people look and see each other from afar. Then Jia returns to the Gen corner and goes to Yi, and Bing returns to the Kun corner and goes to Ding. Jia's and Bing's straight walks together are 920 bu; Jia's returning and going to Yi together are 230 bu; Bing's returning and going to Ding together are 552 bu. The question and answer are the same as before.\end{uniten}

\begin{unitzh}\unitlabel{ceyuan-026.25}法曰：併就數，以減直行共，復以所併就數乘之，為實。併就數減直行共，得數，復加入直行共，為法。得虚弦。\end{unitzh}
\begin{uniten}Rule: Add the "going-to" numbers (jiu shu) and subtract them from the total of the straight walks; again multiply by the added "going-to" numbers, as the dividend. Add the "going-to" numbers and subtract them from the total of the straight walks; add the number obtained again to the total of the straight walks, as the divisor. One obtains the void hypotenuse.\end{uniten}

\begin{unitzh}\unitlabel{ceyuan-026.26}草曰：别得甲丙直行共為大和也，甲還就乙步為小差勾股共也，丙還就丁步為大差勾股共也。以大差勾股共減於大股，餘即虚勾也；以小差勾股共減於大勾，餘即虚股也。二數相併，得□，為大弦虚弦共也；二數相減，餘□，為通差及大虚勾股差共也。又併二數而半之，得□，為太極弦虚弦共，又為太極勾股共也。立天元一為虚弦，先以二共數減於大和，餘□，為虚勾虚股和，於上；次以虚弦減於二共數，餘□□，為大弦，以乘上位，得下□□{\zhsmall 寄左}。然後以天元乘大和，得□，為同數。與左相消，得□□，上法下實，得一百二步，即虚弦也。加入虚和，得二百四十步，即城徑也。合問。\end{unitzh}
\begin{uniten}Working: By distinguishing, one obtains that the total of Jia's and Bing's straight walks is the great sum; Jia's bu returning and going to Yi are the xiaocha ("small difference") base and altitude together; and Bing's bu returning and going to Ding are the dacha ("great difference") base and altitude together. Subtract the dacha base and altitude together from the great altitude; the remainder is the void base. Subtract the xiaocha base and altitude together from the great base; the remainder is the void altitude. Add the two numbers together, obtaining □ as the great hypotenuse and void hypotenuse together. Subtract the two numbers from each other; the remainder, □, is the tong ("through") difference (tong cha) and the base–altitude difference of the great void (da xu) together. Again add the two numbers and halve, obtaining □ as the taiji ("supreme ultimate") hypotenuse and void hypotenuse together; it is also the taiji base and altitude together. Set up the celestial unknown as the void hypotenuse. First subtract the two totals from the great sum; the remainder, □, is the sum of the void base and void altitude, [placed] above. Next subtract the void hypotenuse from the two totals; the remainder, □□, is the great hypotenuse. Multiply the upper position by it, obtaining below □□ [Note: set aside on the left]. Then multiply the great sum by the celestial unknown, obtaining □ as the equal quantity. Cancel against the left, obtaining □□; with the upper as divisor and the lower as dividend, one obtains 102 bu, which is the void hypotenuse. Add in the void sum, obtaining 240 bu, which is the city diameter. This agrees with the question.\end{uniten}

\begin{unitzh}\unitlabel{ceyuan-026.27}又法：併云數，減大和，復以二數相減乘之，為實。併云數減大和，得數，復加入大和，為法。得虚差。\end{unitzh}
\begin{uniten}Another rule: Add the stated numbers and subtract them from the great sum; again multiply by the difference of the two numbers, as the dividend. Add the stated numbers and subtract them from the great sum; add the number obtained again to the great sum, as the divisor. One obtains the void difference (xu cha).\end{uniten}

\begin{unitzh}\unitlabel{ceyuan-027.1}草曰：立天元一為虚較。先以併云數減大和，餘□為虚和，於上；次以天元減於二就步較□，得□□，為通差，以乘之，得□□{\zhsmall 寄左}。然後以天元乘大和，得□，為同數，與左相消，得□□，上法下實，得四十二步，即虚差也。副置虚和為二位：上加虚差而半之，得九十，即虚股也；下減虚差而半之，得四十八，即虚勾也。勾冪、股冪相併，得□，開平方，得一百二步，即虚弦也。加入虚和，得二百四十步，即城徑也。合問。\end{unitzh}
\begin{uniten}Working: Set up the celestial unknown (tianyuan) as the void (xu) difference (jiao). First subtract the sum of the stated numbers (yunshu) from the great sum (dahe); the remainder □ is the void sum, placed above. Next subtract the celestial unknown from the difference of the two approach paces (jiubu), □, obtaining □□, which is the overall (tong) difference; multiply by it, obtaining □□ [Note: deposit on the left (jizuo)]. Then multiply the great sum by the celestial unknown, obtaining □, as the equal quantity (tongshu); eliminate (xiangxiao) it against the left, obtaining □□; the upper is the divisor (fa), the lower the dividend (shi); obtaining 42 bu (步), which is the void difference (cha). Set out the void sum in duplicate (fuzhi) in two positions: to the upper add the void difference and halve it, obtaining 90, which is the void altitude (gu); from the lower subtract the void difference and halve it, obtaining 48, which is the void base (gou). Sum the area (mi) of the base and the area of the altitude, obtaining □; extract the square root, obtaining 102 bu, which is the void hypotenuse (xian). Add the void sum, obtaining 240 bu, which is the diameter of the city. This agrees with the question.\end{uniten}

\begin{unitzh}\unitlabel{ceyuan-027.2}或問：依前見大和，只云股圓差上勾弦差二百一十六，勾圓差上股弦差二十步。問答同前。\end{unitzh}
\begin{uniten}Someone asks: According to the foregoing, the great sum is known; it only says that the base–hypotenuse difference on the altitude–circle difference (guyuancha) is 216, and the altitude–hypotenuse difference on the base–circle difference (gouyuancha) is 20 bu. The question and answer are the same as before.\end{uniten}

\begin{unitzh}\unitlabel{ceyuan-027.3}法曰：以云數二十步減通和，復以二十步乘之，於上；以云數二百一十六減九百步{\zhsmall 按：即併二差以減大和}而半之，乘上位，為立實。三因二十步以減通和，得八百六十；以二百一十六減通和而半之，得二百四十二；二數相乗訖，内減二十之九百步；又以三百四十二及二百一十六共得五百五十八，又以之以減之，為從方{\zhsmall 按：取從方内語有誤，當云：三因小差減大和，併二差減大和，半之，相乘於上；三因大和，加大差，減三之小差，半之，以小差乘之，得數減上位，為從方}。以二百一十六減通和，又以三之二十步減通和，相併於上；以二之五百五十八内却減二十步，餘以減上位，為益亷{\zhsmall 按：取益亷内語亦有誤，當云：三因大和減六之小差，為益亷}。四步常法，得小差股。\end{unitzh}
\begin{uniten}Rule: Subtract the stated number 20 bu from the overall sum (tonghe), then multiply it by 20 bu, placed above. Subtract the stated number 216 from 900 bu [Siku editors: That is, sum the two differences and subtract them from the great sum] and halve it; multiply the upper position by it, as the cubic dividend (lishi). Triple 20 bu and subtract it from the overall sum, obtaining 860; subtract 216 from the overall sum and halve it, obtaining 242; multiply the two numbers together, and when this is done subtract from it 20 times 900 bu; further take 342 and 216, together obtaining 558, and further use it to subtract from it, as the following square (congfang) [Siku editors: The wording in obtaining the following square is in error. It should say: triple the small difference (xiaocha) and subtract it from the great sum; sum the two differences, subtract them from the great sum and halve; multiply these together, placed above; triple the great sum, add the great difference (dacha), subtract 3 times the small difference, halve, and multiply by the small difference; subtract the resulting number from the upper position, as the following square]. Subtract 216 from the overall sum, and further subtract 3 times 20 bu from the overall sum; sum these, placed above; from 2 times 558 then subtract 20 bu, and with the remainder subtract from the upper position, as the augmenting edge (yilian) [Siku editors: The wording in obtaining the augmenting edge is also in error. It should say: triple the great sum and subtract 6 times the small difference, as the augmenting edge]. 4 bu is the constant divisor (changfa). Obtain the small-difference altitude.\end{uniten}

\begin{unitzh}\unitlabel{ceyuan-027.4}草曰：别得小差上股弦差□，加二股，為大勾也；大差上勾弦差□，加二勾，為大股也。立天元一為小差股，加□，得□□，為小差弦也。小差弦上又加天元，得□□，為通勾，以減於和步，得□□，為通股也。通股内減大差上勾弦差□，得□□，半之，得下式□□，即大差之勾也。大差勾上又加勾弦差□，得□□，為大差弦也。再置通股，以小差弦乘之，得□□□，以天元除之，得□□□，為一个大弦也{\zhsmall 泛寄}。再置通勾，以大差弦乘之，得□□□，合以大差勾除，不除，寄母，便以為大弦{\zhsmall 寄左}。乃以大差勾乘泛寄，得□□□□，為同數，與左相消，得□□□□，益積開立方，得一百五十步，為小差股也。合問。\end{unitzh}
\begin{uniten}Working: Separately we get: the altitude–hypotenuse difference on the small difference, □, with two altitudes added, is the great base; the base–hypotenuse difference on the great difference, □, with two bases added, is the great altitude. Set up the celestial unknown as the small-difference altitude; add □, obtaining □□, which is the small-difference hypotenuse. To the small-difference hypotenuse further add the celestial unknown, obtaining □□, which is the overall base; subtract it from the sum in bu (hebu), obtaining □□, which is the overall altitude. From the overall altitude subtract the base–hypotenuse difference on the great difference, □, obtaining □□; halve it, obtaining the expression below, □□, which is the base of the great difference. To the great-difference base further add the base–hypotenuse difference, □, obtaining □□, which is the great-difference hypotenuse. Again set out the overall altitude and multiply it by the small-difference hypotenuse, obtaining □□□; divide it by the celestial unknown, obtaining □□□, which is one great hypotenuse [Note: provisional deposit (fanji)]. Again set out the overall base and multiply it by the great-difference hypotenuse, obtaining □□□; it ought to be divided by the great-difference base; not dividing, deposit the denominator (jimu), and directly take it as the great hypotenuse [Note: deposit on the left]. Then multiply the provisional deposit by the great-difference base, obtaining □□□□, as the equal quantity; eliminate it against the left, obtaining □□□□; with augmenting volume (yiji) extract the cube root, obtaining 150 bu, which is the small-difference altitude. This agrees with the question.\end{uniten}

\begin{unitzh}\unitlabel{ceyuan-027.5}或問：依見前大和，只云髙弦平弦共得三百九十一步，髙弦平弦相較得一百一十九步。問答同前。\end{unitzh}
\begin{uniten}Someone asks: According to the great sum known before, it only says that the high (gao) hypotenuse and the level (ping) hypotenuse together make 391 bu, and the difference between the high hypotenuse and the level hypotenuse is 119 bu. The question and answer are the same as before.\end{uniten}

\begin{unitzh}\unitlabel{ceyuan-027.6}法曰：以較數冪減於共數冪，又半之，為實；以共數減大和，為益從；一步常法；開平方，得圓徑。\end{unitzh}
\begin{uniten}Rule: Subtract the area (mi) of the difference number from the area of the sum number, and further halve it, as the dividend (shi); subtract the sum number from the great sum, as the augmenting following (yicong); 1 bu is the constant divisor (changfa); extract the square root, obtaining the diameter of the circle.\end{uniten}

\begin{unitzh}\unitlabel{ceyuan-027.7}草曰：别得髙數減於通股，為邊股内減明股也；平弦減於通勾，為邊勾内減明勾也。其共數即大弦内減皇極弦，又為皇極勾股共也；其相較步即皇極差也。二云數相併，即黄廣弦也；二云數相減，餘即黄長弦也。以共數減於大和，餘□為皇極弦與圓徑共。立天元一為圓徑，以減皇極弦與圓徑共，得□□，為皇極弦也。以共數自之，得□，於上；以相較數自之，得□，減上位，餘□，又半之，得□，為兩段皇極積{\zhsmall 寄左}。乃以天元乘皇極弦，得卜□，為同數，與左相消，得下□□□，開平方，得二百四十步，即城徑也。合問。\end{unitzh}
\begin{uniten}Working: Separately we get: the high number subtracted from the overall altitude is the side (bian) altitude minus the bright (ming) altitude; the level hypotenuse subtracted from the overall base is the side base minus the bright base. Their sum is the great hypotenuse minus the imperial-pole (huangji) hypotenuse, and also the sum of the imperial-pole base and altitude; their difference in bu is the imperial-pole difference. The two stated numbers summed are the yellow-broad (huangguang) hypotenuse; the two stated numbers subtracted from each other, the remainder is the yellow-long (huangchang) hypotenuse. Subtract the sum number from the great sum; the remainder □ is the imperial-pole hypotenuse and the circle diameter together. Set up the celestial unknown as the circle diameter; subtract it from the imperial-pole hypotenuse and circle diameter together, obtaining □□, which is the imperial-pole hypotenuse. Multiply the sum number by itself, obtaining □, placed above; multiply the difference number by itself, obtaining □, and subtract it from the upper position, remainder □; further halve it, obtaining □, which is two pieces of the imperial-pole area [Note: deposit on the left]. Then multiply the imperial-pole hypotenuse by the celestial unknown, obtaining 卜□ [Translator's note: 卜 is evidently part of the counting-rod diagram], as the equal quantity; eliminate it against the left, obtaining below □□□; extract the square root, obtaining 240 bu, which is the diameter of the city. This agrees with the question.\end{uniten}

\begin{unitzh}\unitlabel{ceyuan-027.8}或問：依前見大和，只云大差弦四百八步，小差弦一百七十步。問答同前。\end{unitzh}
\begin{uniten}Someone asks: According to the foregoing, the great sum is known; it only says that the great-difference hypotenuse is 408 bu and the small-difference hypotenuse is 170 bu. The question and answer are the same as before.\end{uniten}

\begin{unitzh}\unitlabel{ceyuan-027.9}法曰：以併云數減大和，復以乘大和，又倍之，為平實；三之通和，於上，又以併云數減大和，加上位，為從；二步虚法；得圓徑。\end{unitzh}
\begin{uniten}Rule: Subtract the sum of the stated numbers from the great sum, then multiply the great sum by it, and further double it, as the square dividend (pingshi). Triple the overall sum, placed above; further subtract the sum of the stated numbers from the great sum and add it to the upper position, as the following (cong). 2 bu is the empty divisor (xufa). Obtain the diameter of the circle.\end{uniten}

\begin{unitzh}\unitlabel{ceyuan-027.10}草曰：大差弦減和步，餘□為大勾大差勾共也；以小差弦減大和，餘□為大股小差股共也。云數相併□，即大弦内減虚弦也；云數相減，得□，為虚弦平弦共也{\zhsmall 按：此二語因數偶合而誤，見前}。以相併數減於大和，餘□為大差勾小差股共，又為圓徑虚弦共也。立天元一為圓徑，減於□，得□□，為虚弦也；返以減於圓徑，得□□，為小和也。以天元減大和，得□□，為大弦，以乘小和，得□□□{\zhsmall 寄左}。乃再置虚弦，以通和乘之，得□□，與左相消，得□□□，開平方，得二百四十步，即城徑也。合問。\end{unitzh}
\begin{uniten}Working: Subtract the great-difference hypotenuse from the sum in bu; the remainder □ is the great base and the great-difference base together. Subtract the small-difference hypotenuse from the great sum; the remainder □ is the great altitude and the small-difference altitude together. The stated numbers summed, □, are the great hypotenuse minus the void hypotenuse; the stated numbers subtracted from each other give □, which is the void hypotenuse and the level hypotenuse together [Siku editors: These two statements are erroneous because the numbers happen to coincide; see above]. Subtract the summed number from the great sum; the remainder □ is the great-difference base and the small-difference altitude together, and also the circle diameter and the void hypotenuse together. Set up the celestial unknown as the circle diameter; subtract it from □, obtaining □□, which is the void hypotenuse; in turn subtract this from the circle diameter, obtaining □□, which is the small sum (xiaohe). Subtract the celestial unknown from the great sum, obtaining □□, which is the great hypotenuse; with it multiply the small sum, obtaining □□□ [Note: deposit on the left]. Then again set out the void hypotenuse and multiply it by the overall sum, obtaining □□; eliminate it against the left, obtaining □□□; extract the square root, obtaining 240 bu, which is the diameter of the city. This agrees with the question.\end{uniten}

\begin{unitzh}\unitlabel{ceyuan-027.11}或問：依前見大和，只云黄廣弦五百一十步，黄長弦二百七十二步。問答同前。\end{unitzh}
\begin{uniten}Someone asks: According to the foregoing, the great sum is known; it only says that the yellow-broad hypotenuse is 510 bu and the yellow-long hypotenuse is 272 bu. The question and answer are the same as before.\end{uniten}

\begin{unitzh}\unitlabel{ceyuan-027.12}法曰：云數相併減大和，復以相併數乘之，為實；云數相併減大和，得數，復以加大和，為法；得虚弦。\end{unitzh}
\begin{uniten}Rule: Sum the stated numbers and subtract from the great sum; then multiply it by the summed number, as the dividend (shi). Sum the stated numbers and subtract from the great sum; again add the resulting number to the great sum, as the divisor (fa). Obtain the void hypotenuse.\end{uniten}

\begin{unitzh}\unitlabel{ceyuan-027.13}草曰：别得黄廣弦又為大差弦虚弦共，又為邊股叀股共也；黄長弦又為小差弦虚弦共，又為底勾明勾共也。以黄廣弦減於大股，餘即虚股；以黄長弦減於大勾，餘即虚勾。故併數以減於大和，餘□為虚和也。以虚和減徑□□，即虚弦也。二云數相併，得□，為大弦虚弦共也；云數相減，餘□，為虚弦平弦共{\zhsmall 按：此句誤同上}。立天元一為虚弦，以減於七百八十二，得□□，為大弦也；以小和乘之，得□□{\zhsmall 寄左}。乃以天元虚弦乘大和，得□呔，為同數，與左相消，得□□，上法下實，得一百二步，即虚弦也。合問。\end{unitzh}
\begin{uniten}Working: Separately we get: the yellow-broad hypotenuse is also the great-difference hypotenuse and the void hypotenuse together, and also the side altitude and the zhuan (叀) altitude together; the yellow-long hypotenuse is also the small-difference hypotenuse and the void hypotenuse together, and also the bottom (di) base and the bright base together. Subtract the yellow-broad hypotenuse from the great altitude: the remainder is the void altitude; subtract the yellow-long hypotenuse from the great base: the remainder is the void base. Therefore subtract the summed number from the great sum; the remainder □ is the void sum. Subtract the void sum from the diameter, □□: this is the void hypotenuse. The two stated numbers summed give □, which is the great hypotenuse and the void hypotenuse together; the stated numbers subtracted from each other, the remainder □ is the void hypotenuse and the level hypotenuse together [Siku editors: The error in this sentence is the same as above]. Set up the celestial unknown as the void hypotenuse; subtract it from 782, obtaining □□, which is the great hypotenuse; multiply it by the small sum, obtaining □□ [Note: deposit on the left]. Then multiply the great sum by the celestial unknown, the void hypotenuse, obtaining □呔 [Translator's note: 呔 is evidently part of the counting-rod diagram], as the equal quantity; eliminate it against the left, obtaining □□; the upper is the divisor, the lower the dividend; obtaining 102 bu, which is the void hypotenuse. This agrees with the question.\end{uniten}

\begin{unitzh}\unitlabel{ceyuan-027.14}或問：依前見大和，只云邊弦五百四十四步，底弦四百二十五。問答同前。\end{unitzh}
\begin{uniten}Someone asks: According to the foregoing, the great sum is known; it only says that the side hypotenuse is 544 bu and the bottom hypotenuse 425. The question and answer are the same as before.\end{uniten}

\begin{unitzh}\unitlabel{ceyuan-027.15}法曰：云數相減，自之，為實；以大和減併數，為法；得皇極弦。\end{unitzh}
\begin{uniten}Rule: Subtract the stated numbers from each other and multiply by itself, as the dividend (shi); subtract the great sum from the summed number, as the divisor (fa); obtain the imperial-pole hypotenuse.\end{uniten}

\begin{unitzh}\unitlabel{ceyuan-027.16}草曰：别得以邊弦減大股，餘為半徑内減平勾，又為平弦内減勾圓差也；以底弦減於大勾，餘為髙股内少半徑，又為股圓差内少髙股也。二云數相併，得九百六十九，為大弦皇極弦共也；二云數相減□，為皇極勾股差也。併數内減通和，餘□為皇極弦内減圓徑也。立天元一為皇極弦，以自之，於上；以一百一十九自之，減上位，得丨□□，為二皇積{\zhsmall 寄左}。復置天元，内減四十九，得下式□□，為黄方；復以天元乘之，得丨□，與左相消，得□□，上法下實，得二百八十九步，即皇極弦也。内減四十九，餘即城徑也。合問。\end{unitzh}
\begin{uniten}Working: Separately we get: subtracting the side hypotenuse from the great altitude, the remainder is the half-diameter (banjing) minus the level base, and also the level hypotenuse minus the base–circle difference; subtracting the bottom hypotenuse from the great base, the remainder is the high altitude less the half-diameter, and also the altitude–circle difference less the high altitude. The two stated numbers summed give 969, which is the great hypotenuse and the imperial-pole hypotenuse together; the two stated numbers subtracted from each other, □, is the base–altitude difference of the imperial pole. From the summed number subtract the overall sum; the remainder □ is the imperial-pole hypotenuse minus the circle diameter. Set up the celestial unknown as the imperial-pole hypotenuse and multiply it by itself, placed above; multiply 119 by itself and subtract it from the upper position, obtaining 丨□□ [Translator's note: 丨, here and below, is evidently part of the counting-rod diagram], which is two imperial areas [Note: deposit on the left]. Again set out the celestial unknown and subtract 49 from it, obtaining the expression below, □□, which is the yellow square (huangfang); again multiply it by the celestial unknown, obtaining 丨□; eliminate it against the left, obtaining □□; the upper is the divisor, the lower the dividend; obtaining 289 bu, which is the imperial-pole hypotenuse. Subtract 49 from it; the remainder is the diameter of the city. This agrees with the question.\end{uniten}

\begin{unitzh}\unitlabel{ceyuan-027.17}按：右大和八問，每問於大和外復設二數，然多有大和外設一數即可求者。細考其法草所載，皆三數並用，婉轉求之，蓋意在發明三數取用之理，非不知其可省也。\end{unitzh}
\begin{uniten}Siku editors' note: Of the eight great-sum problems above, each problem sets two further numbers besides the great sum; yet there are many which could be solved by setting one number besides the great sum. Examining closely what is recorded in their rules and workings, all of them use the three numbers together and solve in a roundabout way; presumably the intention lies in elucidating the principle of taking and using three numbers; it is not that he did not know that this could be dispensed with.\end{uniten}

\begin{unitzh}\unitlabel{ceyuan-027.18}測圓海鏡卷九\end{unitzh}
\begin{uniten}Sea Mirror of Circle Measurements, juan 9\end{uniten}

\section*{卷10}\addcontentsline{toc}{section}{卷10}

\begin{unitzh}\unitlabel{ceyuan-028.1}欽定四庫全書\end{unitzh}
\begin{uniten}Imperially Commissioned Complete Library of the Four Treasuries (Qinding Siku quanshu)\end{uniten}

\begin{unitzh}\unitlabel{ceyuan-028.2}測圓海鏡卷十\end{unitzh}
\begin{uniten}Sea Mirror of Circle Measurements, juan 10\end{uniten}

\begin{unitzh}\unitlabel{ceyuan-028.3}元 李冶 撰\end{unitzh}
\begin{uniten}Composed by Li Ye of the Yuan\end{uniten}

\begin{unitzh}\unitlabel{ceyuan-028.4}三事和八問\end{unitzh}
\begin{uniten}8 questions on the sum of the three things (sanshi he)\end{uniten}

\begin{unitzh}\unitlabel{ceyuan-028.5}或問：甲、乙同立於乾隅，乙向東行，不知步數而立；甲向南直行，多於乙步，望見乙，復就東北斜行，與乙相㑹。二人共行了一千六百步。又云：甲南行不及斜行八十步。問答同前。\end{unitzh}
\begin{uniten}Someone asks: Jia and Yi stand together at the northwest corner (qian yu). Yi walks east an unknown number of bu (步) and stands still. Jia walks straight south, more bu than Yi, catches sight of Yi, and then walks obliquely toward the northeast and meets Yi. The two together walked 1600 bu. It is also said: Jia's walk south falls short of the oblique walk by 80 bu. The question and answer are the same as before.\end{uniten}

\begin{unitzh}\unitlabel{ceyuan-028.6}法曰：共步内減四之小差，復以自之，於上；以十八个小差冪減於上，為實。四之共步，内減十六个小差，於上；却以十八小差加上，為益從。四步常法。開平方，得中差。\end{unitzh}
\begin{uniten}Rule: From the combined bu (gong bu) subtract 4 times the small difference (xiao cha), and then multiply it by itself; [place it] above. Subtract 18 areas (mi) of the small difference from the above, making the dividend (shi). [Take] 4 times the combined bu, from it subtract 16 small differences, [and place it] above; then add 18 small differences to the above, making the subtractive joining term (yi cong). 4 bu as the constant divisor (chang fa). Extract the square root, obtaining the middle difference (zhong cha).\end{uniten}

\begin{unitzh}\unitlabel{ceyuan-028.7}草曰：别得共步為三事和也，不及步即小差也。立天元一為中差，加二之小差，得□□，為大小差併。以加入三事和，得□□，為三弦也。倍三事，得三千二百，内去大小差併，得□□，為三和也。内減三弦，餘□□，為三个黄方。以自之，得□□□，為九段黄方冪{\zhsmall 寄左}。再置天元中差，加小差，得□□，為大差。以小差□乘之，得□□，為半个黄方冪。就一十八之，得□□，為同數。與左相消，得□□□。開平方，得二百八十步，即中差也。其餘各依法求之，合問。\end{unitzh}
\begin{uniten}Working: By distinction (bie) [one] obtains that the combined bu is the sum of the three things, and that the bu by which [it] falls short (buji bu) is the small difference. Set up the celestial unknown (tianyuan) as the middle difference; add twice the small difference, obtaining □□ [Translator's note: each □ stands for an entry of a counting-rod array not reproduced in the transcription; 丨 and 川 are rod numerals.] as the large and small differences combined. Add this into the sum of the three things, obtaining □□ as 3 hypotenuses (xian). Double the three things, obtaining 3200; from it remove the large and small differences combined, obtaining □□ as 3 sums (he). From it subtract the 3 hypotenuses; the remainder, □□, is 3 yellow squares (huang fang). Multiply it by itself, obtaining □□□ as 9 pieces of the yellow-square area (mi) [Note: store on the left (ji zuo)]. Again set up the celestial unknown, the middle difference, and add the small difference, obtaining □□ as the large difference (da cha). Multiply it by the small difference □, obtaining □□ as half a yellow-square area. Then multiply it by 18, obtaining □□ as the equal quantity (tong shu). Eliminate it against the left (xiang xiao), obtaining □□□. Extract the square root, obtaining 280 bu, which is the middle difference. For the rest, find each according to the rule. This fits the question.\end{uniten}

\begin{unitzh}\unitlabel{ceyuan-028.8}或問：以前三事和，又云：大差三百六十步。問答同前。\end{unitzh}
\begin{uniten}Someone asks: [Take] the preceding sum of the three things; it is also said: the large difference is 360 bu. The question and answer are the same as before.\end{uniten}

\begin{unitzh}\unitlabel{ceyuan-028.9}法曰：倍云數，以云數乘之，又九之，於上。倍云數，加三事和，為前數；倍云數，減二之三事和，為後數。二數又相減，餘一百六十，為泛率。以自乘，減上位，為平實。十八之云數，内又加四之泛率，為從。四常法。得中差。\end{unitzh}
\begin{uniten}Rule: Double the stated number (yun shu), multiply it by the stated number, and again multiply by 9; [place it] above. Double the stated number and add the sum of the three things, making the former number (qian shu); double the stated number and subtract it from twice the sum of the three things, making the latter number (hou shu). The two numbers are again subtracted from each other; the remainder, 160, makes the provisional rate (fan lü). Multiply it by itself and subtract from the upper position, making the square dividend (ping shi). 18 times the stated number, to which again add 4 times the provisional rate, makes the joining term (cong). 4 as the constant divisor. Obtain the middle difference.\end{uniten}

\begin{unitzh}\unitlabel{ceyuan-028.10}草曰：立天元一為中差。置云步，倍之，内減天元，得□□，為大小差共數。加於三事和，得□□，為三弦也。倍三事，内減大小差共數，得下式□□，為三和也。内減三弦，得□□，為三个黄方靣也。以自之，得□{\zhsmall □□}□，為九段黄方冪{\zhsmall 寄左}。再以天元減大差，得下式□□，為小差。又倍之，得□□。以云數乘之，得下式□□。又就分九之，得下式□□。與左相消，得下式□□□。開平方，得二百八十步，即中差也。合問。\end{unitzh}
\begin{uniten}Working: Set up the celestial unknown as the middle difference. Set up the stated bu, double it, and from it subtract the celestial unknown, obtaining □□ as the combined number of the large and small differences. Add it to the sum of the three things, obtaining □□ as 3 hypotenuses. Double the three things and from it subtract the combined number of the large and small differences, obtaining the expression below (xia shi) □□ as 3 sums. From it subtract the 3 hypotenuses, obtaining □□ as 3 sides of the yellow square. Multiply it by itself, obtaining □[□□]□ as 9 pieces of the yellow-square area [Note: store on the left]. Again subtract the celestial unknown from the large difference, obtaining the expression below □□ as the small difference. Again double it, obtaining □□. Multiply it by the stated number, obtaining the expression below □□. Again, by parts (jiu fen), multiply it by 9, obtaining the expression below □□. Eliminate it against the left, obtaining the expression below □□□. Extract the square root, obtaining 280 bu, which is the middle difference. This fits the question.\end{uniten}

\begin{unitzh}\unitlabel{ceyuan-028.11}或問：依前見三事和，又云：中差二百八十步。問答同前。\end{unitzh}
\begin{uniten}Someone asks: As before, the sum of the three things is given (jian); it is also said: the middle difference is 280 bu. The question and answer are the same as before.\end{uniten}

\begin{unitzh}\unitlabel{ceyuan-028.12}法曰：和步加差步，以自乘，於上；又和步内減差步，以自乘，加上位，為平實。四之和步，為從。二步益隅。得大弦。\end{unitzh}
\begin{uniten}Rule: Add the bu of the difference to the bu of the sum and multiply it by itself; [place it] above. Again, from the bu of the sum subtract the bu of the difference, multiply it by itself, and add it to the upper position, making the square dividend. 4 times the bu of the sum makes the joining term. 2 bu as the subtractive corner (yi yu). Obtain the large hypotenuse (da xian).\end{uniten}

\begin{unitzh}\unitlabel{ceyuan-028.13}草曰：立天元一為大弦，減共步，得□□，為和。副置之：上位減差步，得□□，為二勾，以自之，得丨□□，為四段勾冪也；下位加差步，得□□，為二股，以自之，得丨□□，為四段股冪也。二位相併，得□□□，為四段弦冪{\zhsmall 寄左}。然後以天元自之，又四之，得□□，為同數。與左相消，得□□□。開平方，得六百八十步，即大弦也。倍之，以減於三事和，餘即城徑也。合問。\end{unitzh}
\begin{uniten}Working: Set up the celestial unknown as the large hypotenuse; subtract it from the combined bu, obtaining □□ as the sum. Set it up in two places (fu zhi): from the upper position subtract the bu of the difference, obtaining □□ as 2 bases (gou); multiply it by itself, obtaining 丨□□ as 4 pieces of the base area (gou mi). To the lower position add the bu of the difference, obtaining □□ as 2 altitudes (gu); multiply it by itself, obtaining 丨□□ as 4 pieces of the altitude area (gu mi). Combine the two positions, obtaining □□□ as 4 pieces of the hypotenuse area (xian mi) [Note: store on the left]. Afterwards multiply the celestial unknown by itself and again by 4, obtaining □□ as the equal quantity. Eliminate it against the left, obtaining □□□. Extract the square root, obtaining 680 bu, which is the large hypotenuse. Double it and subtract it from the sum of the three things; the remainder is the diameter of the city (cheng jing). This fits the question.\end{uniten}

\begin{unitzh}\unitlabel{ceyuan-028.14}或問：依前見三事和，又云：小差、大差併四百四十步。問答同前。\end{unitzh}
\begin{uniten}Someone asks: As before, the sum of the three things is given; it is also said: the small difference and the large difference combined are 440 bu. The question and answer are the same as before.\end{uniten}

\begin{unitzh}\unitlabel{ceyuan-028.15}法曰：併前後二數，三而一，為弦；反以減共步，得數；又以減弦，得城徑。\end{unitzh}
\begin{uniten}Rule: Combine the former and latter two numbers and divide by 3, making the hypotenuse; in turn subtract it from the combined bu, obtaining a number; again subtract it from the hypotenuse, obtaining the diameter of the city.\end{uniten}

\begin{unitzh}\unitlabel{ceyuan-028.16}草曰：二數相併，得□，三而一，得□，即弦也。以弦減三事和，得□，即和也。弦、和又相減，餘二百四十步，即城徑也。合問。\end{unitzh}
\begin{uniten}Working: Combine the two numbers, obtaining □; divide by 3, obtaining □, which is the hypotenuse. Subtract the hypotenuse from the sum of the three things, obtaining □, which is the sum. The hypotenuse and the sum are again subtracted from each other; the remainder, 240 bu, is the diameter of the city. This fits the question.\end{uniten}

\begin{unitzh}\unitlabel{ceyuan-028.17}或問：依前見三事和，又云：小差、中差、大差共七百二十。問答同前。\end{unitzh}
\begin{uniten}Someone asks: As before, the sum of the three things is given; it is also said: the small difference, the middle difference and the large difference together are 720. The question and answer are the same as before.\end{uniten}

\begin{unitzh}\unitlabel{ceyuan-028.18}法曰：半云數，自之，又三之，於上；以三事減上位，為平實{\zhsmall 按：「以三事減上位」有誤，此係偶合三事之數耳。當云：「加半段三事冪；又倍三事和，加大差，復以大差乘之，減上位，為平實。」}。倍三事，於上；半云數而五之，加上位，為益從。二常法。得小差。\end{unitzh}
\begin{uniten}Rule: Halve the stated number, multiply it by itself, and again multiply by 3; [place it] above. Subtract the three things from the upper position, making the square dividend [Siku editors: "Subtract the three things from the upper position" is in error; this merely happens to coincide with the number of the three things. It should say: "Add half a piece of the area (mi) of the three things; also double the sum of the three things, add the large difference, then multiply by the large difference, and subtract from the upper position, making the square dividend."]. Double the three things, [place it] above; halve the stated number and multiply by 5, and add to the upper position, making the subtractive joining term. 2 as the constant divisor. Obtain the small difference.\end{uniten}

\begin{unitzh}\unitlabel{ceyuan-028.19}草曰：别得三差共為二大差也。立天元一為小差，併大差，加入三事和，得□□，為三弦也。以自之，得丨□□，為十八積九較冪{\zhsmall 寄起}。又以共三事步自之，得□方，於上；又以天元小差乘大差，倍之，得□，加於上，得□□，為十二積四較冪。又加五{\zhsmall 按：即三因二歸}，得□□，為十八个直積六个較冪。以減寄起，餘得丨□□，為三个較冪{\zhsmall 寄左}。然後以天元小差減大差，得□□，為中差。以自之，得丨□□，又三之，得下式川□□，為同數。與寄左相消，得□□□。平方而一，得八十步，即小差也。餘各依數求之，合問。\end{unitzh}
\begin{uniten}Working: By distinction [one] obtains that the three differences together make 2 large differences. Set up the celestial unknown as the small difference; combine it with the large difference and add it into the sum of the three things, obtaining □□ as 3 hypotenuses. Multiply it by itself, obtaining 丨□□ as 18 areas (ji) and 9 areas of the difference (jiao mi) [Note: store it (ji qi)]. Again multiply the combined bu of the three things by itself, obtaining □ square (fang), [placing it] above; again multiply the large difference by the celestial unknown, the small difference, and double it, obtaining □; add it to the above, obtaining □□ as 12 areas and 4 areas of the difference. Again "add five" (jia wu) [Siku editors: that is, multiply by 3 and divide by 2 (san yin er gui)], obtaining □□ as 18 rectangle areas (zhi ji) and 6 areas of the difference. Subtract it from what was stored (ji qi); the remainder obtained, 丨□□, is 3 areas of the difference [Note: store on the left]. Afterwards subtract the celestial unknown, the small difference, from the large difference, obtaining □□ as the middle difference. Multiply it by itself, obtaining 丨□□; again multiply by 3, obtaining the expression below 川□□ as the equal quantity. Eliminate it against what is stored on the left, obtaining □□□. Dividing by square-root extraction (pingfang er yi), obtain 80 bu, which is the small difference. For the rest, find each according to the numbers. This fits the question.\end{uniten}

\begin{unitzh}\unitlabel{ceyuan-028.20}或問：依前見三事和，又云：明黄方、叀黄方共六十六。問答同前。\end{unitzh}
\begin{uniten}Someone asks: As before, the sum of the three things is given; it is also said: the ming yellow square and the zhuan (叀) yellow square together are 66. The question and answer are the same as before.\end{uniten}

\begin{unitzh}\unitlabel{ceyuan-028.21}法曰：二事内加二之共步，復以二之共步乘之，於上位；三事内減二之共步，復以二之共步乘之，得數，減上位，為平實。三事内加二之共步，又倍之，於上；又三{\zhsmall 按：三當作六}之共步，加上位，為泛寄。三事内減二之共步，又四之，於上；又三{\zhsmall 按：三亦當作六}之共步，減上位，得數，以減泛寄，為從。作十八段虚平方開之，得虚黄方。\end{unitzh}
\begin{uniten}Rule: To the two things [Translator's note: 二事 "two things" is presumably an error for 三事 "three things".] add twice the combined bu, then multiply it by twice the combined bu; [place it] in the upper position. From the three things subtract twice the combined bu, then multiply it by twice the combined bu; subtract the number obtained from the upper position, making the square dividend. To the three things add twice the combined bu, and again double it; [place it] above; again [take] 3 [Siku editors: "3" should be "6"] times the combined bu and add it to the upper position, making the provisional stored quantity (fan ji). From the three things subtract twice the combined bu, and again multiply it by 4; [place it] above; again [take] 3 [Siku editors: "3" also should be "6"] times the combined bu and subtract it from the upper position; subtract the number obtained from the provisional stored quantity, making the joining term. Taking 18 pieces as the empty square (xu pingfang), extract it, obtaining the xu yellow square.\end{uniten}

\begin{unitzh}\unitlabel{ceyuan-028.22}草曰：别得共步即虚大小差也。立天元一為虚黄方，以三之，加入倍之共步，得□□，為圓徑也。以圓徑加三事，得□□，為二通和；以圓徑減三事，得□□，為二通弦。又副置圓徑：上加天元，得□□，為二虚和；下減天元，得□□，為二虚弦。乃置二大和，以二小弦乘之，得下□□□{\zhsmall 寄左}。然後置二大弦，以二小和乘之，得下式□□□。與左相消，得□□□。平方開之，得三十六步，即虚黄方也。其餘各依法求之，合問。\end{unitzh}
\begin{uniten}Working: By distinction [one] obtains that the combined bu is the xu large and small differences. Set up the celestial unknown as the xu yellow square; multiply it by 3 and add it into twice the combined bu, obtaining □□ as the diameter of the circle (yuan jing). Add the diameter of the circle to the three things, obtaining □□ as 2 through sums (tong he); subtract the diameter of the circle from the three things, obtaining □□ as 2 through hypotenuses (tong xian). Again set up the diameter of the circle in two places: to the upper add the celestial unknown, obtaining □□ as 2 xu sums; from the lower subtract the celestial unknown, obtaining □□ as 2 xu hypotenuses. Then set up the 2 large sums and multiply them by the 2 small hypotenuses, obtaining the following □□□ [Note: store on the left]. Afterwards set up the 2 large hypotenuses and multiply them by the 2 small sums, obtaining the expression below □□□. Eliminate it against the left, obtaining □□□. Extract its square root, obtaining 36 bu, which is the xu yellow square. For the rest, find each according to the rule. This fits the question.\end{uniten}

\begin{unitzh}\unitlabel{ceyuan-028.23}或問：依前見三事和，又云：皇極弦二百八十九步。問答同前。\end{unitzh}
\begin{uniten}Someone asks: As before, the sum of the three things is given; it is also said: the huangji hypotenuse is 289 bu. The question and answer are the same as before.\end{uniten}

\begin{unitzh}\unitlabel{ceyuan-028.24}法曰：二數相乘，為實。從空。一益隅。得大弦。\end{unitzh}
\begin{uniten}Rule: Multiply the two numbers together, making the dividend (shi). The joining term is empty (cong kong). 1 as the subtractive corner. Obtain the large hypotenuse.\end{uniten}

\begin{unitzh}\unitlabel{ceyuan-028.25}草曰：立天元一為通弦，内減皇弦，餘□□，為皇極勾股和。以自之，得丨□□，於上；以皇極弦冪減上位，得丨□，為二直積。合於皇極除之，不除，寄為母，便以此為城徑{\zhsmall 寄左}。乃以二之天元弦減共步，得□□，為黄方面。以皇弦通之，得□□。與左相消，得丨□□。開平方，得六百八十步，即大弦也。合問。\end{unitzh}
\begin{uniten}Working: Set up the celestial unknown as the through hypotenuse (tong xian); from it subtract the huang hypotenuse; the remainder, □□, is the sum of the huangji base and altitude. Multiply it by itself, obtaining 丨□□, [placing it] above; subtract the area (mi) of the huangji hypotenuse from the upper position, obtaining 丨□ as 2 rectangle areas. This ought to be divided by the huangji; not dividing, store it as the denominator (mu), and then take this as the diameter of the city [Note: store on the left]. Then subtract twice the celestial-unknown hypotenuse from the combined bu, obtaining □□ as the side of the yellow square. Multiply it through (tong) by the huang hypotenuse, obtaining □□. Eliminate it against the left, obtaining 丨□□. Extract the square root, obtaining 680 bu, which is the large hypotenuse. This fits the question.\end{uniten}

\begin{unitzh}\unitlabel{ceyuan-028.26}或問：依前見三事和，又云：見太虚弦一百二步。問答同前。\end{unitzh}
\begin{uniten}Someone asks: As before, the sum of the three things is given; it is also said: the given (jian) taixu hypotenuse is 102 bu. The question and answer are the same as before.\end{uniten}

\begin{unitzh}\unitlabel{ceyuan-028.27}法曰：半虚弦，乘三事，為實。三事為從。四虚隅。翻開之，得半大弦。\end{unitzh}
\begin{uniten}Rule: Halve the xu hypotenuse and multiply it by the three things, making the dividend (shi). The three things make the joining term. 4 as the empty corner (xu yu). Extract it by reversal (fan kai), obtaining half the large hypotenuse.\end{uniten}

\begin{unitzh}\unitlabel{ceyuan-028.28}草曰：識别得：以虚弦減大弦，半之，為皇極弦；以虚弦加大弦，半之，為皇極勾股共也。立天元一為半大弦，以二之，内減虚弦，得□□；折半，得□□，為皇極弦也。又以虚弦加大弦而半之，得□□，為皇極和也。和自之，得丨□□，於上；又以弦自之，得丨□□，減上位，餘得下□，為二直積。合以皇極弦除之，不除，寄為分母，便以此為城徑{\zhsmall 寄左}。然後以四之天元減三事共，餘□□；又以皇極弦分母通之，得□□□，為同數。與左相消，得□□□。倒積開，得三百四十步，倍之，即大弦也。合問。\end{unitzh}
\begin{uniten}Working: By recognition and distinction (shibie) [one] obtains: subtract the xu hypotenuse from the large hypotenuse and halve it, making the huangji hypotenuse; add the xu hypotenuse to the large hypotenuse and halve it, making the huangji base and altitude combined. Set up the celestial unknown as half the large hypotenuse; double it, and from it subtract the xu hypotenuse, obtaining □□; halve it, obtaining □□ as the huangji hypotenuse. Again add the xu hypotenuse to the large hypotenuse and halve it, obtaining □□ as the huangji sum. Multiply the sum by itself, obtaining 丨□□, [placing it] above; again multiply the hypotenuse by itself, obtaining 丨□□, and subtract it from the upper position; the remainder obtained, the following □, is 2 rectangle areas. This ought to be divided by the huangji hypotenuse; not dividing, store it as the denominator (fen mu), and then take this as the diameter of the city [Note: store on the left]. Afterwards subtract 4 times the celestial unknown from the three things combined; the remainder is □□; again multiply it through (tong) by the huangji-hypotenuse denominator, obtaining □□□ as the equal quantity. Eliminate it against the left, obtaining □□□. Extract with inverted area (dao ji kai), obtaining 340 bu; double it, and that is the large hypotenuse. This fits the question.\end{uniten}

\begin{unitzh}\unitlabel{ceyuan-028.29}測圓海鏡卷十\end{unitzh}
\begin{uniten}Sea Mirror of Circle Measurements, juan 10\end{uniten}

\section*{卷11}\addcontentsline{toc}{section}{卷11}

\begin{unitzh}\unitlabel{ceyuan-029.1}欽定四庫全書\end{unitzh}
\begin{uniten}Imperially Commissioned Complete Library of the Four Treasuries (Qinding Siku quanshu)\end{uniten}

\begin{unitzh}\unitlabel{ceyuan-029.2}測圓海鏡卷十一\end{unitzh}
\begin{uniten}Sea Mirror of Circle Measurements (Ceyuan haijing), juan 11\end{uniten}

\begin{unitzh}\unitlabel{ceyuan-029.3}元 李冶 撰\end{unitzh}
\begin{uniten}Composed by Li Ye of the Yuan\end{uniten}

\begin{unitzh}\unitlabel{ceyuan-029.4}雜糅一十八問\end{unitzh}
\begin{uniten}Miscellaneous Mixture (zarou): 18 problems\end{uniten}

\begin{unitzh}\unitlabel{ceyuan-029.5}或問：城南有槐樹一株，城東有柳樹一株。甲出北門東行，丙出西門南行，甲、丙、槐、柳悉與城叅相直。既而丙就柳行五百四十四步，至柳樹下；甲就槐行四百二十五步，至槐樹下。問答同前。\end{unitzh}
\begin{uniten}Someone asks: South of the city there is one locust tree; east of the city there is one willow tree. Jia goes out of the north gate and walks east; Bing goes out of the west gate and walks south; Jia, Bing, the locust and the willow are all mutually in line with the city. Thereupon Bing, approaching the willow, walks 544 bu (步) and arrives beneath the willow tree; Jia, approaching the locust, walks 425 bu and arrives beneath the locust tree. The question and answer are the same as before.\end{uniten}

\begin{unitzh}\unitlabel{ceyuan-029.6}法曰：甲就步自之於上；以二行相減數自之，減上位，為實；二之二行相減數，併入二之甲就步，為從；一步常法。得平弦。\end{unitzh}
\begin{uniten}Rule: Multiply the bu of Jia's approach by itself, in the upper position; multiply the difference of the two walks by itself and subtract it from the upper position, making the dividend (shi); double the difference of the two walks and add it into twice the bu of Jia's approach, making the "following" (cong); 1 bu as the constant divisor (changfa). One obtains the level hypotenuse (pingxian).\end{uniten}

\begin{unitzh}\unitlabel{ceyuan-029.7}草曰：别得丙就步為邊弦也，甲就步為底弦也。邊弦即皇弦、髙弦共也；底弦即皇弦、平弦共也。二行相併，即大弦、皇弦共也；二行相減，即皇極勾股較也。倍皇弦以減於大弦，餘即虚弦也；倍皇弦内減邊弦，餘即叀弦也；倍皇弦内減底弦，餘即明弦也。皇極弦加一差{\zhsmall 按：一差即皇極勾股較。}，則大差弦也；内減一差，則小差弦也。立天元一為平弦，加一皇極勾股差，得□□，即髙弦也。髙弦自之，得丨□□；内加天元冪，得□□□，為皇弦冪{\zhsmall 寄左}。然後以天元減底弦，得下式□□；自之，得丨□□，為同數，與左相消，得丨□□。開平方，得一百三十六步，即平弦也。餘各依法求之。合問。\end{unitzh}
\begin{uniten}Working: Separately one finds that the bu of Bing's approach is the side hypotenuse (bianxian), and the bu of Jia's approach is the bottom hypotenuse (dixian). The side hypotenuse is the august hypotenuse (huangxian) and the high hypotenuse (gaoxian) together; the bottom hypotenuse is the august hypotenuse and the level hypotenuse together. The sum of the two walks is the great hypotenuse (daxian) and the august hypotenuse together; the difference of the two walks is the base–altitude difference (gougu jiao) of the august pole (huangji). Double the august hypotenuse and subtract it from the great hypotenuse: the remainder is the void hypotenuse (xuxian). From double the august hypotenuse subtract the side hypotenuse: the remainder is the exclusive hypotenuse (zhuanxian). From double the august hypotenuse subtract the bottom hypotenuse: the remainder is the bright hypotenuse (mingxian). The august-pole hypotenuse plus one difference (cha) [Note: Observe that the one difference is the base–altitude difference of the august pole.] is then the great-difference hypotenuse (dacha xian); with one difference subtracted from it, it is then the small-difference hypotenuse (xiaocha xian). Set up the celestial unknown (tianyuan) as the level hypotenuse; add one base–altitude difference (gougu cha) of the august pole, obtaining □□ [Translator's note: □ marks counting-rod numbers and arrays not reproduced in the transcription; 丨 and ○ are the rod numerals 1 and 0.], which is the high hypotenuse. Multiply the high hypotenuse by itself, obtaining 丨□□; add into it the area (mi) of the celestial unknown, obtaining □□□, as the area of the august hypotenuse [Note: lodged on the left (jizuo)]. Then subtract the celestial unknown from the bottom hypotenuse, obtaining the form below (xiashi) □□; multiply it by itself, obtaining 丨□□, as the equal quantity (tongshu); cancel it against the left (xiangxiao), obtaining 丨□□. Extract the square root, obtaining 136 bu, which is the level hypotenuse. The rest are each found according to the rules. This agrees with the question.\end{uniten}

\begin{unitzh}\unitlabel{ceyuan-029.8}或問：出南門東行有槐樹一株。甲出北門東行，斜望槐樹，與城相直，就槐樹行二百七十二步。出東門南行有柳樹一株。丙出西門南行，斜望柳樹，與城相直，就柳樹行五百一十步。問答同前。\end{unitzh}
\begin{uniten}Someone asks: Going out of the south gate and walking east, there is one locust tree. Jia goes out of the north gate and walks east, looks obliquely at the locust tree in line with the city, and, approaching the locust tree, walks 272 bu. Going out of the east gate and walking south, there is one willow tree. Bing goes out of the west gate and walks south, looks obliquely at the willow tree in line with the city, and, approaching the willow tree, walks 510 bu. The question and answer are the same as before.\end{uniten}

\begin{unitzh}\unitlabel{ceyuan-029.9}法曰：云數相併而半之，以自乘於上；半丙斜行以為冪，半甲斜行以為冪，併二冪，減上位，為實；併云數為益從；一步平隅。得虚弦。\end{unitzh}
\begin{uniten}Rule: Add the stated numbers (yunshu) together and halve; multiply it by itself, in the upper position. Halve Bing's oblique walk and make it an area; halve Jia's oblique walk and make it an area; add the two areas and subtract from the upper position, making the dividend (shi). Add the stated numbers to make the augmenting "following" (yicong); 1 bu as the level corner (pingyu). One obtains the void hypotenuse.\end{uniten}

\begin{unitzh}\unitlabel{ceyuan-029.10}草曰：别得丙斜行為黄廣弦也，亦為兩个髙弦也，此勾則城徑也；甲斜行即黄長弦也，亦為兩个平弦也，此股則城徑也。二數相併，得□，即大弦、虚弦共也；二數相減，餘□，即兩个皇極差也；二數相併而半之，得□，即皇極和也。立天元一為虚弦，以減于皇極和，得□□，即皇極弦也。以自之，得丨□□，為皇弦冪{\zhsmall 寄左}。然後以髙弦自之，得□；以平弦自之，得□；二自乘數相併，得□，與左相消，得□□□。開平方，得一百二，即虚弦也。合問。\end{unitzh}
\begin{uniten}Working: Separately one finds that Bing's oblique walk is the yellow-breadth hypotenuse (huangguang xian), and is also two high hypotenuses; its base (gou) is then the city diameter (chengjing). Jia's oblique walk is the yellow-length hypotenuse (huangchang xian), and is also two level hypotenuses; its altitude (gu) is then the city diameter. Adding the two numbers gives □, which is the great hypotenuse and the void hypotenuse together; subtracting the two numbers leaves □, which is two august-pole differences (huangji cha); adding the two numbers and halving gives □, which is the august-pole sum (huangji he). Set up the celestial unknown as the void hypotenuse and subtract it from the august-pole sum, obtaining □□, which is the august-pole hypotenuse. Multiply it by itself, obtaining 丨□□, as the area of the august hypotenuse [Note: lodged on the left]. Then multiply the high hypotenuse by itself, obtaining □; multiply the level hypotenuse by itself, obtaining □; add the two self-products together, obtaining □, and cancel against the left, obtaining □□□. Extract the square root, obtaining 102, which is the void hypotenuse. This agrees with the question.\end{uniten}

\begin{unitzh}\unitlabel{ceyuan-029.11}或問：甲從坤隅南行，不知步數而立；乙從艮隅南行一百五十步，望見甲，復斜行五百一十步，與甲相㑹。問答同前。\end{unitzh}
\begin{uniten}Someone asks: Jia walks south from the kun corner, the number of bu not being known, and stands; Yi walks south from the gen corner 150 bu, catches sight of Jia, and then walks obliquely 510 bu and meets Jia. The question and answer are the same as before.\end{uniten}

\begin{unitzh}\unitlabel{ceyuan-029.12}法曰：斜行自之於上；倍南行，減斜，餘自之，以減上，為實；倍南行，減斜，又四之，為從；八步常法。平方，得半徑。\end{unitzh}
\begin{uniten}Rule: Multiply the oblique walk by itself, in the upper position. Double the southward walk and subtract the oblique walk; multiply the remainder by itself and subtract it from the upper, making the dividend (shi). Double the southward walk, subtract the oblique walk, and further multiply by 4, making the "following"; 8 bu as the constant divisor. Extracting the square root, one obtains the semi-diameter (banjing).\end{uniten}

\begin{unitzh}\unitlabel{ceyuan-029.13}草曰：别得南行即小差股，斜行即黄廣弦也。小差股内減半徑，餘即半个黄廣積上股弦差也；全徑即其勾也。立天元一為半城徑，減于乙南行，倍之，得□□，即一个黄廣即上股弦差也。以減于斜行步，餘□□，即股也。自之，得□□□，為股冪也。又倍天元以自之，得□□，為大勾冪；加入大股冪，得□□□{\zhsmall 寄左}。然後以斜行冪□與寄左相消，得下式□□□。開平方，得一百二十步，即半徑也。合問。\end{unitzh}
\begin{uniten}Working: Separately one finds that the southward walk is the small-difference altitude (xiaocha gu), and the oblique walk is the yellow-breadth hypotenuse. From the small-difference altitude subtract the semi-diameter: the remainder is half of the altitude–hypotenuse difference (gu xian cha) on the yellow-breadth area; the whole diameter is its base. Set up the celestial unknown as half the city diameter; subtract it from Yi's southward walk and double, obtaining □□, which is one yellow-breadth, that is, the upper altitude–hypotenuse difference. Subtract it from the bu of the oblique walk; the remainder, □□, is the altitude. Multiply it by itself, obtaining □□□, as the area of the altitude. Further, double the celestial unknown and multiply it by itself, obtaining □□, as the area of the great base (dagou); add in the area of the great altitude (dagu), obtaining □□□ [Note: lodged on the left]. Then cancel the area of the oblique walk, □, against what is lodged on the left, obtaining the form below □□□. Extract the square root, obtaining 120 bu, which is the semi-diameter. This agrees with the question.\end{uniten}

\begin{unitzh}\unitlabel{ceyuan-029.14}或問：乙從艮隅東行，不知逺近而止；甲從坤隅東行一百九十二步，望見乙，復斜行二百七十二步，與乙相㑹。問答同前。\end{unitzh}
\begin{uniten}Someone asks: Yi walks east from the gen corner, how far not being known, and stops; Jia walks east from the kun corner 192 bu, catches sight of Yi, and then walks obliquely 272 bu and meets Yi. The question and answer are the same as before.\end{uniten}

\begin{unitzh}\unitlabel{ceyuan-029.15}法曰：倍東行，減斜行，得數自為冪，以減于斜行冪，為平實；倍東行，減斜行，又四之，為從；八益隅。翻法開平方，得半徑。\end{unitzh}
\begin{uniten}Rule: Double the eastward walk and subtract the oblique walk; make the resulting number into an area by itself, and subtract it from the area of the oblique walk, making the level dividend (pingshi). Double the eastward walk, subtract the oblique walk, and further multiply by 4, making the "following"; 8 as the augmenting corner (yiyu). Extract the square root by the reversed method (fanfa), obtaining the semi-diameter.\end{uniten}

\begin{unitzh}\unitlabel{ceyuan-029.16}草曰：别得甲東行即大差勾也，斜行則黄長弦也。大差勾内減半徑，餘即半个黄長積上勾弦差也；全徑即其股也。立天元一為半徑，減於東行，倍之，得□□，即一个黄長積上勾弦差也。以減于斜行步，得□□，即黄長勾也。以自之，得□□□，為勾冪于上；倍天元以自之，得□□，加上位，得下式□□□，為弦冪{\zhsmall 寄左}。然後以斜行冪□為同數，與左相消，得□□□。平開，得一百二十步，即半城徑也。合問。\end{unitzh}
\begin{uniten}Working: Separately one finds that Jia's eastward walk is the great-difference base (dacha gou), and the oblique walk is the yellow-length hypotenuse. From the great-difference base subtract the semi-diameter: the remainder is half of the base–hypotenuse difference (gou xian cha) on the yellow-length area; the whole diameter is its altitude. Set up the celestial unknown as the semi-diameter; subtract it from the eastward walk and double, obtaining □□, which is one base–hypotenuse difference on the yellow-length area. Subtract it from the bu of the oblique walk, obtaining □□, which is the yellow-length base. Multiply it by itself, obtaining □□□, as the area of the base, in the upper position; double the celestial unknown and multiply it by itself, obtaining □□; add it to the upper position, obtaining the form below □□□, as the area of the hypotenuse [Note: lodged on the left]. Then take the area of the oblique walk, □, as the equal quantity, and cancel against the left, obtaining □□□. Extract the square root, obtaining 120 bu, which is half the city diameter. This agrees with the question.\end{uniten}

\begin{unitzh}\unitlabel{ceyuan-029.17}或問：甲從坤東行一百九十二步，丙從艮南行一百五十步，望見之。問答同前。\end{unitzh}
\begin{uniten}Someone asks: Jia walks east from kun 192 bu; Bing walks south from gen 150 bu and catches sight of him. The question and answer are the same as before.\end{uniten}

\begin{unitzh}\unitlabel{ceyuan-029.18}法曰：二行相乘，倍之，為平實。如法，得圓徑。\end{unitzh}
\begin{uniten}Rule: Multiply the two walks together and double, making the level dividend (pingshi). Proceeding according to the rule, one obtains the circle's diameter (yuanjing).\end{uniten}

\begin{unitzh}\unitlabel{ceyuan-029.19}草曰：别得甲行即大差勾，丙行即小差股。此二數相乘，恰與大小差相乘正同。如法相乘訖，倍之，得□，為圓徑冪{\zhsmall 寄左}。然然立天元為圓徑，以自之，與左相消，得丨□□。開平方，得二百四十步，即城徑也。合問。\end{unitzh}
\begin{uniten}Working: Separately one finds that Jia's walk is the great-difference base and Bing's walk is the small-difference altitude. The product of these two numbers is exactly the same as the product of the great and small differences. Having multiplied according to the rule, double, obtaining □, as the area of the circle's diameter [Note: lodged on the left]. Then [Translator's note: the text reads 然然, evidently for 然後 "then".] set up the celestial unknown as the circle's diameter, multiply it by itself, and cancel against the left, obtaining 丨□□. Extract the square root, obtaining 240 bu, which is the city diameter. This agrees with the question.\end{uniten}

\begin{unitzh}\unitlabel{ceyuan-029.20}又法：以二行相減數減于二行相併數，餘者半之于上；復以二行相減數加于上，即城徑。\end{unitzh}
\begin{uniten}Another rule: Subtract the difference of the two walks from the sum of the two walks; halve the remainder, in the upper position; again add the difference of the two walks to the upper: this is the city diameter.\end{uniten}

\begin{unitzh}\unitlabel{ceyuan-029.21}草曰：别得甲東行減于徑，為虚勾也；丙南行減于徑，為虚股也。二行共為一徑、一虚弦共也。二行相減，即虚和也。以相併數、相減數又相減，即兩个虚弦也。如法求得虚和□、虚弦□，相併，得□，即城徑也。合問。\end{unitzh}
\begin{uniten}Working: Separately one finds that Jia's eastward walk subtracted from the diameter makes the void base (xugou), and Bing's southward walk subtracted from the diameter makes the void altitude (xugu). The two walks together make one diameter and one void hypotenuse together. The difference of the two walks is the void sum (xuhe). Subtract the sum and the difference from each other again: this is two void hypotenuses. Having found by the rule the void sum, □, and the void hypotenuse, □, add them, obtaining □, which is the city diameter. This agrees with the question.\end{uniten}

\begin{unitzh}\unitlabel{ceyuan-029.22}按：又法未合。蓋以二行相減為虚較，而草中誤以為虚和也。其義甚淺，非難知者，是殆偶爾之遺忘；然亦可以決其為當日未定之稿矣。\end{unitzh}
\begin{uniten}Note: The other rule does not agree. For the difference of the two walks makes the void difference (xujiao), but in the working it is mistakenly taken as the void sum. Its sense is very shallow and not something hard to know; this is probably an occasional lapse of memory; yet from it one can also determine that this was a draft not yet settled at that time.\end{uniten}

\begin{unitzh}\unitlabel{ceyuan-029.23}或問：出西門南行二百二十五步有塔；出北門東行六十四步，望塔，正當城徑之半。問答同前。\end{unitzh}
\begin{uniten}Someone asks: Going out of the west gate and walking south 225 bu, there is a pagoda; going out of the north gate and walking east 64 bu, one looks at the pagoda exactly at half of the city's diameter. The question and answer are the same as before.\end{uniten}

\begin{unitzh}\unitlabel{ceyuan-029.24}法曰：二行相乘，為平實；一步常法。得半徑。\end{unitzh}
\begin{uniten}Rule: Multiply the two walks together, making the level dividend; 1 bu as the constant divisor. One obtains the semi-diameter.\end{uniten}

\begin{unitzh}\unitlabel{ceyuan-029.25}草曰：别得二百二十五步為髙股，此乃半徑為勾之股也；其六十四步為平勾，此乃半徑為股之勾也。二數相併，即太極弦也；二數相減，即中差内去皇極差也。又别得二行相乘，恰是半徑冪一段，此與半梯頭相乘，其意正同。今且以弦上容圓取之：立天元一為半徑，副之；上加南行，得□□，為股也；下加東行步，得□□，為勾也。勾股相乘，得丨□□，為大直積；以天元半徑除之，得□□□，為勾股和{\zhsmall 寄左}。然後併勾股，得□□，與左相消，得丨○□。開平方，得一百二十步，即半徑也。合問。\end{unitzh}
\begin{uniten}Working: Separately one finds that the 225 bu is the high altitude (gaogu): this is the altitude for which the semi-diameter is the base; the 64 bu is the level base (pinggou): this is the base for which the semi-diameter is the altitude. The sum of the two numbers is the great-pole hypotenuse (taijixian); the difference of the two numbers is the middle difference (zhongcha) with the august-pole difference taken away. Further, separately one finds that the product of the two walks is exactly one piece (duan) of the area of the semi-diameter; this is exactly the same in meaning as multiplying the half trapezoid-heads (ban titou). Now let us take it by the circle contained on the hypotenuse (xianshang rongyuan): set up the celestial unknown as the semi-diameter and duplicate it (fu); to the upper add the southward walk, obtaining □□, as the altitude; to the lower add the bu of the eastward walk, obtaining □□, as the base. Multiply base and altitude together, obtaining 丨□□, as the great rectangular area (zhiji); divide it by the celestial-unknown semi-diameter, obtaining □□□, as the base–altitude sum (gougu he) [Note: lodged on the left]. Then add base and altitude, obtaining □□, and cancel against the left, obtaining 丨○□. Extract the square root, obtaining 120 bu, which is the semi-diameter. This agrees with the question.\end{uniten}

\begin{unitzh}\unitlabel{ceyuan-029.26}或問：丙從乾隅南行，丁從艮隅亦南行；甲從乾隅東行，乙從坤隅亦東行；各不知步數，四人悉與城相直。只云：丙行内減丁行，餘四百五十步；甲行内減乙行，餘一百二十八步。問答同前。\end{unitzh}
\begin{uniten}Someone asks: Bing walks south from the qian corner, and Ding also walks south from the gen corner; Jia walks east from the qian corner, and Yi also walks east from the kun corner; none of their numbers of bu is known; the four persons are all in line with the city. It is only said: Bing's walk minus Ding's walk leaves 450 bu; Jia's walk minus Yi's walk leaves 128 bu. The question and answer are the same as before.\end{uniten}

\begin{unitzh}\unitlabel{ceyuan-029.27}法曰：二行相乘，為實；一步常法。得城徑。\end{unitzh}
\begin{uniten}Rule: Multiply the two walks together, making the dividend (shi); 1 bu as the constant divisor. One obtains the city diameter.\end{uniten}

\begin{unitzh}\unitlabel{ceyuan-029.28}草曰：别得丙行即大股，丁行即小差之股也；甲行即大勾，乙行即大差之勾也。其□即黄廣股，其□即黄長之勾也。立天元一為城徑。先置黄廣股□為股方差，以□為勾方差，以乘之，得□，為城徑冪{\zhsmall 寄左}。然後以天元冪與左相消，得下式丨□□。開平方，得二百四步。合問。\end{unitzh}
\begin{uniten}Working: Separately one finds that Bing's walk is the great altitude and Ding's walk is the altitude of the small difference; Jia's walk is the great base and Yi's walk is the base of the great difference. The □ is the yellow-breadth altitude; the □ is the base of the yellow length. Set up the celestial unknown as the city diameter. First set up the yellow-breadth altitude, □, as the altitude square-difference (gu fangcha); take □ as the base square-difference (gou fangcha); multiply them together, obtaining □, as the area of the city diameter [Note: lodged on the left]. Then cancel the area of the celestial unknown against the left, obtaining the form below 丨□□. Extract the square root, obtaining 204 bu [Translator's note: so the text; elsewhere in this chunk the city diameter is 240 bu.]. This agrees with the question.\end{uniten}

\begin{unitzh}\unitlabel{ceyuan-029.29}或問：出南門東行有槐樹一株，出東門南行有柳樹一株。丙、丁二人同立于坤隅，甲、乙二人同立于艮隅。丁直東行，至槐而止；乙直南行，至柳而止；丙直南行，甲直東行，四人遥相望見。只云：丙行多于丁行一百六十八步，乙行多于甲行七十步。問答同前。\end{unitzh}
\begin{uniten}Someone asks: Going out of the south gate and walking east, there is one locust tree; going out of the east gate and walking south, there is one willow tree. Bing and Ding, two persons, stand together at the kun corner; Jia and Yi, two persons, stand together at the gen corner. Ding walks straight east, reaches the locust and stops; Yi walks straight south, reaches the willow and stops; Bing walks straight south, Jia walks straight east, and the four persons catch sight of one another from afar. It is only said: Bing's walk exceeds Ding's walk by 168 bu, and Yi's walk exceeds Jia's walk by 70 bu. The question and answer are the same as before.\end{uniten}

\begin{unitzh}\unitlabel{ceyuan-029.30}法曰：云數相乘，為實；二數相減，又半之，為法。得城徑。\end{unitzh}
\begin{uniten}Rule: Multiply the stated numbers together, making the dividend (shi); subtract the two numbers from each other and further halve, making the divisor (fa). One obtains the city diameter.\end{uniten}

\begin{unitzh}\unitlabel{ceyuan-029.31}草曰：别得□即大差勾股較也，其□即小差上勾股較也。二數相併，為大差弦内減小差弦也；二數相較又半之，皇極弦與城徑差也；二數相併而半之，即皇極差也。立天元一為圓徑。二云相減數又半之，加天元，得□□，為極弦也。併二數而半之，得□，為極差也。副置極弦：上位加極差，得□□，為弦較和也；下位内減極差，得□□，為弦較較也。上下相乘，得丨□□，為二直積{\zhsmall 寄左}。然後以天元一乘極弦，得下式丨□，為同數，與左相消，得□□。上法下實而一，得二百四十步，即城徑也。合問。\end{unitzh}
\begin{uniten}Working: Separately one finds that the □ is the base–altitude difference of the great difference, and the □ is the base–altitude difference on the small difference. Adding the two numbers makes the great-difference hypotenuse with the small-difference hypotenuse subtracted from it; the difference of the two numbers, further halved, is the difference between the august-pole hypotenuse and the city diameter; the two numbers added and halved are the august-pole difference. Set up the celestial unknown as the circle's diameter. Halve the difference of the two stated numbers and add the celestial unknown, obtaining □□, as the pole hypotenuse (jixian). Add the two numbers and halve, obtaining □, as the pole difference (jicha). Set up the pole hypotenuse in duplicate: to the upper position add the pole difference, obtaining □□, as the hypotenuse–difference sum (xianjiao he); from the lower position subtract the pole difference, obtaining □□, as the hypotenuse–difference difference (xianjiao jiao). Multiply upper and lower together, obtaining 丨□□, as two rectangular areas [Note: lodged on the left]. Then multiply the pole hypotenuse by the celestial unknown, obtaining the form below 丨□, as the equal quantity; cancel against the left, obtaining □□. With the upper as divisor (fa) and the lower as dividend (shi), divide, obtaining 240 bu, which is the city diameter. This agrees with the question.\end{uniten}

\begin{unitzh}\unitlabel{ceyuan-029.32}或問：甲從坤東行，丙從艮南行，適相見，斜行一百二步，甲、丙相㑹。丙云：「我南行不及汝四十二步。」問答同前。\end{unitzh}
\begin{uniten}Someone asks: Jia walks east from kun and Bing walks south from gen; just then they see each other, and, walking obliquely 102 bu, Jia and Bing meet. Bing says: "My southward walk falls short of yours by 42 bu." The question and answer are the same as before.\end{uniten}

\begin{unitzh}\unitlabel{ceyuan-029.33}法曰：二數相併，以斜行乘於上；二數相併而半之，以乘相併數，減上位，為平實；不及步為從；一步常法。得虚勾。\end{unitzh}
\begin{uniten}Rule: Add the two numbers together and multiply by the oblique walk, in the upper position; add the two numbers and halve, multiply the sum by it, and subtract from the upper position, making the level dividend; the bu of the shortfall as the "following"; 1 bu as the constant divisor. One obtains the void base.\end{uniten}

\begin{unitzh}\unitlabel{ceyuan-029.34}草曰：别得一百二步即虚弦，四十二步即虚較也。又斜行得虚股為乙東行，此便為大差勾也；斜行步得虚勾為丙東行，此便是小差股也。立天元一為虚勾，加斜行步，得□□，為小差股也。以不及步加于小差股，得下式□□，為大差勾也。勾股相乘，得丨□□，為半段黄方冪{\zhsmall 寄左}。然後再置虚勾，加不及步，得□□，為虚股；又加入天元，得□□，為虚和；又加入虚弦，得□□，為圓徑。以自之，得□□□；又半之，得□□□，與寄左相消，得丨□□。平方開，得四十八步，即虚勾也。合問。\end{unitzh}
\begin{uniten}Working: Separately one finds that the 102 bu is the void hypotenuse and the 42 bu is the void difference. Further, the oblique walk getting the void altitude makes Yi's eastward walk; this is then the great-difference base; the bu of the oblique walk getting the void base makes Bing's eastward walk; this is then the small-difference altitude. Set up the celestial unknown as the void base; add the bu of the oblique walk, obtaining □□, as the small-difference altitude. Add the bu of the shortfall to the small-difference altitude, obtaining the form below □□, as the great-difference base. Multiply base and altitude together, obtaining 丨□□, as half a piece of the area of the yellow square (huangfang) [Note: lodged on the left]. Then again set up the void base and add the bu of the shortfall, obtaining □□, as the void altitude; further add in the celestial unknown, obtaining □□, as the void sum; further add in the void hypotenuse, obtaining □□, as the circle's diameter. Multiply it by itself, obtaining □□□; further halve it, obtaining □□□; cancel it against what is lodged on the left, obtaining 丨□□. Extract the square root, obtaining 48 bu, which is the void base. This agrees with the question.\end{uniten}

\begin{unitzh}\unitlabel{ceyuan-029.35}或問：甲從城心東行，丙從城心南行，庚從巽隅西行，壬從巽隅北行，四人遥相望見，各不知步數。只云：甲、丙共行了三百九十一，庚、壬共行了一百三十八。問答同前。\end{unitzh}
\begin{uniten}Someone asks: Jia walks east from the center of the city; Bing walks south from the center of the city; Geng walks west from the xun corner; Ren walks north from the xun corner; the four persons catch sight of one another from afar; none of their numbers of bu is known. It is only said: Jia and Bing together walked 391; Geng and Ren together walked 138. The question and answer are the same as before.\end{uniten}

\begin{unitzh}\unitlabel{ceyuan-029.36}法曰：云數相乘，為實；相併，為法。得虚弦。\end{unitzh}
\begin{uniten}Rule: Multiply the stated numbers together, making the dividend (shi); add them together, making the divisor (fa). One obtains the void hypotenuse.\end{uniten}

\begin{unitzh}\unitlabel{ceyuan-030.1}草曰：别得甲丙共為皇極和也，又為極弦、極黄共；庚壬共為太虚和也，又為虚弦、虚黄共。立天元一為皇極黄方靣{\zhsmall 亦為虚弦也}，減于甲丙共，得□□，即極弦也。又以天元減于庚壬共，得□□，即太虚黄方靣也。以太虚黄方靣乘極弦，得丨□□{\zhsmall 寄左}。然後以天元冪與左相消，得□□，上法下實，如法得一百二步，即皇極黄方靣也。合問。{\zhsmall 按此亦係相消後得一邊之二數者}\end{unitzh}
\begin{uniten}Working: Separately it is found that the joint total of Jia and Bing is the sum (he) of the Imperial Pole (huangji), and is also the Pole hypotenuse and the Pole yellow (jihuang) together; the joint total of Geng and Ren is the sum of the Great Void (taixu), and is also the Void hypotenuse and the Void yellow together. Set up the celestial unknown as the side of the yellow square (huangfang mian) of the Imperial Pole [Note: it is also the Void hypotenuse]; subtract it from the joint total of Jia and Bing, obtaining □□ [Translator's note: □ marks counting-rod numerals and arrays not reproduced in the transcription.], which is the Pole hypotenuse. Again subtract the celestial unknown from the joint total of Geng and Ren, obtaining □□, which is the side of the yellow square of the Great Void. Multiply the Pole hypotenuse by the side of the yellow square of the Great Void, obtaining 丨□□ [Note: set aside on the left (jizuo)]. Then eliminate (xiangxiao) the area (mi) of the celestial unknown against the left, obtaining □□; the upper is the divisor (fa), the lower the dividend (shi); dividing by the divisor gives 102 bu (步), which is the side of the yellow square of the Imperial Pole. This agrees with the question. [Siku editors: This too is one of the cases in which, after the elimination, the two numbers of one side are obtained.]\end{uniten}

\begin{unitzh}\unitlabel{ceyuan-030.2}或問：甲從乾隅東行，不知步數而止；丙向南行，亦不知步數，望見甲，就甲斜行七百八十步，與甲相㑹。甲云：「我行地雖少于汝，以我東行步為法，除汝南行步，則汝止得二步四分。」問答同前。\end{unitzh}
\begin{uniten}Someone asks: Jia walks east from the northwest corner (qian yu) and stops, the number of bu not being known; Bing walks south, the number of bu also not being known; he sees Jia and, walking obliquely toward Jia for 780 bu, meets Jia. Jia says: "Although the ground I have walked is less than yours, if my eastward bu are taken as the divisor to divide your southward bu, then you get only 2 bu 4 fen (分)." The question and answer are the same as before.\end{uniten}

\begin{unitzh}\unitlabel{ceyuan-030.3}法曰：斜步自之為平實，除步自之，又加一步為隅，得甲東行。\end{unitzh}
\begin{uniten}Rule: Multiply the oblique bu by itself to make the square dividend (pingshi); multiply the bu of the division by itself, and further add 1 bu, to make the corner (yu); this gives Jia's eastward walk.\end{uniten}

\begin{unitzh}\unitlabel{ceyuan-030.4}草曰：此問所求城徑與諸問並同，其勾股則與前後諸率不同。今特為此草者，欲使後學有以考較諸率當否也。立天元一為甲東行{\zhsmall 即大勾}，以乗二步四分，得□為長；以自之，得□□為股冪；又併入天元冪，得□□為弦冪{\zhsmall 寄左}。乃以斜行自之，得□為同數，與左相消，得□□□，開平方得三百，即甲東行也。以二步四分乘之，得七百二十步，即丙南行也。倍丙東行，以甲東行乘之，得四十三萬二千為實；以三事和一千八百為法除之，得二百四十步，即城徑也。合問。\end{unitzh}
\begin{uniten}Working: The diameter of the city (chengjing) sought in this question is the same as in all the questions, but its base (gou) and altitude (gu) differ from the rates (lü) before and after. That I now specially make this working is because I wish to give later students the means to check whether the rates are correct or not. Set up the celestial unknown as Jia's eastward walk [Note: this is the great base (da gou)]; multiply it by 2 bu 4 fen, obtaining □ as the length (chang); multiply it by itself, obtaining □□ as the area of the altitude; further add in the area of the celestial unknown, obtaining □□ as the area of the hypotenuse [Note: set aside on the left]. Then multiply the oblique walk by itself, obtaining □ as the equal number (tongshu); eliminate it against the left, obtaining □□□; extracting the square root gives 300, which is Jia's eastward walk. Multiply it by 2 bu 4 fen, obtaining 720 bu, which is Bing's southward walk. Double Bing's eastward walk [Translator's note: 東 "eastward" here evidently stands for 南 "southward".] and multiply it by Jia's eastward walk, obtaining 432000 as the dividend (shi); take the sum of the three things (sanshi he), 1800, as the divisor (fa) and divide by it, obtaining 240 bu, which is the diameter of the city. This agrees with the question.\end{uniten}

\begin{unitzh}\unitlabel{ceyuan-030.5}或問：小差黄方靣少于大差黄方靣八十四步，太虚黄方靣少于皇極黄方靣六十六步。問答同前。\end{unitzh}
\begin{uniten}Someone asks: The side of the yellow square of the Small Difference (xiaocha) is less than the side of the yellow square of the Great Difference (dacha) by 84 bu; the side of the yellow square of the Great Void is less than the side of the yellow square of the Imperial Pole by 66 bu. The question and answer are the same as before.\end{uniten}

\begin{unitzh}\unitlabel{ceyuan-030.6}法：半八十四為中差；以中差減六十六為二小差，半之為小差；又中、小差相併為大差。乃以小差乘大差為平實，半步常法，得虚黄。\end{unitzh}
\begin{uniten}Rule: Halve 84 to make the middle difference (zhongcha); subtract the middle difference from 66 to make two small differences; halve it to make the small difference; further add the middle and small differences together to make the great difference. Then multiply the great difference by the small difference to make the square dividend; 1/2 bu as the constant divisor (changfa); this gives the Void yellow.\end{uniten}

\begin{unitzh}\unitlabel{ceyuan-030.7}草曰：别得八十四為兩个虚積中差，其六十六為虚積大小差併。半八十四，得□，為虚中差也。以中差減六十六，餘二十四，半之，得□，即虚小差也。以小差反減六十六，餘□，即虚大差也。又别得小差黄方為兩叀股，大差黄方為兩明勾也。立天元一為虚黄方，置三位：上加小差，得□□，為虚勾也；中加大差，得下□□，為虚股也；下加大小差併，得□□，為虚弦也。三位併之，得□□，即城徑也。倍虚勾減城徑，得□□，為大差黄方靣也；又倍虚股減城徑，得□□，為小差黄方靣也。半小差黄方靣，得□□，以乘大差黄方，得□□□，為一个虚直積{\zhsmall 寄左}。乃以虚勾、虚股相乘，得丨□□，為同數，與左相消，得□□□，平方開得三十六步，即虚黄方靣也。其餘依法求之。合問。據此問既别得大小差正數，自可以求得黄方靣也。諸如此數，實不湏草，今特為細草者，庶使後學知其來歴。\end{unitzh}
\begin{uniten}Working: Separately it is found that 84 is two middle differences of the Void area (xuji), and the 66 is the great and small differences of the Void area combined. Halve 84, obtaining □, as the Void middle difference. Subtract the middle difference from 66, leaving 24; halve it, obtaining □, which is the Void small difference. Conversely subtract the small difference from 66, leaving □, which is the Void great difference. Further, separately it is found that the yellow square of the Small Difference is two Special (zhuan) altitudes, and the yellow square of the Great Difference is two Bright (ming) bases. Set up the celestial unknown as the Void yellow square, and place it in three positions: to the upper add the small difference, obtaining □□, as the Void base; to the middle add the great difference, obtaining the following □□, as the Void altitude; to the lower add the great and small differences combined, obtaining □□, as the Void hypotenuse. Add the three positions together, obtaining □□, which is the diameter of the city. Double the Void base and subtract it from the diameter of the city, obtaining □□, as the side of the yellow square of the Great Difference; further double the Void altitude and subtract it from the diameter of the city, obtaining □□, as the side of the yellow square of the Small Difference. Halve the side of the yellow square of the Small Difference, obtaining □□, and multiply the yellow square of the Great Difference by it, obtaining □□□, as one Void rectangle area (zhiji) [Note: set aside on the left]. Then multiply the Void base and the Void altitude together, obtaining 丨□□, as the equal number; eliminate it against the left, obtaining □□□; extracting the square root gives 36 bu, which is the side of the Void yellow square. The rest is to be found according to the rule. This agrees with the question. According to this question, since the true numbers of the great and small differences have been separately obtained, the side of the yellow square can naturally be found. For all numbers such as these, really no working is needed; that I now specially make a detailed working (xicao) is in the hope that later students may know where it comes from.\end{uniten}

\begin{unitzh}\unitlabel{ceyuan-030.8}或問：大差弦較較減皇極弦，餘四十九步；小差弦較和減太虚弦，餘一百三十八步；又皇極差一百一十九步。問答同前。\end{unitzh}
\begin{uniten}Someone asks: The hypotenuse–difference difference (xianjiao jiao) of the Great Difference, subtracted from the hypotenuse of the Imperial Pole, leaves 49 bu; the hypotenuse–difference sum (xianjiao he) of the Small Difference, diminished by the hypotenuse of the Great Void, leaves 138 bu; further, the difference of the Imperial Pole is 119 bu. The question and answer are the same as before.\end{uniten}

\begin{unitzh}\unitlabel{ceyuan-030.9}法曰：併前二數為冪，内減極差冪為平實，從空，二益隅，得虚弦。\end{unitzh}
\begin{uniten}Rule: Add the first two numbers together to make an area (mi); from it subtract the area of the Pole difference to make the square dividend; the joining term (cong) is empty; 2 as the subtractive corner (yiyu); this gives the Void hypotenuse.\end{uniten}

\begin{unitzh}\unitlabel{ceyuan-030.10}草曰：别得大差弦較較與小差弦較和，皆同為圓徑也。又二數相併，得□，為明弦、叀弦共，又為極和内少兩个虚弦也。其一百三十八即虚和也，□則旁差也。立天元一為虚弦，加入一百三十八，得□□，為圓徑也；又加入□，得□□，為極弦。以自之，得丨□□；又倍之，得□□□；内却減極差冪□，得下式□□□，為和冪{\zhsmall 寄左}。乃倍天元加併數，得□□，為極和，以自增乘，得□□□，為同數，與左相消，得□□□，開平方得一百二步，即虚弦也。加入一百三十八，得二百四十步，為圓徑。合問。{\zhsmall 前二數相併加虚弦便是極弦}\end{unitzh}
\begin{uniten}Working: Separately it is found that the hypotenuse–difference difference of the Great Difference and the hypotenuse–difference sum of the Small Difference are both equal to the diameter of the circle (yuanjing). Further, add the two numbers together, obtaining □, which is the Bright hypotenuse and the Special hypotenuse together, and is also the Pole sum less two Void hypotenuses. The 138 is the Void sum, and □ is the side difference (pangcha). Set up the celestial unknown as the Void hypotenuse; add in 138, obtaining □□, as the diameter of the circle; further add in □, obtaining □□, as the Pole hypotenuse. Multiply it by itself, obtaining 丨□□; further double it, obtaining □□□; from it then subtract the area of the Pole difference, □, obtaining the following form □□□, as the area of the sum [Note: set aside on the left]. Then double the celestial unknown and add the combined number, obtaining □□, as the Pole sum; multiply it by itself, obtaining □□□, as the equal number; eliminate it against the left, obtaining □□□; extracting the square root gives 102 bu, which is the Void hypotenuse. Add in 138, obtaining 240 bu, as the diameter of the circle. This agrees with the question. [Note: The first two numbers added together, with the Void hypotenuse added, are precisely the Pole hypotenuse.]\end{uniten}

\begin{unitzh}\unitlabel{ceyuan-030.11}或問：小差不及平弦五十六步，髙弦不及大差一百五步。問答同前。\end{unitzh}
\begin{uniten}Someone asks: The Small Difference falls short of the Level (ping) hypotenuse by 56 bu; the High (gao) hypotenuse falls short of the Great Difference by 105 bu. The question and answer are the same as before.\end{uniten}

\begin{unitzh}\unitlabel{ceyuan-030.12}法曰：以前數自之為實，二數相減為法，得平勾。\end{unitzh}
\begin{uniten}Rule: Multiply the first number by itself to make the dividend; subtract the two numbers from each other to make the divisor; this gives the Level base.\end{uniten}

\begin{unitzh}\unitlabel{ceyuan-030.13}草曰：别得云數相併，得□，為平勾不及髙股也。此數得極差，則通差也；此數内減虚差，則極差也。云數相減，餘□，即城徑不及極弦也。以前數減于半徑，餘即平勾；以後數加于半徑，即髙股也。倍前數加小差，則為股圓差之勾也，此與前數加平弦同；倍後數減于大差，則為勾圓差之股也，此與後數減于髙弦同。立天元一為平勾，加相併數，得□□，即髙股也；又加天元，得□□，即極弦也；内減二云數差，得□□，為城徑也。半之，得□□，以自之，得丨□□，為半徑冪{\zhsmall 寄左}。然後以天元乘髙股，得丨□，為同數，與左相消，得□□，上法下實，得六十四步，即平勾也。合問。\end{unitzh}
\begin{uniten}Working: Separately it is found that the stated numbers (yunshu) added together, giving □, are the amount by which the Level base falls short of the High altitude. This number, gaining the Pole difference, is then the Comprehensive (tong) difference; this number, with the Void difference subtracted from it, is then the Pole difference. The stated numbers subtracted from each other leave □, which is the amount by which the diameter of the city falls short of the Pole hypotenuse. Subtract the first number from the radius (banjing), and the remainder is the Level base; add the latter number to the radius, and it is the High altitude. Double the first number and add the Small Difference, and it is the base of the altitude–circle difference (gu yuan cha); this is the same as the first number added to the Level hypotenuse. Double the latter number and subtract it from the Great Difference, and it is the altitude of the base–circle difference (gou yuan cha); this is the same as the latter number subtracted from the High hypotenuse. Set up the celestial unknown as the Level base; add the combined number, obtaining □□, which is the High altitude; further add the celestial unknown, obtaining □□, which is the Pole hypotenuse; from it subtract the difference of the two stated numbers, obtaining □□, as the diameter of the city. Halve it, obtaining □□; multiply it by itself, obtaining 丨□□, as the area of the radius [Note: set aside on the left]. Then multiply the High altitude by the celestial unknown, obtaining 丨□, as the equal number; eliminate it against the left, obtaining □□; the upper is the divisor, the lower the dividend; this gives 64 bu, which is the Level base. This agrees with the question.\end{uniten}

\begin{unitzh}\unitlabel{ceyuan-030.14}又法：云數相得為實，相減為法，得半徑。\end{unitzh}
\begin{uniten}Another rule: Multiply the stated numbers together (xiangde) to make the dividend; subtract them from each other to make the divisor; this gives the radius.\end{uniten}

\begin{unitzh}\unitlabel{ceyuan-030.15}草曰：立天元為半徑，副之：上内減五十六，得□□，為平勾；下加一百五，得□□，為髙股。上下相乘，得丨□□，為半徑冪{\zhsmall 寄左}。以天元冪與左相消，得下式□□，上法下實，得一百二十步，即半徑也。合問。\end{unitzh}
\begin{uniten}Working: Set up the celestial unknown as the radius, and duplicate it (fu): from the upper subtract 56, obtaining □□, as the Level base; to the lower add 105, obtaining □□, as the High altitude. Multiply the upper and the lower together, obtaining 丨□□, as the area of the radius [Note: set aside on the left]. Eliminate the area of the celestial unknown against the left, obtaining the following form □□; the upper is the divisor, the lower the dividend; this gives 120 bu, which is the radius. This agrees with the question.\end{uniten}

\begin{unitzh}\unitlabel{ceyuan-030.16}或問：通勾、通弦共一千步，大差、小差共得四百四十步。問答同前。\end{unitzh}
\begin{uniten}Someone asks: The Comprehensive base and the Comprehensive hypotenuse together are 1000 bu; the Great Difference and the Small Difference together make 440 bu. The question and answer are the same as before.\end{uniten}

\begin{unitzh}\unitlabel{ceyuan-030.17}法曰：以二差共減于一千，又半之，以自乘為平實；以二差共減于一千，又半之，加入二之前數為縱{\zhsmall 前數謂一千也。按此語有誤，應加入二之後數，後數謂大小差共也}；二步二分五釐益隅；得勾圓差。\end{unitzh}
\begin{uniten}Rule: Subtract the two differences combined from 1000, further halve it, and multiply it by itself to make the square dividend; subtract the two differences combined from 1000, further halve it, and add in twice the first number to make the joining term (zong) [Note: "The first number" means 1000.] [Siku editors: This wording is in error; it should be "add in twice the latter number", the latter number meaning the Great and Small Differences together.]; 2 bu 2 fen 5 li (釐) as the subtractive corner; this gives the base–circle difference.\end{uniten}

\begin{unitzh}\unitlabel{ceyuan-030.18}草曰：立天元一為小差數，加入後數，得□□；却以減于前數，得□□；折半，得□□，為一个圓徑也。以自之，得下式□□□{\zhsmall 寄左}。然後以天元減後數，得□□，為大差；以天元乘之，又倍之，得□□，與左相消，得□□□，開平方得八十步，即勾圓差也。\end{unitzh}
\begin{uniten}Working: Set up the celestial unknown as the number of the Small Difference; add in the latter number, obtaining □□; then subtract it from the first number, obtaining □□; halve it, obtaining □□, as one diameter of the circle. Multiply it by itself, obtaining the following form □□□ [Note: set aside on the left]. Then subtract the celestial unknown from the latter number, obtaining □□, as the Great Difference; multiply it by the celestial unknown, and further double it, obtaining □□; eliminate it against the left, obtaining □□□; extracting the square root gives 80 bu, which is the base–circle difference.\end{uniten}

\begin{unitzh}\unitlabel{ceyuan-030.19}或問：皇極三事和六百八十步，太虚弦和較三十六。問答同前。\end{unitzh}
\begin{uniten}Someone asks: The sum of the three things of the Imperial Pole is 680 bu; the hypotenuse–sum difference (xianhe jiao) of the Great Void is 36. The question and answer are the same as before.\end{uniten}

\begin{unitzh}\unitlabel{ceyuan-030.20}法曰：二數相得為實，半之後數為益從，五分常法，平開得城徑。\end{unitzh}
\begin{uniten}Rule: Multiply the two numbers together to make the dividend; half of the latter number as the subtractive joining term (yicong); 5 fen as the constant divisor; extracting the square root gives the diameter of the city.\end{uniten}

\begin{unitzh}\unitlabel{ceyuan-030.21}草曰：别得皇極三事和即大弦也。立天元一為城徑，減三个後數□而半之，得□□，為太虚大小差併也；却加入兩个後數□，得下□□，為虚和也；又以虚和減天元，得下□□，為虚弦也。置通弦{\zhsmall 即皇極三事和也}，内加天元，得下式□□，即通和也。乃置通和，以虚弦乘之，得下式□□□{\zhsmall 寄左}。再置虚和，以通弦乘之，得下□□，為同數，與左相消，得□□□，開平方得二百四十步，即城徑也。合問。\end{unitzh}
\begin{uniten}Working: Separately it is found that the sum of the three things of the Imperial Pole is the great hypotenuse (da xian). Set up the celestial unknown as the diameter of the city; subtract three of the latter number, □, and halve it, obtaining □□, as the great and small differences of the Great Void combined; then add in two of the latter number, □, obtaining the following □□, as the Void sum; further subtract the Void sum from the celestial unknown, obtaining the following □□, as the Void hypotenuse. Set up the Comprehensive hypotenuse [Note: this is the sum of the three things of the Imperial Pole], add the celestial unknown to it, obtaining the following form □□, which is the Comprehensive sum. Then set up the Comprehensive sum and multiply it by the Void hypotenuse, obtaining the following form □□□ [Note: set aside on the left]. Again set up the Void sum and multiply it by the Comprehensive hypotenuse, obtaining the following □□, as the equal number; eliminate it against the left, obtaining □□□; extracting the square root gives 240 bu, which is the diameter of the city. This agrees with the question.\end{uniten}

\begin{unitzh}\unitlabel{ceyuan-030.22}或問：出南門行一百三十五步有樹；出北門行一十五步，折而東行二百八步，望見。問答同前。\end{unitzh}
\begin{uniten}Someone asks: Going out of the south gate and walking 135 bu, there is a tree; going out of the north gate and walking 15 bu, then turning and walking east 208 bu, one sees it. The question and answer are the same as before.\end{uniten}

\begin{unitzh}\unitlabel{ceyuan-030.23}法曰：以東行步乘南行步，得數，又自乘為實；以東行步自乘，乘南行步，又倍之為從；東行步自乘于上，併南北二行步以減于東行步，餘數自之為冪，以減上，再寄位；又併南北二行步，以東行步乘而倍之，内減再寄，為第一益亷；四之東行步于上，又併南北二行步減于東行步，又四之，減上位，為第二益亷；四步虚隅；開三乘方，得半徑。\end{unitzh}
\begin{uniten}Rule: Multiply the southward bu by the eastward bu, obtaining a number; further multiply it by itself to make the dividend. Multiply the eastward bu by itself, multiply by the southward bu, and further double it to make the joining term. [Place] the eastward bu multiplied by itself in the upper [position]; add together the two walks, southward and northward, and subtract them from the eastward bu; multiply the remaining number by itself to make an area; subtract it from the upper, and set [the result] aside again in a position (zai jiwei). Further add together the two walks, southward and northward, multiply by the eastward bu and double it; from this subtract what was set aside again, to make the first subtractive side term (yilian). [Place] 4 times the eastward bu in the upper [position]; further add together the two walks, southward and northward, and subtract them from the eastward bu, further multiply by 4, and subtract it from the upper position, to make the second subtractive side term. 4 bu as the empty corner (xuyu). Extract the root of the thrice-multiplied square (sancheng fang), obtaining the radius.\end{uniten}

\begin{unitzh}\unitlabel{ceyuan-030.24}草曰：立天元一為半徑{\zhsmall 即髙勾也}。置南行步，加天元，得□□，為髙弦也。置大勾□，以髙弦乘之，得□□；復以髙勾除之，得下式□□，為大弦也。令之自乘，得□□□{\zhsmall 寄左}。又置二之天元，加南北行併，得□□，為大股；復用大勾二百八減之，得□□，為較也。以自乘，得□□□，為較冪；以減寄左，得□□□□□，為二直積{\zhsmall 寄左}。再置大股□□，以大勾□乘之，得□□，為直積；又倍之，得□□，為同數，與左相消，得□□□□□，翻法開三乘方，得一百二十步，即城徑之半也。合問。\end{unitzh}
\begin{uniten}Working: Set up the celestial unknown as the radius [Note: this is the High base]. Set up the southward bu and add the celestial unknown, obtaining □□, as the High hypotenuse. Set up the great base, □, and multiply it by the High hypotenuse, obtaining □□; again divide it by the High base, obtaining the following form □□, as the great hypotenuse. Let it be multiplied by itself, obtaining □□□ [Note: set aside on the left]. Further set up twice the celestial unknown and add the southward and northward walks combined, obtaining □□, as the great altitude; again subtract the great base, 208, from it, obtaining □□, as the difference (jiao). Multiply it by itself, obtaining □□□, as the area of the difference; subtract it from what was set aside on the left, obtaining □□□□□, as two rectangle areas [Note: set aside on the left]. Again set up the great altitude, □□, and multiply it by the great base, □, obtaining □□, as the rectangle area; further double it, obtaining □□, as the equal number; eliminate it against the left, obtaining □□□□□; reversing the method (fanfa), extract the root of the thrice-multiplied square, obtaining 120 bu, which is half the diameter of the city. This agrees with the question.\end{uniten}

\begin{unitzh}\unitlabel{ceyuan-030.25}或問：出北門一十五步，折而東行二百八步有樹；出西門八步，折而南行四百九十五步見之。問答同前。\end{unitzh}
\begin{uniten}Someone asks: Going out of the north gate 15 bu and turning to walk east 208 bu, there is a tree; going out of the west gate 8 bu and turning to walk south 495 bu, one sees it. The question and answer are the same as before.\end{uniten}

\begin{unitzh}\unitlabel{ceyuan-030.26}法曰：先置南行步，内減一東二西併步，餘二百七十一，為前泛率；次併一南二北，内減東行步，餘三百一十七，為中泛率；次併東西步，以南行步乘之于上位，又以西行乘南北併，得數減上位，餘一十萬二千八百四十，為後泛率。乃以後泛率自乘，得一百五億七千六百六萬五千六百，為三乘方實；以前中二泛相減，餘四十六，以乘後法數為從；前中二泛相乘，得八萬五千九百七，加入二之後泛數，共得二十九萬一千五百八十七于上位；又併東西行以乘南北併，得二十二萬三百二十，加上位，通得五十一萬一千九百七，為第一亷；二之前泛數，加入四之東西併，得一千四百五十二于上位；又以前中二泛相減于四十六，減上位，餘一千四百六，為第二亷；一步常法；得半徑。{\zhsmall 按此法乃取於又法草中，其求第二亷云「二之前泛數」句誤，當云「二之四數併」。若二之前泛數加入四之東西併，便得第二亷一千四百零六，更不待再減，然原文之意不如是也}\end{unitzh}
\begin{uniten}Rule: First set up the southward bu; from it subtract one eastward and two westward [walks] combined; the remainder, 271, is the first provisional rate (fanlü). Next combine one southward and two northward [walks]; from it subtract the eastward bu; the remainder, 317, is the middle provisional rate. Next combine the eastward and westward bu and multiply them by the southward bu, [placing the result] in the upper position; further multiply the southward and northward combined by the westward walk, and subtract the number obtained from the upper position; the remainder, 102840, is the latter provisional rate. Then multiply the latter provisional rate by itself, obtaining 10576065600, as the dividend of the thrice-multiplied square. Subtract the first and middle provisional [rates] from each other, leaving 46, and multiply the latter fa number [Translator's note: 法 here is perhaps an error for 泛, "provisional".] by it to make the joining term. Multiply the first and middle provisional [rates] together, obtaining 85907; add in twice the latter provisional number, obtaining altogether 291587, in the upper position; further combine the eastward and westward walks and multiply by the southward and northward combined, obtaining 220320; add it to the upper position, obtaining in all 511907, as the first side term. Twice the first provisional number, with four times the eastward and westward combined added in, gives 1452, in the upper position; further subtract the first and middle provisional [rates] from each other, at 46, and subtract it from the upper position, leaving 1406, as the second side term. 1 bu as the constant divisor. This gives the radius. [Siku editors: This rule is taken from the working of the alternative rule. Where it seeks the second side term, the phrase "twice the first provisional number" is in error; it should say "twice the four numbers combined". If twice the first provisional number is added to four times the eastward and westward combined, one obtains at once the second side term, 1406, with no need for a further subtraction; but the intention of the original text is not like this.]\end{uniten}

\begin{unitzh}\unitlabel{ceyuan-031.1}草曰：立天元一為半城徑，加入東行、西行併，得□□，為大勾也。又置天元，加入南行、北行併，得□□，為大股也。置西行八步，以大股乘之，得下式□□，合以大勾除之，不除，寄為母，便以此為股尖也。置南行四百九十五步，減天元，得□□，用分母大勾乘之，乘訖得下式□□□，内減了股尖，餘□□□，為小股也。{\zhsmall 内帶大勾分母。}置小股，合以大勾乘了，復以大股除之，為小勾；今為小股内已有大勾為母，更不湏乘，只以小股□□□便為小勾也。{\zhsmall 内帶大股為母。}小勾、小股相乘，得數為一个小勾股相乘直積，内帶大勾股相乘直積為分母也。乃以半城徑{\zhsmall 即天元也}除之，為一个弦較和也：丨□□□□。此法本取勾外容圓，合以弦較和除二積，為勾外所容之圓；今用天元半徑除一个積，則却得一个弦較和也。内依舊帶大積分母也。{\zhsmall 寄左。}\end{unitzh}
\begin{uniten}Working: Set up the celestial unknown (tianyuan) as half the city diameter, and add it to the sum of the east walk and the west walk, obtaining □□ [Translator's note: □ marks counting-rod arrays not reproduced in the transcription.], which is the large base (gou). Again put down the celestial unknown and add it to the sum of the south walk and the north walk, obtaining □□, which is the large altitude (gu). Put down the west walk, 8 bu (步), and multiply it by the large altitude, obtaining the following expression (shi) □□. It should be divided by the large base; do not divide, but lodge it as the denominator, and take this as the altitude tip (gujian). Put down the south walk, 495 bu, and subtract the celestial unknown, obtaining □□. Multiply it by the denominator, the large base; when the multiplication is finished, one obtains the following expression □□□. From it subtract the altitude tip, leaving □□□, which is the small altitude. [Note: It contains the large base as denominator.] Put down the small altitude; it should be multiplied by the large base and then divided by the large altitude to make the small base. Now, since the small altitude already contains the large base as denominator, there is no further need to multiply: just take the small altitude □□□ as the small base. [Note: It contains the large altitude as denominator.] Multiply the small base and the small altitude together; the number obtained is one rectangle area (zhiji) of the small base and altitude multiplied together, which contains as denominator the rectangle area of the large base and altitude multiplied together. Then divide it by half the city diameter [Note: that is, the celestial unknown], making one hypotenuse–difference sum (xianjiaohe): 丨□□□□. This method is originally taken from the circle contained outside the base (gouwai rongyuan): one should divide 2 areas by the hypotenuse–difference sum to make the circle contained outside the base; now, dividing one area by the celestial-unknown radius, one instead obtains one hypotenuse–difference sum. Inside, it still carries the large area as denominator, as before. [Note: Lodge it on the left (jizuo).]\end{uniten}

\begin{unitzh}\unitlabel{ceyuan-031.2}然後再置小股□□□，合用大積乘之，縁内已帶大勾分母，今只用大股□□乘之，得□□□□，為大積所乘小股，于上。再置小勾，合用大積乘之，縁内已帶大股分母，合只用大勾□□乘之，得□□□□，為大積所乘之小勾也。以此小勾減上小股，得□□□，即帶分小較也。又二因小較，得下式□□□，為帶分二較也。又以大勾股直積丨□□乘二之天元半徑，得□□□，為一个帶分弦較較也。{\zhsmall 弦較較乘弦較和為二直積，既以圓徑除二百積為弦較和，則是圓徑為弦較較也；今又為半天元圓徑除一積為弦較和，故倍天元半徑作一个弦較較也。}遂将此弦較較加入前二較，得□□□□，亦為一个弦較和也。與寄左相消，得下式丨□□□□，開三乘方，得一百二十步，即城半徑也。合問。\end{unitzh}
\begin{uniten}Afterwards, again put down the small altitude □□□. It should be multiplied by the large area, but because it already contains the large base as denominator, now just multiply it by the large altitude □□, obtaining □□□□, which is the small altitude multiplied by the large area, placed above. Again put down the small base. It should be multiplied by the large area, but because it already contains the large altitude as denominator, one should just multiply it by the large base □□, obtaining □□□□, which is the small base multiplied by the large area. Subtract this small base from the small altitude above, obtaining □□□, which is the small difference (xiaojiao) carrying the denominator. Again multiply (yin) the small difference by 2, obtaining the following expression □□□, which is the 2 differences carrying the denominator. Again, with the rectangle area of the large base and altitude, 丨□□, multiply twice the celestial-unknown radius, obtaining □□□, which is one hypotenuse–difference difference (xianjiaojiao) carrying the denominator. [Note: The hypotenuse–difference difference multiplied by the hypotenuse–difference sum makes 2 rectangle areas. Since dividing the 200 areas [Translator's note: 二百積, "200 areas", is evidently an error for 二直積, "2 rectangle areas".] by the circle diameter makes the hypotenuse–difference sum, the circle diameter is the hypotenuse–difference difference. Now, moreover, because dividing one area by half the circle diameter, the celestial unknown, makes the hypotenuse–difference sum, therefore double the celestial-unknown radius to make one hypotenuse–difference difference.] Then add this hypotenuse–difference difference to the preceding 2 differences, obtaining □□□□, which is also one hypotenuse–difference sum. Cancel it against what was lodged on the left (xiangxiao), obtaining the following expression 丨□□□□. Extract the root of the thrice-multiplied square (sancheng fang), obtaining 120 bu, which is the city radius. This agrees with the question.\end{uniten}

\begin{unitzh}\unitlabel{ceyuan-031.3}又法：此問係是洞淵測圓門第一十三，前答亦依洞淵細草，用勾外容圓術以入于弦較和，然其數煩碎，宛轉費力。今别草一法，其亷從與前不殊，而中間段絡逕捷明白，方之前術，極為省易，學者當自知也。\end{unitzh}
\begin{uniten}Another method: This question is the 13th of Dongyuan's Circle-Measuring Section (Ceyuan men). The preceding answer also followed Dongyuan's detailed working (xicao), using the method of the circle contained outside the base in order to enter into the hypotenuse–difference sum. But its numbers are tedious and fragmentary, and it is roundabout and laborious. Now I separately draft one method: its edges (lian) and joining term (cong) do not differ from the preceding, but the sections and links (duanluo) in between are direct, quick, clear and plain; compared with the preceding method it is extremely economical and easy. Learners should come to know this for themselves.\end{uniten}

\begin{unitzh}\unitlabel{ceyuan-031.4}立天元一為半徑，副之：上併加東西行，得□□，為通勾率；下併加南北行，得□□，為通股率。乃置西行八步，以通股乘之，得下□□，合通勾除，不除，寄為母，便以此為南小股也。又置南行四百九十五步，内減天元，得□□，用通勾乘之，得□□□，内減了南小股，下式卜□□為股圓差也，内帶通勾分母。又置北行一十五步，以通勾乘之，得□□，合通股除，不除，寄為母，便以此為北小勾也。又置東行二百八步，内減天元，得□□，用通股乘之，得□□□，内減了北小勾，餘□□□，為勾圓差也。{\zhsmall 内帶通股分母。}乃以二差相乘，得下式丨□{\zhsmall □□}□{\zhsmall □□}，為半段圓徑冪也，内帶通積為母。{\zhsmall 寄左。}然後以通勾、通股相乘，得丨□□，以天元冪乘之，得丨□□□，又倍之，得下式□□□□，為同數，與左相消，得亷從一與前同。合問。\end{unitzh}
\begin{uniten}Set up the celestial unknown as the radius and duplicate it (fu): to the upper one add the east and west walks together, obtaining □□, which is the rate (lü) of the total base (tonggou); to the lower one add the south and north walks together, obtaining □□, which is the rate of the total altitude (tonggu). Then put down the west walk, 8 bu, and multiply it by the total altitude, obtaining the following □□. It should be divided by the total base; do not divide, but lodge it as the denominator, and take this as the southern small altitude. Again put down the south walk, 495 bu; from it subtract the celestial unknown, obtaining □□; multiply it by the total base, obtaining □□□; from it subtract the southern small altitude; the following expression 卜□□ is the altitude–circle difference (guyuan cha); it contains the total base as denominator. Again put down the north walk, 15 bu, and multiply it by the total base, obtaining □□. It should be divided by the total altitude; do not divide, but lodge it as the denominator, and take this as the northern small base. Again put down the east walk, 208 bu; from it subtract the celestial unknown, obtaining □□; multiply it by the total altitude, obtaining □□□; from it subtract the northern small base, leaving □□□, which is the base–circle difference (gouyuan cha). [Note: It contains the total altitude as denominator.] Then multiply the 2 differences together, obtaining the following expression 丨□[□□]□[□□], which is a half portion of the area (mi) of the circle diameter; it contains the total area (tongji) as denominator. [Note: Lodge it on the left.] Afterwards multiply the total base and the total altitude together, obtaining 丨□□; multiply it by the area (mi) of the celestial unknown, obtaining 丨□□□; again double it, obtaining the following expression □□□□, which is the same number (tongshu). Cancel it against the left, obtaining edges and joining term entirely the same as before. This agrees with the question.\end{uniten}

\begin{unitzh}\unitlabel{ceyuan-031.5}按：洞淵疑為古之精於算者。序中謂「老大以來，得洞淵九容之說」，而於此問又明其為洞淵測圓門第十三題，前答亦依其細草。大抵是書之作，皆師其意而演之者也。今洞淵之為人與書雖不可考，而即此一草觀之，其取徑遙深，而惟變所適，亦可見文豹之一班矣。至謂「其數煩碎宛轉費力」，特為初學難易而言，讀者宜善㑹也。\end{unitzh}
\begin{uniten}Note: Dongyuan is suspected to have been one of the ancients who were expert in computation. The preface says, "Since growing old I have obtained Dongyuan's theory of the nine containments (jiurong)", and in this question it is further made clear that it is the 13th problem of Dongyuan's Circle-Measuring Section, and that the preceding answer also followed his detailed working. On the whole, the making of this book consisted entirely in taking his ideas as a model and developing them. Now, although Dongyuan as a person and his book can no longer be investigated, yet, looking at this one working alone, his choice of path is far-reaching and deep, and he goes only where change leads; here too one can see one spot of the patterned leopard. As for saying that "its numbers are tedious and fragmentary, roundabout and laborious", this was said merely with regard to difficulty and ease for beginners; readers should understand this well.\end{uniten}

\begin{unitzh}\unitlabel{ceyuan-031.6}測圓海鏡卷十一\end{unitzh}
\begin{uniten}Sea Mirror of Circle Measurements (Ceyuan haijing), juan 11\end{uniten}

\section*{卷12}\addcontentsline{toc}{section}{卷12}

\begin{unitzh}\unitlabel{ceyuan-032.1}欽定四庫全書 測圓海鏡卷十二 元 李冶 撰\end{unitzh}
\begin{uniten}Imperially Commissioned Complete Library of the Four Treasuries (Qinding Siku quanshu). Sea Mirror of Circle Measurements, juan 12. Composed by Li Ye of the Yuan.\end{uniten}

\begin{unitzh}\unitlabel{ceyuan-032.2}之分一十四問\end{unitzh}
\begin{uniten}Fractional parts (zhi fen): 14 problems.\end{uniten}

\begin{unitzh}\unitlabel{ceyuan-032.3}或問：甲、乙二人俱在西北隅。乙向直東行，不知步數而止；甲向直南行，望見乙，復向乙斜行。甲告乙云：「我直行、斜行共一千二百八十步，汝東行步居我南行步十五分之八。」\end{unitzh}
\begin{uniten}Someone asks: Jia and Yi, two people, are both at the northwest corner. Yi walks due east, not knowing the number of paces (bu 步), and stops; Jia walks due south, catches sight of Yi, and then walks obliquely toward Yi. Jia tells Yi: "My straight walk and oblique walk together are 1280 bu; your eastward paces amount to 8/15 of my southward paces."\end{uniten}

\begin{unitzh}\unitlabel{ceyuan-032.4}法曰：十六之共步冪為實，二百五十七之共步為益從，一十六步常法，得勾圓差。\end{unitzh}
\begin{uniten}Rule: 16 times the area (mi) of the total paces is the dividend (shi); 257 times the total paces is the subtractive linear term (yi cong); 16 bu is the constant divisor (chang fa); one obtains the base-circle difference (gou yuan cha).\end{uniten}

\begin{unitzh}\unitlabel{ceyuan-032.5}草曰：别得共步即股弦共也。立天元一為小差，以乘共步，得□，為勾冪；就分以二百二十五通之，得□，為二百二十五段勾冪{\zhsmall 寄左}。然後再置共步，内減小差，得□□，為二股；就分四之，得□□，為一十五勾；以自之，得□□□，為同數。與左相消，得□□□，平方開之，得八十步，即小差也。既得小差，加共步而半之，得六百八十步，即弦也；若以減共步而半之，得六百步，即股也。以股冪減弦冪，餘一十萬二千四百步，開平方，得三百二十步，即勾也。勾股相乘，倍之，得三十八萬四千步，為實；以弦和和一千六百步為法；實如法而一，得二百四十步，即城徑也。合問。\end{unitzh}
\begin{uniten}Working: It is separately obtained that the total paces are the altitude and the hypotenuse combined (gu xian gong). Set up the celestial unknown (tianyuan) as the small difference (xiao cha); multiply the total paces by it, obtaining □ [Translator's note: here and below, □ marks part of a counting-rod array not reproduced in the transcription.] as the area of the base. According to the parts (jiu fen), multiply it through (tong) by 225, obtaining □ as 225 pieces (duan) of the base area [Note: set aside on the left (ji zuo).] Then again set up the total paces and subtract the small difference from it, obtaining □□ as 2 altitudes; according to the parts, multiply by 4, obtaining □□ as 15 bases; square it, obtaining □□□ as the equal quantity (tong shu). Eliminate it against the left (xiang xiao), obtaining □□□; extract the square root, obtaining 80 bu, which is the small difference. Having obtained the small difference, add it to the total paces and halve, obtaining 680 bu, which is the hypotenuse; if one subtracts it from the total paces and halves, one obtains 600 bu, which is the altitude. Subtract the altitude area from the hypotenuse area; the remainder is 102,400 bu; extract the square root, obtaining 320 bu, which is the base. Multiply base and altitude together and double, obtaining 384,000 bu as the dividend; take the sum of the hypotenuse and the sum (xian he he), 1600 bu, as the divisor (fa); divide the dividend by the divisor, obtaining 240 bu, which is the city's diameter. This agrees with the question.\end{uniten}

\begin{unitzh}\unitlabel{ceyuan-032.6}或問：甲、乙二人俱在西北隅。乙直南行，不知步數而立；甲直東行，望見乙，復向乙斜行，與乙相㑹。甲云：「我共行了一千步。」又云：「我東行步居汝南行步十五分之八。」\end{unitzh}
\begin{uniten}Someone asks: Jia and Yi, two people, are both at the northwest corner. Yi walks due south, not knowing the number of paces, and stands; Jia walks due east, catches sight of Yi, then walks obliquely toward Yi and meets Yi. Jia says: "Altogether I have walked 1000 bu." He further says: "My eastward paces amount to 8/15 of your southward paces."\end{uniten}

\begin{unitzh}\unitlabel{ceyuan-032.7}法曰：二百二十五段共步冪為實，七百六之共步為益從，二百二十五步常法，得股圓差。\end{unitzh}
\begin{uniten}Rule: 225 pieces of the area of the total paces are the dividend; 706 times the total paces is the subtractive linear term; 225 bu is the constant divisor; one obtains the altitude-circle difference (gu yuan cha).\end{uniten}

\begin{unitzh}\unitlabel{ceyuan-032.8}草曰：别得共步即勾弦共也。立天元一為大差，以乘共步，得□；又就分以二百五十六通之，得□，為二百五十六个股冪{\zhsmall 寄左}。然後再置共步，内減天元大差，得□□，為二勾；就分以一十五之，得□□，為十六个股也；以自之，得□□{\zhsmall □□}，為同數。與左相消，得□□□，開平方，得三百六十，即大差也。副置共步：上位減大差而半之，得三百二十步，即勾也；下位加大差而半之，得六百八十步，即弦也。餘數各依法求之。合問。\end{unitzh}
\begin{uniten}Working: It is separately obtained that the total paces are the base and the hypotenuse combined (gou xian gong). Set up the celestial unknown as the large difference (da cha); multiply the total paces by it, obtaining □; further, according to the parts, multiply it through by 256, obtaining □ as 256 altitude areas [Note: set aside on the left.] Then again set up the total paces and subtract the celestial-unknown large difference from it, obtaining □□ as 2 bases; according to the parts, multiply by 15, obtaining □□ as 16 altitudes; square it, obtaining □□ [Note: □□] as the equal quantity. Eliminate it against the left, obtaining □□□; extract the square root, obtaining 360, which is the large difference. Set up the total paces in two positions (fu zhi): from the upper position subtract the large difference and halve, obtaining 320 bu, which is the base; to the lower position add the large difference and halve, obtaining 680 bu, which is the hypotenuse. The remaining numbers are each found according to the methods. This agrees with the question.\end{uniten}

\begin{unitzh}\unitlabel{ceyuan-032.9}或問：甲、乙俱在城西北隅。甲南行，不知步數而立；乙東行，亦不知步數，望見甲，就甲斜行，與之相㑹。乙云：「我東步少于城周九分之五。」甲云：「我南行却多于汝東行二百八十步。」問答同前。\end{unitzh}
\begin{uniten}Someone asks: Jia and Yi are both at the northwest corner of the city. Jia walks south, not knowing the number of paces, and stands; Yi walks east, also not knowing the number of paces, catches sight of Jia, walks obliquely to Jia and meets him. Yi says: "My eastward paces are less than the city's circumference by 5/9." Jia says: "My southward walk, however, is 280 bu more than your eastward walk." The question and the answer are the same as before.\end{uniten}

\begin{unitzh}\unitlabel{ceyuan-032.10}法曰：别得周居九分，徑居三分，乙東行居四分。{\zhsmall 按：此法未詳，當加「倍較步為實，徑分數自之，内減二分數為法，得數三之，即城徑」二十四字。}\end{unitzh}
\begin{uniten}Rule: It is separately obtained that the circumference amounts to 9 parts, the diameter amounts to 3 parts, and Yi's eastward walk amounts to 4 parts. [Siku editors: This rule is not detailed. One should add the 24 characters: "Double the difference paces (jiao bu) as the dividend; square the number of parts of the diameter, subtract from it the number for 2 parts, as the divisor; triple the number obtained, and that is the city's diameter."]\end{uniten}

\begin{unitzh}\unitlabel{ceyuan-032.11}草曰：立天元一為一分之數，以三之，得□，為徑；以四之，得□，為勾。以徑減勾，餘□，為小差{\zhsmall 只天元便是小差}。再置小差，加入甲多步，得□□，為大差。倍大差，以天元乘之，得□□，為一段圓徑冪{\zhsmall 寄左}。再置城徑，以自之，得下式□□，為同數。與左相消，得□□，上法下實，得八十步，即一分之數也。以三之，得二百四十步，即城徑也。合問。\end{unitzh}
\begin{uniten}Working: Set up the celestial unknown as the number of one part; multiply it by 3, obtaining □ as the diameter; multiply it by 4, obtaining □ as the base. Subtract the diameter from the base; the remainder, □, is the small difference [Note: the celestial unknown by itself is just the small difference.] Again set up the small difference and add in Jia's "more" paces, obtaining □□ as the large difference. Double the large difference and multiply it by the celestial unknown, obtaining □□ as one piece of the area of the circle's diameter [Note: set aside on the left.] Again set up the city's diameter and square it, obtaining the form below (xia shi), □□, as the equal quantity. Eliminate it against the left, obtaining □□: the upper is the divisor, the lower the dividend; one obtains 80 bu, which is the number of one part. Multiply it by 3, obtaining 240 bu, which is the city's diameter. This agrees with the question.\end{uniten}

\begin{unitzh}\unitlabel{ceyuan-032.12}或問：甲出西門南行，不知步數而立；乙出北門東行，望見甲。既而乙云：「我所行居城徑六分之五。」甲云：「然則我所行却多于汝二百八十步。」問答同前。\end{unitzh}
\begin{uniten}Someone asks: Jia goes out of the west gate and walks south, not knowing the number of paces, and stands; Yi goes out of the north gate and walks east and catches sight of Jia. Thereupon Yi says: "What I have walked amounts to 5/6 of the city's diameter." Jia says: "In that case, what I have walked is, however, 280 bu more than yours." The question and the answer are the same as before.\end{uniten}

\begin{unitzh}\unitlabel{ceyuan-032.13}法曰：四之却多步為實；分自之于上，半分母減子，得數倍之，又以減數乘之，減上位，為法；得一分之數。\end{unitzh}
\begin{uniten}Rule: 4 times the "more" paces (que duo bu) is the dividend. Square the parts, putting this on top; subtract half the denominator from the numerator, double the number obtained, further multiply it by the number from the subtraction, and subtract this from the upper position, as the divisor; one obtains the number of one part.\end{uniten}

\begin{unitzh}\unitlabel{ceyuan-032.14}草曰：别得却多步即勾股差也。乃立天元一為一分數，以六之為城徑，以五之為乙行。置乙行，内減半城徑，得□，為小差也；又加入却多步，得□□；又二之，得□□，為二大差；又以小差乘之，得□□，為徑冪{\zhsmall 寄左}。然後以徑冪□□與左相消，得下□□，上法下實，得四十步，即一分之數也。六之則為城徑，五之則為乙行，又以却多步加乙行，即甲行步也。合問。\end{unitzh}
\begin{uniten}Working: It is separately obtained that the "more" paces are the difference between base and altitude (gou gu cha). Then set up the celestial unknown as the number of one part; multiply it by 6 as the city's diameter, and by 5 as Yi's walk. Set up Yi's walk and subtract half the city's diameter from it, obtaining □ as the small difference; further add in the "more" paces, obtaining □□; further multiply by 2, obtaining □□ as 2 large differences; further multiply by the small difference, obtaining □□ as the area of the diameter [Note: set aside on the left.] Then eliminate the area of the diameter, □□, against the left, obtaining the following □□: the upper is the divisor, the lower the dividend; one obtains 40 bu, which is the number of one part. Multiplied by 6, it is then the city's diameter; multiplied by 5, it is then Yi's walk; further add the "more" paces to Yi's walk, and that is Jia's paces walked. This agrees with the question.\end{uniten}

\begin{unitzh}\unitlabel{ceyuan-032.15}或問：甲、丙二人俱在西北隅。甲向東行，不知步數而立；丙向南行，望見甲，與之相㑹。丙語甲云：「我行既多于汝，又城徑少于我四十分之十六。」{\zhsmall 按：四十為股分，十六為徑，當云「徑少於我為四十分之十六」，原文脫「為」字，似十六為股圓差分矣。}甲云：「然則吾二人共行了九百二十步。」問答同前。\end{unitzh}
\begin{uniten}Someone asks: Jia and Bing, two people, are both at the northwest corner. Jia heads east, not knowing the number of paces, and stands; Bing heads south, catches sight of Jia, and meets him. Bing says to Jia: "My walk is more than yours, and moreover the city's diameter is less than mine by 16/40." [Siku editors: 40 is the parts of the altitude and 16 is the diameter. It should say: "the diameter, less than mine, is 16/40". The original text has dropped the character 為 (wei, "is"), so that it seems that 16 were the parts of the altitude-circle difference.] Jia says: "In that case we two have walked 920 bu together." The question and the answer are the same as before.\end{uniten}

\begin{unitzh}\unitlabel{ceyuan-032.16}法曰：倍子以減倍母，又乘共行步，為實；倍子減倍母，以乘子母併數，于上，又以子冪加上位，為法；如法，得一十五步，即一分之數也。\end{unitzh}
\begin{uniten}Rule: Subtract double the numerator from double the denominator, and further multiply by the total paces walked, as the dividend. Subtract double the numerator from double the denominator and multiply the sum of the numerator and denominator by it, putting this on top; further add the area of the numerator to the upper position, as the divisor. Dividing by the divisor, one obtains 15 bu, which is the number of one part.\end{uniten}

\begin{unitzh}\unitlabel{ceyuan-032.17}草曰：别得共行步即通和也。又别得四十分之十六，或作二十分之八，或作十分之四，亦得，但所得分數不同耳。乃立天元一為一分之數，以十六之為城徑，以四十之為丙行。丙行減和步，得□□，為通勾；勾内減徑，餘得□□，為小差，于上。以分母、分子相減，餘□；又倍之，得□，為兩个大差；以乘上位，得□□，為圓徑冪{\zhsmall 寄左}。然後以分子十六分自之，得下□□，與左相消，得□□，上法下實，得一十五步，即一分之數也。以十六之，得二百四十步，即城徑也。合問。\end{unitzh}
\begin{uniten}Working: It is separately obtained that the total paces walked are the overall sum (tong he). It is further separately obtained that 16/40 may also be made 8/20, or made 4/10, and it also works; only the numbers of parts obtained differ. Then set up the celestial unknown as the number of one part; multiply it by 16 as the city's diameter, and by 40 as Bing's walk. Subtract Bing's walk from the sum paces, obtaining □□ as the overall base (tong gou); subtract the diameter from the base; the remainder obtained, □□, is the small difference, on top. Subtract the denominator and the numerator of the parts from each other, the remainder □; further double it, obtaining □ as two large differences; multiply the upper position by it, obtaining □□ as the area of the circle's diameter [Note: set aside on the left.] Then square the 16 parts of the numerator, obtaining the following □□; eliminate it against the left, obtaining □□: the upper is the divisor, the lower the dividend; one obtains 15 bu, which is the number of one part. Multiply it by 16, obtaining 240 bu, which is the city's diameter. This agrees with the question.\end{uniten}

\begin{unitzh}\unitlabel{ceyuan-032.18}或問：甲、乙俱立于城中心。乙出東門直行，不知步數而立；甲出南門直行，亦不知步數，望見乙，向乙斜行，與之相㑹。乙云：「我居汝南行十五分之八。」又云：「斜行步内若減甲直行，餘三十四步；若減乙直行，餘一百五十三步。」問答同前。\end{unitzh}
\begin{uniten}Someone asks: Jia and Yi both stand at the center of the city. Yi goes out of the east gate and walks straight on, not knowing the number of paces, and stands; Jia goes out of the south gate and walks straight on, also not knowing the number of paces, catches sight of Yi, walks obliquely toward Yi and meets him. Yi says: "Mine amounts to 8/15 of your southward walk." He further says: "If from the oblique paces one subtracts Jia's straight walk, 34 bu remain; if one subtracts Yi's straight walk, 153 bu remain." The question and the answer are the same as before.\end{uniten}

\begin{unitzh}\unitlabel{ceyuan-032.19}法曰：以云數二減步為小差、大差，以相乘，倍之，開平方，加入大小差併，以自之於上；又以大小差相較數以自之，減上位，為實；甲行分、乙行分相乘，又倍之，為隅法；得一分之數。\end{unitzh}
\begin{uniten}Rule: Take the stated numbers, the two paces from the subtractions, as the small difference and the large difference; multiply them together, double, extract the square root, add in the sum of the large and small differences, and square it, putting this on top. Further take the difference between the large and small differences, square it, and subtract it from the upper position, as the dividend. Multiply the parts of Jia's walk and the parts of Yi's walk together and further double, as the corner divisor (yu fa); one obtains the number of one part.\end{uniten}

\begin{unitzh}\unitlabel{ceyuan-032.20}草曰：别得云步相併，得一百八十七，是于皇極弦内少一个皇極黄方靣也。又别得三十四步是个小勾圓差，其一百五十三步是一个小股圓差，此二差又相減，餘一百一十九，即中差也。乃立天元一為一分之數，以八之，得□，為乙東行數；以十五之，得□，為甲南行數；以二數相乘，又倍之，得□□，為二直積，于上{\zhsmall 寄左}。然後以云步三十四乘一百五十三，得五千二百二；又倍之，得一萬四百四，為平方實，開之，得一百二步，即小黄方也；加入相併數一百八十七，得二百八十九，為小弦也；以自之，得八萬三千五百二十一，為弦冪，于上；以中差冪一萬四千一百六十一減上位，餘□，與左相消，得□□□，平方開之，得一十七步，即一分之數也。副置一分之數：上位以八之，得一百三十六，即乙東行也；下位以十五之，得二百五十五，即甲東行也。二位相乘，得三萬四千六百八十；又倍之，得六萬九千三百六十，為實；以弦二百八十九為法；如法，得二百四十步，即城徑也。合問。\end{unitzh}
\begin{uniten}Working: It is separately obtained that the stated paces added together give 187; this is the hypotenuse of the August Pole (huangji) less one side (mian) of the August Pole yellow square (huang fang). It is further separately obtained that the 34 bu is a small base-circle difference, and that the 153 bu is one small altitude-circle difference; these two differences, further subtracted from each other, leave 119, which is the middle difference (zhong cha). Then set up the celestial unknown as the number of one part; multiply it by 8, obtaining □ as the number of Yi's eastward walk; multiply it by 15, obtaining □ as the number of Jia's southward walk; multiply the two numbers together and further double, obtaining □□ as 2 rectangle areas (zhi ji), on top [Note: set aside on the left.] Then multiply 153 by the stated paces 34, obtaining 5202; further double it, obtaining 10,404 as the dividend for the square root (ping fang shi); extract it, obtaining 102 bu, which is the small yellow square; add in the combined number 187, obtaining 289 as the small hypotenuse; square it, obtaining 83,521 as the hypotenuse area, on top; subtract the area of the middle difference, 14,161, from the upper position, the remainder □; eliminate it against the left, obtaining □□□; extract the square root, obtaining 17 bu, which is the number of one part. Set up the number of one part in two positions: the upper position multiplied by 8 gives 136, which is Yi's eastward walk; the lower position multiplied by 15 gives 255, which is Jia's eastward walk [Translator's note: so the text, although Jia walked south]. Multiply the two positions together, obtaining 34,680; further double it, obtaining 69,360 as the dividend; take the hypotenuse, 289, as the divisor; dividing by the divisor, one obtains 240 bu, which is the city's diameter. This agrees with the question.\end{uniten}

\begin{unitzh}\unitlabel{ceyuan-032.21}或問：甲出西門南行，乙出北門東行，各不知逺近，兩相望見，復相斜行，各行了三百四十步相㑹。甲云：「城徑居我南行二分之一。」乙云：「我東行居城徑六分之五。」問答同前。\end{unitzh}
\begin{uniten}Someone asks: Jia goes out of the west gate and walks south; Yi goes out of the north gate and walks east; each does not know how far. The two catch sight of each other and then walk obliquely toward each other; each walks 340 bu and they meet. Jia says: "The city's diameter amounts to 1/2 of my southward walk." Yi says: "My eastward walk amounts to 5/6 of the city's diameter." The question and the answer are the same as before.\end{uniten}

\begin{unitzh}\unitlabel{ceyuan-032.22}法曰：以二之斜行步自之為實；以各行分數自之為冪{\zhsmall 按：此語未詳，當云「以城徑六分乘甲南行二分，得十二分，加半城徑三分，得十五分，為大股分；乙東行五分，加半城徑三分，得八分，為大勾分；各自之為冪」。}又相併為隅法；開平方，得一分之數。\end{unitzh}
\begin{uniten}Rule: Square twice the oblique paces as the dividend; square the numbers of parts of each walk as areas [Siku editors: This wording is not detailed. It should say: "Multiply the 2 parts of Jia's southward walk by the 6 parts of the city's diameter, obtaining 12 parts; add the 3 parts of half the city's diameter, obtaining 15 parts, as the parts of the large altitude (da gu); to the 5 parts of Yi's eastward walk add the 3 parts of half the city's diameter, obtaining 8 parts, as the parts of the large base (da gou); square each as areas."] and further add them together as the corner divisor; extract the square root, obtaining the number of one part.\end{uniten}

\begin{unitzh}\unitlabel{ceyuan-032.23}草曰：别得倍斜行為大弦。又别得乙行五分，城徑六分，甲行十二分。乃立天元一為一分之數，以六之，得□，為城徑；以五之，得□，為乙行分；以十二之，得□，為甲行分。乃副置半城徑：上位加甲行步，得□，以自之，得□□，為甲行冪；下位加乙行步，得□，以自之，得□□，為乙行冪。二冪又相併，得□□，為大弦冪{\zhsmall 寄左}。然後置大弦六百八十步，以自之，得□，與左相消，得□□□，平方開之，得四十步，即一分之數也。以六之，得二百四十步，即城徑也。合問。\end{unitzh}
\begin{uniten}Working: It is separately obtained that double the oblique walk is the large hypotenuse (da xian). It is further separately obtained that Yi's walk is 5 parts, the city's diameter 6 parts, and Jia's walk 12 parts. Then set up the celestial unknown as the number of one part; multiply it by 6, obtaining □ as the city's diameter; multiply it by 5, obtaining □ as the parts of Yi's walk; multiply it by 12, obtaining □ as the parts of Jia's walk. Then set up half the city's diameter in two positions: to the upper position add Jia's walk paces, obtaining □; square it, obtaining □□ as the area of Jia's walk; to the lower position add Yi's walk paces, obtaining □; square it, obtaining □□ as the area of Yi's walk. The two areas, further added together, give □□ as the area of the large hypotenuse [Note: set aside on the left.] Then set up the large hypotenuse, 680 bu, and square it, obtaining □; eliminate it against the left, obtaining □□□; extract the square root, obtaining 40 bu, which is the number of one part. Multiply it by 6, obtaining 240 bu, which is the city's diameter. This agrees with the question.\end{uniten}

\begin{unitzh}\unitlabel{ceyuan-032.24}或問：甲出西門南行，不知步數而立；乙出北門東行，見之，乙斜行與甲相㑹。甲、乙二人共行了一千三百六十步，其甲南行居斜十七分之十二，其乙東行居斜十七分之五。問答同前。\end{unitzh}
\begin{uniten}Someone asks: Jia goes out of the west gate and walks south, not knowing the number of paces, and stands; Yi goes out of the north gate and walks east, sees him, and Yi walks obliquely to meet Jia. Jia and Yi, the two people, have walked 1360 bu together; Jia's southward walk amounts to 12/17 of the oblique, and Yi's eastward walk amounts to 5/17 of the oblique. The question and the answer are the same as before.\end{uniten}

\begin{unitzh}\unitlabel{ceyuan-032.25}法曰：别得共步即二弦也。半共步，得六百八十步，副置：上位以五之，得三千四百，以十七而一，得二百步，即乙東行也；下位以十二之，得八萬一千六百，以十七而一，得四百八十，即甲南行也。二行相減，餘二百八十，即勾股差也。其餘各依法求之。合問。\end{unitzh}
\begin{uniten}Rule: It is separately obtained that the total paces are 2 hypotenuses. Halve the total paces, obtaining 680 bu; set it up in two positions: the upper position multiplied by 5 gives 3400; divided by 17, this gives 200 bu, which is Yi's eastward walk; the lower position multiplied by 12 gives 81,600 [Translator's note: 680 × 12 is 8160]; divided by 17, this gives 480, which is Jia's southward walk. Subtract the two walks from each other; the remainder, 280, is the difference between base and altitude. The rest are each found according to the methods. This agrees with the question.\end{uniten}

\begin{unitzh}\unitlabel{ceyuan-033.1}或問：甲出西門南行，不知步數而立；乙出北門東行，望見之。既而乙謂甲云：「我取汝六分之五，得六百步。」甲謂乙云：「我取汝五分之三，亦得六百步。」問答同前。\end{unitzh}
\begin{uniten}Someone asks: Jia goes out of the west gate and walks south an unknown number of bu (步), and stands; Yi goes out of the north gate, walks east, and sees him. Then Yi says to Jia: "Taking 5/6 of yours, I obtain 600 bu." Jia says to Yi: "Taking 3/5 of yours, I also obtain 600 bu." The question and answer are the same as before.\end{uniten}

\begin{unitzh}\unitlabel{ceyuan-033.2}法曰：求得各行步{\zhsmall 按見後草}，相併，以自之于上；併甲南行冪、乙東行冪，以減上，為實；併各行為從；半步常法；得全徑。\end{unitzh}
\begin{uniten}Rule: Find the bu walked by each [Siku editors: see the working below]. Add them together and multiply by itself, [placing it] above. Add together the area (mi) of Jia's southward walk and the area of Yi's eastward walk, and subtract this from the upper; this is the dividend (shi). Add together the walks of each as the follower (cong). Half a bu is the constant divisor (chang fa). Obtain the full diameter (quan jing).\end{uniten}

\begin{unitzh}\unitlabel{ceyuan-033.3}草曰：置{\zhsmall 乙取甲六分之五，六百步；甲取乙五分之三，六百步}。以上六分、五分各自直乗步數訖，得人{\zhsmall 六分之五，三千六百步；五分之三，三千步}。别得左行三千六百步，為六乙行、五甲行也；右行三千步，為五甲行、三乙行也。以方程法入之，乃再置{\zhsmall 五甲行，六乙行，三千六百步；五甲行，三乙行，三千步}。先以左行直減右行：右上空，中餘三乙行，下餘六百步，上法下實，得二百步，即乙行也。却以今右行減于元左行：上餘五甲行，空中，下餘二千四百步，上法下實，得四百八十步，即甲行也。既得此數，乃立天元一為城徑，以半之，副置二位：上以加甲行，得□□，為通股；以自之，得□□□，為大股冪。下位加乙行，得□□，為通勾；以自之，得□□□，為大勾冪。二冪相併，得□□□，為大弦冪{\zhsmall 寄左}。乃併甲行、乙行，以自乗，得下式□，亦為大弦冪。與左相消，得下□□□，開平方，得二百四十步，即城徑也。合問。\end{unitzh}
\begin{uniten}Working: Set out [Note: Yi taking 5/6 of Jia, 600 bu; Jia taking 3/5 of Yi, 600 bu]. With the 6 parts and 5 parts above, multiply each directly into the numbers of bu. When this is done, obtain ren (人) [Translator's note: 人 "person" here appears to be corrupt.] [Note: 5/6, 3600 bu; 3/5, 3000 bu]. Separately it is obtained that the left column (hang), 3600 bu, is 6 Yi-walks and 5 Jia-walks, and the right column, 3000 bu, is 5 Jia-walks and 3 Yi-walks. Enter this by the method of rectangular arrays (fangcheng), and set out again [Note: 5 Jia-walks, 6 Yi-walks, 3600 bu; 5 Jia-walks, 3 Yi-walks, 3000 bu]. First directly subtract the left column from the right column. On the right the top is empty, the middle leaves 3 Yi-walks, and the bottom leaves 600 bu. The upper is the divisor (fa) and the lower the dividend, giving 200 bu, which is Yi's walk. Then subtract the present right column from the original left column. The top leaves 5 Jia-walks, the middle is empty, and the bottom leaves 2400 bu. The upper is the divisor and the lower the dividend, giving 480 bu, which is Jia's walk. Having obtained these numbers, set up the celestial unknown (tianyuan) as the city diameter (cheng jing). Halve it, and set it out in duplicate in 2 places. To the upper add Jia's walk, obtaining □□ [Translator's note: □ evidently stands for a counting-rod array not reproduced in the transcription; so throughout.] as the through-altitude (tong gu). Multiply it by itself, obtaining □□□ as the area of the great altitude (da gu). To the lower place add Yi's walk, obtaining □□ as the through-base (tong gou). Multiply it by itself, obtaining □□□ as the area of the great base (da gou). Add the 2 areas together, obtaining □□□ as the area of the great hypotenuse (da xian) [Note: set aside on the left]. Then add Jia's walk and Yi's walk together and multiply by itself, obtaining the following expression □, which is also the area of the great hypotenuse. Eliminate it against the left (xiang xiao), obtaining the following □□□. Extract the square root, obtaining 240 bu, which is the city diameter. This agrees with the question.\end{uniten}

\begin{unitzh}\unitlabel{ceyuan-033.4}或問：甲從坤隅南行，不知步數而立；乙從艮隅東行，望見之。既而乙謂甲云：「我所行取汝所行三分之一，得二百步。」甲謂乙云：「我所行内減汝所行四分之三，得三百步。」問答同前。\end{unitzh}
\begin{uniten}Someone asks: Jia walks south from the Kun (坤) corner an unknown number of bu, and stands; Yi walks east from the Gen (艮) corner and sees him. Then Yi says to Jia: "What I walked, taking 1/3 of what you walked, makes 200 bu." Jia says to Yi: "What I walked, subtracting from it 3/4 of what you walked, makes 300 bu." The question and answer are the same as before.\end{uniten}

\begin{unitzh}\unitlabel{ceyuan-033.5}法曰：如法求得各行{\zhsmall 按見後草}，以相乗，又二之，開平方，得全徑。\end{unitzh}
\begin{uniten}Rule: Find each walk according to the method [Siku editors: see the working below]. Multiply them together, double it, and extract the square root, obtaining the full diameter.\end{uniten}

\begin{unitzh}\unitlabel{ceyuan-033.6}草曰：置{\zhsmall 乙取甲三分之一，二百步；甲減乙四分之三，三百步}。以上三分、四分置乗步數訖，得{\zhsmall 三分之一，六百步；四分之三，一千二百步}。别得右行六百步，為三乙行、一甲行也；左行一千二百步，為四甲行内少三之乙行步也。以方程法入之，乃再置{\zhsmall 一甲行，三乙行，六百步；四甲行，三乙行負，一千二百步}。先以左行直加右行：右上得五甲行，中空，下一千八百步，上法下實，得三百六十步，即甲行也。次以一甲行減元右行六百步，餘二百四十步，以中三除之，得八十步，即乙行步也。甲行、乙行二數相乘，得數，又倍之，開平方，即城徑也。合問。\end{unitzh}
\begin{uniten}Working: Set out [Note: Yi taking 1/3 of Jia, 200 bu; Jia subtracting 3/4 of Yi, 300 bu]. With the 3 parts and 4 parts above, set out and multiply the numbers of bu. When this is done, obtain [Note: 1/3, 600 bu; 3/4, 1200 bu]. Separately it is obtained that the right column, 600 bu, is 3 Yi-walks and 1 Jia-walk, and the left column, 1200 bu, is 4 Jia-walks less 3 times Yi's walk in bu. Enter this by the method of rectangular arrays, and set out again [Note: 1 Jia-walk, 3 Yi-walks, 600 bu; 4 Jia-walks, 3 Yi-walks negative, 1200 bu]. First directly add the left column to the right column. On the right the top gets 5 Jia-walks, the middle is empty, and the bottom is 1800 bu. The upper is the divisor and the lower the dividend, giving 360 bu, which is Jia's walk. Next subtract 1 Jia-walk from the 600 bu of the original right column, leaving 240 bu. Divide it by the 3 in the middle, obtaining 80 bu, which is Yi's walk in bu. Multiply the 2 numbers, Jia's walk and Yi's walk, together to obtain a number. Double it and extract the square root: this is the city diameter. This agrees with the question.\end{uniten}

\begin{unitzh}\unitlabel{ceyuan-033.7}或問：股圓差如股五分之三，勾圓差如勾四分之一。又云：其大、小差相減，餘二百八十步。問答同前。\end{unitzh}
\begin{uniten}Someone asks: the altitude–circle difference (gu yuan cha) is as 3/5 of the altitude (gu), and the base–circle difference (gou yuan cha) is as 1/4 of the base (gou). It is also said: subtracting the great and small differences (da cha, xiao cha) from each other leaves 280 bu. The question and answer are the same as before.\end{uniten}

\begin{unitzh}\unitlabel{ceyuan-033.8}法曰：二之中差為實；置股子，以勾母乗之，内減股母，為法；得小差。\end{unitzh}
\begin{uniten}Rule: Twice the middle difference (zhong cha) is the dividend. Set out the altitude numerator (gu zi), multiply it by the base denominator (gou mu), and subtract the altitude denominator (gu mu) from it; this is the divisor. Obtain the small difference.\end{uniten}

\begin{unitzh}\unitlabel{ceyuan-033.9}草曰：别得勾圓差即小差，股圓差即大差，云步即中差。乃立天元一為小差，以四之，得□，為勾；勾上加中差，得□□，為股；又三之，得□□，為五个大差也；内減五个天元，得□□，為五个中差也{\zhsmall 寄左}。乃以五之相減步□，與左相消，得□□，上法下實，得八十步，即小差也。合問。\end{unitzh}
\begin{uniten}Working: Separately it is obtained that the base–circle difference is the small difference, the altitude–circle difference is the great difference, and the stated bu are the middle difference. Then set up the celestial unknown as the small difference. Multiply it by 4, obtaining □ as the base. To the base add the middle difference, obtaining □□ as the altitude. Multiply further by 3, obtaining □□ as 5 great differences. Subtract 5 celestial unknowns from it, obtaining □□ as 5 middle differences [Note: set aside on the left]. Then take 5 times the bu of the subtraction, □, and eliminate it against the left, obtaining □□. The upper is the divisor and the lower the dividend, giving 80 bu, which is the small difference. This agrees with the question.\end{uniten}

\begin{unitzh}\unitlabel{ceyuan-033.10}或問：股圓差如股五分之三，勾圓差如勾四分之一。又云：勾母每分少于股母每分四十步。問答同前。\end{unitzh}
\begin{uniten}Someone asks: the altitude–circle difference is as 3/5 of the altitude, and the base–circle difference is as 1/4 of the base. It is also said: each part of the base denominator is less than each part of the altitude denominator by 40 bu. The question and answer are the same as before.\end{uniten}

\begin{unitzh}\unitlabel{ceyuan-033.11}法曰：二之少步，實；以股子母相減數減勾子母相減數，為法；如法得小差。\end{unitzh}
\begin{uniten}Rule: Twice the bu by which it is less is the dividend. Subtract the difference between the altitude numerator and denominator from the difference between the base numerator (gou zi) and denominator; this is the divisor. Dividing by the divisor, obtain the small difference.\end{uniten}

\begin{unitzh}\unitlabel{ceyuan-033.12}草曰：立天元一為勾圓差，便為勾母每分數。以天元加四十步，得□□，為股母每分數，于上。乃以股子減股母，餘二分，以乘上位，得□□，為城徑{\zhsmall 寄左}。再置天元在地，以勾子減勾母，餘三分，以乗之，得□□，為同數。與左相消，得下丨□，上法下實，得八十步，即勾圓差也。合問。\end{unitzh}
\begin{uniten}Working: Set up the celestial unknown as the base–circle difference, which is then the number of each part of the base denominator. Add 40 bu to the celestial unknown, obtaining □□ as the number of each part of the altitude denominator, [placed] above. Then subtract the altitude numerator from the altitude denominator, leaving 2 parts. Multiply the upper place by it, obtaining □□ as the city diameter [Note: set aside on the left]. Again set out the celestial unknown at the earth (di) [position]. Subtract the base numerator from the base denominator, leaving 3 parts, and multiply by it, obtaining □□ as the equal number (tong shu). Eliminate against the left, obtaining the following 丨□. The upper is the divisor and the lower the dividend, giving 80 bu, which is the base–circle difference. This agrees with the question.\end{uniten}

\begin{unitzh}\unitlabel{ceyuan-033.13}或問：甲出南門直行；乙出東門直行，望見甲，斜行與甲相㑹。甲云：「我行不及股圓差二十四分之十五。」乙云：「我行不及勾圓差五分之四。」又云：甲行多于乙行一百一十九，股圓差多于勾圓差二百八十。問答同前。\end{unitzh}
\begin{uniten}Someone asks: Jia goes out of the south gate and walks straight; Yi goes out of the east gate and walks straight, sees Jia, and walks obliquely to meet Jia. Jia says: "My walk falls short of the altitude–circle difference by 15/24." Yi says: "My walk falls short of the base–circle difference by 4/5." It is also said: Jia's walk exceeds Yi's walk by 119, and the altitude–circle difference exceeds the base–circle difference by 280. The question and answer are the same as before.\end{uniten}

\begin{unitzh}\unitlabel{ceyuan-033.14}法曰：以大差母分二十四，以乘甲多一百一十九，得數；倍小差母五，得一十，以乘之，于上。以小差母五乗二之二差相較數，又九之，減上位，為實。倍小差母，得一十，却以小差乗之，又九之，于上；倍甲分母，以小差母乗之，得數，減上位，以為法。得小差一分之數。\end{unitzh}
\begin{uniten}Rule: Take the parts of the great-difference denominator, 24, and with them multiply Jia's excess, 119, obtaining a number. Double the small-difference denominator, 5, obtaining 10, and multiply by it, [placing it] above. With the small-difference denominator, 5, multiply twice the difference between the 2 differences, and multiply further by 9. Subtract it from the upper place; this is the dividend. Double the small-difference denominator, obtaining 10. Then multiply it by the small difference, and further by 9, [placing it] above. Double Jia's denominator and multiply it by the small-difference denominator, obtaining a number. Subtract it from the upper place to make the divisor. Obtain the number of 1 part of the small difference.\end{uniten}

\begin{unitzh}\unitlabel{ceyuan-033.15}草曰：立天元一為小差一分之數{\zhsmall 此一分之數便是乙直行之數也}，以五之，得□，為小差；加二百八十，得下□□，為大差；又倍之，得□□；以小差乗之，得下式□□，為一个圓徑冪；又九之，得□□{\zhsmall 寄左}。乃又置乙行步，加一百一十九，□□，即甲行步也；以二十四之，得□□，為九个大差也；倍小差母，得□，以乘之，得□□，為同數。與左相消，得□□，上法下實，得一十六步，即小差一分之數也。既得此數，餘各如法求之。合問。\end{unitzh}
\begin{uniten}Working: Set up the celestial unknown as the number of 1 part of the small difference [Note: this number of 1 part is precisely the number of Yi's straight walk]. Multiply it by 5, obtaining □ as the small difference. Add 280, obtaining the following □□ as the great difference. Double it, obtaining □□. Multiply it by the small difference, obtaining the following expression □□ as 1 area of the circle diameter (yuan jing mi). Multiply further by 9, obtaining □□ [Note: set aside on the left]. Then again set out Yi's walk in bu and add 119: □□, which is Jia's walk in bu. Multiply it by 24, obtaining □□ as 9 great differences. Double the small-difference denominator, obtaining □, and multiply by it, obtaining □□ as the equal number. Eliminate against the left, obtaining □□. The upper is the divisor and the lower the dividend, giving 16 bu, which is the number of 1 part of the small difference. Having obtained this number, find the rest each according to the method. This agrees with the question.\end{uniten}

\begin{unitzh}\unitlabel{ceyuan-033.16}或問：大勾、大股、大弦三事和一千六百步；以明勾除大股，得八步三分之一；以□股除大勾，得一十步三分之二；以虚勾、明勾相減，餘二十四步；以虚股、□股相減，餘六十步。問答同前。\end{unitzh}
\begin{uniten}Someone asks: the sum of the three things (san shi he), the great base, great altitude and great hypotenuse, is 1600 bu. Dividing the great altitude by the bright base (ming gou) gives 8 1/3 bu. Dividing the great base by the [illegible] altitude gives 10 2/3 bu. Subtracting the void base (xu gou) and the bright base from each other leaves 24 bu. Subtracting the void altitude (xu gu) and the [illegible] altitude from each other leaves 60 bu. The question and answer are the same as before.\end{uniten}

\begin{unitzh}\unitlabel{ceyuan-033.17}法曰：六十步加入大三事和，又三之，二而一，為實；併二云數分母分子，内減六步，為法；如法得□股。\end{unitzh}
\begin{uniten}Rule: Add 60 bu to the great sum of the three things (da san shi he), multiply further by 3, and divide by 2; this is the dividend. Add the denominators and numerators of the 2 stated numbers, and subtract 6 bu from it; this is the divisor. Dividing by the divisor, obtain the [illegible] altitude.\end{uniten}

\begin{unitzh}\unitlabel{ceyuan-033.18}草曰：别得六十步與二十四步二數相併而半之，得□，即明勾、□股差也，又為虚勾、虚股差也。若以二數直相減，即虚黄方也。其二十四步得二虚勾，即半徑也；其六十步得二□股，亦為半徑也。立天元一為□股，加差步，得□□，為明勾也；以乗八步三分之一，得□□，為大股也。以天元乗一十步三分之二，得□，為大勾也。勾股相併，得下□□，為大和也{\zhsmall 寄左}。然後四之天元，加入二之六十步，得□□，為小三事和；以小三事和加入大三事和，得□□，為二个大和也。合折半為大和了，又就三分之為前數；今不折半三因，但身外加五，得□□，為同數。與左相消，得□□，上法下實，得三十步，即□股也。四之□股，加入二之六十步，得二百四十步，即城徑也。合問。\end{unitzh}
\begin{uniten}Working: Separately it is obtained that the 2 numbers 60 bu and 24 bu, added together and halved, give □, which is the difference of the bright base and the [illegible] altitude, and is also the difference of the void base and the void altitude. If the 2 numbers are directly subtracted from each other, it is the void yellow square (xu huang fang). The 24 bu, gaining 2 void bases, is the half-diameter (ban jing); the 60 bu, gaining 2 [illegible] altitudes, is also the half-diameter. Set up the celestial unknown as the [illegible] altitude. Add the bu of the difference, obtaining □□ as the bright base. Multiply it by 8 1/3 bu, obtaining □□ as the great altitude. Multiply the celestial unknown by 10 2/3 bu, obtaining □ as the great base. Add base and altitude together, obtaining the following □□ as the great sum (da he) [Note: set aside on the left]. Then multiply the celestial unknown by 4 and add twice the 60 bu, obtaining □□ as the small sum of the three things (xiao san shi he). Add the small sum of the three things to the great sum of the three things, obtaining □□ as 2 great sums. It ought to be halved to make the great sum, and then turned into thirds to make the previous number. Now, without halving and tripling (san yin), simply "add five beyond the body" (shen wai jia wu) [Translator's note: i.e. increase the number by half], obtaining □□ as the equal number. Eliminate against the left, obtaining □□. The upper is the divisor and the lower the dividend, giving 30 bu, which is the [illegible] altitude. Multiply the [illegible] altitude by 4 and add twice the 60 bu, obtaining 240 bu, which is the city diameter. This agrees with the question.\end{uniten}

\begin{unitzh}\unitlabel{ceyuan-033.19}按：之分即通分也。張邱建謂學者不患乘除之為難，而患通分之為難；又謂夏侯陽之方倉、孫子之蕩杯，皆未盡其妙；於是作為《算經》三卷，以發其義。是書末設十四問，皆以立天元一之法御之，尤為簡妙，殆所以明立天元一之法其用無不周也。又按：問中兩言「以方程入之」，《張邱建算經》内數問亦然。蓋有通分而乗除不窮，有方程而通分益便；此又因通分及之，非立天元一本法也。秦九韶謂時人誤以大衍法為方程者，蓋此類也。\end{unitzh}
\begin{uniten}Note: "Parts of" (zhi fen) is precisely the connecting of parts (tongfen). Zhang Qiujian said that students do not worry that multiplication and division are difficult, but worry that the connecting of parts is difficult. He also said that Xiahou Yang's "square granary" and Sunzi's "washing cups" both failed to exhaust its subtlety, and therefore he composed the Suanjing (Mathematical Classic) in 3 juan to bring out its meaning. This book sets 14 problems at its end and masters all of them by the method of setting up the celestial unknown, which is especially concise and subtle. This is probably meant to make clear that there is nothing the method of setting up the celestial unknown does not cover in its use. A further note: twice in the problems it says "enter it by rectangular arrays"; several problems in the Zhang Qiujian suanjing are also so. For with the connecting of parts, multiplication and division are inexhaustible, and with rectangular arrays the connecting of parts becomes all the more convenient. Here, again, it comes to them by way of the connecting of parts; it is not the basic method of setting up the celestial unknown. When Qin Jiushao says that people of his time mistakenly took the dayan (大衍) method for rectangular arrays, it is probably of this kind.\end{uniten}

\begin{unitzh}\unitlabel{ceyuan-033.20}按：右書十二卷，皆為立天元一法而作也。其法神明變化，不可端倪。今略舉數端言之：如諸法中有求之不可得者，此法求之可得；若此法求之不可得者，則必不可求矣。又諸法中有難求者，雖强探力索，毫釐未至，則不可得；此法但知大意，不待深思，加以步算，即可得矣。又諸法中有所求，或先得彼而後得此者，不能移易；此法任其所求，或先得此，或先得彼，無不如志。又諸法有數始可求，一數不具，則不可求；此法數不具亦可求，且有無數即可求者。又諸法遇甚繁甚密者，湏次第步算，或累日累月，其功不能再省；此法有經年步算可約之頃刻而得者。凡此皆尋常智慮所不能及，要皆自然之理數，易知易從；然自不習者觀之，蓋有茫然莫解其故者矣。是書之作，殆深憂𫝊習者難其人，而其法遂泯於後世也。其謄寫魯魚、算式舛訛，今悉正之。\end{unitzh}
\begin{uniten}Note: The 12 juan of the book above were all composed for the method of setting up the celestial unknown. Its method is divine and luminous in its transformations, and its bounds cannot be discerned. Now let us briefly raise a few points and speak of them. For instance, if among the various methods there are things that cannot be obtained by seeking, sought by this method they can be obtained; and if something cannot be obtained when sought by this method, then it certainly cannot be sought. Again, among the various methods there are things that are hard to find: although one probes forcefully and searches strenuously, if one falls short by a hao or a li (毫釐), one cannot obtain them. With this method one need only know the general idea, without waiting for deep thought, and by applying reckoning (bu suan) one can obtain them. Again, in the various methods, where in seeking something one first obtains that and afterwards obtains this, it cannot be shifted or altered. This method lets one seek as one pleases, first obtaining this or first obtaining that, and nothing fails to go as intended. Again, in the various methods there must be numbers before one can seek, and if 1 number is not provided, one cannot seek. With this method one can seek even when the numbers are not provided, and there are even things that can be sought with no numbers at all. Again, when the various methods meet something extremely complicated and extremely intricate, one must reckon in order, perhaps for days and months on end, and the labour cannot be reduced any further. With this method there are cases where a whole year of reckoning can be condensed and obtained in a moment. All these are beyond the reach of ordinary intelligence and deliberation; yet in essence they are all natural principles and numbers, easy to know and easy to follow. But seen by those who are not practised in it, there are probably some who are at a loss and cannot understand the reason. This book was probably made out of a deep worry that those who would transmit and practise it would be hard to find, and that its method would then vanish in later ages. Its copying errors ("lu" for "yu") and the mistakes in its computational expressions (suan shi) have now all been corrected.\end{uniten}

\begin{unitzh}\unitlabel{ceyuan-033.21}《測圓海鏡》卷十二\end{unitzh}
\begin{uniten}Sea Mirror of Circle Measurements, juan 12\end{uniten}

\begin{unitzh}\unitlabel{ceyuan-033.22}後序\end{unitzh}
\begin{uniten}Postface\end{uniten}

\begin{unitzh}\unitlabel{ceyuan-033.23}敬齋先生病且革，語其子克修曰：「吾平生著述，死後可盡燔去；獨《測圓海鏡》一書，雖九九小數，吾嘗精思致力焉，後世必有知者，庶可布廣垂永乎！」先生于六藝百家靡不貫串，文集近數百卷，常謙謙不自伐；惟于此書，不忘稱異于易簀之間，想有𤣥妙内得于心者。予以先生與先人同牓之故，素常兄事克修，克修兄命予重為序之。予不敢詭論艷藻，刻畫無鹽，唐突西子，直以所聞語意載之于後。至元二十四年春三月朔，翰林修撰、承直郎、廣平王德淵後序。\end{unitzh}
\begin{uniten}When Master Jingzhai was ill and near the end, he said to his son Kexiu: "The writings of my whole life may all be burned after my death. Only the one book, the Sea Mirror of Circle Measurements, is different: although it is the petty numbers of the nine-nines, I once applied refined thought and effort to it. Later ages will surely have those who understand it; may it, I hope, be spread widely and handed down forever!" The Master had thoroughly penetrated all of the Six Arts and the Hundred Schools, and his collected writings came to nearly several hundred juan, yet he was always modest and did not boast of himself. Only of this book did he not forget to speak as remarkable at the time of changing his mat [Translator's note: i.e. on his deathbed]; I suppose there was something profound and subtle that he had grasped inwardly in his mind. Because the Master and my late father were on the same [examination] list, I have always served Kexiu as an elder brother, and Brother Kexiu bade me write a preface for it anew. I do not dare to use outlandish argument and gaudy phrasing, carving and painting Wuyan and affronting Xizi; I simply set down after it the purport of the words I heard. On the 1st day of the 3rd month of spring in the 24th year of Zhiyuan, postface by Wang Deyuan of Guangping, Compiler of the Hanlin Academy and Gentleman for Rendering Service (chengzhi lang).\end{uniten}
\clearpage\section*{Sources}\addcontentsline{toc}{section}{Sources}
The Chinese text is Siku quanshu (四庫全書) edition, as transcribed on Chinese Wikisource. The text of the work is in the public domain; the Wikisource transcription is credited here with the pages and revisions used. The punctuation, the division into units and the English translation are AI-produced and dedicated to the public domain (CC0 1.0).

\begin{itemize}
\item \href{https://zh.wikisource.org/wiki/\%E6\%B8\%AC\%E5\%9C\%93\%E6\%B5\%B7\%E9\%8F\%A1_\%28\%E5\%9B\%9B\%E5\%BA\%AB\%E5\%85\%A8\%E6\%9B\%B8\%E6\%9C\%AC\%29}{測圓海鏡}, revision 619836
\item \href{https://zh.wikisource.org/wiki/\%E6\%B8\%AC\%E5\%9C\%93\%E6\%B5\%B7\%E9\%8F\%A1_\%28\%E5\%9B\%9B\%E5\%BA\%AB\%E5\%85\%A8\%E6\%9B\%B8\%E6\%9C\%AC\%29/\%E5\%8D\%B701}{卷01}, revision 788292
\item \href{https://zh.wikisource.org/wiki/\%E6\%B8\%AC\%E5\%9C\%93\%E6\%B5\%B7\%E9\%8F\%A1_\%28\%E5\%9B\%9B\%E5\%BA\%AB\%E5\%85\%A8\%E6\%9B\%B8\%E6\%9C\%AC\%29/\%E5\%8D\%B702}{卷02}, revision 788293
\item \href{https://zh.wikisource.org/wiki/\%E6\%B8\%AC\%E5\%9C\%93\%E6\%B5\%B7\%E9\%8F\%A1_\%28\%E5\%9B\%9B\%E5\%BA\%AB\%E5\%85\%A8\%E6\%9B\%B8\%E6\%9C\%AC\%29/\%E5\%8D\%B703}{卷03}, revision 788296
\item \href{https://zh.wikisource.org/wiki/\%E6\%B8\%AC\%E5\%9C\%93\%E6\%B5\%B7\%E9\%8F\%A1_\%28\%E5\%9B\%9B\%E5\%BA\%AB\%E5\%85\%A8\%E6\%9B\%B8\%E6\%9C\%AC\%29/\%E5\%8D\%B704}{卷04}, revision 788297
\item \href{https://zh.wikisource.org/wiki/\%E6\%B8\%AC\%E5\%9C\%93\%E6\%B5\%B7\%E9\%8F\%A1_\%28\%E5\%9B\%9B\%E5\%BA\%AB\%E5\%85\%A8\%E6\%9B\%B8\%E6\%9C\%AC\%29/\%E5\%8D\%B705}{卷05}, revision 788301
\item \href{https://zh.wikisource.org/wiki/\%E6\%B8\%AC\%E5\%9C\%93\%E6\%B5\%B7\%E9\%8F\%A1_\%28\%E5\%9B\%9B\%E5\%BA\%AB\%E5\%85\%A8\%E6\%9B\%B8\%E6\%9C\%AC\%29/\%E5\%8D\%B706}{卷06}, revision 788303
\item \href{https://zh.wikisource.org/wiki/\%E6\%B8\%AC\%E5\%9C\%93\%E6\%B5\%B7\%E9\%8F\%A1_\%28\%E5\%9B\%9B\%E5\%BA\%AB\%E5\%85\%A8\%E6\%9B\%B8\%E6\%9C\%AC\%29/\%E5\%8D\%B707}{卷07}, revision 788306
\item \href{https://zh.wikisource.org/wiki/\%E6\%B8\%AC\%E5\%9C\%93\%E6\%B5\%B7\%E9\%8F\%A1_\%28\%E5\%9B\%9B\%E5\%BA\%AB\%E5\%85\%A8\%E6\%9B\%B8\%E6\%9C\%AC\%29/\%E5\%8D\%B708}{卷08}, revision 788308
\item \href{https://zh.wikisource.org/wiki/\%E6\%B8\%AC\%E5\%9C\%93\%E6\%B5\%B7\%E9\%8F\%A1_\%28\%E5\%9B\%9B\%E5\%BA\%AB\%E5\%85\%A8\%E6\%9B\%B8\%E6\%9C\%AC\%29/\%E5\%8D\%B709}{卷09}, revision 788309
\item \href{https://zh.wikisource.org/wiki/\%E6\%B8\%AC\%E5\%9C\%93\%E6\%B5\%B7\%E9\%8F\%A1_\%28\%E5\%9B\%9B\%E5\%BA\%AB\%E5\%85\%A8\%E6\%9B\%B8\%E6\%9C\%AC\%29/\%E5\%8D\%B710}{卷10}, revision 788311
\item \href{https://zh.wikisource.org/wiki/\%E6\%B8\%AC\%E5\%9C\%93\%E6\%B5\%B7\%E9\%8F\%A1_\%28\%E5\%9B\%9B\%E5\%BA\%AB\%E5\%85\%A8\%E6\%9B\%B8\%E6\%9C\%AC\%29/\%E5\%8D\%B711}{卷11}, revision 788313
\item \href{https://zh.wikisource.org/wiki/\%E6\%B8\%AC\%E5\%9C\%93\%E6\%B5\%B7\%E9\%8F\%A1_\%28\%E5\%9B\%9B\%E5\%BA\%AB\%E5\%85\%A8\%E6\%9B\%B8\%E6\%9C\%AC\%29/\%E5\%8D\%B712}{卷12}, revision 788315
\end{itemize}
\end{document}
